
\documentclass[11pt]{article}

\usepackage[margin=1in]{geometry}
\usepackage{amsmath,amssymb,amsthm,mathtools,mathrsfs,bm}
\usepackage{enumitem}
\usepackage{booktabs,array,longtable,tabularx}
\usepackage{xcolor}
\usepackage{microtype}
\usepackage[colorlinks=true,linkcolor=blue,citecolor=blue,urlcolor=blue,hypertexnames=false]{hyperref}
\hypersetup{
  pdftitle={The Arithmetic Renewal Gran Tensor: Conditional Closure Mechanisms for GRH and Affine Prime Correlations},
  pdfauthor={A. P\'elissier},
  pdfsubject={Arithmetic renewal tensors, Dirichlet channels, and affine prime correlations},
  pdfkeywords={Renewal Gran tensor, generalized Riemann hypothesis, affine primes, twin primes, Mobius inversion}
}
\usepackage[nameinlink,capitalise,noabbrev]{cleveref}
\allowdisplaybreaks[3]
\setlength{\emergencystretch}{4em}
\raggedbottom
\hbadness=10000

\newtheorem{theorem}{Theorem}[section]
\newtheorem{proposition}[theorem]{Proposition}
\newtheorem{lemma}[theorem]{Lemma}
\newtheorem{corollary}[theorem]{Corollary}
\newtheorem{criterion}[theorem]{Criterion}
\newtheorem{obstruction}[theorem]{Obstruction}
\theoremstyle{definition}
\newtheorem{definition}[theorem]{Definition}
\newtheorem{construction}[theorem]{Construction}
\newtheorem{assumption}[theorem]{Assumption}
\newtheorem{principle}[theorem]{Principle}
\theoremstyle{remark}
\newtheorem{remark}[theorem]{Remark}

\newcommand{\N}{\mathbb N}
\newcommand{\Z}{\mathbb Z}
\newcommand{\R}{\mathbb R}
\newcommand{\C}{\mathbb C}
\newcommand{\F}{\mathbb F}
\newcommand{\T}{\mathbb T}
\newcommand{\Torus}{\mathbb T}
\newcommand{\one}{\mathbf 1}
\newcommand{\id}{\mathrm{id}}
\newcommand{\im}{\operatorname{Im}}
\newcommand{\E}{\mathbb E}
\newcommand{\Prob}{\mathbb P}
\newcommand{\Tr}{\operatorname{Tr}}
\newcommand{\rank}{\operatorname{rank}}
\newcommand{\Ker}{\operatorname{Ker}}
\newcommand{\Ran}{\operatorname{Ran}}
\newcommand{\Span}{\operatorname{span}}
\newcommand{\supp}{\operatorname{supp}}
\newcommand{\spec}{\operatorname{spec}}
\newcommand{\rad}{\operatorname{rad}}
\newcommand{\ord}{\operatorname{ord}}
\newcommand{\lcm}{\operatorname{lcm}}
\newcommand{\diag}{\operatorname{diag}}
\newcommand{\dist}{\operatorname{dist}}
\newcommand{\KL}{\operatorname{KL}}
\newcommand{\TV}{\operatorname{TV}}
\newcommand{\conv}{\operatorname{conv}}
\newcommand{\relint}{\operatorname{relint}}
\newcommand{\dd}{\mathrm d}
\newcommand{\ee}[1]{\exp(2\pi\mathrm i #1)}
\newcommand{\op}{\mathrm{op}}
\newcommand{\norm}[1]{\left\lVert#1\right\rVert}
\newcommand{\abs}[1]{\left\lvert#1\right\rvert}
\newcommand{\ip}[2]{\left\langle#1,#2\right\rangle}
\newcommand{\Aff}{\mathrm{Aff}}
\newcommand{\loc}{\mathrm{loc}}
\newcommand{\conn}{\mathrm{conn}}
\newcommand{\occ}{\mathrm{occ}}
\newcommand{\prodsec}{\mathrm{prod}}
\newcommand{\exc}{\mathrm{exc}}
\newcommand{\amp}{\mathrm{amp}}
\newcommand{\fib}{\mathrm{fib}}
\newcommand{\HL}{\mathrm{HL}}
\newcommand{\PP}{\mathrm{PP}}
\newcommand{\Short}{\operatorname{Short}}
\newcommand{\Phase}{\mathrm{ph}}
\newcommand{\Hard}{\mathrm{hard}}
\newcommand{\Rel}{\mathrm{rel}}
\newcommand{\Act}{\mathrm{act}}
\newcommand{\PhiRate}{\Phi}
\newcommand{\Gran}{\mathfrak G}
\newcommand{\Fesh}{\operatorname{Fesh}}
\newcommand{\Schur}{\operatorname{Schur}}
\newcommand{\Christ}{\operatorname{Chr}}
\newcolumntype{L}[1]{>{\raggedright\arraybackslash}p{#1}}

\title{\textbf{The Arithmetic Renewal Gran Tensor}\\[0.35em]
\Large Conditional Closure Mechanisms for the Generalized Riemann Hypothesis\\
and Affine Prime Correlations}
\author{A.~P\'elissier}
\date{12 September 2026}

\begin{document}
\maketitle

\begin{abstract}
We construct a finite, cutoff-compatible arithmetic realization of the
Renewal Gran tensor and use it to place conditional closures for Dirichlet
GRH and affine prime correlations on one common source geometry.  The tensor
retains finite same-history cylinders, complete mixed source maps,
simultaneous nuisance shorting, product and occurrence sectors, Feshbach
reduction, accepted/action likelihoods, and the exact old/detail calculus
which separates signed retuning from nonnegative carrier birth.  A
state-preserving score predictor controls the upper relative-metric defect
through an explicitly acquired mixed-Gram residual.  Divisibility
incidence gives the finite zeta operator, M\"obius inversion gives its inverse,
and $\Lambda=\mu*\log$ supplies the universal logarithmic compiler.  An exact
expansion--pullback theorem shows that factorization, determinant, and
provenance coordinates do not by themselves create a new endpoint source;
analytic input must therefore be applied only after the reader appropriate to
that theorem has been restored.

For a primitive Dirichlet character, a contracted finite phase action is
shown to be equivalent at the extensive scale to square-root cancellation in
twisted Mertens sums, so a source-fixed entropic descent with subpower action
implies GRH.  The outer M\"obius--log compiler gives a second route: the
coefficient of the complementary quotient $q=1$ is $\log1=0$, the direct
terminal Mertens coordinate disappears, the nonprincipal divisor head closes
at square-root scale, and conductor periodicity removes the terminal quotient
range.  A third exact entrance applies two-threshold M\"obius descent directly
to the reciprocal Dirichlet series.  Its high reader, full same-product
diagonal, and short product collisions below the cubic frontier close
unconditionally.  Equality-preserving double completion then resolves the
remaining four-M\"obius determinant current into a nonzero-frequency
reciprocal-phase channel and an exact zero-frequency sinc sheet; inverse-class
product compression and a primitive one-dimensional fibre kernel remove two
further coordinate artifacts.  A fourth, cutoff-independent entrance uses the
completed primitive theta kernel, its de Bruijn heat flow, and finite Jensen
polynomials.  Failure of GRH then returns either a finite multiple-root
collision with a signed heat curvature or a cofinal escape through increasing
Jensen degree or shift.  A fifth completed Euler--Weil entrance gives strict
local Schur cells, finite Hermite--Biehler spectra, and an exact
parity-renormalized archimedean form.  Its universal logarithmic core is
diagonal in the Legendre basis with harmonic-number eigenvalues, so every
fixed localized history has a coercive high-mode complement and only a finite
protected low-mode Schur obstruction.  Exact Legendre acquisition gives a
finite-data bracket for that obstruction, while Gaussian averaging of the
complete Weil form yields a scalar heat threshold whose vanishing is
equivalent to GRH.  Exhaustive Euler-region Pick tests detect every hidden
pole, but no fixed finite test controls the continuation.  Normalized Mellin
banks similarly restore spectral visibility without removing the conditioning
cost of translation tomography.  On the finite Jensen side, a universal
Hermite heat limit leads to sharp one-pair padding bounds, an algebraic
quartic optimum, and a conjugation-preserving Appell criterion for mixed
blocks.  Exact pairing expansions expose the nonmultiplicative
couplings, and equispaced growing Gaussian symbols have a logarithmically
translating critical point.  For actual spectral filters, finite
interpolation preserves a positive tail but incurs a target-normalized
coefficient cost.  A completed two-channel identity forces all uncancelled
terms to compensate a hypothetical off-line zero.  Exact residue formulas,
last-zero bounds, and rational sign certificates distinguish finite
admission from the unresolved growing-order transfer and cofinal signed
compensation.

For primitive affine pairs we prove the exact local root formula, singular
series convergence, Chinese-remainder tensorization, degree-two occurrence
closure, the one-dimensional neutral parity quotient, an exact $W$-fibre
amplitude decomposition, and prime-power cleanup.  For twins, endpoint fusion
turns the outer compiler into a target-free M\"obius--prime packet, while a
source-independent complement pressure has zero carrier birth for every
positive pressure.  Its prime-cleaned odd--odd current conditionally closes
the Hardy--Littlewood asymptotic.  Restoring the reflected prime reader yields
an additional analytic reduction: Maynard's fixed-residue theorem removes
absolutely every balanced critical modulus with a divisor in an explicit
convenient interval.  Each surviving hard modulus has a canonical prefix
below that interval and a rough cofactor of at most ten prime factors; the
prime-modulus base layer is a literal bilinear prime--prime occurrence.

The two branches are linked exactly rather than heuristically.  They share
logarithmic annihilation of the direct quotient, factorization-to-endpoint
fusion, the birth--retuning identity, and the requirement that a signed
protected read be controlled by an independently declared source rather than
by an ambient norm fitted after observing the endpoint.  They separate again
at their first analytic readers: fixed-character Mertens, reciprocal-phase,
completed-theta, or localized Weil currents for GRH, and prime-progression or
pressure currents for affine pairs.  All finite identities and all displayed
implications are proved; the remaining conjectural content is isolated in
explicit signed-current, protected-amplitude, reciprocal-phase, sinc-sheet,
low-mode Schur, backward relative-heat, first-prime continuation,
normalized spectral compensation, collision, or compactness certificates.
\end{abstract}

\tableofcontents
\clearpage


% =============================================================================
\section{Introduction and statement of the conditional closures}
\label{sec:introduction}
% =============================================================================

\subsection{Purpose and scope}

The Generalized Riemann Hypothesis and the prime-pair conjectures are usually
presented through different analytic languages.  The former is encoded by
zeros of Dirichlet $L$-functions and square-root cancellation in character
sums; the latter is encoded by singular series, parity, bilinear forms, and
asymptotic sieve estimates; standard background may be found in \cite{Davenport,IK,MV}.  The aim of this manuscript is not to identify
these conjectures, nor to claim that one implies the other.  It is to prove
that their strongest finite reductions admit a common arithmetic tensor
architecture and to state the remaining analytic estimates in that common
language.

The phrase \emph{Renewal Gran tensor} will mean the complete finite relational
object obtained by retaining, at every cutoff,
\begin{enumerate}[label=\textup{(G\arabic*)},leftmargin=3.2em]
\item the actual finite history cylinder and its positive reference law;
\item every jointly visible source map and every mixed Gram block;
\item simultaneous shorting of nuisance coordinates;
\item the product occurrence and the orthogonal occurrence innovation;
\item the accepted/action likelihood record;
\item the old/detail refinement map and its complete Schur data.
\end{enumerate}
The adjective \emph{complete} is essential.  Diagonal source energies do not
determine a protected signed coefficient, and marginal old/detail norms do not
determine carrier-birth observability.

The arithmetic loading is fixed before a terminal Mertens sum or prime-pair
residual is evaluated.  This order excludes two circular constructions:
choosing a positive action whose kernel is the desired endpoint, and adjoining
an endpoint coordinate only after its numerical value is known.

\subsection{Exact, standard-input, and conditional statements}

Three proof statuses occur.
\begin{enumerate}[label=\textup{(S\arabic*)},leftmargin=3.2em]
\item \emph{Exact finite statement}: a finite-dimensional identity, positivity
      assertion, tensor decomposition, or finite arithmetic rearrangement
      proved in the manuscript.
\item \emph{Standard analytic input}: a classical zero-free-region,
      Bombieri--Vinogradov, Weil-criterion, Tauberian, or published dispersion
      estimate used with its hypotheses stated explicitly.
\item \emph{Conditional closure}: a proved implication from a displayed
      quantitative certificate to GRH, a Hardy--Littlewood asymptotic, or a
      prime-pair existence statement.  The certificate itself is not asserted
      unless proved separately.
\end{enumerate}
No local factor, formal Schur complement, spectral gap, or positivity statement
is promoted into a conjectural asymptotic without the corresponding
quantitative hypothesis.

Throughout, when a nonnegative quantity occurring in an upper-bound hypothesis
is written as $F(X)=X^{\alpha+o(1)}$, this means that for every $\varepsilon>0$
one has $F(X)\le X^{\alpha+\varepsilon}$ for all sufficiently large $X$.
Two-sided order is denoted by $\asymp$; an asymptotic formula is always written
with an explicit main term and an $o(\cdot)$ or $O(\cdot)$ remainder.

\subsection{The GRH closures}

For a primitive character $\chi$, write
\[
 M_\chi(X)=\sum_{n\le X}\mu(n)\chi(n),
 \qquad
 \psi_\chi(X)=\sum_{n\le X}\chi(n)\Lambda(n),
\]
and let $\Psi_\chi(X)=\psi_\chi(X)-\delta_\chi X$, where
$\delta_\chi$ is the principal-channel indicator with the fixed missing Euler
factors incorporated in the local main term.

The first closure is action-theoretic.

\begin{theorem}[Source-fixed entropic GRH closure]
\label{thm:intro-entropic-GRH}
Suppose that for every primitive Dirichlet character $\chi$ and every dyadic
shell $(Y,2Y]$ there is a source-preserving descent from the ratio-complete
phase reference to the hard-counting phase law, constructed on one compatible
Renewal Gran cylinder, whose extensive relative-entropy action satisfies
\[
 A_{\chi,Y}\,\mathfrak D_{\chi,Y}=Y^{o(1)}.
\]
Then, for every $\varepsilon>0$,
\[
 M_\chi(X)=O_{\chi,\varepsilon}(X^{1/2+\varepsilon}),
\]
and GRH holds for every primitive Dirichlet $L$-function.
\end{theorem}

The proof is given in \Cref{sec:hard-phase,sec:entropic-descent}: entropy
contraction bounds the exact phase rate function, dyadic summation gives the
Mertens estimate, and the reciprocal Dirichlet series excludes zeros to the
right of the critical line.

The second closure is a concrete target-free compiler.  Fix $B>2$ and put
\[
 D_X=\left\lfloor X^{1/2}(\log X)^{-B}\right\rfloor,
 \qquad
 Q_X=\left\lfloor\frac{X}{D_X+1}\right\rfloor.
\]
For a nonprincipal primitive $\chi$, choose $E_X\to\infty$ with
$E_X=X^{o(1)}$ and set
\[
 R_X=\left\lfloor\frac{X^{1/2}}{E_X}\right\rfloor.
\]

\begin{theorem}[Target-free short-quotient GRH closure]
\label{thm:intro-short-GRH}
For every fixed nonprincipal primitive character,
\[
 \Psi_\chi(X)
 =\sum_{2\le q<R_X}\chi(q)\log q
   \bigl[M_\chi(\lfloor X/q\rfloor)-M_\chi(D_X)\bigr]
  +X^{1/2+o(1)}.
\]
Consequently, the estimate
\[
 \left|
 \sum_{2\le q<R_X}\chi(q)\log q
   \bigl[M_\chi(\lfloor X/q\rfloor)-M_\chi(D_X)\bigr]
 \right|
 \le X^{1/2+o(1)}
\]
implies GRH for $L(s,\chi)$.
\end{theorem}

The direct coordinate $M_\chi(X)$ does not occur: its complementary quotient
is $q=1$ and its coefficient is $\log1=0$.  For the principal channel, the
outer identity remains exact but the divisor head has an explicit profile;
full GRH from this route requires a square-root estimate for that profile or
an independent principal-channel closure.

The third closure is an exact two-M\"obius descent.  Let $K\ge2$, let
$E_K\to\infty$ with $E_K=K^{o(1)}$, and put
\[
 R_K=\left\lfloor K^{1/2}/E_K\right\rfloor,
 \qquad
 Y_K=\left\lfloor K/(R_K+1)\right\rfloor.
\]
For a nonprincipal primitive $\chi$, define
\[
 \mathcal Q_\chi(K;R_K,Y_K)
 =\sum_{m\le R_K}\sum_{n\le Y_K}
  \mu(m)\mu(n)\chi(mn)
  \sum_{k\le K/(mn)}\chi(k).
\]

\begin{theorem}[Protected rectangular GRH closure]
\label{thm:intro-rectangular-GRH}
One has
\[
 M_\chi(K)=M_\chi(R_K)+M_\chi(Y_K)
 -\mathcal Q_\chi(K;R_K,Y_K).
\]
The two linear terms are $K^{1/2+o(1)}$.  Hence
\[
 |\mathcal Q_\chi(K;R_K,Y_K)|\le K^{1/2+o(1)}
\]
implies GRH for $L(s,\chi)$.  Unconditionally, the high reader states, the
complete same-product diagonal, and short product collisions below the cubic
product scale are removed; the remaining certificate is the connected
four-M\"obius determinant current of \Cref{sec:grh-rectangular}.
\end{theorem}

The determinant frontier has a further exact source resolution.  The
unsquared low-reader corner is already
$-M_\chi(K)+K^{1/2+o(1)}$ and is therefore quarantined as a circular final
scalar.  On the protected square, equality-preserving completion introduces
an independent frequency $\nu$ and gives
\[
 \mathcal O_{\chi,K}
 =\mathcal O_{\chi,K}^{(0)}
  +\mathcal O_{\chi,K}^{(\ne0)}.
\]
The nonzero channel retains reciprocal-residue phases; the zero channel has
an exact sinc kernel with resolution width comparable to the larger reciprocal
modulus.

\begin{theorem}[Frequency-resolved rectangular GRH closure]
\label{thm:intro-frequency-GRH}
After the proved high-reader, complete same-product diagonal, and short
product-collision removals, suppose that
\[
 |\mathcal O_{\chi,K}^{(\ne0)}|
 +|\mathcal O_{\chi,K}^{(0)}|
 \le K^{1+o(1)}.
\]
Then $M_\chi(K)=O(K^{1/2+o(1)})$ and GRH holds for $L(s,\chi)$.  The
nonzero term is the legitimate landing point for source-normalized
Kloosterman-fraction estimates, while the zero term is the phase-free
sinc-sheet current of \Cref{sec:grh-frequency-resolved}.
\end{theorem}

The fourth closure is cutoff independent.  The completed primitive theta
kernel defines a real entire function $\mathscr X_{\chi,t}$ under backward
heat flow and a de Bruijn--Newman constant $\Lambda_{\rm DN}(\chi)\ge0$.
Its Taylor moments generate finite Jensen packets
$J_{\chi,t}^{d,n}$.  GRH is equivalent to hyperbolicity of all packets at
$t=0$.

\begin{theorem}[Completed theta--Jensen GRH closure]
\label{thm:intro-theta-GRH}
Suppose that, in every primitive nonprincipal channel, a positive first heat
boundary cannot escape cofinally through $d+n\to\infty$, every fixed first
collision is a generic double collision, and its signed Jensen curvature is
nonpositive.  Then $\Lambda_{\rm DN}(\chi)=0$ and GRH holds in that channel.
The classical pole-removed zeta completion gives the corresponding principal
criterion.  Thus these conditions in all primitive channels imply full GRH.
\end{theorem}

The theorem is the contrapositive of the finite collision/cofinal escape
alternative and the positive-curvature law proved in
\Cref{sec:grh-theta-heat}.  It is independent of the outer quotient compiler:
the theta source is fixed by completion data before any terminal cutoff or
zero is read.

The fifth closure is a completed spectral positivity criterion.  Finite
unramified Euler banks form strict Schur cascades and finite
Hermite--Biehler systems, but the completed scattering ratio is the constant
root number and a scalar-normalized Euler cascade has no nonzero entire
limit.  Completion must therefore be installed before positivity is tested.
After exact archimedean resummation, every compact support wall has a universal
logarithmic core whose Legendre eigenvalues are the harmonic numbers.

\begin{theorem}[Completed Weil--Schur GRH closure]
\label{thm:intro-Weil-Schur-GRH}
For every primitive channel $\chi$ and every support wall $A>0$, let
$\mathsf H_{\chi,A}^{[r],\rm sc}$ be the completed finite-history operator of
\Cref{sec:grh-weil-schur}.  Its high Legendre complement is strictly positive
for all sufficiently large endpoint-independent cutoffs $N$ determined from the completed active history.  If, for every
$A$, the corresponding finite protected Schur complement is nonnegative for
one such $N$, then the localized Weil form is nonnegative for every compactly
supported test function and GRH holds in that channel.  The pole-removed zeta
completion gives the principal analogue.
\end{theorem}

This criterion does not infer low-mode positivity from the harmonic gap.  It
reduces each fixed infinite-dimensional sign problem to a finite matrix and
leaves the cofinal evolution of those protected matrices as the analytic
frontier.

The finite and cofinal parts of these closures can now be separated more
sharply.  A normalized modulation bank restores visibility throughout a
growing Mellin strip, but monomial translation tomography has an exponential
proper-arc conditioning obstruction.  The completed source geometry also
admits a state-preserving score predictor whose upper comparison defect is
exactly controlled by one acquired mixed-Gram residual.  Neither result
supplies the missing cofinal lower comparison.

Within the theta entrance, a fixed coefficient block has an exact Hermite
heat limit.  A conjugate pair with $p\ge1$ arbitrary real padding factors
can reach a real-rooted heat image only if $y^2\le(p+2)K$, strictly for even
$p$; the quartic optimum is the algebraic constant
$c_\star\simeq3.814511867$.  Multi-pair blocks have exact total-height and
first-collision budgets.  Their remaining real-axis test is the absence of
a common real zero of two consecutive Appell polynomials, not a presumed
positivity of a Gaussian weight.  These are finite-block constraints, not
new zero-exclusion assertions.  Pairing expansions make the interactions
between conjugate blocks explicit: a sufficient one-pair determinant
condition is not multiplicative.  For a genuinely growing equispaced
Gaussian symbol, the last critical point is at least
$(\log M)/s-O(1)$; transferring this statement to a beta-kernel filter
with $2M$ derivatives remains a separate problem.

Within the Weil entrance, every low matrix entry is an explicit Legendre
translation coefficient, wall integral, or gamma moment.  The exact protected
Schur matrix admits the computable bracket
$L_N-d_N^{-1}R_N\preceq F_N\preceq L_N$.
A complementary Gaussian average of the complete form gives a real scalar
heat symbol and the following additional closure formulation.

\begin{theorem}[Completed Gaussian heat and exhaustive Pick closure]
\label{thm:intro-heat-pick-GRH}
For each primitive nonprincipal channel, Gaussian averaging of the complete
Weil form has a finite positivity threshold $\vartheta_\chi\ge0$, and
GRH in that channel is equivalent to $\vartheta_\chi=0$.
The relative differential inequality in
\Cref{crit:gauss-relative-flow} is sufficient to force this vanishing.
Equivalently for the conjugate pair, every finite Pick matrix of the actual
completed logarithmic derivatives on a fixed dense Euler-region sample must
be positive semidefinite.  A hidden pole forces a negative matrix at some
finite stage, whereas any prescribed finite collection of positive matrices
does not control the continuation.
\end{theorem}

The proof in \Cref{sec:gaussian-weil-pick} keeps the whole signed arithmetic
form under the Gaussian semigroup.  Eventual positivity, contact-jet
identities, and forward dissipation do not imply positivity at time zero.
In the coefficient-sensitive Pick chart, the gamma and higher-prime-power
layers are holomorphic throughout the target half-plane; all continuation
work is localized in the first-prime layer.

A direct radial-filter formulation in
\Cref{sec:admissible-spectral-filters} identifies the corresponding
normalization barrier.  Any finite compatible spectral table admits an
exact interpolant positive beyond a fixed wall.  Nevertheless, if its pole
response vanishes and its target response is $T>0$, then its coefficients
obey $\sum_j|b_j|\ge k!T/L^k$.  Arithmetic positivity simultaneously
forces the actual completed remainder to be at least $m_{\rho_*}T$.
Thus the remaining task is not finite cancellation, but a cofinal signed
estimate that contradicts this normalized compensation while preserving
admission on prime-power logarithms.  Last-zero and negative-interval
results restrict particular half-line-positive filter architectures;
they do not by themselves exclude a spectral zero.

\subsection{The affine-prime closures}

For a primitive nondegenerate affine pair
$L_i(n)=a_i n+b_i$, write
\[
 \mathcal C_{\bm L}(X)
 =\sum_{n\in I_X(\bm L)}\Lambda(L_1(n))\Lambda(L_2(n)),
\]
and let $\mathfrak S(\bm L)$ be its Hardy--Littlewood singular series
\cite{HardyLittlewood}.  The exact $W$-fibre decomposition isolates a
protected normalized amplitude $\mathfrak a_{\bm L}(X;W)$ and a nonnegative residue
dispersion.

\begin{theorem}[Fixed-packet affine conditional closure]
\label{thm:intro-affine-amplitude}
Let $\bm L$ be a fixed admissible primitive affine pair and let
$W_X=X^{o(1)}$ be a squarefree local regulator tending to infinity.  If
\[
 \mathfrak a_{\bm L}(X;W_X)=1+o(1)
\]
and the exceptional prime-power term is $o(X)$, then
\[
 \mathcal C_{\bm L}(X)
 =\mathfrak S(\bm L)|I_X(\bm L)|+o(X).
\]
If instead $\mathfrak a_{\bm L}(X;W_X)\ge\eta>0$, then the interval contains genuine
prime pairs for all sufficiently large $X$ after the proved prime-power
cleanup.  The analogous regulator-uniform statement along the reflection
family $(n,N-n)$ gives the strong Goldbach conclusion.
\end{theorem}

For the twin packet, the endpoint and pressure reductions yield a sharper
single-current certificate.  Put
\[
 \mathcal C_X(2)=\sum_{X<n\le2X}\Lambda(n)\Lambda(n+2),
 \qquad
 R_X(2)=\mathcal C_X(2)-\mathfrak S(2)X.
\]

\begin{theorem}[Complement-pressure twin closure]
\label{thm:intro-pressure-twin}
Let $\mathfrak J_X^{\rm oo,pr}(t)$ be the prime-cleaned odd--odd pressure
current constructed in \Cref{sec:affine-pressure}, and put
$t_{0,X}=(\log X)^{-20}$.  If
\[
 \sup_{t_{0,X}\le t\le1}
 \abs{\mathfrak J_X^{\rm oo,pr}(t)}=o(X),
\]
or merely
\[
 \int_{t_{0,X}}^1\mathfrak J_X^{\rm oo,pr}(t)\,\dd t=o(X),
\]
then $R_X(2)=o(X)$.  Hence the Hardy--Littlewood twin asymptotic holds, and
there are infinitely many twin primes.
\end{theorem}

A complementary analytic reduction restores the prime reader before opening
a second M\"obius--log factor.  In a smooth balanced block, the modulus and
quotient scales are
\[
 M=X^{17/33}(\log X)^{-O(1)},
 \qquad Q\asymp X/M=X^{16/33}(\log X)^{O(1)},
 \qquad m\le X^{17/33}\text{ on the support}.
\]
Let
\[
 L_X=X^{1/33+2/10000},
 \qquad
 U_X=X^{13/165-2/10000}.
\]

\begin{theorem}[Convenient-divisor and hard-modulus reduction]
\label{thm:intro-hard-modulus}
For either reflected prime-reader orientation, every odd balanced modulus
$m\asymp M$ having a divisor in $[L_X,U_X]$ contributes
$O_A(X(\log X)^{-A})$ in total, by Maynard's fixed-residue theorem.  Every
surviving squarefree modulus has a canonical factorization
\[
 m=r(m)s(m),
 \qquad r(m)<L_X,
\]
where $P^-(n)$ denotes the least prime factor of $n>1$, with
\[
 P^-(s(m))>X^{8/165-4/10000},
 \qquad 1\le\omega(s(m))\le10.
\]
Thus a source-independent smooth decomposition whose nonbalanced remainder is
$o(X)$ closes the twin asymptotic once the finite-depth hard prime-progression
currents are $o(X)$.  The depth-one base layer contains prime moduli $p$ and
the literal prime--prime reader $\Lambda(pq\pm2)$.
\end{theorem}

This theorem is proved in \Cref{sec:affine-hard-modulus}.  It is an absolute
pruning theorem for one balanced analytic frontier, not a replacement for the
global pressure or protected-amplitude conditions.  Proper prime powers
contribute only $O(X^{1/2}(\log X)^2)$ to the weighted correlation.

\subsection{The common mechanism}

The exact shared outer diagram is
\[
 \boxed{
 \begin{gathered}
 \Lambda=\mu*\log
 \quad\Longrightarrow\quad
 \text{outer divisor/complement compiler}\\[0.2em]
 \quad\Longrightarrow\quad
 q=1\text{ annihilated by }\log1=0\\[0.2em]
 \quad\Longrightarrow\quad
 \text{target-free signed quotient current}\\[0.2em]
 \quad\Longrightarrow\quad
 \text{endpoint fusion and birth--retuning decomposition}\\[0.2em]
 \quad\Longrightarrow\quad
 \begin{cases}
 \text{conductor-block short quotient for GRH},\\
 \text{endpoint/pressure or restored-prime current for twins}.
 \end{cases}
 \end{gathered}}
\]
The expansion--pullback theorem supplies an additional common audit: an exact
factorization or determinant expansion of a protected endpoint is not an
independent source until a genuinely new action or occurrence row has been
loaded.  This forces reader restoration before a theorem about complete
primes is used, and forces factor preservation through dispersion when a
convolution theorem is used.

The two branches also possess source-native entrances outside the common
outer diagram.  Direct two-threshold M\"obius descent creates the protected
rectangular GRH current; equality-preserving completion then separates its
reciprocal-phase and sinc-sheet channels.  The completed theta kernel creates
a cutoff-independent heat/Jensen current, while completed Euler cells and the
localized Weil form create a spectral-renewal chart with an unconditional
harmonic high-mode complement.  On the affine side, the reflected outer
orientation restores a varying fixed-residue prime reader and admits
Maynard's convenient-divisor pruning.  That pruning has no source-preserving
transfer to a fixed character, because the latter has neither a varying
progression modulus nor a second prime endpoint.

The score predictor in \Cref{subsec:gran-predictor}, the normalized Mellin
bank in \Cref{sec:cofinal-mellin}, and the explicit Weil acquisition in
\Cref{subsec:weil-acquisition} quantify three further parts of this
source-restoration requirement.  Their contraction, visibility, and finite
matrix identities do not identify the protected currents of the two branches.

The shared tensor therefore explains which reductions and no-go statements
are exact.  It does not supply the final reciprocal-phase or sinc-sheet
cancellation, low-mode Schur positivity, collision curvature, or cofinal
compactness estimate.

\subsection{Organization}

\Cref{sec:gran-cylinder,sec:gran-sources,sec:gran-occurrence,sec:gran-refinement}
build the finite tensor, its complete mixed-Gram prediction estimates, and
its cutoff calculus.  \Cref{sec:arithmetic-loading} installs arithmetic,
proves endpoint pullback, and states the reader-restoration rule.
\Cref{sec:hard-phase,sec:grh-outer,sec:grh-hyperbola} develop the phase and
target-free Chebyshev routes.  \Cref{sec:grh-rectangular,sec:grh-frequency-resolved}
give the two-M\"obius descent and its completion-frequency resolution.
\Cref{sec:cofinal-mellin} separates spectral visibility from conditioned
cofinal reconstruction.  \Cref{sec:grh-theta-heat,sec:finite-heat-blocks}
develop the completed theta--Jensen entrance, finite heat-block
constraints, pair couplings, and growing Gaussian symbols.  \Cref{sec:grh-weil-schur} gives completed Euler--Weil renewal,
exact low-mode acquisition, and protected Schur brackets;
\Cref{sec:gaussian-weil-pick} gives the Gaussian heat threshold and the
coefficient-sensitive Pick formulation.
\Cref{sec:admissible-spectral-filters} gives finite interpolation,
normalized completed compensation, exact filter residues, and the
remaining growing-order admission problem.
\Cref{sec:local-geometry,sec:affine-outer,sec:affine-endpoint,sec:affine-pressure}
establish the local, endpoint, and pressure affine reductions, while
\Cref{sec:affine-hard-modulus} restores the prime reader and isolates the
finite-depth hard-modulus frontier.  \Cref{sec:common-spine} gives the
formal comparison and the combined closure theorem.

% =============================================================================
\part{The arithmetic Renewal Gran tensor}
% =============================================================================

\section{Finite renewal cylinders}
\label{sec:gran-cylinder}

\subsection{Histories, cutoffs, and renewal}

\begin{definition}[Finite same-history cylinder]
Let $X\ge1$.  A finite same-history cylinder is a finite set $\Omega_X$
together with a strictly positive probability law $\rho_X$.  Its Hilbert space
is
\[
 \mathcal H_X=L^2(\Omega_X,\rho_X),
 \qquad
 \ip{f}{g}_{\rho_X}=\sum_{\omega\in\Omega_X}
 \overline{f(\omega)}g(\omega)\rho_X(\omega).
\]
A history coordinate is any function on $\Omega_X$.  A coordinate is
\emph{physical} only when it is declared as part of the cylinder before the
protected endpoint is evaluated.
\end{definition}

\begin{definition}[Cutoff-compatible renewal system]
A renewal system is a sequence
\[
 (\Omega_X,\rho_X,\pi_{Y,X})_{1\le X\le Y},
\]
where $\pi_{Y,X}:\Omega_Y\to\Omega_X$ is a surjection,
$\pi_{Z,X}=\pi_{Y,X}\circ\pi_{Z,Y}$, and
$(\pi_{Y,X})_\#\rho_Y=\rho_X$.  The pullback
\[
 U_{Y,X}f=f\circ\pi_{Y,X}
\]
is then an isometry from $\mathcal H_X$ to $\mathcal H_Y$.
\end{definition}

\begin{lemma}[Isometric renewal embedding]
\label{lem:renewal-isometry}
For every $X\le Y$,
\[
 U_{Y,X}^*U_{Y,X}=I_{\mathcal H_X},
 \qquad
 U_{Z,X}=U_{Z,Y}U_{Y,X}.
\]
The orthogonal projection onto the old sector is
$P^{\rm old}_{Y,X}=U_{Y,X}U_{Y,X}^*$.
\end{lemma}

\begin{proof}
The pushforward identity gives
\[
 \norm{U_{Y,X}f}_{\rho_Y}^2
 =\sum_{\omega_Y}|f(\pi_{Y,X}\omega_Y)|^2\rho_Y(\omega_Y)
 =\sum_{\omega_X}|f(\omega_X)|^2\rho_X(\omega_X).
\]
The composition law is immediate.  An isometry times its adjoint is the
orthogonal projection onto its range.
\end{proof}

\begin{definition}[Old and detail sectors]
For adjacent cutoffs, write
\[
 \mathcal H_{X+1}
 =U_{X+1,X}\mathcal H_X\oplus\mathcal D_{X+1},
 \qquad
 \mathcal D_{X+1}=\Ker U_{X+1,X}^*.
\]
The first summand is the transported old carrier and the second is the
new detail carrier.
\end{definition}

\begin{proposition}[Martingale interpretation]
Let $f\in\mathcal H_{X+1}$.  Then
\[
 U_{X+1,X}U_{X+1,X}^*f
 =\E_{\rho_{X+1}}[f\mid\pi_{X+1,X}],
\]
and the detail $f-P^{\rm old}_{X+1,X}f$ has conditional mean zero on every
old fibre.
\end{proposition}

\begin{proof}
The adjoint averages over each fibre with its conditional probability.  The
stated projection is therefore the finite conditional expectation.  The
orthogonal complement has zero inner product with every old-measurable
function, which is equivalent to fibrewise zero conditional mean.
\end{proof}

\subsection{Relation records and tensor completion}

A one-time history space is not sufficient when the protected observable
depends on relative source orientation.  The cylinder must contain the
same-history relation record before either endpoint is aggregated.

\begin{definition}[Relational completion]
Let $E_X$ be a finite set of admissible directed history relations.  A
relational completion of $\Omega_X$ is a finite set
$\widehat\Omega_X$ with maps
\[
 s,t:\widehat\Omega_X\to\Omega_X,
 \qquad
 \ell:\widehat\Omega_X\to E_X,
\]
and a positive reference law $\widehat\rho_X$.  The maps record source,
target, and relation label on one history.  Its Hilbert space is
$\widehat{\mathcal H}_X=L^2(\widehat\Omega_X,\widehat\rho_X)$.
\end{definition}

\begin{construction}[Gran tensor at cutoff $X$]
\label{con:gran-tensor}
The Renewal Gran tensor $\Gran_X$ consists of
\[
 \Gran_X=
 \bigl(
 \widehat{\mathcal H}_X,
 \{S_X^{(a)}\}_{a\in\mathcal A_X},
 N_X,
 U_{2,X},
 \mathcal A_X^{\rm acc},
 U_{X+1,X},
 \widehat\rho_X
 \bigr),
\]
where the entries are respectively the relational history carrier, the
complete source family, the nuisance map, the two-occurrence fusion map, the
accepted/action record, the cutoff renewal embedding, and the fixed positive
reference.  Completeness means that all mixed inner products among jointly
visible source columns are retained.
\end{construction}

\begin{remark}[Why the object is a tensor]
Each source bank may itself be a product of chronology, residue, divisor,
character, affine, occurrence, and action coordinates.  The complete carrier
retains their multilinear mixed moments before any quotient is taken.  The
word tensor refers to this actual joint occurrence, not to the formal
replacement of every occurrence by an abstract tensor product.
\end{remark}

\section{Complete source Grams and nuisance shorting}
\label{sec:gran-sources}

\subsection{Source maps}

\begin{definition}[Source map and complete mixed Gram]
Let $\mathcal E_X$ be a finite-dimensional coefficient space and let
\[
 S_X:\mathcal E_X\longrightarrow\widehat{\mathcal H}_X
\]
be a linear source map.  Its complete Gram is
\[
 G_X=S_X^*S_X\succeq0.
\]
For a family $S_X^{(1)},\ldots,S_X^{(k)}$, the complete joint map is
$S_X=[S_X^{(1)}\ \cdots\ S_X^{(k)}]$; hence
\[
 G_X^{ab}=(S_X^{(a)})^*S_X^{(b)}.
\]
\end{definition}

\begin{lemma}[Source realization of a positive form]
Every finite positive semidefinite matrix $G$ is the Gram of a source map.
One may take $S=G^{1/2}$ on $\overline{\Ran G}$.
\end{lemma}

\begin{proof}
The spectral theorem defines the nonnegative square root and gives
$S^*S=G$.
\end{proof}

\begin{definition}[Protected influence]
For $G\succeq0$ and a coefficient vector $e$, define
\[
 \Lambda_G(e)
 :=\inf\{\Lambda\ge0:ee^*\preceq\Lambda G\}.
\]
Equivalently,
\[
 \Lambda_G(e)=
 \begin{cases}
 \ip{e}{G^\dagger e},&e\in\Ran G,\\
 +\infty,&e\notin\Ran G.
 \end{cases}
\]
\end{definition}

\begin{proof}[Proof of the equivalence]
On $\overline{\Ran G}$ write $u=G^{\dagger/2}e$.  Then
$ee^*=G^{1/2}uu^*G^{1/2}\preceq\norm{u}^2G$, and equality of the optimal
constant follows by testing against $G^\dagger e$.  If $e\notin\Ran G$, a
vector in $\Ker G$ pairs nontrivially with $e$, so no finite domination is
possible.
\end{proof}

\begin{proposition}[Riesz protection inequality]
\label{prop:gran-riesz}
If $e\in\Ran G$, then for every $x$,
\[
 \abs{\ip{e}{x}}^2
 \le \ip{x}{Gx}\,\Lambda_G(e).
\]
\end{proposition}

\begin{proof}
Write $e=G^{1/2}G^{\dagger/2}e$ and apply Cauchy--Schwarz.
\end{proof}

\subsection{Nuisance maps and simultaneous shorting}

\begin{definition}[Nuisance short]
Let $N:\mathcal N\to\widehat{\mathcal H}$ be a nuisance source map.  Its
orthogonal short is
\[
 Q_N=I-N(N^*N)^\dagger N^*.
\]
For a protected source map $S$, the shorted source and Gram are
\[
 S^{\rm sh}=Q_NS,
 \qquad
 G^{\rm sh}=S^*Q_NS.
\]
\end{definition}

\begin{lemma}[Exact nuisance projection]
$Q_N$ is the orthogonal projection onto $\Ker N^*$, and
\[
 \norm{Q_Nf}^2
 =\inf_{u\in\mathcal N}\norm{f-Nu}^2.
\]
\end{lemma}

\begin{proof}
$N(N^*N)^\dagger N^*$ is the orthogonal projection onto $\Ran N$.
The Pythagorean theorem gives the variational formula.
\end{proof}

\begin{theorem}[Order independence of simultaneous shorting]
\label{thm:simultaneous-short}
Let $N=[N_1\ \cdots\ N_k]$.  The projection $Q_N$ depends only on the joint
nuisance span.  If the corresponding orthogonal nuisance projections commute
pairwise, sequential shorting in any order gives the same joint short.
For two projections this commutation condition is also necessary.  For a
larger family no converse based on pairwise commutation is asserted.
In general, sequential one-source shorting may erase mixed protected
orientation, whereas the joint short does not depend on an arbitrary order.
\end{theorem}

\begin{proof}
The first statement is uniqueness of the orthogonal projection onto
$(\sum_j\Ran N_j)^\perp$.  A product of pairwise commuting orthogonal
projections is the orthogonal projection onto the intersection of their
ranges.  For two projections, self-adjointness of the product forces the
two orders to coincide, proving necessity.  With three or more projections,
noncommuting products may degenerate to the intersection projection, so
pairwise commutation is only used as a sufficient hypothesis here.
\end{proof}

\begin{principle}[Assembly before shorting]
Whenever two arithmetic rows may contribute to the same protected signed
coefficient, their mixed Gram is assembled before either row is shorted.  The
protected sector is declared only after all prescribed nuisance columns are
present.
\end{principle}


\subsection{State-preserving prediction and a complete-Gram bridge}
\label{subsec:gran-predictor}

A comparison between two score descriptions can be made quantitative
without declaring them identical.  The following finite-algebra
construction provides a contractive predictor and identifies exactly
which mixed source row measures its failure to reproduce the actual
arithmetic score.

\begin{theorem}[Canonical state-preserving score predictor]
\label{thm:gran-petz-predictor}
Let $\mathcal M$ be a finite-dimensional $C^*$-algebra with faithful
trace $\tau$, let $\mathcal N\subseteq\mathcal M$ be a unital
subalgebra, and let $\mathcal E_{\mathcal N}$ be the trace-preserving
conditional expectation.  For $\rho\succ0$, $\tau(\rho)=1$, put
$\varphi(x)=\tau(\rho x)$ and
$\rho_{\mathcal N}=\mathcal E_{\mathcal N}(\rho)$.
The map
\begin{equation}
 \Pi_{\rho,\mathcal N}(x)=\rho_{\mathcal N}^{-1/2}
 \mathcal E_{\mathcal N}(\rho^{1/2}x\rho^{1/2})
 \rho_{\mathcal N}^{-1/2}
 \label{eq:gran-petz-predictor}
\end{equation}
is unital, completely positive, state preserving, and contractive for
$\|x\|_{2,\rho}^2=\varphi(x^*x)$.  The contraction holds at every
matrix amplification.  No modular-invariance assumption is needed for
these assertions; without additional structure the map need not fix
$\mathcal N$ and is not asserted to be a conditional expectation.
\end{theorem}

\begin{proof}
Positivity of $\rho$ implies invertibility of
$\rho_{\mathcal N}$.  Every factor in
\eqref{eq:gran-petz-predictor} is completely positive, and substitution
of $x=I$ proves unitality.  For $y\in\mathcal N$ one has
$\tau(\rho y)=\tau(\rho_{\mathcal N}y)$, by trace orthogonality of
$\mathcal E_{\mathcal N}$.  Applying this to the output of the map,
then using cyclicity and trace preservation, gives
$\varphi(\Pi_{\rho,\mathcal N}(x))=\tau(\rho^{1/2}x\rho^{1/2})
=\varphi(x)$.
The Schwarz inequality for a unital completely positive map now gives
\[
 \|\Pi_{\rho,\mathcal N}(x)\|_{2,\rho}^2
 \le\varphi(\Pi_{\rho,\mathcal N}(x^*x))
 =\|x\|_{2,\rho}^2.
\]
The same proof applies with the matrix trace at every amplification.
\end{proof}

Let $A,S$ be two score syntheses on a common coefficient space, embedded
in the same physical $L^2(\rho)$ carrier, and let $Q$ be the common
orthogonal nuisance short.  Put
\begin{equation}
 \widehat A=QA,\quad \widehat S=QS,\quad
 M=\widehat A^*\widehat A,\quad E=\widehat S^*\widehat S,
 \quad T=Q\Pi_{\rho,\mathcal N}\widehat A,\quad
 R=\widehat S-T.
 \label{eq:gran-predictor-data}
\end{equation}
All comparisons below are on $\supp M$; extension to the whole
coefficient space requires $\Ker M\subseteq\Ker E$.

\begin{theorem}[Residual control of the relative score metric]
\label{thm:gran-predictor-bridge}
Define
$q=\|TM^{-1/2}\|\le1$,
$\varepsilon=\|RM^{-1/2}\|$, and $H=M^{-1/2}EM^{-1/2}$.
Then
\begin{equation}
 T^*T\preceq M,\qquad
 \|H\|\le(q+\varepsilon)^2\le(1+\varepsilon)^2,
 \qquad (\|H\|-1)_+\le2\varepsilon+\varepsilon^2.
 \label{eq:gran-predictor-bridge}
\end{equation}
For every $\theta>0$,
\begin{equation}
 E\preceq(1+\theta)M+(1+\theta^{-1})R^*R.
 \label{eq:gran-predictor-order}
\end{equation}
Most importantly, with $P=T^*T$ and $C=T^*\widehat S$,
\begin{equation}
 R^*R=E+P-C-C^*,\qquad
 \begin{pmatrix}P&C\\C^*&E\end{pmatrix}\succeq0.
 \label{eq:gran-predictor-mixed-gram}
\end{equation}
\end{theorem}

\begin{proof}
Contractivity of the predictor and of $Q$ gives $T^*T\preceq M$.
The identity $\widehat S=T+R$ and the triangle inequality after
whitening give \eqref{eq:gran-predictor-bridge}.
Young's inequality applied to $\|Tc+Rc\|^2$ for each coefficient
vector proves \eqref{eq:gran-predictor-order}.
Expansion of $(\widehat S-T)^*(\widehat S-T)$ gives the first
identity in \eqref{eq:gran-predictor-mixed-gram}; the second is the
joint Gram of $T$ and $\widehat S$.
\end{proof}

Thus $\varepsilon_X\to0$ on an exhaustive source family, together
with the kernel condition, makes the upper relative-metric defect tend
to zero.  This does not supply a restoring lower comparison, nor does
it establish the final arithmetic descent.  The residual is not
recoverable from marginal score norms: the maps
$Tc=ce_1$, $S_0c=ce_1$, and $S_1c=ce_2$ have identical scalar
marginal Grams, but residual squares zero and two, respectively.
The mixed block $C$ must therefore be acquired from the same declared
operational record, not selected after the protected endpoint is known.


\section{Product occurrence, innovation, and Feshbach reduction}
\label{sec:gran-occurrence}

\subsection{Actual two-occurrence carrier}

\begin{definition}[Product fusion and occurrence innovation]
Let $\mathcal H_1$ be a one-occurrence carrier and $\mathcal H_2$ an actual
two-occurrence carrier.  A product fusion is an isometry
\[
 U_2:\mathcal H_1^{\otimes2}\longrightarrow\mathcal H_2.
\]
The product projection and occurrence innovation are
\[
 P_{\prodsec}=U_2U_2^*,
 \qquad
 P_{\occ}=I-U_2U_2^*.
\]
\end{definition}

\begin{theorem}[Product--occurrence Pythagoras]
For every $f\in\mathcal H_2$,
\[
 f=P_{\prodsec}f+P_{\occ}f,
 \qquad
 \norm f^2=\norm{P_{\prodsec}f}^2+\norm{P_{\occ}f}^2.
\]
The product model is complete if and only if $P_{\occ}=0$.
\end{theorem}

\begin{proof}
$U_2U_2^*$ is an orthogonal projection because $U_2$ is an isometry.  The
claims follow from orthogonal decomposition.
\end{proof}

\begin{remark}[Occurrence innovation is not higher form degree]
A nonzero $P_{\occ}$ records two-event information not generated by the
product occurrence.  It need not create a third affine form or a third
prime-coordinate degree.  This distinction is central in the affine-pair
analysis.
\end{remark}

\subsection{Feshbach reduction}

\begin{theorem}[Finite Feshbach formula]
\label{thm:gran-feshbach}
Let $\mathcal H=\mathcal H_P\oplus\mathcal H_Q$ and let
\[
 A=\begin{pmatrix}A_{PP}&A_{PQ}\\A_{QP}&A_{QQ}\end{pmatrix}\succeq0.
\]
Assume first that $A_{QQ}$ is invertible.  Then
\[
 \ip{x_P\oplus x_Q}{A(x_P\oplus x_Q)}
 =\ip{x_P}{\Fesh_Q(A)x_P}
 +\ip{x_Q+A_{QQ}^{-1}A_{QP}x_P}
          {A_{QQ}(x_Q+A_{QQ}^{-1}A_{QP}x_P)},
\]
where
\[
 \Fesh_Q(A)=A_{PP}-A_{PQ}A_{QQ}^{-1}A_{QP}\succeq0.
\]
Consequently,
\[
 \inf_{x_Q}\ip{x_P\oplus x_Q}{A(x_P\oplus x_Q)}
 =\ip{x_P}{\Fesh_Q(A)x_P}.
\]
With the usual range condition, the same statements hold with
$A_{QQ}^{-1}$ replaced by $A_{QQ}^\dagger$.
\end{theorem}

\begin{proof}
Expand the square on the right.  The cross terms are
$\ip{x_P}{A_{PQ}x_Q}+\ip{x_Q}{A_{QP}x_P}$ and the remaining correction is
$\ip{x_P}{A_{PQ}A_{QQ}^{-1}A_{QP}x_P}$.  Positivity follows by minimizing a
positive form.  The pseudoinverse version follows after restricting to
$\overline{\Ran A_{QQ}}$.
\end{proof}

\begin{corollary}[Occurrence Feshbach reduction]
If $\mathcal H_P=\Ran U_2$ and $\mathcal H_Q=\Ran P_{\occ}$, then every
positive two-occurrence action admits an exact effective product action by
shorting the occurrence innovation.  The mixed product--occurrence block must
be retained; deleting it replaces the Schur complement by the generally
larger principal block.
\end{corollary}

\section{Accepted/action geometry}
\label{sec:gran-action}

\subsection{Finite relative entropy}

\begin{definition}[Accepted/action record]
Let $R$ and $P$ be probability laws on the same finite relational carrier,
with $P\ll R$.  The accepted/action likelihood record is
\[
 \ell=\log\frac{\dd P}{\dd R},
 \qquad
 \mathfrak D(P\|R)=\E_P\ell=\KL(P\|R).
\]
The reference $R$ is fixed from the source data before the terminal law $P$ is
read.
\end{definition}

\begin{theorem}[Data processing]
Let $\pi$ be any finite measurable quotient.  Then
\[
 \KL(\pi_\#P\|\pi_\#R)\le\KL(P\|R).
\]
\end{theorem}

\begin{proof}
On each quotient fibre, apply the log-sum inequality
\[
 \sum_i a_i\log\frac{a_i}{b_i}
 \ge\left(\sum_i a_i\right)
 \log\frac{\sum_i a_i}{\sum_i b_i}.
\]
Sum over fibres.
\end{proof}

\begin{theorem}[Finite path-space chain rule]
Let $P$ and $R$ be laws on a finite path
$(Z_0,\ldots,Z_k)$, and suppose $P\ll R$.  Then
\[
 \KL(P\|R)
 =\KL(P_0\|R_0)
 +\sum_{j=1}^k
 \E_P\KL\bigl(P(Z_j\mid Z_{<j})\|R(Z_j\mid Z_{<j})\bigr).
\]
Every term is nonnegative.
\end{theorem}

\begin{proof}
Factor both densities into an initial density and conditional transition
densities, take the logarithm of the ratio, and average under $P$.
\end{proof}

\begin{remark}[Local obstruction output]
A large global action returns a large nonnegative conditional increment at
some chronology, relation, destination, provenance, or within-fibre stage.
This is a finite obstruction, not an appeal to an unspecified global failure.
\end{remark}

\section{Cutoff refinement: carrier birth and signed retuning}
\label{sec:gran-refinement}

The most delicate finite operation is enlargement of a source carrier.  An
adjacent cutoff can both change the source on the old coordinates and create a
new detail coordinate.  These contributions have opposite signs in the
protected influence ledger and must not be conflated.

\begin{definition}[Old/detail source block]
Let $\mathcal E$ be the old coefficient space and $\mathcal D$ a new detail
space.  A complete refined source Gram has the block form
\[
 \mathbb G^+
 =\begin{pmatrix}A&B\\B^*&D\end{pmatrix}\succeq0,
 \qquad A\succ0.
\]
Its detail Schur complement is
\[
 S=D-B^*A^{-1}B\succeq0.
\]
For a protected read $f=e\oplus r$, define the prediction residual
\[
 q=r-B^*A^{-1}e.
\]
\end{definition}

\begin{theorem}[Exact old/detail influence decomposition]
\label{thm:gran-old-detail}
One has
\[
 f\in\Ran\mathbb G^+
 \quad\Longleftrightarrow\quad
 q\in\Ran S.
\]
On this range,
\[
 \Lambda_{\mathbb G^+}(f)
 =\ip{e}{A^{-1}e}+\ip{q}{S^\dagger q}.
\]
The second term is nonnegative and is the carrier-birth innovation.
\end{theorem}

\begin{proof}
Use the block factorization
\[
 \mathbb G^+
 =\begin{pmatrix}I&0\\B^*A^{-1}&I\end{pmatrix}
  \begin{pmatrix}A&0\\0&S\end{pmatrix}
  \begin{pmatrix}I&A^{-1}B\\0&I\end{pmatrix}.
\]
The inverse coordinate of $e\oplus r$ is
$e\oplus(r-B^*A^{-1}e)=e\oplus q$.  Range membership and the pseudoinverse
quadratic form therefore split as stated.
\end{proof}

\begin{definition}[Old-sector retuning capture]
Let $G\succ0$ be the previous old Gram and let $A\succ0$ be the old principal
block after refinement.  Define
\[
 \mathcal C^{\rm ret}(e;G,A)
 :=\ip{e}{G^{-1}e}-\ip{e}{A^{-1}e}.
\]
It is signed.  A positive value means that the old-sector retuning lowers the
protected influence.
\end{definition}

\begin{theorem}[Birth--retuning balance]
\label{thm:gran-birth-retuning}
With the preceding notation,
\[
 \Lambda_{\mathbb G^+}(e\oplus r)
 =\Lambda_G(e)
  -\mathcal C^{\rm ret}(e;G,A)
  +\eta^{\rm birth},
\]
where
\[
 \eta^{\rm birth}=\ip{q}{S^\dagger q}\ge0.
\]
Moreover,
\[
 \mathcal C^{\rm ret}(e;G,A)
 =\ip{A^{-1}e}{(A-G)G^{-1}e}.
\]
\end{theorem}

\begin{proof}
The first identity is \Cref{thm:gran-old-detail} with the definition of
$\mathcal C^{\rm ret}$.  The resolvent identity
$G^{-1}-A^{-1}=A^{-1}(A-G)G^{-1}$ gives the second.
\end{proof}

\begin{corollary}[Nested extensions cannot improve influence]
If the old block is unchanged, $A=G$, then
\[
 \Lambda_{\mathbb G^+}(e\oplus r)=\Lambda_G(e)+\eta^{\rm birth}
 \ge\Lambda_G(e).
\]
Thus a pure source extension has no favorable chronology; any decrease must
come from signed retuning of already visible rows.
\end{corollary}

\begin{corollary}[Cumulative cutoff ledger]
For a compatible sequence of cutoffs,
\[
 \Lambda_X
 =\Lambda_{X_0}
  -\sum_{j=X_0}^{X-1}\mathcal C_j^{\rm ret}
  +\sum_{j=X_0}^{X-1}\eta_j^{\rm birth},
\]
provided every range condition holds.
\end{corollary}

\begin{proof}
Apply \Cref{thm:gran-birth-retuning} at each adjacent cutoff and telescope.
\end{proof}

\begin{proposition}[Mixed old/detail data are indispensable]
There exist two positive block Grams with the same old block, the same detail
block, and the same protected marginals, but with different birth innovation.
Consequently, diagonal energies do not determine the cutoff ledger.
\end{proposition}

\begin{proof}
Take scalar old and detail spaces, $A=D=2$, $e=r=1$, and
\[
 \mathbb G_0=\begin{pmatrix}2&0\\0&2\end{pmatrix},
 \qquad
 \mathbb G_1=\begin{pmatrix}2&1\\1&2\end{pmatrix}.
\]
For $\mathbb G_0$, $S=2$ and $q=1$, so the birth term is $1/2$.
For $\mathbb G_1$, $S=3/2$ and $q=1/2$, so the birth term is $1/6$.
The marginals coincide while the mixed block changes the prediction residual.
\end{proof}

\begin{principle}[No terminally fitted source]
A source map, nuisance short, pressure family, or cutoff retuning used in a
closure theorem is declared independently of the numerical value of the
protected Mertens or prime-pair endpoint.  The only permitted dependence is on
public arithmetic data such as the cutoff, conductor, affine coefficients,
residue class, and prescribed local model.
\end{principle}


% =============================================================================
\section{Canonical arithmetic loading}
\label{sec:arithmetic-loading}
% =============================================================================

\subsection{Divisibility incidence and M\"obius inversion}

Let $\mathcal A_X=\C^{\{1,\ldots,X\}}$.  Dirichlet convolution on the finite
cutoff is
\[
 (f*g)(n)=\sum_{d\mid n}f(d)g(n/d).
\]

\begin{definition}[Finite zeta and M\"obius operators]
Define
\[
 (Z_Xf)(n)=\sum_{d\mid n}f(d),
 \qquad
 (M_Xf)(n)=\sum_{d\mid n}\mu(d)f(n/d).
\]
The corresponding incidence matrices are
\[
 Z_X(n,d)=\one_{d\mid n},
 \qquad
 M_X(n,d)=\mu(n/d)\one_{d\mid n}.
\]
\end{definition}

\begin{theorem}[Finite M\"obius inversion]
\label{thm:finite-mobius}
For every cutoff,
\[
 M_XZ_X=Z_XM_X=I.
\]
Equivalently, if $F(n)=\sum_{d\mid n}f(d)$, then
$f(n)=\sum_{d\mid n}\mu(d)F(n/d)$.
\end{theorem}

\begin{proof}
The coefficient of $f(e)$ in $(M_XZ_Xf)(n)$ is
\[
 \sum_{e\mid d\mid n}\mu(n/d)
 =\sum_{r\mid n/e}\mu(r)
 =\one_{e=n}.
\]
The other product is identical because Dirichlet convolution is commutative.
\end{proof}

\begin{corollary}[Canonical von Mangoldt loading]
\label{cor:lambda-mu-log}
Let $\log(n)$ denote the arithmetic logarithm and let $\one(n)=1$.  Then
\[
 \log=\one*\Lambda,
 \qquad
 \boxed{\Lambda=\mu*\log},
\]
so that
\[
 \Lambda(n)=\sum_{d\mid n}\mu(d)\log\frac nd.
\]
\end{corollary}

\begin{proof}
Both sides of $\log n=\sum_{d\mid n}\Lambda(d)$ equal
$\sum_{p^k\parallel n}k\log p$.  Apply \Cref{thm:finite-mobius}.
\end{proof}

\subsection{Character, residue, affine, and factor-depth coordinates}

\begin{definition}[Dirichlet loading]
For a Dirichlet character $\chi$ modulo $q$, define the diagonal character
operator
\[
 (C_\chi f)(n)=\chi(n)f(n).
\]
The twisted M\"obius and Chebyshev coordinates are
\[
 \mu_\chi(n)=\mu(n)\chi(n),
 \qquad
 M_\chi(X)=\sum_{n\le X}\mu_\chi(n),
\]
\[
 \psi_\chi(X)=\sum_{n\le X}\chi(n)\Lambda(n).
\]
\end{definition}

\begin{definition}[Residue and affine pullbacks]
For a modulus $W$ and residue $r$, put
\[
 \mathbf r_{W,r}(n)=\one_{n\equiv r\pmod W}.
\]
For an affine form $L(n)=an+b$, define the pullback
\[
 L^*f(n)=f(an+b).
\]
For a pair $\bm L=(L_1,L_2)$, the two-occurrence arithmetic row is
\[
 n\longmapsto (L_1^*\Lambda)(n)(L_2^*\Lambda)(n).
\]
\end{definition}

\begin{definition}[Factor-depth and Liouville loading]
Let $\Omega(n)$ be the total number of prime factors counted with
multiplicity.  The Liouville coordinate is
\[
 \lambda(n)=(-1)^{\Omega(n)}.
\]
On squarefree support, the factor-depth sign is $(-1)^{\omega(n)}=\mu(n)$.
The one- and two-form signs, their product, and their residue restrictions are
loaded as joint source columns before parity shorting.
\end{definition}

\begin{remark}[Why the factor-depth product is retained]
Shorting the two one-form signs separately does not remove their neutral
product.  The affine parity square below proves that this product is the only
selector-stable parity coordinate remaining after constants and one-form
signs are removed.
\end{remark}

\subsection{The universal outer M\"obius--log compiler}

\begin{theorem}[Finite outer compiler]
\label{thm:universal-outer}
Let $b$ be any arithmetic function and let $I$ be a finite set of positive
integers.  Then
\[
 \sum_{m\in I}b(m)\Lambda(m)
 =\sum_{d\ge1}\mu(d)
   \sum_{\substack{q\ge1\\dq\in I}}(\log q)b(dq).
\]
All sums are finite.
\end{theorem}

\begin{proof}
Insert \Cref{cor:lambda-mu-log} and write $m=dq$.
\end{proof}

\begin{corollary}[Logarithmic annihilation of the terminal quotient]
\label{cor:q1-annihilation}
In \Cref{thm:universal-outer}, the section $q=1$ vanishes identically.  Thus
any protected coordinate carried only by the direct factorization $m=d$ is
absent from the outer readout.
\end{corollary}

\begin{proof}
Its coefficient is $\log1=0$.
\end{proof}

\begin{theorem}[Character outer compiler]
\label{thm:character-outer}
For every primitive character $\chi$ and integer $X\ge2$,
\[
 \psi_\chi(X)
 =\sum_{d\le X}\mu(d)\chi(d)
   \sum_{q\le X/d}\chi(q)\log q
 =\sum_{q=2}^{X}\chi(q)\log q\,M_\chi(\lfloor X/q\rfloor).
\]
\end{theorem}

\begin{proof}
Use complete multiplicativity on the supported integers in
\Cref{thm:universal-outer}, then interchange $d$ and $q$.
\end{proof}

\begin{theorem}[Affine outer compiler]
\label{thm:affine-universal-outer}
Let $h$ be fixed and let $I$ be a finite interval on which $m-h\ge2$.
Then
\[
 \sum_{m\in I}\Lambda(m-h)\Lambda(m)
 =\sum_{d\ge1}\mu(d)
   \sum_{\substack{q\ge1\\dq\in I}}
      (\log q)\Lambda(dq-h).
\]
Again the direct section $q=1$ is absent.
\end{theorem}

\begin{proof}
Apply \Cref{thm:universal-outer} with $b(m)=\Lambda(m-h)$.
\end{proof}

\subsection{Source realization of the outer compiler}

\begin{construction}[Divisor--quotient occurrence carrier]
At cutoff $X$, let
\[
 \mathcal O_X=
 \{(d,q,m):dq=m,\ m\in I_X,\ q\ge1\}.
\]
Equip it with any fixed positive reference law depending only on the cutoff
and the declared interval.  The divisor, quotient, endpoint, character,
M\"obius, logarithmic, local-model, and prime-occurrence columns are functions
on $\mathcal O_X$.  Their complete source map is the column operator into
$L^2(\mathcal O_X)$.
\end{construction}

\begin{proposition}[Outer coefficients are mixed source pairings]
On the divisor--quotient carrier, both the character packet
\[
 \sum_{d,q}\mu(d)\chi(d)\chi(q)(\log q)\one_{dq\le X}
\]
and the affine packet
\[
 \sum_{d,q}\mu(d)(\log q)
 \bigl[\Lambda(dq-h)-\gamma_h(d)\bigr]
\]
are mixed pairings of predeclared source columns.  Their signs are therefore
encoded in cross-Gram entries, not in the separate diagonal energies.
\end{proposition}

\begin{proof}
Choose one source vector with entries proportional to the M\"obius divisor
column and the other with entries proportional to the logarithmic completed
quotient row.  Their ordinary inner product is the displayed coefficient.
The complete Gram contains this mixed entry by construction.
\end{proof}

\begin{proposition}[Endpoint fusion is a Gran-tensor short]
Let $F:\ell^2(\mathcal O_X)\to\ell^2(I_X)$ sum all factorizations with common
endpoint $m$.  Then
\[
 P_{\rm end}=F^*(FF^*)^\dagger F
\]
is the orthogonal projection onto endpoint-readable factorization profiles.
The complementary projection $I-P_{\rm end}$ is read-null factorization
innovation.  For every endpoint read $c$ and factorization vector $z$,
\[
 \ip{F^*c}{z}=\ip{F^*c}{P_{\rm end}z}.
\]
\end{proposition}

\begin{proof}
The formula is the orthogonal projection onto $\Ran F^*$.  Since
$F^*c\in\Ran F^*$, it is orthogonal to $\Ker F=\Ran(I-P_{\rm end})$.
\end{proof}


\subsection{Exact pullback of endpoint forms and reader restoration}
\label{subsec:arithmetic-pullback}

An arithmetic expansion may introduce many factorization, determinant, or
provenance coordinates without introducing a new endpoint source.  The
following finite statement distinguishes a genuine source enlargement from a
coordinate pullback of an already protected read.

\begin{theorem}[Exact expansion--pullback theorem]
\label{thm:arithmetic-expansion-pullback}
Let $\mathcal N$ be a finite endpoint set, let $\mathcal A_n$ be a finite
allocation fibre over each $n\in\mathcal N$, and put
\[
 \mathcal H=\ell^2(\mathcal N),
 \qquad
 \widetilde{\mathcal H}
 =\bigoplus_{n\in\mathcal N}\ell^2(\mathcal A_n).
\]
Let
\[
 (Fy)(n)=\sum_{\alpha\in\mathcal A_n}y(n,\alpha)
\]
be endpoint fusion, and let $e:\mathcal H\to\C$ be any fixed endpoint
functional.  The completely expanded positive form is
\[
 \boxed{\widetilde A=(eF)^*(eF)=F^*e^*eF.}
 \tag{L.1}
\]
For every $x\in\Ran F$,
\[
 \boxed{
 \inf_{Fy=x}\ip{y}{\widetilde Ay}=|ex|^2.}
 \tag{L.2}
\]
Thus the sharp fusion quotient of the expanded form is exactly the original
endpoint form $e^*e$.
\end{theorem}

\begin{proof}
If $Fy=x$, then
\[
 \ip{y}{\widetilde Ay}=|eFy|^2=|ex|^2,
\]
independently of the representative in the allocation fibre.  Taking the
infimum proves the claim.
\end{proof}

\begin{corollary}[Factorization and determinant coordinates do not create
source innovation]
\label{cor:arithmetic-factorization-pullback}
Suppose an endpoint coefficient has an exact expansion
\[
 x(n)=\sum_{\alpha\in\mathcal A_n}c(n,\alpha).
\]
Squaring the completely expanded sum, applying fibrewise unitary transforms,
or grouping cross terms by a determinant, congruence, or factorization label
only decomposes the pullback $F^*e^*eF$.  After endpoint fusion its sharp
protected quotient remains $e^*e$.
\end{corollary}

\begin{proof}
Every operation described is either a change of coordinates inside a fusion
fibre or a partition of the matrix entries of the same quadratic form.  Apply
\Cref{thm:arithmetic-expansion-pullback}.
\end{proof}

A particularly transparent instance is the signed two-scale Fej\'er packet.
For $L\ge1$, let $\mathcal H_L=\ell^2\{1,\ldots,L\}$, extend vectors by zero
to $\Z$, and define
\[
 (W_Hx)(t)=\frac1{\sqrt H}\sum_{j=0}^{H-1}x(t+j),
 \qquad F_H=W_H^*W_H.
\]

\begin{proposition}[Rank-one collapse of the signed Fej\'er chronology]
\label{prop:arithmetic-Fejer-collapse}
If $e_L=(1,\ldots,1)\in\mathcal H_L$, then
\[
 \boxed{2F_{2L}-F_L=e_Le_L^*.}
 \tag{L.3}
\]
Equivalently,
\[
 2\norm{W_{2L}x}^2-\norm{W_Lx}^2
 =\abs{\ip{e_L}{x}}^2
 \qquad(x\in\mathcal H_L).
\]
Consequently the signed two-scale packet is exactly the protected terminal
rank-one form; the two positive window energies may not be spent separately
without retaining a strictly stronger ambient action.
\end{proposition}

\begin{proof}
The matrix of $F_H$ is
\[
 F_H(m,n)=\left(1-\frac{|m-n|}{H}\right)_+.
\]
For $1\le m,n\le L$, one has $|m-n|<L$, and hence
\[
 2\left(1-\frac{|m-n|}{2L}\right)
 -\left(1-\frac{|m-n|}{L}\right)=1.
\]
Thus every matrix entry of $2F_{2L}-F_L$ equals one.  This is $e_Le_L^*$.
\end{proof}

\begin{principle}[Reader restoration before analytic input]
\label{prin:reader-restoration}
When a von Mangoldt, character, or endpoint reader has been opened by an exact
arithmetic identity and a later theorem applies to the complete reader, the
allocation fibres are fused before that theorem is invoked.  An estimate for
the restored reader is not silently assigned to one selected factorization
slice.  Conversely, a factorization theorem is useful only when its factors
remain jointly visible through the first operation which creates the required
averaging or dispersion geometry.
\end{principle}

\subsection{Classical analytic inputs used below}

The following inputs are standard.  They are stated here so that every later
conditional statement has an explicit dependence.

\begin{theorem}[Maximal Bombieri--Vinogradov input]
\label{thm:weighted-BV-input}
For every $A>0$ there is $B_A>0$ such that, uniformly for
$D\le X^{1/2}(\log X)^{-B_A}$,
\[
 \sum_{d\le D}
 \max_{Y\le2X+2}
 \max_{(a,d)=1}
 \left|
 \sum_{\substack{n\le Y\\ n\equiv a\pmod d}}\Lambda(n)
 -\frac{Y}{\varphi(d)}
 \right|
 \ll_A\frac{X}{(\log X)^A},
\]
with the standard equivalent variants for dyadic intervals and moving
endpoints.
\end{theorem}

\begin{remark}
This is the standard maximal form obtained from Bombieri--Vinogradov \cite{Bombieri1965,Vinogradov,IK} by dyadic
and partial-summation arguments.  It is used only for the affine divisor head.
\end{remark}

\begin{theorem}[Zero-free-region summatory input]
\label{thm:zero-free-input}
Let $H(s)$ be represented by an absolutely convergent Dirichlet series for
$\Re s>1$ and suppose that it extends holomorphically, with at most polynomial
vertical growth, to a de la Vall\'ee Poussin region
\[
 \Re s\ge 1-\frac{c_0}{\log(3+|\Im s|)}.
\]
If
\[
 F(s)=-\frac{\zeta'(s)}{\zeta(s)}H(s)
      =\sum_{n\ge1}\frac{a_F(n)}{n^s}
 \qquad(\Re s>1),
\]
then, after possibly decreasing $c_0$, there is $c>0$ such that
\[
 \sum_{n\le x}a_F(n)
 =H(1)x+O_H\!\left(xe^{-c\sqrt{\log x}}\right).
\]
Under the same continuation and growth hypotheses,
$G(s)=H(s)/\zeta(s)=\sum a_G(n)n^{-s}$ satisfies
\[
 \sum_{n\le x}a_G(n)
 =O_H\!\left(xe^{-c\sqrt{\log x}}\right).
\]
\end{theorem}

\begin{remark}
Only the displayed consequences of the classical zero-free region \cite{Davenport,MV} are used;
no zero-free estimate near the critical line is assumed.
\end{remark}

\begin{principle}[Arithmetic loading order]
The incidence, M\"obius, logarithmic, character, residue, affine, factor-depth,
local-model, occurrence, and cutoff columns are loaded before the protected
endpoint is selected.  The outer compiler is then an exact contraction of the
loaded tensor.
\end{principle}


% =============================================================================
\part{Conditional closure of Dirichlet channels}
% =============================================================================

\section{The phase and outer-compiler GRH entrances}
\label{sec:grh-entrance}

The action entrance and the outer-compiler entrance address different parts of
the same problem.  The phase action gives a conceptually complete conditional
criterion for every primitive channel.  The outer compiler gives a more
explicit target-free arithmetic packet, closes the nonprincipal divisor head,
and reveals the square-root quotient barrier.  Neither route permits a source
to be chosen after observing the terminal M\"obius sum.

For the phase sections below, all phase groups, ratio kernels, and transport
laws are finite.  The only asymptotic input is the proposed subpower extensive
action.  For the concrete outer sections, fix $B>2$ and set
\[
 D_X=\left\lfloor X^{1/2}(\log X)^{-B}\right\rfloor,
 \qquad
 Q_X=\left\lfloor\frac{X}{D_X+1}\right\rfloor.
 \tag{G.0}
\]


\section{Hard phase packets and a canonical extensive action}
\label{sec:hard-phase}
% =============================================================================

Fix a primitive Dirichlet character $\chi$ modulo $q$.  Let
\[
 H_\chi:=\chi\bigl((\mathbb Z/q\mathbb Z)^\times\bigr)\subset\T,
 \qquad
 G_\chi:=\langle-1,H_\chi\rangle\subset\T.
 \tag{2.1}
\]
The finite subgroup $G_\chi$ is cyclic and has even cardinality at least two.
Write $u_\chi$ for its uniform law.

\begin{definition}[Dyadic hard-counting phase packet]
For $Y\ge2$ put
\[
 \Omega_{\chi,Y}
 :=\{n\in\N:Y<n\le2Y,\ \mu(n)^2=1,\ (n,q)=1\},
 \qquad
 A_{\chi,Y}:=|\Omega_{\chi,Y}|.
 \tag{2.2}
\]
For $n\in\Omega_{\chi,Y}$ define the protected phase
\[
 \omega_\chi(n):=\mu(n)\chi(n)\in G_\chi.
 \tag{2.3}
\]
When $A_{\chi,Y}>0$, let $\tau_{\chi,Y}$ be the pushforward to $G_\chi$ of the
uniform law on $\Omega_{\chi,Y}$.  Put
\[
 S_\chi(Y):=\sum_{Y<n\le2Y}\mu(n)\chi(n)
 =A_{\chi,Y}\sum_{g\in G_\chi}\tau_{\chi,Y}(g)g.
 \tag{2.4}
\]
When $A_{\chi,Y}=0$, set $S_\chi(Y)=0$ and all actions below equal to zero.
\end{definition}

\begin{lemma}[The source-fixed phase reference is centered]
\label{lem:uniform-phase-centered}
One has
\[
 \boxed{\sum_{g\in G_\chi}u_\chi(g)g=0.}
 \tag{2.5}
\]
\end{lemma}

\begin{proof}
Let $\zeta$ generate the nontrivial finite cyclic group $G_\chi$ of order $m$.
Then
\[
 \sum_{g\in G_\chi}g=\sum_{j=0}^{m-1}\zeta^j
 =\frac{1-\zeta^m}{1-\zeta}=0.
\]
Divide by $m$.
\end{proof}

\begin{lemma}[Positive density of the active hard packet]
\label{lem:active-density}
For fixed $q$ there is a constant
\[
 c_q
 =\frac{\varphi(q)}q
  \sum_{\substack{d\ge1\\(d,q)=1}}\frac{\mu(d)}{d^2}
 =\frac{\varphi(q)}q
  \prod_{p\nmid q}\left(1-\frac1{p^2}\right)>0
 \tag{2.6}
\]
such that
\[
 \boxed{A_{\chi,Y}=c_qY+O_q(Y^{1/2}).}
 \tag{2.7}
\]
In particular $A_{\chi,Y}\asymp_qY$.
\end{lemma}

\begin{proof}
Let
\[
 N_q(x)=\sum_{\substack{n\le x\\(n,q)=1}}\mu(n)^2.
\]
Using $\mu(n)^2=\sum_{d^2\mid n}\mu(d)$ and writing $n=d^2m$ gives
\[
 N_q(x)
 =\sum_{\substack{d\le\sqrt x\\(d,q)=1}}\mu(d)
   \#\{m\le x/d^2:(m,q)=1\}.
\]
For fixed $q$,
\[
 \#\{m\le z:(m,q)=1\}=\frac{\varphi(q)}qz+O_q(1).
\]
Therefore
\[
 N_q(x)
 =\frac{\varphi(q)}q x
  \sum_{\substack{d\le\sqrt x\\(d,q)=1}}\frac{\mu(d)}{d^2}
  +O_q(\sqrt x).
\]
The tail of the absolutely convergent $d^{-2}$ series is $O(x^{-1/2})$,
so
\[
 N_q(x)=c_qx+O_q(\sqrt x).
\]
Subtracting $N_q(Y)$ from $N_q(2Y)$ proves \textup{(2.7)}.  The Euler product
in \textup{(2.6)} is positive.
\end{proof}

% -----------------------------------------------------------------------------
\subsection{The contracted phase action}
% -----------------------------------------------------------------------------

Identify $\C$ with $\R^2$ and put
\[
 V_\chi:=\Span_\R G_\chi\subseteq\C,
 \qquad
 \ip{z}{w}_\R:=\operatorname{Re}(\overline zw).
\]
If $G_\chi=\{\pm1\}$, then $V_\chi=\R$; otherwise $V_\chi=\C$ as a real
space.

\begin{definition}[Uniform phase pressure and contracted action]
For $\theta\in V_\chi$ define
\[
 \Lambda_\chi(\theta)
 :=\log\left(
   \frac1{|G_\chi|}\sum_{g\in G_\chi}
   e^{\ip{\theta}{g}_\R}
 \right).
 \tag{2.8}
\]
For $z\in\conv(G_\chi)$ define the Legendre transform
\[
 \boxed{
 I_\chi(z)
 :=\sup_{\theta\in V_\chi}
 \{\ip{\theta}{z}_\R-\Lambda_\chi(\theta)\}.}
 \tag{2.9}
\]
The extensive dyadic phase action is
\[
 \boxed{
 \mathfrak J_\chi(Y)
 :=A_{\chi,Y}\,
 I_\chi\!\left(\frac{S_\chi(Y)}{A_{\chi,Y}}\right).}
 \tag{2.10}
\]
\end{definition}

\begin{theorem}[Entropy contraction onto the protected phase mean]
\label{thm:phase-entropy-contraction}
For every $z\in\conv(G_\chi)$,
\[
 \boxed{
 I_\chi(z)
 =\min\left\{
  \KL(p\|u_\chi):
  p\in\Prob(G_\chi),\ \sum_gp(g)g=z
 \right\}.}
 \tag{2.11}
\]
The minimizer is unique.  For $z\in\relint\conv(G_\chi)$ it has the
exponential form
\[
 \boxed{
 p_{\theta_z}(g)
 =u_\chi(g)
  e^{\ip{\theta_z}{g}_\R-\Lambda_\chi(\theta_z)},
 \qquad
 \nabla\Lambda_\chi(\theta_z)=z.}
 \tag{2.12}
\]
\end{theorem}

\begin{proof}
For any $\theta\in V_\chi$ and any probability $p$ on $G_\chi$,
\[
 \begin{aligned}
 \KL(p\|p_\theta)
 &=\sum_gp(g)\log\frac{p(g)}{u_\chi(g)}
   -\ip{\theta}{\sum_gp(g)g}_\R
   +\Lambda_\chi(\theta)\\
 &=\KL(p\|u_\chi)-\ip{\theta}{\E_p g}_\R
   +\Lambda_\chi(\theta).
 \end{aligned}
\]
If $\E_pg=z$, nonnegativity gives
\[
 \KL(p\|u_\chi)
 \ge\ip{\theta}{z}_\R-\Lambda_\chi(\theta).
\]
Taking the supremum in $\theta$ proves that the right side of
\textup{(2.11)} is at least $I_\chi(z)$.

The feasible set is a nonempty compact convex polytope and relative entropy is
strictly convex, so it has a unique minimizer.  In the relative interior, the
Lagrange equations for normalization and the real mean constraint give
\textup{(2.12)}.  Substitution into the preceding identity gives equality with
$I_\chi(z)$.  Boundary points follow by taking a sequence from the relative
interior and using compactness and lower semicontinuity.
\end{proof}

\begin{theorem}[Explicit quadratic bounds]
\label{thm:phase-action-quadratic}
For every $z\in\conv(G_\chi)$,
\[
 \boxed{I_\chi(z)\ge\frac12|z|^2.}
 \tag{2.13}
\]
For $|z|\le1/2$,
\[
 \boxed{I_\chi(z)\le2|z|^2.}
 \tag{2.14}
\]
If $G_\chi=\{\pm1\}$, the constant two in \textup{(2.14)} may be replaced
by one.
\end{theorem}

\begin{proof}
Let $p$ have mean $z$.  By \Cref{lem:uniform-phase-centered},
\[
 |z|
 =\left|\sum_g(p(g)-u_\chi(g))g\right|
 \le\norm{p-u_\chi}_1.
\]
Pinsker's inequality gives
$\norm{p-u_\chi}_1^2\le2\KL(p\|u_\chi)$.  Minimize over $p$ and use
\Cref{thm:phase-entropy-contraction} to obtain \textup{(2.13)}.

Suppose first that $|G_\chi|\ge4$.  Then
\[
 \E_{u_\chi}g=0,
 \qquad
 \E_{u_\chi}g^2=0,
 \qquad
 \E_{u_\chi}|g|^2=1.
\]
For $|z|\le1/2$ define
\[
 p_z(g)=u_\chi(g)\bigl(1+2\operatorname{Re}(\overline zg)\bigr).
\]
This is nonnegative and normalized.  Moreover
\[
 \E_{p_z}g
 =\E_u\bigl[g+g(\overline zg+z\overline g)\bigr]=z.
\]
Since $\log x\le x-1$,
\[
 \KL(p_z\|u_\chi)
 \le\chi^2(p_z,u_\chi)
 =\E_u\bigl[2\operatorname{Re}(\overline zg)\bigr]^2
 =2|z|^2.
\]
Apply \Cref{thm:phase-entropy-contraction}.  If
$G_\chi=\{\pm1\}$, $z\in[-1,1]$ and the law
$p_z(\pm1)=(1\pm z)/2$ has mean $z$ and chi-square divergence $z^2$.
\end{proof}

\begin{theorem}[The phase action is exactly the square-root Mertens target]
\label{thm:action-Mertens-equivalence}
For fixed primitive $\chi$, the following two dyadic estimates are equivalent:
\[
 \boxed{\mathfrak J_\chi(Y)=Y^{o(1)},}
 \tag{2.15}
\]
\[
 \boxed{|S_\chi(Y)|^2=Y^{1+o(1)}.}
 \tag{2.16}
\]
More quantitatively,
\[
 \boxed{|S_\chi(Y)|^2
 \le2A_{\chi,Y}\mathfrak J_\chi(Y)
 \le2Y\mathfrak J_\chi(Y).}
 \tag{2.17}
\]
Conversely, once $|S_\chi(Y)|/A_{\chi,Y}\le1/2$,
\[
 \boxed{\mathfrak J_\chi(Y)
 \le\frac{2|S_\chi(Y)|^2}{A_{\chi,Y}}.}
 \tag{2.18}
\]
\end{theorem}

\begin{proof}
Apply \textup{(2.13)} to $z=S_\chi(Y)/A_{\chi,Y}$ and multiply by
$A_{\chi,Y}$ to obtain the first inequality in \textup{(2.17)}.  The second
uses $A_{\chi,Y}\le Y$.  Hence \textup{(2.15)} implies
\textup{(2.16)}.

Conversely, \Cref{lem:active-density} and \textup{(2.16)} give
$|S_\chi(Y)|/A_{\chi,Y}=Y^{-1/2+o(1)}$, so the ratio is eventually at most
$1/2$.  Apply \textup{(2.14)} and multiply by $A_{\chi,Y}$ to get
\textup{(2.18)}.  The density lower bound then turns \textup{(2.16)} into
\textup{(2.15)}.
\end{proof}

\begin{theorem}[Dyadic phase-action criterion for GRH]
\label{thm:phase-action-GRH}
Assume that for every primitive Dirichlet character $\chi$ and every dyadic
shell,
\[
 \mathfrak J_\chi(Y)=Y^{o(1)}.
 \tag{2.19}
\]
Then, for every $\varepsilon>0$,
\[
 \boxed{
 M_\chi(X):=\sum_{n\le X}\mu(n)\chi(n)
 =O_{\chi,\varepsilon}(X^{1/2+\varepsilon}).}
 \tag{2.20}
\]
Consequently every primitive Dirichlet $L$-function has all nontrivial zeros
on $\operatorname{Re}s=1/2$; hence GRH holds.
\end{theorem}

\begin{proof}
By \Cref{thm:action-Mertens-equivalence}, for every $\varepsilon>0$ and all
large $Y$,
\[
 |S_\chi(Y)|\le C_{\chi,\varepsilon}Y^{1/2+\varepsilon}.
\]
Partition $(1,X]$ into the shells
$(X/2^{j+1},X/2^j]$.  Their bounds form a convergent geometric series, so,
after absorbing the finitely many small shells,
\[
 |M_\chi(X)|
 \le C'_{\chi,\varepsilon}X^{1/2+\varepsilon}.
\]

Let $a_n=\mu(n)\chi(n)$.  For $s=\sigma+it$ and $N\ge1$, partial summation
gives
\[
 \sum_{n\le N}\frac{a_n}{n^s}
 =\frac{M_\chi(N)}{N^s}
  +s\int_1^NM_\chi(x)x^{-s-1}\,\dd x.
\]
If $\sigma>1/2$, choose
$0<\varepsilon<\sigma-1/2$.  The boundary term tends to zero and the integral
converges absolutely.  Thus
\[
 F_\chi(s):=\sum_{n\ge1}\frac{\mu(n)\chi(n)}{n^s}
\]
is analytic on $\operatorname{Re}s>1/2$.  On $\operatorname{Re}s>1$ the
Euler product gives $F_\chi(s)L(s,\chi)=1$.  By analytic continuation this
identity holds throughout $\operatorname{Re}s>1/2$ (meromorphically at the
principal pole), so $L(s,\chi)$ has no zero there.  The functional equation
reflects every nontrivial zero $\rho$ to $1-\overline\rho$.  Hence there is no
zero with real part below $1/2$ either, and all nontrivial zeros lie on the
critical line.
\end{proof}

\begin{remark}[Why the action orientation matters]
The accepted/action calculus naturally produces
$\KL(\text{hard}\|\text{reference})$.  This is the correct orientation for a
source-fixed comparator and has an exact path chain rule.  Although Pinsker's
inequality is available in either orientation, the hard-to-reference divergence
alone decomposes canonically into local likelihood costs without fitting the
reference to the terminal phase counts.
\end{remark}

% =============================================================================
\section{The ratio-complete Cayley source and its source-fixed reference}
\label{sec:Cayley-reference}
% =============================================================================

Put
\[
 R_\chi:=-H_\chi\subseteq G_\chi.
 \tag{3.1}
\]
The previously constructed ratio-complete replacement bank contains every ratio in
$R_\chi$ with equal outgoing source weight before terminal aggregation.  The
following finite calculation is included for completeness.

\begin{definition}[Uniform ratio operator]
On $\ell^2(G_\chi,u_\chi)$ define
\[
 (P_\chi f)(g)
 =\frac1{|H_\chi|}\sum_{\eta\in R_\chi}f(\eta g),
 \qquad
 L_\chi^{\Phase}:=I-P_\chi.
 \tag{3.2}
\]
\end{definition}

\begin{theorem}[Exact Cayley spectrum and phase gap]
\label{thm:Cayley-spectrum}
The operator $P_\chi$ is self-adjoint and Markov.  For
$\psi\in\widehat G_\chi$,
\[
 P_\chi\psi=\widehat u_{R_\chi}(\psi)\psi,
\]
where
\[
 \boxed{
 \widehat u_{R_\chi}(\psi)
 =\begin{cases}
 1,&\psi=1,\\
 -1,&\psi\ne1\text{ and }\psi|_{H_\chi}=1,\\
 0,&\psi|_{H_\chi}\ne1.
 \end{cases}}
 \tag{3.3}
\]
(The middle case is absent when $G_\chi=H_\chi$.)  Consequently
\[
 \boxed{\operatorname{spec}(L_\chi^{\Phase})\subseteq\{0,1,2\},
 \qquad
 L_\chi^{\Phase}\succeq P_{\one^\perp}.}
 \tag{3.4}
\]
\end{theorem}

\begin{proof}
The set $R_\chi=-H_\chi$ is invariant under inversion, so convolution by its
uniform law is self-adjoint.  Characters diagonalize convolution, and
\[
 \widehat u_{R_\chi}(\psi)
 =\psi(-1)\frac1{|H_\chi|}\sum_{h\in H_\chi}\psi(h).
\]
The subgroup sum vanishes unless $\psi$ is trivial on $H_\chi$.  In that
case it equals one for the trivial character and, when the quotient
$G_\chi/H_\chi$ is nontrivial, equals minus one for its unique nontrivial
character.  Subtracting from one proves \textup{(3.4)}.
\end{proof}

\begin{theorem}[Poissonized source-fixed phase kernel]
\label{thm:Poissonized-phase-kernel}
For $t>0$ put
\[
 K_{\chi,t}:=e^{-tL_\chi^{\Phase}}.
 \tag{3.5}
\]
Then $K_{\chi,t}$ is a symmetric strictly positive Markov kernel, its invariant
law is $u_\chi$, and
\[
 \boxed{
 \norm{K_{\chi,t}f}_2\le e^{-t}\norm f_2
 \quad(f\perp\one).}
 \tag{3.6}
\]
The joint law
\[
 \boxed{
 R_{\chi,t}(g,h):=u_\chi(g)K_{\chi,t}(g,h)}
 \tag{3.7}
\]
has both source and target marginals equal to $u_\chi$ and is fixed without
using hard terminal phase counts.
\end{theorem}

\begin{proof}
Since $P_\chi$ is Markov,
\[
 K_{\chi,t}
 =e^{-t}e^{tP_\chi}
 =e^{-t}\sum_{n=0}^\infty\frac{t^n}{n!}P_\chi^n
\]
is Markov and symmetric.  The support $R_\chi$ generates $G_\chi$: if
$G_\chi=H_\chi$ it is the whole group, while in the index-two case odd and
even products cover the two cosets.  Hence for every pair $g,h$ some power
$P_\chi^n(g,h)$ is positive, and the exponential series makes
$K_{\chi,t}(g,h)>0$.  The spectral estimate follows from
\Cref{thm:Cayley-spectrum}.  Symmetry and preservation of constants give both
marginals of \textup{(3.7)}.
\end{proof}

\begin{definition}[Source-fixed relational phase reference]
\label{def:source-fixed-reference}
A finite ratio-labelled relational reference consists of:
\begin{enumerate}[label=\textup{(Q\arabic*)},leftmargin=3em]
\item a finite source record $S$, a finite history record $\mathcal E$, and a
      finite target record $T$, with maps
      $s:\mathcal E\to S$ and $t:\mathcal E\to T$;
\item source and target phase maps
      $\phi_S:S\to G_\chi$ and $\phi_T:T\to G_\chi$;
\item a protected ratio mark $\eta:\mathcal E\to G_\chi$ satisfying
      \[
       \phi_T(t(e))=\eta(e)\phi_S(s(e));
      \]
\item a faithful reference law $R$ on $\mathcal E$ whose source phase is
      uniform and whose conditional ratio law is the source-fixed
      Poissonized Cayley kernel, possibly with additional normalized
      provenance fibres independent of the phase quotient.
\end{enumerate}
The reference is \emph{source fixed}: its law is determined by the local ratio
instrument and the chosen physical phase time $t>0$, before hard terminal
counts are read.
\end{definition}

\begin{proposition}[Phase neutrality of the relational reference]
\label{prop:reference-phase-neutral}
Every reference of \Cref{def:source-fixed-reference} has
\[
 \boxed{(\phi_T\circ t)_*R=u_\chi.}
 \tag{3.8}
\]
\end{proposition}

\begin{proof}
After the normalized provenance fibres are summed out, the source--target phase
law is \textup{(3.7)}.  Its target marginal is uniform by
\Cref{thm:Poissonized-phase-kernel}.
\end{proof}

\begin{remark}[No terminal fitting]
Replacing $u_\chi$ in \textup{(3.7)} by the observed hard phase law makes the
comparison cost vanish identically and has no arithmetic content.  The
reference time, ratio weights, source phase law, and provenance decoder must be
fixed from the local relational instrument before the hard target marginal is
used.
\end{remark}

% =============================================================================
\section{Reference-preserving relational descent as entropic transport}
\label{sec:entropic-descent}
% =============================================================================

Fix one dyadic packet and abbreviate
\[
 \Omega=\Omega_{\chi,Y},\qquad A=A_{\chi,Y},
 \qquad G=G_\chi,\qquad u=u_\chi,\qquad\tau=\tau_{\chi,Y}.
\]
Let $\beta$ be the uniform hard law on $\Omega$.  A relational history may have
more than one provenance realization of the same source--target pair; these
histories remain separate elements of $\mathcal E$.

\begin{definition}[Phase and hard descent polytopes]
Let $R$ be a source-fixed relational reference with source marginal
$\alpha=s_*R$.  Define
\[
 \Pi_{\Phase}(\alpha,\tau)
 :=\{\pi\in\Prob(\mathcal E):
      s_*\pi=\alpha,
      (\phi_T\circ t)_*\pi=\tau\},
 \tag{4.1}
\]
\[
 \Pi_{\Hard}(\alpha,\beta)
 :=\{\pi\in\Prob(\mathcal E):
      s_*\pi=\alpha,
      t_*\pi=\beta\}.
 \tag{4.2}
\]
Their entropic descent actions are
\[
 \boxed{
 \mathfrak D_{\Phase}(Y)
 :=\inf_{\pi\in\Pi_{\Phase}(\alpha,\tau)}\KL(\pi\|R),}
 \tag{4.3}
\]
\[
 \boxed{
 \mathfrak D_{\Hard}(Y)
 :=\inf_{\pi\in\Pi_{\Hard}(\alpha,\beta)}\KL(\pi\|R),}
 \tag{4.4}
\]
with value $+\infty$ when the corresponding polytope is empty.
\end{definition}

\begin{remark}
The phase action simultaneously shorts all within-phase target identities.  It
is the load-bearing protected quotient.  The hard action is stronger: it asks
for the actual uniform integer target law on the same physical history graph.
Every hard descent is a phase descent, so
$\mathfrak D_{\Phase}\le\mathfrak D_{\Hard}$.
\end{remark}

\begin{theorem}[Relational action dominates the exact GRH action]
\label{thm:relational-dominates-phase}
For every dyadic packet,
\[
 \boxed{
 \mathfrak J_\chi(Y)
 \le A\KL(\tau\|u)
 \le A\mathfrak D_{\Phase}(Y)
 \le A\mathfrak D_{\Hard}(Y).}
 \tag{4.5}
\]
Consequently either of the estimates
\[
 A\mathfrak D_{\Phase}(Y)=Y^{o(1)},
 \qquad\text{or}\qquad
 A\mathfrak D_{\Hard}(Y)=Y^{o(1)}
 \tag{4.6}
\]
for every primitive $\chi$ and every dyadic shell proves GRH.
\end{theorem}

\begin{proof}
The entropy contraction theorem gives
\[
 I_\chi\!\left(\sum_g\tau(g)g\right)\le\KL(\tau\|u).
\]
For any $\pi\in\Pi_{\Phase}(\alpha,\tau)$, data processing under the terminal
phase map gives
\[
 \KL(\tau\|u)
 \le\KL(\pi\|R),
\]
because the terminal phase pushforward of $R$ is $u$ by
\Cref{prop:reference-phase-neutral}.  Take the infimum and use the inclusion of
hard descents into phase descents.  Multiply by $A$.  The final assertion is
\Cref{thm:phase-action-GRH}.
\end{proof}

\begin{remark}[The new load-bearing quantity]
The conjecture is reduced to the extensive, source-fixed cost
\[
 \boxed{A_{\chi,Y}\mathfrak D_{\Phase}(Y).}
\]
This is not a terminal statistic inserted by hand.  It is the minimum
accepted/action likelihood needed to keep the independently fixed source law
while producing the hard phase marginal through actual ratio-labelled
histories.  The full hard descent action is a stronger physical certificate;
the phase action is the exact protected quotient relevant to GRH.
\end{remark}

% -----------------------------------------------------------------------------
\subsection{Weighted Hall feasibility}
% -----------------------------------------------------------------------------

The next theorem applies both to the target $Z=G$ of the phase quotient and to
the full hard target $Z=\Omega$.  Let $z:\mathcal E\to Z$ be the relevant
target map, and let $\nu$ be the desired target law.

\begin{definition}[Neighbour set]
For $U\subseteq S$ put
\[
 N(U):=\{v\in Z:\exists e\in\mathcal E,
          \ s(e)\in U,\ z(e)=v\}.
 \tag{4.7}
\]
\end{definition}

\begin{theorem}[Exact weighted Hall compiler]
\label{thm:weighted-Hall}
There exists a probability $\pi$ on $\mathcal E$ satisfying
\[
 s_*\pi=\alpha,
 \qquad
 z_*\pi=\nu
 \tag{4.8}
\]
if and only if
\[
 \boxed{\alpha(U)\le\nu(N(U))
 \quad\text{for every }U\subseteq S.}
 \tag{4.9}
\]
Failure returns a concrete source set $U$ whose mass cannot be transported
through the actual relational histories.
\end{theorem}

\begin{proof}
Necessity follows because all mass emitted by $U$ must arrive in $N(U)$.
For sufficiency, build a finite network with a supersource joined to each
$s\in S$ by capacity $\alpha(s)$, an infinite-capacity edge from $s$ to $v$
whenever an admitted history connects them, and an edge from each $v\in Z$ to
the sink of capacity $\nu(v)$.  A cut determined by $U\subseteq S$ and its
neighbour set has capacity
\[
 1-\alpha(U)+\nu(N(U)).
\]
Condition \textup{(4.9)} makes every cut have capacity at least one.  The
max-flow/min-cut theorem gives a unit flow.  Split the flow among parallel
relational histories to obtain $\pi$.
\end{proof}

\begin{corollary}[Degree--congestion sufficient condition]
\label{cor:degree-Hall}
Suppose $S$ and $Z$ carry uniform laws, every source has at least $d$ admitted
neighbours, every target has at most $\Delta$ source neighbours, and
\[
 \boxed{d|S|\ge\Delta|Z|.}
 \tag{4.10}
\]
Then the Hall condition holds.
\end{corollary}

\begin{proof}
For $U\subseteq S$, count incident edges in two ways:
\[
 d|U|\le |E(U,N(U))|\le\Delta|N(U)|.
\]
Thus
\[
 \frac{|N(U)|}{|Z|}
 \ge\frac d\Delta\frac{|U|}{|Z|}
 \ge\frac{|U|}{|S|},
\]
which is \textup{(4.9)} for uniform laws.
\end{proof}

\begin{remark}[Interpretation]
The previously constructed ratio-availability theorem controls the outgoing side of this
criterion.  The genuinely new arithmetic quantity is target congestion after
all-strata aggregation.  A failed descent is therefore not an unspecified
lack of mixing: it is a weighted Hall cut or an excessive target indegree on
the actual relational graph.
\end{remark}

% -----------------------------------------------------------------------------
\subsection{The entropic minimizer and pressure dual}
% -----------------------------------------------------------------------------

Let $R(e)>0$ on every admitted history.  For two probability laws on
$\mathcal E$, use
\[
 \KL(\pi\|R)
 =\sum_e\pi(e)\log\frac{\pi(e)}{R(e)}
 =\sum_eR(e)h\!\left(\frac{\pi(e)}{R(e)}\right),
 \quad
 h(x)=x\log x-x+1.
 \tag{4.11}
\]
The second equality uses equality of the total masses.

\begin{theorem}[Unique entropic descent and pressure duality]
\label{thm:entropic-dual}
Assume the Hall condition, and let
\[
 \Pi(\alpha,\nu)
 =\{\pi\in\Prob(\mathcal E):s_*\pi=\alpha,
                                  z_*\pi=\nu\}.
\]
Then:
\begin{enumerate}[label=\textup{(E\arabic*)},leftmargin=3em]
\item the minimum
      \[
       \mathfrak D(\alpha,\nu\mid R)
       :=\min_{\pi\in\Pi(\alpha,\nu)}\KL(\pi\|R)
      \]
      exists and is unique;
\item it has the exact dual representation
      \[
       \boxed{
       \mathfrak D(\alpha,\nu\mid R)
       =\sup_{a\in\R^S,\,b\in\R^Z}
       \left\{
        \sum_s\alpha(s)a_s+\sum_v\nu(v)b_v
        -\sum_{e\in\mathcal E}R(e)
         \bigl(e^{a_{s(e)}+b_{z(e)}}-1\bigr)
       \right\};}
       \tag{4.12}
      \]
\item if the transport polytope contains a law positive on every admitted
      history, then the minimizing law has the Sinkhorn--pressure form
      \[
       \boxed{
       \pi_*(e)=R(e)e^{a_{s(e)}+b_{z(e)}}}
       \tag{4.13}
      \]
      for dual optimizers $a,b$, unique up to
      $a\mapsto a+c$, $b\mapsto b-c$.
\end{enumerate}
Every trial pair $(a,b)$ in \textup{(4.12)} is a rigorous lower-bound
certificate for the descent action.
\end{theorem}

\begin{proof}
The feasible set is a nonempty compact polytope.  Relative entropy is lower
semicontinuous and strictly convex on the probability simplex because
$x\mapsto x\log x$ is strictly convex.  This proves existence and uniqueness.

Introduce multipliers $a_s,b_v$ for the two marginal constraints.  Since the
convex conjugate of $h$ is
\[
 h^*(y)=e^y-1,
\]
minimization over each nonnegative coordinate $\pi(e)$ gives
\[
 \inf_{x\ge0}R(e)h(x/R(e))-(a_{s(e)}+b_{z(e)})x
 =-R(e)(e^{a_{s(e)}+b_{z(e)}}-1).
\]
Summing and adding the marginal terms gives the right side of
\textup{(4.12)}.  Finite-dimensional convex duality is exact after restricting
to the minimal face of the feasible transport polytope; forced-zero edges are
obtained as limits of dual potentials.  This proves the dual formula.

If a strictly positive feasible point exists, the minimizer lies in the
relative interior and the coordinate stationarity equations are
\[
 \log\frac{\pi_*(e)}{R(e)}=a_{s(e)}+b_{z(e)},
\]
which is \textup{(4.13)}.  The one-dimensional gauge freedom follows because
the two marginal totals agree.
\end{proof}

\begin{remark}[The dual is a witness, not a fitted proof]
The target pressure $b$ may be used to certify that the action is large.  It
may not be chosen from the terminal Mertens sum and then advertised as an
independent source law.  To prove a small action one must construct a feasible
relational flow or derive a uniform upper bound from the source graph,
capacities, and common action.
\end{remark}

% =============================================================================
\section{Same-history path action and the local-to-global bridge}
\label{sec:path-action}
% =============================================================================

The persistent relation packet retains the source state, intermediate ratio
and provenance marks, accepted/action record, and target state on one history.
This makes the following chain rule a table-native finite identity.

\begin{theorem}[Discrete relational path-action chain rule]
\label{thm:path-chain-rule}
Let $P$ and $R$ be faithful laws on a finite path
$X_0,X_1,\ldots,X_m$.  Write their conditional kernels, with the complete past
retained, as
\[
 p_j(\cdot\mid X_{<j}),
 \qquad
 r_j(\cdot\mid X_{<j}).
\]
Then
\[
 \boxed{
 \KL(P\|R)
 =\KL(P_0\|R_0)
  +\sum_{j=1}^m
   \E_P\KL\bigl(
    p_j(\cdot\mid X_{<j})
    \|r_j(\cdot\mid X_{<j})
   \bigr).}
 \tag{5.1}
\]
If $P_0=R_0$, the initial term vanishes.  If
$C_j=C_j(X_j)$ is a coarse protected coordinate, then every summand further
splits as
\[
 \boxed{
 \begin{aligned}
 &\KL(p_j\|r_j)\\
 &=\KL(p_j^{C_j}\|r_j^{C_j})
  +\sum_c p_j^{C_j}(c)
    \KL\bigl(p_j(\cdot\mid c)\|r_j(\cdot\mid c)\bigr).
 \end{aligned}}
 \tag{5.2}
\]
Thus phase-ratio mismatch and hidden provenance/target mismatch are distinct
nonnegative action coordinates.
\end{theorem}

\begin{proof}
The path likelihood ratio factors as
\[
 \frac{P(X_0,\ldots,X_m)}{R(X_0,\ldots,X_m)}
 =\frac{P_0(X_0)}{R_0(X_0)}
  \prod_{j=1}^m
  \frac{p_j(X_j\mid X_{<j})}{r_j(X_j\mid X_{<j})}.
\]
Take logarithms and expectation under $P$ to obtain \textup{(5.1)}.  Formula
\textup{(5.2)} is the ordinary relative-entropy chain rule under the map
$C_j$.
\end{proof}

\begin{corollary}[Continuous-time accepted/action formula]
\label{cor:continuous-path-action}
Let $(p,L)$ and $(p,\widetilde L)$ be finite continuous-time Markov models with
the same initial law and common directed support.  Put
\[
 \PhiRate(a,b):=a\log(a/b)-a+b.
\]
Then on $[0,T]$,
\[
 \boxed{
 \KL(\mathbb P_{p,L}^T\|\mathbb P_{p,\widetilde L}^T)
 =\int_0^T\sum_up_t(u)
   \sum_{v\ne u}\PhiRate(L(u,v),\widetilde L(u,v))\,\dd t.}
 \tag{5.3}
\]
Occupation times and directed jump counts are a finite sufficient record for
this likelihood.
\end{corollary}

\begin{proof}
The likelihood of a resolved path is the product of the initial mass, all jump
rates, and all exponential holding factors.  Their logarithmic ratio is
\[
 \sum_{u\ne v}N_{uv}\log\frac{L(u,v)}{\widetilde L(u,v)}
 -\sum_u\Theta_u(\lambda_L(u)-\lambda_{\widetilde L}(u)).
\]
Taking expectation and using
$\E[N_{uv}(\dd t)]=p_t(u)L(u,v)\dd t$ gives \textup{(5.3)}.
\end{proof}

\begin{theorem}[Local accepted/action criterion for GRH]
\label{thm:local-action-GRH}
For every primitive $\chi$ and dyadic shell $Y$, suppose a source-fixed
relational reference and a feasible phase descent path law $P_{\chi,Y}$ are
constructed on the same protected cylinder.  Let
\[
 \mathcal A_{\chi,Y}
 :=\KL(P_{\chi,Y}\|R_{\chi,Y}),
 \tag{5.4}
\]
or equivalently use the sum of local conditional actions in
\Cref{thm:path-chain-rule}.  If
\[
 \boxed{
 A_{\chi,Y}\mathcal A_{\chi,Y}=Y^{o(1)}}
 \tag{5.5}
\]
for every primitive $\chi$, then GRH holds.
\end{theorem}

\begin{proof}
The constructed path law is an admissible element of
$\Pi_{\Phase}(\alpha,\tau)$, so
\[
 \mathfrak D_{\Phase}(Y)
 \le\mathcal A_{\chi,Y}.
\]
Apply \Cref{thm:relational-dominates-phase,thm:phase-action-GRH}.
\end{proof}

\begin{corollary}[A sufficient local increment budget]
\label{cor:local-increment-budget}
In the setting of \Cref{thm:path-chain-rule}, suppose $P_0=R_0$ and
\[
 \sum_{j=1}^{m_Y}
 \E_P\KL(p_j(\cdot\mid X_{<j})\|r_j(\cdot\mid X_{<j}))
 \le A_{\chi,Y}^{-1}Y^{o(1)}.
 \tag{5.6}
\]
Then the shell passes the GRH action criterion.  In particular, if
$m_Y=Y^{o(1)}$, it is enough that every local conditional action be
$A_{\chi,Y}^{-1}Y^{o(1)}/m_Y$ uniformly.
\end{corollary}

\begin{remark}[Correct scale of the estimate]
The normalized divergence must be of order $A^{-1}$ up to subpower factors.
This is not an artifact: a square-root fluctuation of a packet of size $A$
has mean bias $A^{-1/2}$ and relative-entropy cost of order $A^{-1}$.  The
natural quantity is therefore the \emph{extensive} action
$A\mathcal A$, not the unscaled one-step divergence.
\end{remark}

% =============================================================================
\section{The source-shorted influence Laplacian}
\label{sec:influence}
% =============================================================================

The entropic problem is nonlinear.  Near the source-fixed reference it has an
exact quadratic compiler obtained by shorting all source-marginal directions
before measuring the target discrepancy.

Let $\mathcal E$ be a finite history set, let $R(e)>0$ be a reference
probability, and let
\[
 s:\mathcal E\to S,
 \qquad
 z:\mathcal E\to Z.
\]
Write
\[
 \alpha=s_*R,
 \qquad
 \rho=z_*R.
\]
On $\R^{\mathcal E}$ use the weighted inner product
\[
 \ip{h}{k}_{R^{-1}}
 :=\sum_{e\in\mathcal E}\frac{h(e)k(e)}{R(e)}.
 \tag{6.1}
\]
Let
\[
 (Ah)(s)=\sum_{e:s(e)=s}h(e),
 \qquad
 (Bh)(v)=\sum_{e:z(e)=v}h(e).
 \tag{6.2}
\]
Euclidean inner products are used on $\R^S$ and $\R^Z$.

\begin{lemma}[Weighted adjoints and source short]
\label{lem:weighted-adjoints}
The weighted adjoints are
\[
 (A^*f)(e)=R(e)f(s(e)),
 \qquad
 (B^*g)(e)=R(e)g(z(e)).
 \tag{6.3}
\]
Moreover
\[
 AA^*=\diag(\alpha),
 \qquad
 BB^*=\diag(\rho).
 \tag{6.4}
\]
The orthogonal projection onto $\Ker A$ is
\[
 \boxed{
 P_S=I-A^*\diag(\alpha)^{-1}A.}
 \tag{6.5}
\]
\end{lemma}

\begin{proof}
For example,
\[
 \ip{h}{A^*f}_{R^{-1}}
 =\sum_eh(e)f(s(e))
 =\ip{Ah}{f}.
\]
The remaining identities follow by direct summation.  Since
$A^*(AA^*)^{-1}A$ is the orthogonal projection onto $\Ran A^*$,
\textup{(6.5)} projects onto its orthogonal complement $\Ker A$.
\end{proof}

\begin{definition}[Descent influence Laplacian]
The source-shorted target influence is
\[
 \boxed{
 \mathcal L_R:=BP_SB^*\succeq0.}
 \tag{6.6}
\]
If $R_{sv}$ denotes the aggregate reference mass of histories connecting
$s$ to $v$, then
\[
 \boxed{
 \mathcal L_R
 =\diag(\rho)-R^\mathsf T\diag(\alpha)^{-1}R.}
 \tag{6.7}
\]
\end{definition}

\begin{theorem}[Exact Thomson energy for source-preserving descent]
\label{thm:Thomson-influence}
Let $d\in\R^Z$ satisfy $\sum_vd(v)=0$.  The linear correction problem
\[
 Ah=0,
 \qquad
 Bh=d
 \tag{6.8}
\]
is solvable exactly when $d\in\Ran\mathcal L_R$.  On that branch,
\[
 \boxed{
 \min_{Ah=0,\,Bh=d}\norm h_{R^{-1}}^2
 =\ip{d}{\mathcal L_R^\dagger d}.}
 \tag{6.9}
\]
The unique minimum-norm correction is
\[
 \boxed{
 h_*=P_SB^*\mathcal L_R^\dagger d.}
 \tag{6.10}
\]
\end{theorem}

\begin{proof}
Restrict $B$ to the Hilbert subspace $\Ker A$.  Its adjoint there is
$P_SB^*$, and its Gram on the target space is
$BP_SB^*=\mathcal L_R$.  The Moore--Penrose minimum-norm theorem gives
solvability precisely on the Gram range, the solution \textup{(6.10)}, and
\[
 \norm{h_*}_{R^{-1}}^2
 =\ip{\mathcal L_R^\dagger d}
        {\mathcal L_R\mathcal L_R^\dagger d}
 =\ip d{\mathcal L_R^\dagger d}.
\]
\end{proof}

\begin{theorem}[Quadratic influence controls entropic descent]
\label{thm:quadratic-to-entropy}
Assume $d\in\Ran\mathcal L_R$ and put $h_*$ as in
\Cref{thm:Thomson-influence}.  Define the leverage
\[
 \boxed{
 \vartheta_R(d)
 :=\max_{e\in\mathcal E}\frac{|h_*(e)|}{R(e)}.}
 \tag{6.11}
\]
If $\vartheta_R(d)\le1/2$, then $\pi_*=R+h_*$ is a positive probability with
source marginal $\alpha$ and target marginal $\rho+d$, and
\[
 \boxed{
 \frac13\ip d{\mathcal L_R^\dagger d}
 \le\KL(\pi_*\|R)
 \le\ip d{\mathcal L_R^\dagger d}.}
 \tag{6.12}
\]
Consequently
\[
 \boxed{
 \mathfrak D(\alpha,\rho+d\mid R)
 \le\ip d{\mathcal L_R^\dagger d}.}
 \tag{6.13}
\]
\end{theorem}

\begin{proof}
The two linear constraints imply that $R+h_*$ has the desired marginals and
total mass one.  The leverage condition gives positivity.  Put
$x_e=h_*(e)/R(e)$.  Since $\sum_eR(e)x_e=0$,
\[
 \KL(R+h_*\|R)
 =\sum_eR(e)\bigl[(1+x_e)\log(1+x_e)-x_e\bigr].
\]
For $|x|\le1/2$, the second derivative of the bracket is $1/(1+x)$, which
lies between $2/3$ and $2$.  Since the function and its first derivative
vanish at zero,
\[
 \frac13x^2
 \le(1+x)\log(1+x)-x
 \le x^2.
\]
Sum over $e$ and apply \textup{(6.9)}.  The minimum entropic action is no
larger than the action of this feasible correction.
\end{proof}

\begin{corollary}[A concrete subpower flow criterion]
\label{cor:influence-GRH}
For every dyadic packet, suppose the hard or phase target discrepancy $d_Y$
from the source-fixed reference satisfies
\[
 \vartheta_{R_Y}(d_Y)\le\frac12,
 \qquad
 A_{\chi,Y}\ip{d_Y}{\mathcal L_{R_Y}^\dagger d_Y}=Y^{o(1)}.
 \tag{6.14}
\]
Then the corresponding entropic descent action is subpower and the GRH
criterion passes.
\end{corollary}

\begin{proof}
Use \Cref{thm:quadratic-to-entropy} and then
\Cref{thm:relational-dominates-phase}.
\end{proof}

% -----------------------------------------------------------------------------
\subsection{Exact phase influence of the Poissonized Cayley reference}
% -----------------------------------------------------------------------------

Let $m_\chi=|G_\chi|$ and use the phase reference
$R_{\chi,t}(g,h)=m_\chi^{-1}K_{\chi,t}(g,h)$ from
\textup{(3.7)}.

\begin{theorem}[Uniform phase influence gap]
\label{thm:phase-influence-gap}
The source-shorted target influence on $\R^{G_\chi}$ is
\[
 \boxed{
 \mathcal L_{\chi,t}^{\Phase}
 =\frac1{m_\chi}\bigl(I-K_{\chi,t}^2\bigr).}
 \tag{6.15}
\]
Its kernel is exactly the constants and
\[
 \boxed{
 \mathcal L_{\chi,t}^{\Phase}
 \succeq
 \frac{1-e^{-2t}}{m_\chi}P_{\one^\perp}.}
 \tag{6.16}
\]
Hence for every mean-zero phase discrepancy $d$,
\[
 \boxed{
 \ip d{(\mathcal L_{\chi,t}^{\Phase})^\dagger d}
 \le\frac{m_\chi}{1-e^{-2t}}\norm d_2^2.}
 \tag{6.17}
\]
There is a finite constant $C_{\chi,t}$ such that the minimum-energy leverage
obeys
\[
 \boxed{
 \vartheta_{R_{\chi,t}}(d)
 \le C_{\chi,t}\norm d_2.}
 \tag{6.18}
\]
Thus the phase descent action is quadratically controlled in a fixed
neighbourhood of the uniform law.
\end{theorem}

\begin{proof}
Both source and target marginals are $u=m_\chi^{-1}\one$.  Formula
\textup{(6.7)} gives
\[
 \begin{aligned}
 \mathcal L_{\chi,t}^{\Phase}
 &=\frac1{m_\chi}I
  -\left(\frac1{m_\chi}K_{\chi,t}\right)^\mathsf T
   (m_\chi I)
   \left(\frac1{m_\chi}K_{\chi,t}\right)\\
 &=\frac1{m_\chi}(I-K_{\chi,t}^2),
 \end{aligned}
\]
using symmetry.  On a nonconstant character,
$L_\chi^{\Phase}$ has eigenvalue $\lambda\in\{1,2\}$ and
$K_{\chi,t}$ has eigenvalue $e^{-t\lambda}$.  The corresponding influence
eigenvalue is
$m_\chi^{-1}(1-e^{-2t\lambda})$, which is at least the right side of
\textup{(6.16)}.  This proves \textup{(6.17)}.

The map
$d\mapsto h_*=P_SB^*(\mathcal L_{\chi,t}^{\Phase})^\dagger d$ is linear on
the fixed finite-dimensional mean-zero phase space.  Every reference entry is
strictly positive by \Cref{thm:Poissonized-phase-kernel}.  Taking the maximum
of the finitely many weighted coordinate norms gives a finite constant
$C_{\chi,t}$ and proves \textup{(6.18)}.
\end{proof}

\begin{remark}[What the phase gap does and does not prove]
The uniform phase influence gap removes a local controllability obstruction:
every sufficiently small phase discrepancy can be produced by a
source-preserving perturbation with quadratic action.  It does not show that
the hard discrepancy is small.  The missing theorem is precisely that the
actual hard-counting descent occupies the central, subpower extensive-action
regime.
\end{remark}

% =============================================================================

\section{Complete carrier-birth and retuning calculus}
\label{sec:grh-birth-retuning}
\subsection{Universal old--detail endpoint formula}
\label{subsec:birth-block-formula}

Let \(\mathcal E\) be the transported old carrier and \(\mathcal D\) the new
detail carrier.  Let
\[
 \mathbb G^+
 =\begin{pmatrix}A&B\\ B^*&D\end{pmatrix}\succeq0,
 \qquad A\succ0,
 \tag{T.3}
\]
and put
\[
 \boxed{S=D-B^*A^{-1}B\succeq0.}
 \tag{T.4}
\]
For an endpoint readout \(f=e\oplus r\), define
\[
 \boxed{q=r-B^*A^{-1}e.}
 \tag{T.5}
\]
For a positive form \(H\), write
\[
 \Lambda_H(h)=\inf\{\Lambda\ge0:hh^*\preceq\Lambda H\}.
 \tag{T.6}
\]
This is \(\langle h,H^\dagger h\rangle\) on the faithful range and infinity
otherwise.

\begin{theorem}[Exact carrier-birth influence]
\label{thm:birth-carrier-birth}
One has
\[
 \boxed{f\in\Ran\mathbb G^+\iff q\in\Ran S.}
 \tag{T.7}
\]
On this branch,
\[
 \boxed{
 \Lambda_{\mathbb G^+}(f)
 =\langle e,A^{-1}e\rangle
  +\langle q,S^\dagger q\rangle.}
 \tag{T.8}
\]
If \(q\notin\Ran S\), then \(\Lambda_{\mathbb G^+}(f)=+\infty\).  The
nonnegative carrier-birth innovation is therefore
\[
 \boxed{\eta^{\rm birth}=\langle q,S^\dagger q\rangle.}
 \tag{T.9}
\]
\end{theorem}

\begin{proof}
Use the congruence factorization
\[
 \mathbb G^+
 =\begin{pmatrix}I&0\\B^*A^{-1}&I\end{pmatrix}
  \begin{pmatrix}A&0\\0&S\end{pmatrix}
  \begin{pmatrix}I&A^{-1}B\\0&I\end{pmatrix}.
 \tag{T.10}
\]
The inverse left factor sends \(e\oplus r\) to \(e\oplus q\).  The range
condition and the Moore--Penrose value then follow from the diagonal middle
block.
\end{proof}

\begin{corollary}[Scalar detail trichotomy]
\label{cor:birth-scalar-detail}
If \(\mathcal D=\C\), let \(s=D-B^*A^{-1}B\) and
\(q=r-B^*A^{-1}e\).  Then
\[
 \boxed{
 \eta^{\rm birth}=|q|^2/s\quad(s>0).}
 \tag{T.11}
\]
If \(s=0\), the innovation is zero when \(q=0\), while \(q\ne0\) gives
infinite influence.  Thus the zero-Schur branch is either exact prediction or
an endpoint-visible source null.
\end{corollary}

% -----------------------------------------------------------------------------
\subsection{Source-native prediction and Christoffel writer}
\label{subsec:birth-source-native}

Suppose the fine form is realized on one source carrier by
\[
 \mathscr S:\mathcal E\to\mathcal K,
 \qquad
 \mathscr Z:\mathcal D\to\mathcal K,
 \tag{T.12}
\]
with \(A=\mathscr S^*\mathscr S\),
\(B=\mathscr S^*\mathscr Z\), and \(D=\mathscr Z^*\mathscr Z\).  Put
\[
 P_{\mathscr S}=\mathscr SA^{-1}\mathscr S^*,
 \qquad
 \mathscr T=(I-P_{\mathscr S})\mathscr Z,
 \qquad
 d_{\rm old}=\mathscr SA^{-1}e.
 \tag{T.13}
\]

\begin{theorem}[Source-native cutoff writer]
\label{thm:birth-source-birth}
The Schur and endpoint residuals are
\[
 \boxed{S=\mathscr T^*\mathscr T,}
 \qquad
 \boxed{q=r-\mathscr Z^*d_{\rm old}.}
 \tag{T.14}
\]
Moreover,
\[
 \boxed{
 \eta^{\rm birth}
 =\min\{\|h\|^2:\mathscr T^*h=q\}.}
 \tag{T.15}
\]
Thus carrier birth is exactly the minimum source cost needed to reproduce the
part of the new endpoint readout not predicted by the old minimum decoder.
\end{theorem}

\begin{proof}
The first formula is the Gram form of the Anderson--Trapp short.  The second
uses \(B^*A^{-1}e=\mathscr Z^*d_{\rm old}\).  The last formula is the
Moore--Penrose minimum-norm theorem.
\end{proof}

\begin{definition}[Cutoff Christoffel packet]
\label{def:birth-Christoffel}
The source-fixed packet is
\[
 \boxed{
 q^{\rm cut}=r-\mathscr Z^*d_{\rm old},
 \qquad
 S^{\rm cut}=\mathscr Z^*(I-P_{\mathscr S})\mathscr Z,}
 \tag{T.16}
\]
with sharp cost
\[
 \boxed{
 \eta^{\rm cut}
 =\langle q^{\rm cut},(S^{\rm cut})^\dagger q^{\rm cut}\rangle.}
 \tag{T.17}
\]
\end{definition}

\begin{proposition}[Protected-readout inequality]
\label{prop:birth-protected-readout}
For a protected vector \(x\oplus\xi\),
\[
 \boxed{
 \begin{aligned}
 \left\langle\binom{x}{\xi},\mathbb G^+\binom{x}{\xi}\right\rangle
 &=(x+A^{-1}B\xi)^*A(x+A^{-1}B\xi)
   +\langle\xi,S\xi\rangle,\\
 &\ge\langle\xi,S\xi\rangle.
 \end{aligned}}
 \tag{T.18}
\]
On the faithful branch,
\[
 \boxed{
 |\langle q,\xi\rangle|^2
 \le\eta^{\rm birth}\langle\xi,S\xi\rangle.}
 \tag{T.19}
\]
\end{proposition}

\begin{proof}
Complete the square and apply Cauchy--Schwarz in the supported Schur form.
\end{proof}

% -----------------------------------------------------------------------------
\subsection{Signed old-sector retuning}
\label{subsec:birth-retuning}

Let \(G\succ0\) be the transported old-cutoff form and let \(A\succ0\) be the
old--old block seen at the fine cutoff.  Set
\[
 g=G^{-1}e,
 \qquad
 \widehat g=A^{-1}e,
 \tag{T.20}
\]
and define
\[
 \boxed{
 \mathcal C^{\rm ret}
 =\langle e,G^{-1}e\rangle-\langle e,A^{-1}e\rangle.}
 \tag{T.21}
\]

\begin{theorem}[Exact retuning and birth balance]
\label{thm:birth-retuning-balance}
The retuning capture is
\[
 \boxed{
 \mathcal C^{\rm ret}
 =\langle\widehat g,(A-G)g\rangle.}
 \tag{T.22}
\]
The complete fine-cutoff influence is
\[
 \boxed{
 \Lambda_{\mathbb G^+}(e\oplus r)
 =\Lambda_G(e)-\mathcal C^{\rm ret}+\eta^{\rm birth}.}
 \tag{T.23}
\]
Hence the favorable net decrement is
\[
 \boxed{
 \mathcal C^{\rm net}
 =\mathcal C^{\rm ret}-\eta^{\rm birth}.}
 \tag{T.24}
\]
\end{theorem}

\begin{proof}
The resolvent identity
\[
 G^{-1}-A^{-1}=A^{-1}(A-G)G^{-1}
\]
gives \textup{(T.22)}.  Combine it with
\Cref{thm:birth-carrier-birth}.
\end{proof}

On one common source carrier write
\[
 \boxed{
 \mathscr S^+
 =\bigl[\mathscr S_0+\mathscr D_{\rm old},\ \mathscr Z\bigr].}
 \tag{T.25}
\]
Then
\[
 \boxed{
 \begin{aligned}
 A-G
 &={\mathscr S_0}^*\mathscr D_{\rm old}
   +\mathscr D_{\rm old}^*\mathscr S_0
   +\mathscr D_{\rm old}^*\mathscr D_{\rm old},\\
 B&=(\mathscr S_0+\mathscr D_{\rm old})^*\mathscr Z,\\
 D&=\mathscr Z^*\mathscr Z.
 \end{aligned}}
 \tag{T.26}
\]

\begin{corollary}[Source form of retuning]
\label{cor:birth-source-retuning}
The signed retuning current is
\[
 \boxed{
 \begin{aligned}
 \mathcal C^{\rm ret}
 ={}&\langle\mathscr S_0\widehat g,\mathscr D_{\rm old}g\rangle
   +\langle\mathscr D_{\rm old}\widehat g,\mathscr S_0g\rangle\\
  &+\langle\mathscr D_{\rm old}\widehat g,
              \mathscr D_{\rm old}g\rangle.
 \end{aligned}}
 \tag{T.27}
\]
The sum is real, but its terms need not be separately positive or real.
\end{corollary}

\begin{proof}
Insert \textup{(T.26)} into \textup{(T.22)}.
\end{proof}

Thus
\[
 \boxed{
 \begin{aligned}
 \mathscr D_{\rm old}
 &\longrightarrow\text{the only potentially favorable signed capture},\\
 \mathscr Z
 &\longrightarrow\text{the nonnegative carrier-birth payment}.
 \end{aligned}}
 \tag{T.28}
\]

% -----------------------------------------------------------------------------
\subsection{Nested-source and uniform-retuning no-go theorems}
\label{subsec:birth-no-go}

\begin{proposition}[Pure source extension cannot lower influence]
\label{cth:birth-nested-source}
If the old source is strictly compatible,
\[
 \boxed{A=G,}
 \tag{T.29}
\]
then
\[
 \boxed{
 \Lambda_{\mathbb G^+}(e\oplus r)
 =\Lambda_G(e)+\eta^{\rm birth}\ge\Lambda_G(e).}
 \tag{T.30}
\]
Equality holds exactly when the new endpoint readout is perfectly predicted
on the faithful detail quotient.
\end{proposition}

\begin{proof}
The retuning current vanishes in \textup{(T.23)}.
\end{proof}

\begin{corollary}[Cumulative birth--retuning telescope]
\label{cor:birth-telescope}
Along a transported cutoff chain,
\[
 \boxed{
 \Lambda_M
 =\Lambda_{X_0}
  -\sum_{j=X_0}^{M-1}\mathcal C_j^{\rm ret}
  +\sum_{j=X_0}^{M-1}\eta_j^{\rm birth}.}
 \tag{T.31}
\]
A subpower terminal influence therefore requires
\[
 \boxed{
 \sum_{j=X_0}^{M-1}\mathcal C_j^{\rm ret}
 =\Lambda_{X_0}
  +\sum_{j=X_0}^{M-1}\eta_j^{\rm birth}
  -M^{o(1)}.}
 \tag{T.32}
\]
\end{corollary}

\begin{proof}
Apply \textup{(T.23)} at every adjacent cutoff and telescope.
\end{proof}

\begin{proposition}[Orthogonal-birth barrier]
\label{cth:birth-orthogonal-birth}
Let
\(G_N=\diag(s_1,\ldots,s_N)\), \(s_j>0\), let
\(e_N=(1,\ldots,1)\), and let \(|x_j|=1\).  Then
\[
 \boxed{
 \Lambda_N=\sum_js_j^{-1},
 \qquad
 \mathfrak A_N=\sum_js_j,}
 \tag{T.33}
\]
and
\[
 \boxed{\Lambda_N\mathfrak A_N\ge N^2.}
 \tag{T.34}
\]
Thus a near-linear protected action forces polynomial endpoint influence on a
pure orthogonal-birth chronology, independently of the protected phases.
\end{proposition}

\begin{proof}
The formulas are diagonal, and Cauchy--Schwarz gives the product bound.
\end{proof}

\begin{proposition}[Subpower uniform-retuning no-go]
\label{prop:birth-uniform-retuning}
If
\[
 \boxed{A\preceq\kappa G,}
 \tag{T.35}
\]
then
\[
 \boxed{
 \langle e,A^{-1}e\rangle
 \ge\kappa^{-1}\langle e,G^{-1}e\rangle.}
 \tag{T.36}
\]
A retuning with \(\kappa=X^{o(1)}\) cannot turn a polynomial endpoint
influence into a subpower one.
\end{proposition}

\begin{proof}
Use operator monotonicity of inversion.
\end{proof}

\begin{proposition}[Scalar retuning preserves the Cauchy product]
\label{prop:birth-scalar-retuning}
If \(A=\lambda G\), then for every protected old vector \(x\),
\[
 \boxed{
 \langle e,A^{-1}e\rangle\langle x,Ax\rangle
 =\langle e,G^{-1}e\rangle\langle x,Gx\rangle.}
 \tag{T.37}
\]
Useful retuning must therefore be selective: decoder heavy and protected-light.
\end{proposition}

\begin{proof}
The two factors scale by \(\lambda^{-1}\) and \(\lambda\).
\end{proof}

% -----------------------------------------------------------------------------
\subsection{Conditional handoff with complete birth accounting}
\label{subsec:birth-handoff}

Let \(x_X\) be the protected squarefree M\"obius history,
\(\mathfrak A_X=\langle x_X,G_Xx_X\rangle\), and
\(\Lambda_X=\langle e_X,G_X^{-1}e_X\rangle\).

\begin{theorem}[Birth-corrected endpoint criterion]
\label{thm:birth-GRH}
If
\[
 \boxed{\mathfrak A_X\le X^{1+o(1)}}
 \tag{T.42}
\]
and
\[
 \boxed{
 \Lambda_{X_0}
 -\sum_{j=X_0}^{X-1}\mathcal C_j^{\rm ret}
 +\sum_{j=X_0}^{X-1}\eta_j^{\rm birth}
 =X^{o(1)},}
 \tag{T.43}
\]
then
\[
 \boxed{|M_\chi(X)|^2\le X^{1+o(1)}.}
 \tag{T.44}
\]
For every fixed primitive character, this implies GRH.
\end{theorem}

\begin{proof}
By \Cref{cor:birth-telescope}, the expression in
\textup{(T.43)} is \(\Lambda_X\).  Apply Cauchy--Schwarz in \(G_X\):
\[
 |\langle e_X,x_X\rangle|^2\le\Lambda_X\mathfrak A_X.
\]
The square-root Mertens criterion and the functional equation give GRH.
\end{proof}

The estimates \textup{(T.42)}--\textup{(T.43)} are not proved here.

% -----------------------------------------------------------------------------


% =============================================================================
\section{Target-free M\"obius--log compilation in Dirichlet channels}
\label{sec:grh-outer}
% =============================================================================

\subsection{The calibrated survivor compiler}
\label{subsec:grhouter-survivor-compiler}

Let
\[
 \Omega=\{\ell_1,\ldots,\ell_J,\infty\}
\]
be an ordered first-hit alphabet and let \(\nu_0,\nu_1\) be probability laws
on \(\Omega\) with \(\nu_1\ll\nu_0\).  Put
\[
 a_{\star,j}=\nu_\star(\ell_j),
 \qquad
 S_{\star,j}=\nu_\star\{\ell_{j+1},\ldots,\ell_J,\infty\},
 \qquad \star\in\{0,1\},
\]
with \(S_{\star,0}=1\), and assume \(S_{0,j}>0\).  Define
\[
 Q_j=\frac{S_{1,j}}{S_{0,j}}
\]
and the crossing determinant
\[
 \boxed{
 \mathfrak d_j=S_{1,j}a_{0,j}-S_{0,j}a_{1,j}.}
 \tag{G.2}
\]

\begin{theorem}[Exact survivor-current telescope]
\label{thm:grhouter-survivor-current}
For every \(1\le j\le J\),
\[
 \boxed{
 Q_j-Q_{j-1}
 =\frac{\mathfrak d_j}{S_{0,j-1}S_{0,j}}.}
 \tag{G.3}
\]
Consequently,
\[
 \boxed{
 \frac{\nu_1(\infty)}{\nu_0(\infty)}-1
 =\sum_{j=1}^{J}
   \frac{\mathfrak d_j}{S_{0,j-1}S_{0,j}}.}
 \tag{G.4}
\]
The determinant also has the hazard form
\[
 \boxed{
 \mathfrak d_j
 =S_{1,j-1}S_{0,j-1}
  \left(
   \frac{a_{0,j}}{S_{0,j-1}}
   -\frac{a_{1,j}}{S_{1,j-1}}
  \right).}
 \tag{G.5}
\]
\end{theorem}

\begin{proof}
Use \(S_{\star,j-1}=a_{\star,j}+S_{\star,j}\).  Direct subtraction gives
\[
 Q_j-Q_{j-1}
 =\frac{S_{1,j}S_{0,j-1}-S_{1,j-1}S_{0,j}}
       {S_{0,j-1}S_{0,j}},
\]
and the numerator is \(\mathfrak d_j\).  Summation proves the telescope.
Factoring the previous survivor masses gives the hazard formula.
\end{proof}

Let \(\lambda=\dd\nu_1/\dd\nu_0\), and let \(\mathscr F_j\) reveal the
first \(j\) singleton hit cells and the remaining survivor cell.

\begin{theorem}[Doob compiler and positive fallback action]
\label{thm:grhouter-survivor-Pythagoras}
The conditional expectations
\[
 M_j=\mathbb E_{\nu_0}[\lambda\mid\mathscr F_j]
\]
form a martingale.  On the survivor cell their value is \(Q_j\), and
\[
 \boxed{
 \sum_{\xi\in\Omega}\nu_0(\xi)|\lambda(\xi)-1|^2
 =\sum_{j=1}^{J}
   \frac{|\mathfrak d_j|^2}
        {a_{0,j}S_{0,j}S_{0,j-1}},}
 \tag{G.6}
\]
with zero hit cells omitted.
\end{theorem}

\begin{proof}
The formula for the conditional expectation is immediate on the revealed
singletons and on the survivor cell.  The martingale-difference spaces are
orthogonal.  Conditional on the previous survivor, the split probabilities
are \(a_{0,j}/S_{0,j-1}\) and \(S_{0,j}/S_{0,j-1}\), while the two likelihood
values differ by \(\mathfrak d_j/(a_{0,j}S_{0,j})\).  Multiplication by the
previous survivor mass yields the displayed summand.
\end{proof}

\begin{remark}[Signed current versus positive Christoffel envelope]
The protected chronology is \textup{(G.4)}, not the positive sum
\textup{(G.6)}.  The latter can be large in orthogonal threshold directions
whose signed crossing determinants cancel in the open current.  It is a valid
fallback action, but it may not be substituted for the signed telescope.
\end{remark}

Let \(\kappa_j\ne0\) be a source-fixed local transport comparator and set
\(K_j=\prod_{i=1}^j\kappa_i\), \(K_0=1\).

\begin{definition}[Euler-centred survivor current]
\label{def:grhouter-Euler-centred}
Define
\[
 \boxed{
 \mathfrak e_j=Q_j-\kappa_jQ_{j-1}.}
 \tag{G.7}
\]
\end{definition}

\begin{proposition}[Exact calibrated telescope]
\label{prop:grhouter-calibrated-telescope}
One has
\[
 \boxed{
 \frac{Q_J}{K_J}-Q_0
 =\sum_{j=1}^{J}\frac{\mathfrak e_j}{K_j}.}
 \tag{G.8}
\]
Equivalently,
\[
 Q_j-Q_{j-1}
 =\mathfrak e_j+(\kappa_j-1)Q_{j-1}.
 \tag{G.9}
\]
Thus the raw increment is the sum of the completed arithmetic current and the
predictable local transport.
\end{proposition}

\begin{proof}
Divide \(Q_j=\kappa_jQ_{j-1}+\mathfrak e_j\) by \(K_j\) and telescope.
\end{proof}

\begin{remark}[Population rule for cutoff retuning]
Whenever the old-sector cutoff retuning is populated by an ordered positive
first-hit panel, its local arithmetic content is the calibrated determinant
\(\mathfrak e_j\).  A raw old-block difference may contain both earlier
transported likelihood and deterministic Euler drift.  Such a difference
must not be inserted directly into \(\mathcal C_X^{\rm ret}\) before the
survivor and local-transport shorts have been performed.  If the two positive
laws or their mixed hit--survive rows are not operationally present, the
missing survivor population is itself the returned source obstruction.
\end{remark}

% -----------------------------------------------------------------------------
\subsection{The outer M\"obius--log switch for a Dirichlet channel}
\label{subsec:grhouter-outer-switch}

Fix a primitive Dirichlet character \(\chi\) of conductor \(q_\chi\).  Put
\[
 \mu_\chi(n)=\mu(n)\chi(n),
 \qquad
 M_\chi(Y)=\sum_{n\le Y}\mu_\chi(n),
\]
\[
 L_\chi(Y)=\sum_{m\le Y}\chi(m)\log m,
\]
and
\[
 \psi_\chi(X)=\sum_{n\le X}\chi(n)\Lambda(n).
\]
Let \(\delta_\chi=1\) for the principal channel (including a fixed finite set of removed Euler factors) and
\(\delta_\chi=0\) otherwise, and define
\[
 \Psi_\chi(X)=\psi_\chi(X)-\delta_\chi X.
\]
All real arguments of \(M_\chi,L_\chi\) are interpreted with the integer
floor.

\begin{theorem}[Exact target-free outer compiler]
\label{thm:grhouter-outer-compiler}
For every integer \(X\ge2\),
\[
 \boxed{
 \psi_\chi(X)
 =\sum_{d\le X}\mu_\chi(d)L_\chi(X/d)
 =\sum_{q=2}^{X}\chi(q)(\log q)M_\chi(X/q).}
 \tag{G.10}
\]
In particular, the direct section \(M_\chi(X)\) is absent because its
coefficient is \(\log1=0\).
\end{theorem}

\begin{proof}
Complete multiplicativity of \(\chi\) on its supported integers and the
finite identity \(\Lambda=\mu*\log\) give
\[
 \chi(n)\Lambda(n)
 =\sum_{dq=n}\mu(d)\chi(d)\chi(q)\log q.
\]
Summing over \(n\le X\) and interchanging the finite sums gives the first
formula.  Interchanging in the other order gives the second.  The term
\(q=1\) is identically zero.
\end{proof}

Choose an integer \(1\le D<X\) and put
\[
 \boxed{
 Q=Q_X(D):=\left\lfloor\frac{X}{D+1}\right\rfloor.}
 \tag{G.11}
\]
Define the divisor head and short-quotient tail by
\[
 \mathfrak H_{\chi,X}(D)
 :=\sum_{d\le D}\mu_\chi(d)L_\chi(X/d),
\]
\[
 \boxed{
 \mathfrak T_{\chi,X}(D)
 :=\sum_{q=2}^{Q}\chi(q)(\log q)
   \bigl[M_\chi(X/q)-M_\chi(D)\bigr].}
 \tag{G.12}
\]

\begin{theorem}[Exact head--tail split]
\label{thm:grhouter-head-tail}
For every \(D<X\),
\[
 \boxed{
 \Psi_\chi(X)
 =\bigl[\mathfrak H_{\chi,X}(D)-\delta_\chi X\bigr]
  +\mathfrak T_{\chi,X}(D).}
 \tag{G.13}
\]
Every Mertens cutoff occurring in the tail is at most \(X/2\).
\end{theorem}

\begin{proof}
Split the first sum in \textup{(G.10)} at \(D\).  In the tail, interchange
\(d\) and \(q\).  For fixed \(q\), the range is
\(D<d\le\lfloor X/q\rfloor\), which is nonempty exactly for \(q\le Q\).
The inner sum is
\(M_\chi(X/q)-M_\chi(D)\).  Since \(q\ge2\), its upper cutoff is at most
\(X/2\).
\end{proof}

\begin{theorem}[Vanishing top-carrier birth section]
\label{thm:grhouter-no-top-birth}
Let the Mertens coefficient carrier at cutoff \(X\) contain the coordinates
\(M_\chi(0),\ldots,M_\chi(X)\), and decompose it as the old subcarrier
through \(X-1\) plus the top coordinate \(M_\chi(X)\).  The tail readout
\(\mathfrak T_{\chi,X}(D)\) factors through the old subcarrier.  Hence its
top-detail prediction residual is zero and
\[
 \boxed{
 q_{X}^{\rm top}=0,
 \qquad
 \eta_{X}^{\rm top,birth}=0.}
 \tag{G.14}
\]
This statement concerns the endpoint compiler; newly visible source rows on
smaller quotient coordinates may still carry their own mixed old--detail
packets.
\end{theorem}

\begin{proof}
By \textup{(G.12)}, the readout uses only
\(M_\chi(\lfloor X/q\rfloor)\) with \(q\ge2\), and the fixed coordinate
\(M_\chi(D)\).  It is therefore independent of \(M_\chi(X)\).  In the carrier-birth
old--detail formula, the top detail readout equals its old-source prediction,
so the residual and its Moore--Penrose Christoffel cost vanish.
\end{proof}

\begin{remark}[What is and is not bypassed]
The outer compiler removes the direct top-coordinate carrier birth from the
Chebyshev endpoint.  It does not remove the need to estimate the old-sector
short-quotient current, nor does it make quotient-level source births
harmless.  The carrier-birth ledger is bypassed only for the circular target
section which would otherwise contain \(M_\chi(X)\) itself.
\end{remark}

% -----------------------------------------------------------------------------
\subsection{The nonprincipal square-root head closes}
\label{subsec:grhouter-nonprincipal-head}

\begin{lemma}[Periodic logarithmic sum]
\label{lem:grhouter-periodic-log}
If \(\chi\) is nonprincipal and fixed, then
\[
 \boxed{|L_\chi(Y)|\ll_\chi\log(2Y).}
 \tag{G.15}
\]
\end{lemma}

\begin{proof}
The partial sums \(S_\chi(t)=\sum_{n\le t}\chi(n)\) are bounded by a
constant depending only on the fixed conductor.  Abel summation gives
\[
 L_\chi(Y)=S_\chi(Y)\log Y-
 \int_1^Y\frac{S_\chi(t)}{t}\,\dd t,
\]
which proves the estimate.
\end{proof}

Fix \(B>2\) and put
\[
 \boxed{
 D_X=\left\lfloor X^{1/2}(\log X)^{-B}\right\rfloor,
 \qquad
 Q_X=Q_X(D_X)=X^{1/2+o(1)}.}
 \tag{G.16}
\]

\begin{theorem}[Nonprincipal head closure]
\label{thm:grhouter-nonprincipal-head}
For every fixed nonprincipal primitive \(\chi\),
\[
 \boxed{
 \mathfrak H_{\chi,X}(D_X)
 \ll_\chi X^{1/2}(\log X)^{1-B}.}
 \tag{G.17}
\]
Consequently,
\[
 \boxed{
 \Psi_\chi(X)
 =\mathfrak T_{\chi,X}(D_X)
  +O_\chi\bigl(X^{1/2}(\log X)^{1-B}\bigr).}
 \tag{G.18}
\]
\end{theorem}

\begin{proof}
Use \(|\mu_\chi(d)|\le1\) and
\Cref{lem:grhouter-periodic-log}:
\[
 |\mathfrak H_{\chi,X}(D_X)|
 \ll_\chi D_X\log X.
\]
Insert \textup{(G.16)} and apply \Cref{thm:grhouter-head-tail} with
\(\delta_\chi=0\).
\end{proof}

This is the first unconditional gain imported from the affine outer switch:
for every nonprincipal fixed conductor, the only GRH-strength quantity left
by the Chebyshev compiler is the target-free quotient tail.

% -----------------------------------------------------------------------------
\subsection{The principal head is a separate explicit obstruction}
\label{subsec:grhouter-principal-head}

Let \(\chi_0\) denote the principal channel with a fixed finite set of removed Euler factors, encoded by a modulus \(q\), and put
\(c_q=\varphi(q)/q\).  The primitive zeta channel is the case \(q=1\).  Define
\[
 A_q(D)=\sum_{\substack{d\le D\\(d,q)=1}}\frac{\mu(d)}d,
 \qquad
 B_q(D)=\sum_{\substack{d\le D\\(d,q)=1}}
         \frac{\mu(d)\log d}{d}.
\]

\begin{proposition}[Exact principal head profile up to the floor error]
\label{prop:grhouter-principal-head}
For fixed \(q\),
\[
 \boxed{
 \begin{aligned}
 \mathfrak H_{\chi_0,X}(D)
 ={}&c_qX\bigl[(\log X-1)A_q(D)-B_q(D)\bigr]
 \\
 &+O_q(D\log(2X)).
 \end{aligned}}
 \tag{G.19}
\]
Moreover, in the convergent zero-free-region sense,
\[
 A_q(D)\longrightarrow0,
 \qquad
 B_q(D)\longrightarrow-\frac1{c_q}.
\]
\end{proposition}

\begin{proof}
Periodic counting and Abel summation give
\[
 L_{\chi_0}(Y)
 =c_q(Y\log Y-Y)+O_q(\log(2Y)).
\]
Replace \(Y=\lfloor X/d\rfloor\) by \(X/d\); the derivative of
\(y\log y-y\) is \(\log y\), so the accumulated floor error is
\(O_q(D\log(2X))\).  Expanding the main term gives the formula.  The limits
follow from the Laurent expansion
\[
 \frac1{L(s,\chi_0)}=\frac{s-1}{c_q}+O_q((s-1)^2)
\]
and standard partial summation in the classical zero-free region.
\end{proof}

\begin{remark}[The principal head is not claimed to be square-root]
Formula \textup{(G.19)} identifies the exact principal head residual after
subtracting \(X\), but the classical zero-free region gives only a
subexponential saving from the scale \(X\).  It does not imply an
\(X^{1/2+o(1)}\) bound.  The principal channel must therefore retain the
explicit head residual as a separate obstruction; the reduction does not hide it
inside the short-quotient packet.
\end{remark}

% -----------------------------------------------------------------------------
\subsection{Complete quotient Gram and protected influence}
\label{subsec:grhouter-quotient-Gram}

For the cutoff \(D_X\), write
\[
 r_{\chi,X}(q)
 =\chi(q)\bigl[M_\chi(X/q)-M_\chi(D_X)\bigr],
 \qquad 2\le q\le Q_X,
\]
\[
 h_X(q)=\log q.
\]
Then
\[
 \boxed{
 \mathfrak T_{\chi,X}(D_X)
 =\langle h_X,r_{\chi,X}\rangle_{\ell^2}.}
 \tag{G.20}
\]

\begin{theorem}[Positive polarization of the target-free packet]
\label{thm:grhouter-positive-polarization}
One has
\[
 \boxed{
 \begin{pmatrix}
  \|h_X\|^2&\mathfrak T_{\chi,X}(D_X)\\
  \overline{\mathfrak T_{\chi,X}(D_X)}&\|r_{\chi,X}\|^2
 \end{pmatrix}
 \succeq0.}
 \tag{G.21}
\]
The mixed coefficient is reconstructed from the four positive polarization
energies of \(h_X\pm r_{\chi,X}\) and
\(h_X\pm i r_{\chi,X}\), but its sign and phase are not determined by the
two marginal norms.
\end{theorem}

\begin{proof}
The matrix is the Gram of the two vectors \(h_X,r_{\chi,X}\).
\end{proof}

For \(d>D_X\), define the quotient-row vector
\[
 u_{d,X}(q)
 =\one_{\{2\le q\le Q_X,\ dq\le X\}}\chi(q).
\]
Then
\[
 r_{\chi,X}
 =\sum_{d>D_X}\mu_\chi(d)u_{d,X}.
\]

\begin{theorem}[Complete long-divisor mixed Gram]
\label{thm:grhouter-complete-Gram}
The quotient energy is exactly
\[
 \boxed{
 \|r_{\chi,X}\|^2
 =\sum_{d,d'>D_X}
   \mu_\chi(d)\overline{\mu_\chi(d')}
   \mathbb G_{\chi,X}^{\rm quot}(d,d'),}
 \tag{G.22}
\]
where
\[
 \boxed{
 \mathbb G_{\chi,X}^{\rm quot}(d,d')
 =\langle u_{d',X},u_{d,X}\rangle
 =\sum_{\substack{2\le q\le Q_X\\dq\le X,\ d'q\le X}}
   |\chi(q)|^2}
 \tag{G.23}
\]
is a positive Gram kernel.  All \(d\)--\(d'\) cross blocks must be retained
before Cauchy--Schwarz, a large sieve, or a dispersion estimate is applied.
\end{theorem}

\begin{proof}
Expand the squared norm of the signed row sum and interchange the finite sums.
\end{proof}

\begin{proposition}[Diagonal quotient energies are not source minimal]
\label{cth:grhouter-diagonal-no-go}
The protected coefficient \(\langle h_X,r_{\chi,X}\rangle\) is not determined
by the per-divisor energies or by the per-quotient magnitudes.  Two nonzero
long-divisor rows can cancel exactly in the protected sum while both diagonal
energies remain positive.
\end{proposition}

\begin{proof}
Take two row vectors \(u\) and \(-u\) with equal protected coefficients.
Their signed sum vanishes, whereas the sum of their squared norms is
\(2\|u\|^2\).
\end{proof}

Since
\[
 \|h_X\|^2
 =Q_X(\log Q_X)^2+O(Q_X\log Q_X),
\]
Cauchy--Schwarz gives the following scale.

\begin{corollary}[Sufficient quotient-dispersion scale]
\label{cor:grhouter-L2-scale}
The estimate
\[
 \boxed{
 \|r_{\chi,X}\|^2
 \le
 \frac{X^{1+o(1)}}{Q_X(\log Q_X)^2}}
 \tag{G.24}
\]
implies
\[
 \boxed{
 |\mathfrak T_{\chi,X}(D_X)|^2\le X^{1+o(1)}.}
 \tag{G.25}
\]
\end{corollary}

Let \(L_X\succeq0\) be an independently populated positive action on the
quotient carrier, fixed before \(r_{\chi,X}\) is read.

\begin{theorem}[Sharp action--influence criterion]
\label{thm:grhouter-action-influence}
If
\[
 h_X\in\Ran L_X
\]
and
\[
 \boxed{
 \langle r_{\chi,X},L_Xr_{\chi,X}\rangle
 \langle h_X,L_X^\dagger h_X\rangle
 \le X^{1+o(1)},}
 \tag{G.26}
\]
then \textup{(G.25)} holds.
\end{theorem}

\begin{proof}
Apply Cauchy--Schwarz in the seminorm generated by \(L_X\):
\[
 |\langle h_X,r_{\chi,X}\rangle|^2
 \le
 \langle r_{\chi,X},L_Xr_{\chi,X}\rangle
 \langle h_X,L_X^\dagger h_X\rangle.
\]
\end{proof}

\begin{remark}[No terminally fitted action]
Choosing \(L_X\) after observing \(r_{\chi,X}\), or choosing it so that the
desired coefficient is a null or an eigenvector by construction, only rewrites
the target.  The action must come from the complete same-history quotient,
completed, Euler, gamma, contour, cutoff, nuisance, and provenance panel.
\end{remark}

% -----------------------------------------------------------------------------
\subsection{Nonprincipal GRH handoff and the birth--retuning balance}
\label{subsec:grhouter-GRH-handoff}

\begin{lemma}[Chebyshev square-root criterion]
\label{lem:grhouter-Chebyshev-criterion}
For a fixed primitive character \(\chi\), if for every \(\varepsilon>0\)
\[
 \Psi_\chi(X)=O_{\chi,\varepsilon}(X^{1/2+\varepsilon}),
\]
then \(L(s,\chi)\) has no zero with \(\Re s>1/2\).  The functional equation
then gives GRH for the channel.
\end{lemma}

\begin{proof}
Partial summation gives, after subtracting the principal pole when present,
\[
 -\frac{L'}{L}(s,\chi)-\frac{\delta_\chi}{s-1}
 =s\int_1^\infty\Psi_\chi(x)x^{-s-1}\,\dd x
  +\delta_\chi.
\]
The identity is first valid for \(\Re s>1\) and its right-hand side gives
an analytic continuation wherever the integral converges.  The assumed bound makes
the integral analytic in every half-plane \(\Re s>1/2+\varepsilon\).  Hence
there are no zeros there.  Let \(\varepsilon\downarrow0\) and use the
functional equation.
\end{proof}

\begin{theorem}[Target-free conditional GRH criterion for nonprincipal channels]
\label{thm:grhouter-nonprincipal-GRH}
Fix a nonprincipal primitive character.  If either
\textup{(G.24)} or \textup{(G.26)} holds along every sufficiently large
cutoff, then
\[
 \boxed{
 |\Psi_\chi(X)|^2\le X^{1+o(1)},}
 \tag{G.27}
\]
and GRH holds for \(L(s,\chi)\).
\end{theorem}

\begin{proof}
The quotient estimate gives
\(|\mathfrak T_{\chi,X}(D_X)|^2\le X^{1+o(1)}\).  The nonprincipal head is
\(X^{1/2+o(1)}\) by \Cref{thm:grhouter-nonprincipal-head}.  Apply
\Cref{thm:grhouter-head-tail,lem:grhouter-Chebyshev-criterion}.
\end{proof}

\begin{theorem}[Head-or-tail obstruction]
\label{thm:grhouter-head-tail-obstruction}
If the Chebyshev square-root estimate fails in a primitive channel, then along
some cofinal sequence at least one of
\[
 \boxed{
 |\mathfrak H_{\chi,X}(D_X)-\delta_\chi X|
 \ge\frac12|\Psi_\chi(X)|}
 \tag{G.28}
\]
or
\[
 \boxed{
 |\mathfrak T_{\chi,X}(D_X)|
 \ge\frac12|\Psi_\chi(X)|}
 \tag{G.29}
\]
holds.  For a fixed nonprincipal character the first branch is only
\(X^{1/2+o(1)}\), so every polynomial GRH failure is returned by the
short-quotient packet.  For the principal channel both branches remain.
\end{theorem}

\begin{proof}
Use the exact split \textup{(G.13)} and the triangle inequality.  The
nonprincipal conclusion is \Cref{thm:grhouter-nonprincipal-head}.
\end{proof}

\begin{remark}[Carrier-birth population rule]
For the Chebyshev compiler, the top-coordinate Christoffel payment vanishes by
\Cref{thm:grhouter-no-top-birth}.  The first arithmetic object is not an arbitrary
new-coordinate writer but the old-sector quotient current
\(\mathfrak T_{\chi,X}(D_X)\).  If this current is decomposed by thresholds,
its raw pieces must first be converted to calibrated survivor determinants by
\Cref{thm:grhouter-survivor-current,prop:grhouter-calibrated-telescope}.  The positive
first-hit or quotient energy is a fallback action only after the signed
coefficient has been assembled.
\end{remark}

The carrier-birth formula
\[
 \mathcal C_X^{\rm net}
 =\mathcal C_X^{\rm ret}-\eta_X^{\rm birth}
\]
remains exact for the complete source.  The outer compiler identifies a quotient geometry on which
its circular top-coordinate contribution satisfies
\[
 \boxed{
 \eta_X^{\rm top,birth}=0,}
 \tag{G.30}
\]
so the first potentially favorable arithmetic row is pure old-sector
retuning.  Quotient-level and nonprime source births must still be populated
and shorted normally.

% -----------------------------------------------------------------------------


% =============================================================================
\section{Hyperbola--Volterra decomposition and terminal quotient removal}
\label{sec:grh-hyperbola}
% =============================================================================
Fix a nonprincipal primitive character $\chi$ of conductor $f$.  Throughout
this section
\[
 D_X=\left\lfloor X^{1/2}(\log X)^{-B}\right\rfloor,\qquad
 Q_X=\left\lfloor\frac{X}{D_X+1}\right\rfloor=X^{1/2+o(1)},\qquad B>2.
 \tag{H.2}
\]
and
\[
 \mathfrak T_{\chi,X}(D_X)
 =\sum_{2\le q\le Q_X}\chi(q)\log q
   \bigl[M_\chi(\lfloor X/q\rfloor)-M_\chi(D_X)\bigr].
 \tag{H.1}
\]

\subsection{Exact hyperbola-cell factorization}
\label{subsec:hyperbola-hyperbola-factorization}

Fix a nonprincipal primitive character \(\chi\) of conductor \(f\).  Put
\[
 a_\chi(d)=\mu(d)\chi(d),
 \qquad
 S_{\chi,X}(n)
 =\sum_{D_X<d\le n}a_\chi(d)
 =M_\chi(n)-M_\chi(D_X).
 \tag{H.3}
\]
For \(D_X<n\le X/2\), define the hyperbola cell
\[
 \mathcal I_X(n)
 =\left\{q\in\mathbb N:
   2\le q\le Q_X,
   \left\lfloor\frac Xq\right\rfloor=n
  \right\},
 \tag{H.4}
\]
and its character-visible multiplicity
\[
 w_{\chi,X}(n)
 =\sum_{q\in\mathcal I_X(n)}|\chi(q)|^2.
 \tag{H.5}
\]
The integer-part convention in \textup{(H.1)} is used throughout this
version.

\begin{theorem}[Exact hyperbola replication identity]
\label{thm:hyperbola-hyperbola-replication}
The target-free quotient vector satisfies
\[
 \boxed{
 r_{\chi,X}(q)
 =\chi(q)S_{\chi,X}\!\left(\left\lfloor\frac Xq\right\rfloor\right),
 \qquad 2\le q\le Q_X.}
 \tag{H.6}
\]
Consequently,
\[
 \boxed{
 \|r_{\chi,X}\|_{\ell^2}^2
 =\sum_{D_X<n\le X/2}
   w_{\chi,X}(n)|S_{\chi,X}(n)|^2.}
 \tag{H.7}
\]
If
\[
 H_{\chi,X}(n)
 =\sum_{q\in\mathcal I_X(n)}\chi(q)\log q,
 \tag{H.8}
\]
then the protected coefficient has both exact forms
\[
 \boxed{
 \begin{aligned}
 \mathfrak T_{\chi,X}(D_X)
 &=\sum_{D_X<n\le X/2}
   H_{\chi,X}(n)S_{\chi,X}(n)\\
 &=\sum_{d>D_X}\mu(d)\chi(d)
   L_\chi\!\left(
    \min\left\{Q_X,\left\lfloor\frac Xd\right\rfloor\right\}
   \right),
 \end{aligned}}
 \tag{H.9}
\]
where \(L_\chi(y)=\sum_{2\le q\le y}\chi(q)\log q\).  No \(q=1\)
coordinate occurs in either representation.
\end{theorem}

\begin{proof}
Equation \textup{(H.6)} is the definition of the quotient vector after
interchanging the finite divisor and quotient sums.  Grouping equal values of
\(\lfloor X/q\rfloor\) gives \textup{(H.7)} and the first line of
\textup{(H.9)}.  Expanding
\(S_{\chi,X}(n)=\sum_{D_X<d\le n}\mu(d)\chi(d)\), interchanging the finite
sums, and observing that
\[
 d\le\left\lfloor\frac Xq\right\rfloor
 \iff q\le\left\lfloor\frac Xd\right\rfloor
\]
gives the second line.  The lower quotient bound is two, so the top Mertens
section remains absent.
\end{proof}

Let \(\mathscr V_X\) be the finite Volterra map
\[
 (\mathscr V_Xa)(n)=\sum_{D_X<d\le n}a(d)
 \tag{H.10}
\]
and let \(\mathscr H_{\chi,X}\) replicate one hyperbola level to its quotient
cell with multiplier \(\chi(q)\).  Then
\[
 \boxed{
 r_{\chi,X}
 =\mathscr H_{\chi,X}\mathscr V_Xa_\chi.}
 \tag{H.11}
\]
The raw quotient norm is therefore the norm of a cumulative Volterra profile,
not the norm of its arithmetic increments.

% -----------------------------------------------------------------------------
\subsection{The unavoidable hyperbola--Volterra energy floor}
\label{subsec:hyperbola-energy-floor}

\begin{lemma}[Visible size of the terminal hyperbola cells]
\label{lem:hyperbola-cell-size}
There is a constant \(c_f>0\), depending only on the conductor, such that for
all sufficiently large \(X\) and every integer
\[
 D_X<n\le2D_X
\]
one has
\[
 \boxed{
 w_{\chi,X}(n)
 \ge c_f\frac{X}{D_X^2}.}
 \tag{H.12}
\]
\end{lemma}

\begin{proof}
For this range the upper quotient endpoint is automatically at most \(Q_X\),
and
\[
 |\mathcal I_X(n)|
 =\left\lfloor\frac Xn\right\rfloor
  -\left\lfloor\frac X{n+1}\right\rfloor
 \ge\frac{X}{n(n+1)}-1
 \ge\frac{X}{2D_X(2D_X+1)}-1.
\]
Because \(X/D_X^2=(\log X)^{2B+o(1)}\to\infty\), this interval contains a
fixed positive proportion of integers coprime to \(f\).  More explicitly,
periodic counting gives
\[
 \#\{q\in I:(q,f)=1\}
 =\frac{\varphi(f)}f|I|+O_f(1).
\]
For a primitive character \(|\chi(q)|^2=\one_{(q,f)=1}\), which proves the
claim.
\end{proof}

\begin{lemma}[Squarefree mass in one dyadic interval]
\label{lem:hyperbola-squarefree-mass}
Put
\[
 \kappa_f
 =\frac1{\zeta(2)}
  \prod_{p\mid f}\frac1{1+1/p}>0.
 \tag{H.13}
\]
Then
\[
 \boxed{
 \sum_{D<n\le2D}|\mu(n)\chi(n)|^2
 =\kappa_fD+O_f(D^{1/2}).}
 \tag{H.14}
\]
\end{lemma}

\begin{proof}
Use
\(\mu(n)^2=\sum_{r^2\mid n}\mu(r)\) and count integers coprime to \(f\) in
arithmetic periods.  The main coefficient is
\[
 \frac{\varphi(f)}f
 \sum_{(r,f)=1}\frac{\mu(r)}{r^2}
 =\frac1{\zeta(2)}
  \prod_{p\mid f}\frac{1-1/p}{1-1/p^2}
 =\kappa_f.
\]
The truncation and periodic counting errors are \(O_f(D^{1/2})\).
\end{proof}

\begin{theorem}[Sharp kinematic lower floor for the raw quotient energy]
\label{thm:hyperbola-raw-floor}
For the outer-compiler cutoff \textup{(H.2)}, there is \(c_\chi>0\) such that
\[
 \boxed{
 \|r_{\chi,X}\|_{\ell^2}^2
 \ge c_\chi\frac{X}{D_X}
 \asymp_\chi Q_X.}
 \tag{H.15}
\]
In particular, no argument can prove
\(\|r_{\chi,X}\|^2=o(Q_X)\).
\end{theorem}

\begin{proof}
Write \(S(n)=S_{\chi,X}(n)\).  Since \(S(D_X)=0\) and
\[
 a_\chi(n)=S(n)-S(n-1),
\]
the elementary difference inequality gives
\[
 \sum_{D_X<n\le2D_X}|a_\chi(n)|^2
 \le4\sum_{D_X<n\le2D_X}|S(n)|^2.
 \tag{H.16}
\]
By \Cref{lem:hyperbola-squarefree-mass}, the left side is at least
\(\frac12\kappa_fD_X\) for large \(X\).  Insert
\Cref{lem:hyperbola-cell-size} into the exact energy identity
\textup{(H.7)} and restrict to this dyadic interval:
\[
 \|r_{\chi,X}\|^2
 \ge c_f\frac{X}{D_X^2}
       \sum_{D_X<n\le2D_X}|S(n)|^2
 \gg_\chi\frac{X}{D_X}.
\]
Finally \(Q_X\asymp X/D_X\).
\end{proof}

\begin{remark}[Diagonal quotient target]
The sufficient Cauchy--Schwarz criterion above is not a small-energy
statement.  At exponent resolution its best plausible scale is
\[
 \boxed{
 \|r_{\chi,X}\|^2=Q_XX^{o(1)},}
 \tag{H.17}
\]
which is only a subpower inflation above the unavoidable floor
\textup{(H.15)}.  The apparent logarithmic factors in \textup{(G.24)} are not a
source of power saving because \(Q_X/D_X\) is itself only polylogarithmic.
Any proof through the raw norm would have to show that the complete cumulative
Mertens profile stays within a subpower factor of its minimum kinematic scale.
\end{remark}

\begin{proposition}[A small terminal endpoint does not control the quotient profile]
\label{cth:hyperbola-endpoint-profile-no-go}
On the same hyperbola geometry there are coefficient histories
\(a(d)\), with \(|a(d)|\le1\), such that
\[
 \sum_{d\le X}a(d)=0
 \tag{H.18}
\]
but the associated quotient profile satisfies
\[
 \boxed{
 \sum_{2\le q\le Q_X}|\chi(q)|^2
 \left|\sum_{D_X<d\le X/q}a(d)\right|^2
 \gg_\chi X^2.}
 \tag{H.19}
\]
Thus neither a final square-root endpoint nor even an identically zero final
endpoint structurally implies the target-free quotient-dispersion estimate.
\end{proposition}

\begin{proof}
Choose one fixed integer \(q_0\ge2\) coprime to \(f\).  For large \(X\), put
\(N=\lfloor X/q_0\rfloor-D_X\), set \(a(d)=1\) on the first \(N\)
coordinates after \(D_X\), set \(a(d)=-1\) on the next \(N\) coordinates,
and set it to zero elsewhere.  There is enough room because
\(X-\lfloor X/q_0\rfloor\ge N\).  The final sum is zero, whereas the
\(q_0\)-coordinate of the quotient profile equals \(N\asymp X\).  Its single
squared contribution proves the claim.
\end{proof}

\begin{remark}[What the floor does and does not say]
The lower bound \textup{(H.15)} is not evidence against GRH.  It says that
raw quotient energy contains a deterministic cumulative-profile cost.  The
protected theorem is the signed mixed coefficient
\(\langle h_X,r_{\chi,X}\rangle\), and that coefficient may cancel while the
raw profile energy remains large.
\end{remark}

% -----------------------------------------------------------------------------
\subsection{Conductor-block calibration and the signed crossing current}
\label{subsec:hyperbola-conductor-blocks}

Let
\[
 \mathcal B_m
 =\{mf+1,\ldots,mf+f\}
 \tag{H.20}
\]
be a complete conductor block contained in \([2,Q_X]\).  On
\(\mathbb C^f\), let
\[
 P_0=\frac1f\one\one^*,
 \qquad
 Q_0=I-P_0.
 \tag{H.21}
\]
Define the block profile and writer
\[
 s_{X,m}(a)
 =S_{\chi,X}\!\left(
   \left\lfloor\frac{X}{mf+a}\right\rfloor
  \right),
 \tag{H.22}
\]
\[
 h_{\chi,m}(a)
 =\overline{\chi(a)}\log(mf+a),
 \qquad 1\le a\le f.
 \tag{H.23}
\]
With the standard complex inner product, the contribution of this block to
\textup{(H.1)} is
\[
 \mathfrak T_{\chi,X}(m)
 =\langle h_{\chi,m},s_{X,m}\rangle.
 \tag{H.24}
\]

\begin{theorem}[Exact conductor-block transport/current split]
\label{thm:hyperbola-block-split}
Every complete conductor block has the orthogonal decomposition
\[
 \boxed{
 \mathfrak T_{\chi,X}(m)
 =\mathfrak E_{\chi,X}^{\rm tr}(m)
  +\mathfrak C_{\chi,X}^{\rm conn}(m),}
 \tag{H.25}
\]
where
\[
 \boxed{
 \mathfrak E_{\chi,X}^{\rm tr}(m)
 =\frac1f
  \left(\sum_{a=1}^f\chi(a)\log(mf+a)\right)
  \left(\sum_{a=1}^fs_{X,m}(a)\right),}
 \tag{H.26}
\]
\[
 \boxed{
 \mathfrak C_{\chi,X}^{\rm conn}(m)
 =\langle Q_0h_{\chi,m},Q_0s_{X,m}\rangle.}
 \tag{H.27}
\]
The positive profile innovation is
\[
 \boxed{
 \mathcal I_{\chi,X}^{\rm conn}(m)
 =\|Q_0s_{X,m}\|^2
 =\frac1{2f}\sum_{a,b=1}^f
  |s_{X,m}(a)-s_{X,m}(b)|^2.}
 \tag{H.28}
\]
It obeys
\[
 \boxed{
 |\mathfrak C_{\chi,X}^{\rm conn}(m)|^2
 \le
 \|Q_0h_{\chi,m}\|^2
 \mathcal I_{\chi,X}^{\rm conn}(m).}
 \tag{H.29}
\]
\end{theorem}

\begin{proof}
Insert \(I=P_0+Q_0\) in both slots of
\(\langle h_{\chi,m},s_{X,m}\rangle\).  The mixed terms vanish by
orthogonality.  The constant--constant term is \textup{(H.26)}, and the
mean-zero term is \textup{(H.27)}.  The variance identity in
\textup{(H.28)} is standard, and \textup{(H.29)} is Cauchy--Schwarz on
the mean-zero child space.
\end{proof}

The first term in \textup{(H.25)} is the deterministic conductor transport
of the block mean.  The second is the signed crossing current.  The analogy
with the affine first-hit determinant is exact at the level of order of
operations:
\[
 \boxed{
 \begin{gathered}
 \text{block profile}
 \longrightarrow
 \text{constant transport short}
 \longrightarrow
 \text{signed mean-zero current}\\
 \longrightarrow
 \text{optional positive innovation}
 \end{gathered}}
 \tag{H.30}
\]

\begin{lemma}[Thin-hyperbola support of the connected profile]
\label{lem:hyperbola-thin-support}
For \(1\le a,b\le f\), put
\(n_a=\lfloor X/(mf+a)\rfloor\).  Then
\[
 \boxed{
 s_{X,m}(a)-s_{X,m}(b)
 =\begin{cases}
  -\displaystyle\sum_{n_a<d\le n_b}\mu(d)\chi(d),&n_a<n_b,\\[0.8em]
  \displaystyle\sum_{n_b<d\le n_a}\mu(d)\chi(d),&n_b<n_a,\\[0.4em]
  0,&n_a=n_b.
 \end{cases}}
 \tag{H.31}
\]
Moreover, for \(t=mf\ge f\),
\[
 \boxed{
 |s_{X,m}(a)-s_{X,m}(b)|
 \le C_f\left(\frac{X}{t^2}+1\right).}
 \tag{H.32}
\]
\end{lemma}

\begin{proof}
The first identity is subtraction of two Mertens prefixes.  Since
\[
 \left|\frac{X}{t+a}-\frac{X}{t+b}\right|
 \le\frac{X|a-b|}{t^2},
\]
the number of integer coefficients in the difference interval is at most
\(C_f(X/t^2+1)\).  Use \(|\mu(d)\chi(d)|\le1\).
\end{proof}

\begin{lemma}[Small deterministic conductor transport]
\label{lem:hyperbola-small-transport}
For every complete block with \(t=mf\ge f\),
\[
 \boxed{
 \left|\sum_{a=1}^f\chi(a)\log(t+a)\right|
 \le\frac{C_f}{t}.}
 \tag{H.33}
\]
\end{lemma}

\begin{proof}
Because \(\chi\) is nonprincipal,
\(\sum_{a=1}^f\chi(a)=0\).  Subtract \(\log t\) and use
\(|\log(t+a)-\log t|\le a/t\).
\end{proof}

\begin{proposition}[Positive block innovations do not determine the signed sum]
\label{cth:hyperbola-block-energy-no-sign}
Two complete blocks can have identical writer norms and identical positive
innovation energies \(\mathcal I^{\rm conn}\), while their connected currents
are equal and opposite.  Hence
\[
 \sum_m\mathcal I_{\chi,X}^{\rm conn}(m)
 \tag{H.34}
\]
does not determine
\(\sum_m\mathfrak C_{\chi,X}^{\rm conn}(m)\).
\end{proposition}

\begin{proof}
On any mean-zero block choose nonzero vectors \(w,z\).  The two profiles
\(z\) and \(-z\) have the same norm, while
\(\langle w,z\rangle\) and \(\langle w,-z\rangle\) have opposite signs.
Take two physical copies with the same writer Gram.  This finite example is
unchanged by adjoining constant transport coordinates.
\end{proof}

\begin{remark}[Source-minimal quotient row]
The first arithmetic quantity is not the raw norm of the cumulative profile.
It is the coherently assembled sum of
\[
 \boxed{
 \mathfrak E_{\chi,X}^{\rm tr}(m)
 +\mathfrak C_{\chi,X}^{\rm conn}(m)}
 \tag{H.35}
\]
over complete conductor blocks, together with the two bounded endpoint
partials.  The positive innovations are fallback magnitudes and must retain
their cross-block Gram if they are used.
\end{remark}

% -----------------------------------------------------------------------------
\subsection{Unconditional closure of the terminal quotient range}
\label{subsec:hyperbola-terminal-range}

For \(2\le R\le Q_X\), define
\[
 \mathfrak T_{\chi,X}^{\ge R}(D_X)
 =\sum_{R\le q\le Q_X}
  \chi(q)(\log q)
  S_{\chi,X}\!\left(\left\lfloor\frac Xq\right\rfloor\right).
 \tag{H.36}
\]

\begin{theorem}[Terminal conductor-block estimate]
\label{thm:hyperbola-terminal-estimate}
For every fixed nonprincipal primitive character and every
\(2f\le R\le Q_X\),
\[
 \boxed{
 |\mathfrak T_{\chi,X}^{\ge R}(D_X)|
 \le C_\chi\left[
  \frac{X\log(2Q_X)}{R}
  +Q_X\log(2Q_X)
 \right].}
 \tag{H.37}
\]
The estimate uses no cancellation of the M\"obius coefficients.
\end{theorem}

\begin{proof}
Discard at most two incomplete conductor blocks at the endpoints.  Their total
absolute contribution is
\[
 O_\chi\left(\frac{X\log(2Q_X)}R\right)
\]
because \(|S_{\chi,X}(\lfloor X/q\rfloor)|\le X/q\).

Consider a complete block with \(t=mf\).  By
\Cref{thm:hyperbola-block-split,lem:hyperbola-small-transport}, its transport term is
\[
 |\mathfrak E_{\chi,X}^{\rm tr}(m)|
 \le C_\chi\frac{X}{t^2}.
\]
Furthermore
\[
 \|Q_0h_{\chi,m}\|_1\le C_\chi\log(2t),
\]
and \Cref{lem:hyperbola-thin-support} gives
\[
 \|Q_0s_{X,m}\|_\infty
 \le C_\chi\left(\frac{X}{t^2}+1\right).
\]
Hence
\[
 |\mathfrak C_{\chi,X}^{\rm conn}(m)|
 \le C_\chi\log(2t)
      \left(\frac{X}{t^2}+1\right).
\]
Summing over \(t\) in an arithmetic progression of fixed step \(f\) yields
\[
 \sum_{t\ge R}\frac{X\log(2t)}{t^2}
 \ll_\chi\frac{X(1+\log R)}R,
 \qquad
 \sum_{t\le Q_X}\log(2t)
 \ll_\chi Q_X\log(2Q_X).
\]
Together with the endpoint partial blocks this proves the stated bound.
\end{proof}

\begin{corollary}[Square-root closure above a subpower-shrunken boundary]
\label{cor:hyperbola-terminal-square-root}
Let \(E_X\to\infty\) satisfy \(E_X=X^{o(1)}\), and put
\[
 R_X=\left\lfloor\frac{X^{1/2}}{E_X}\right\rfloor.
 \tag{H.38}
\]
Then
\[
 \boxed{
 \mathfrak T_{\chi,X}^{\ge R_X}(D_X)
 =X^{1/2+o(1)}.}
 \tag{H.39}
\]
In particular, one may remove an arbitrarily large subpower fraction of the
upper quotient bank without any M\"obius cancellation theorem.
\end{corollary}

\begin{proof}
The first term in \textup{(H.37)} is
\(X^{1/2}E_X\log(2Q_X)=X^{1/2+o(1)}\).  The second is
\(Q_X\log(2Q_X)=X^{1/2+o(1)}\) by \textup{(H.2)}.
\end{proof}

\begin{remark}[The square-root quotient barrier]
If \(R=X^\theta\), the deterministic conductor-block estimate contains
\(X/R=X^{1-\theta}\).  It reaches the square-root exponent precisely at
\(\theta=1/2\).  Moving the boundary to any fixed power below \(X^{1/2}\)
therefore requires genuine cancellation of the thin-hyperbola M\"obius
increments in \textup{(H.31)}.  Character periodicity and local transport
subtraction alone cannot supply that gain.
\end{remark}

% -----------------------------------------------------------------------------
\subsection{The short-quotient head and GRH handoff}
\label{subsec:hyperbola-head-handoff}

With \(R_X\) from \textup{(H.38)}, define
\[
 \boxed{
 \mathfrak T_{\chi,X}^{<R_X}(D_X)
 =\sum_{2\le q<R_X}
  \chi(q)(\log q)
  \bigl[M_\chi(\lfloor X/q\rfloor)-M_\chi(D_X)\bigr].}
 \tag{H.40}
\]

\begin{theorem}[Reduced nonprincipal Chebyshev packet]
\label{thm:hyperbola-reduced-packet}
For every fixed nonprincipal primitive character,
\[
 \boxed{
 \Psi_\chi(X)
 =\mathfrak T_{\chi,X}^{<R_X}(D_X)
  +X^{1/2+o(1)}.}
 \tag{H.41}
\]
The remaining quotient bank has cardinality
\[
 \boxed{
 R_X=X^{1/2}/E_X+O(1).}
 \tag{H.42}
\]
\end{theorem}

\begin{proof}
the preceding outer-compiler reduction closed the divisor head by
\(X^{1/2+o(1)}\).  Split its target-free quotient tail at \(R_X\) and apply
\Cref{cor:hyperbola-terminal-square-root} to the upper part.
\end{proof}

\begin{corollary}[Conditional GRH from the calibrated short-quotient head]
\label{cor:hyperbola-head-GRH}
If
\[
 \boxed{
 |\mathfrak T_{\chi,X}^{<R_X}(D_X)|
 \le X^{1/2+o(1)},}
 \tag{H.43}
\]
then
\[
 |\Psi_\chi(X)|\le X^{1/2+o(1)},
 \tag{H.44}
\]
and GRH holds for \(L(s,\chi)\).
\end{corollary}

\begin{proof}
Use \Cref{thm:hyperbola-reduced-packet} and the Chebyshev square-root criterion of
the preceding outer-compiler reduction.
\end{proof}

The head remains target free.  Its long-divisor form is
\[
 \boxed{
 \mathfrak T_{\chi,X}^{<R_X}(D_X)
 =\sum_{d>D_X}\mu(d)\chi(d)
   L_\chi\!\left(
    \min\left\{R_X-1,
      \left\lfloor\frac Xd\right\rfloor
    \right\}
   \right).}
 \tag{H.45}
\]
The coefficient \(q=1\) is still absent.

\begin{proposition}[Stability of the vanishing top-birth section]
\label{prop:hyperbola-top-birth-stability}
Any source assembly, nuisance short, conductor-block projection, or quotient
localization acting entirely on
\[
 \ell^2\{2,\ldots,Q_X\}
 \tag{H.46}
\]
factors through Mertens coordinates of height at most \(X/2\).  Its prediction residual on the top coordinate \(M_\chi(X)\) is therefore zero.
A nonzero top-coordinate Christoffel birth can reappear only if a later
factorization explicitly adjoins a \(q=1\) section or an equivalent terminal
read.
\end{proposition}

\begin{proof}
Every admitted quotient satisfies \(q\ge2\), hence
\(\lfloor X/q\rfloor\le X/2\).  Linear maps and orthogonal shorts on this
quotient carrier cannot enlarge the set of Mertens coordinates through which
the source factors.
\end{proof}

\begin{remark}[No return to the raw quotient norm]
The estimate \textup{(H.43)} is a signed current theorem.  It must not be
replaced by a sum of conductor-block innovation energies or by a raw profile
norm unless the complete cross-block Gram is retained and the resulting
stronger theorem is proved explicitly.  The deterministic floor
\textup{(H.15)} and the endpoint counterexample
\textup{(H.19)} show why that substitution is not source minimal.
\end{remark}

% -----------------------------------------------------------------------------


% =============================================================================

% =============================================================================
\section{Protected rectangular M\"obius descent}
\label{sec:grh-rectangular}
% =============================================================================

The outer Chebyshev compiler of \Cref{sec:grh-outer,sec:grh-hyperbola}
removes the direct terminal Mertens coordinate and leaves a signed
short-quotient current.  There is a complementary entrance which acts one
step farther upstream, on the reciprocal Dirichlet series itself.  It creates
two M\"obius occurrences before any prime child or positive quotient norm is
selected.

Fix a nonprincipal primitive character $\chi$ of conductor $f$, and put
\[
 M_\chi(x)=\sum_{n\le x}\mu(n)\chi(n),
 \qquad
 S_\chi(t)=\sum_{1\le k\le t}\chi(k).
\]
All real upper limits are interpreted by the integer floor.

\subsection{Exact two-threshold descent}

\begin{definition}[Rectangular descent form]
\label{def:rectangular-descent}
For positive integers $K,N_1,N_2$, define
\[
 \boxed{
 \mathcal Q_\chi(K;N_1,N_2)
 :=\sum_{m\le N_1}\sum_{n\le N_2}
   \mu(m)\mu(n)\chi(mn)
   S_\chi\!\left(\frac{K}{mn}\right).}
 \tag{R.1}
\]
\end{definition}

\begin{theorem}[Two-threshold M\"obius descent]
\label{thm:rectangular-two-threshold}
If
\[
 K<(N_1+1)(N_2+1),
 \tag{R.2}
\]
then
\[
 \boxed{
 M_\chi(K)
 =M_\chi(\min\{N_1,K\})
  +M_\chi(\min\{N_2,K\})
  -\mathcal Q_\chi(K;N_1,N_2).}
 \tag{R.3}
\]
\end{theorem}

\begin{proof}
Consider the finite triple sum
\[
 \mathcal F=\sum_{abk\le K}\mu(a)\mu(b)\chi(abk).
\]
Grouping by $r=abk$ gives
\[
 \mathcal F
 =\sum_{r\le K}\chi(r)(\mu*\mu*\one)(r)
 =\sum_{r\le K}\mu(r)\chi(r)
 =M_\chi(K),
\]
because $\mu*\one$ is the convolution identity.  Condition \textup{(R.2)}
excludes the simultaneous inequalities $a>N_1$ and $b>N_2$ on the support
$abk\le K$.  Hence the support is the union of $a\le N_1$ and $b\le N_2$.
For fixed $a$,
\[
 \sum_{bk\le K/a}\mu(b)\chi(bk)
 =\sum_{r\le K/a}\chi(r)\sum_{b\mid r}\mu(b)=1.
\]
The first one-sided restriction therefore contributes
$M_\chi(\min\{N_1,K\})$, and the second contributes the analogous term with
$N_2$.  Their intersection is \textup{(R.1)}.  Inclusion--exclusion proves
the identity.
\end{proof}

\begin{corollary}[Symmetric square-root descent]
\label{cor:rectangular-symmetric}
Let $N=\lfloor K^{1/2}\rfloor$.  Then
\[
 \boxed{M_\chi(K)=2M_\chi(N)-\mathcal Q_\chi(K;N,N).}
 \tag{R.4}
\]
Consequently a square-root bound for the rectangular form is equivalent, up
to the trivially bounded term $M_\chi(N)$, to the square-root Mertens bound at
$K$.
\end{corollary}

\begin{proof}
Since $K<(N+1)^2$, apply \Cref{thm:rectangular-two-threshold}.  The remaining
term satisfies $|M_\chi(N)|\le N$.
\end{proof}

The descent is pointwise as well as summatory.  If
\[
 C_{N_1,N_2}(r)
 =\sum_{\substack{abk=r\\a\le N_1,\ b\le N_2}}\mu(a)\mu(b),
\]
then the same inclusion--exclusion argument gives
\[
 \boxed{
 C_{N_1,N_2}(r)
 =\mu(r)\bigl(
   \one_{r\le N_1}+\one_{r\le N_2}-1
  \bigr).}
 \tag{R.5}
\]
Thus high M\"obius coefficients are replaced exactly by a signed bank of two
short M\"obius occurrences.  Taking absolute values at this stage would
replace one unknown sign by divisor multiplicity and would lose the purpose
of the descent.

\subsection{Corner-matched thresholds and the protected matrix}

Let $E_K\to\infty$ with $E_K=K^{o(1)}$, and set
\[
 R_K=\left\lfloor\frac{K^{1/2}}{E_K}\right\rfloor,
 \qquad
 Y_K=\left\lfloor\frac{K}{R_K+1}\right\rfloor.
 \tag{R.6}
\]
Then
\[
 R_K=K^{1/2}E_K^{-1}(1+o(1)),
 \qquad
 Y_K=K^{1/2}E_K(1+o(1)),
 \qquad
 R_KY_K=K(1+o(1)),
 \tag{R.7}
\]
and
\[
 R_KY_K<K<(R_K+1)(Y_K+1).
 \tag{R.8}
\]

\begin{theorem}[Corner-matched rectangular closure]
\label{thm:rectangular-corner}
One has the exact identity
\[
 \boxed{
 M_\chi(K)=M_\chi(R_K)+M_\chi(Y_K)
 -\mathcal Q_\chi(K;R_K,Y_K).}
 \tag{R.9}
\]
The two linear terms are $K^{1/2+o(1)}$ by the trivial estimate.  Therefore
\[
 \abs{\mathcal Q_\chi(K;R_K,Y_K)}\le K^{1/2+o(1)}
 \tag{R.10}
\]
implies
\[
 M_\chi(K)=O(K^{1/2+o(1)}),
\]
and hence GRH for $L(s,\chi)$.
\end{theorem}

\begin{proof}
The floor definition in \textup{(R.6)} gives \textup{(R.8)}.  Apply
\Cref{thm:rectangular-two-threshold}.  The implication to GRH is the reciprocal
Dirichlet-series argument of \Cref{thm:phase-action-GRH}.
\end{proof}

Put
\[
 A_{\chi,K}(m,n)
 =S_\chi\!\left(\left\lfloor\frac{K}{mn}\right\rfloor\right),
 \qquad
 1\le m\le R_K,
 \quad
 1\le n\le Y_K.
 \tag{R.11}
\]
Then $\mathcal Q_\chi$ is the bilinear transpose of this matrix against the
two M\"obius--character vectors.  For a complex character this is not a
Hermitian quadratic form.

\begin{lemma}[Finite conductor alphabet]
\label{lem:rectangular-finite-alphabet}
If
\[
 C_\chi=\max_{0\le r<f}|S_\chi(r)|,
\]
then
\[
 S_\chi(t)=S_\chi(t\bmod f),
 \qquad
 |A_{\chi,K}(m,n)|\le C_\chi.
 \tag{R.12}
\]
Thus the reader takes values in one fixed finite set, independently of $K$.
\end{lemma}

\begin{proof}
The sum of a nonprincipal character over every complete conductor block is
zero.  Remove complete blocks from the interval $[1,\lfloor t\rfloor]$.
\end{proof}

The bounded alphabet does not imply a small ambient operator norm.

\begin{theorem}[Macroscopic quotient-one block]
\label{thm:rectangular-ambient-obstruction}
Fix $\theta\in(2^{-1/2},1)$.  For all sufficiently large $K$, the rectangle
\[
 \theta R_K<m\le R_K,
 \qquad
 \theta Y_K<n\le Y_K
\]
satisfies
\[
 \left\lfloor\frac{K}{mn}\right\rfloor=1,
 \qquad A_{\chi,K}(m,n)=1.
\]
Consequently
\[
 \boxed{
 \norm{A_{\chi,K}}_{2\to2}
 \gg_\theta(R_KY_K)^{1/2}\asymp K^{1/2}.}
 \tag{R.13}
\]
In particular, an ambient operator-norm estimate cannot prove
\textup{(R.10)} at the natural $\ell^2$ sizes of the two M\"obius vectors.
\end{theorem}

\begin{proof}
By \textup{(R.7)}, $R_KY_K/K\to1$.  Since $\theta^2>1/2$, the indicated
products satisfy $K/2<mn<K$ for large $K$.  The corresponding matrix block is
an all-ones rectangle of dimensions comparable to $R_K$ and $Y_K$, whose
nonzero singular value is the square root of the product of those dimensions.
\end{proof}

This is a protected-vector obstruction, not evidence against the desired
scalar estimate.  It proves only that a positive action controlling every
ambient direction is nonminimal.

\subsection{High-reader closure and product compression}

Let
\[
 q_K(m,n)=\left\lfloor\frac{K}{mn}\right\rfloor
\]
and split $\mathcal Q_\chi=\mathcal Q_\chi^{\ge R_K}
+\mathcal Q_\chi^{<R_K}$ according as $q_K(m,n)\ge R_K$ or $q_K(m,n)<R_K$.

\begin{lemma}[Hyperbolic pair count]
\label{lem:rectangular-pair-count}
For $U,V,T\ge1$,
\[
 \#\{(m,n):m\le U,\ n\le V,\ mn\le T\}
 \le T\bigl(1+\log(2U)\bigr).
 \tag{R.14}
\]
\end{lemma}

\begin{proof}
The left side is at most
\[
 \sum_{m\le U}\min\{V,\lfloor T/m\rfloor\}
 \le T\sum_{m\le U}\frac1m.
\]
\end{proof}

\begin{theorem}[Absolute closure of the high reader]
\label{thm:rectangular-high-reader}
One has
\[
 \boxed{
 \abs{\mathcal Q_\chi^{\ge R_K}(K)}
 \ll_\chi\frac{K}{R_K}\log(2R_K)
 =K^{1/2+o(1)}.}
 \tag{R.15}
\]
The residual current is exactly
\[
 \boxed{
 \mathcal Q_\chi^{<R_K}(K)
 =\sum_{k=1}^{R_K-1}\chi(k)
   \sum_{\substack{m\le R_K,\ n\le Y_K\\
                   K/R_K<mn\le K/k}}
   \mu(m)\mu(n)\chi(mn).}
 \tag{R.16}
\]
\end{theorem}

\begin{proof}
The high-reader condition implies $mn\le K/R_K$.  Use
\Cref{lem:rectangular-finite-alphabet,lem:rectangular-pair-count}.  For the
second identity, open the partial character sum and interchange finite sums;
$q_K(m,n)<R_K$ is equivalent to $mn>K/R_K$.
\end{proof}

The first row in \textup{(R.16)} must remain assembled with the later rows.
Indeed,
\[
 \sum_{\substack{m\le R_K,\ n\le Y_K\\mn>K/R_K}}
 \mu(m)\mu(n)\chi(mn)
 =M_\chi(R_K)M_\chi(Y_K)+K^{1/2+o(1)},
 \tag{R.17}
\]
because the omitted lower-product region has only $K^{1/2+o(1)}$ pairs.  A
separate estimate for this rank-one row would already require strong Mertens
control and would discard the cancellation created by the exact descent.

Define the truncated product coefficient
\[
 \beta_{R_K,Y_K}(r)
 =\sum_{\substack{mn=r\\m\le R_K,\ n\le Y_K}}\mu(m)\mu(n),
 \qquad
 c_{\chi,K}(r)=\chi(r)\beta_{R_K,Y_K}(r).
 \tag{R.18}
\]
Then $|\beta_{R_K,Y_K}(r)|\le\tau(r)$ and
\[
 \boxed{
 \mathcal Q_\chi(K;R_K,Y_K)
 =\sum_{r\le R_KY_K}c_{\chi,K}(r)
 S_\chi\!\left(\left\lfloor\frac{K}{r}\right\rfloor\right).}
 \tag{R.19}
\]

\begin{lemma}[Divisor-square energy]
\label{lem:rectangular-divisor-square}
For $Z\ge2$,
\[
 \sum_{r\le Z}\tau(r)^2\ll Z(1+\log Z)^3.
 \tag{R.20}
\]
\end{lemma}

\begin{proof}
For every prime power $p^a$, $(a+1)^2\le\binom{a+3}{3}$, so
$\tau(r)^2\le d_4(r)$.  Summing ordered four-factor representations gives
\[
 \sum_{r\le Z}d_4(r)
 \le Z\sum_{a,b,c\le Z}\frac1{abc}
 \ll Z(1+\log Z)^3.
\]
\end{proof}

Let $W_K(r)=S_\chi(\lfloor K/r\rfloor)$.  Squaring the residual product sum
gives
\[
 |\mathcal Q_\chi^{<R_K}(K)|^2
 =\mathcal D_{\chi,K}+\mathcal O_{\chi,K},
 \tag{R.21}
\]
where
\[
 \mathcal D_{\chi,K}
 =\sum_{K/R_K<r\le R_KY_K}|c_{\chi,K}(r)|^2|W_K(r)|^2
\]
and $\mathcal O_{\chi,K}$ is the coherently assembled sum over distinct
products $r\ne s$.

\begin{theorem}[Closure of the complete same-product diagonal]
\label{thm:rectangular-product-diagonal}
One has
\[
 \boxed{
 \mathcal D_{\chi,K}\ll_\chi K(1+\log K)^3=K^{1+o(1)}.}
 \tag{R.22}
\]
Consequently the rectangular GRH certificate is reduced to the protected
connected estimate
\[
 \boxed{\abs{\mathcal O_{\chi,K}}\le K^{1+o(1)}.}
 \tag{R.23}
\]
\end{theorem}

\begin{proof}
Use $|W_K(r)|\le C_\chi$, $|\beta_{R_K,Y_K}(r)|\le\tau(r)$, and
\Cref{lem:rectangular-divisor-square}.  If \textup{(R.23)} holds, then
$|\mathcal Q_\chi^{<R_K}(K)|\le K^{1/2+o(1)}$ by \textup{(R.21)}; combine
with \Cref{thm:rectangular-high-reader,thm:rectangular-corner}.
\end{proof}

The diagonal in this theorem includes all distinct factorizations with the
same product.  Retaining only the literal entry diagonal
$(m,n)=(m',n')$ would leave an artificial factorization-provenance energy.

\subsection{Short product collisions and determinant descent}

On a dyadic block $r,s\asymp P$, write $s=r+h$ and put
\[
 H_P=1+\frac{P^2}{K}.
 \tag{R.24}
\]
The quotient displacement satisfies
\[
 \left\lfloor\frac{K}{r}\right\rfloor
 -\left\lfloor\frac{K}{r+h}\right\rfloor
 =\frac{Kh}{r(r+h)}+O(1).
 \tag{R.25}
\]

\begin{theorem}[Short product shifts close below the cubic scale]
\label{thm:rectangular-short-shifts}
Let $\Omega_K=K^{o(1)}$.  On a near-diagonal block $r,s\asymp P$, the
complete absolute contribution of
\[
 0<|s-r|\le\Omega_KH_P
\]
to the connected square is $K^{1+o(1)}$, provided
\[
 \boxed{P\le K^{2/3}\Omega_K^{-1/3}.}
 \tag{R.26}
\]
Equivalently, all subpower quotient-gap collisions are harmless while the
quotient scale $K/P$ is at least $K^{1/3+o(1)}$.
\end{theorem}

\begin{proof}
For one positive shift $h$, Cauchy--Schwarz and
\Cref{lem:rectangular-divisor-square} give
\[
 \sum_{r\asymp P}
 |\beta_{R_K,Y_K}(r)\beta_{R_K,Y_K}(r+h)|
 \ll P(1+\log K)^3.
\]
Summing over $|h|\le\Omega_KH_P$ yields
\[
 \ll\Omega_K\left(P+\frac{P^3}{K}\right)(1+\log K)^3,
\]
which is $K^{1+o(1)}$ under \textup{(R.26)}.  Character and reader factors
are bounded.
\end{proof}

The remaining off-product current admits an exact four-M\"obius determinant
lift.  For nonzero $h$,
\[
 \boxed{
 \begin{aligned}
 &\sum_r c_{\chi,K}(r)\overline{c_{\chi,K}(r+h)}
 W_K(r)\overline{W_K(r+h)}\\
 &\quad=
 \sum_{\substack{m,a\le R_K,\ n,b\le Y_K\\ab-mn=h}}
 \mu(m)\mu(n)\mu(a)\mu(b)
 \chi(mn)\overline{\chi(ab)}
 W_K(mn)\overline{W_K(ab)}.
 \end{aligned}}
 \tag{R.27}
\]

\begin{theorem}[Primitive short-gcd descent]
\label{thm:rectangular-short-gcd}
For a nonzero term in \textup{(R.27)}, put
\[
 g=(m,a),\qquad m=gv,\qquad a=gu,\qquad (u,v)=1.
\]
Then $g\mid h$ and, with $\ell=h/g$,
\[
 \boxed{ub-vn=\ell.}
 \tag{R.28}
\]
The common factor cancels exactly from the arithmetic phase:
\[
 \boxed{
 \mu(m)\mu(a)\chi(m)\overline{\chi(a)}
 =\mu(u)\mu(v)\chi(v)\overline{\chi(u)}.}
 \tag{R.29}
\]
Moreover,
\[
 \boxed{
 b\equiv\ell\overline u\pmod v,
 \qquad
 n\equiv-\ell\overline v\pmod u.}
 \tag{R.30}
\]
\end{theorem}

\begin{proof}
Nonzero M\"obius weights force $m$ and $a$ to be squarefree.  Thus $g,u,v$
are pairwise coprime and
\[
 \mu(gv)\mu(gu)=\mu(g)^2\mu(u)\mu(v)=\mu(u)\mu(v).
\]
Nonzero character weights force $(g,f)=1$, so the common character factor has
modulus one and cancels.  Substitution into $ab-mn=h$ gives
$g(ub-vn)=h$.  Reduction modulo $v$ and modulo $u$ gives the reciprocal
congruences.
\end{proof}

\begin{criterion}[Protected rectangular determinant criterion]
\label{crit:rectangular-determinant}
For every fixed nonprincipal primitive $\chi$, suppose the coherently
assembled remote product-collision current remaining after
\Cref{thm:rectangular-high-reader,thm:rectangular-product-diagonal,thm:rectangular-short-shifts}
satisfies \textup{(R.23)}.  Then
\[
 M_\chi(K)=O(K^{1/2+o(1)}),
\]
and GRH holds for $L(s,\chi)$.
\end{criterion}

The first unresolved source in this route is therefore a finite-state,
four-M\"obius determinant correlation on \textup{(R.28)}, concentrated in
remote product collisions and the low-quotient range below $K^{1/3+o(1)}$.
The fixed conductor supplies a bounded reader but cannot by itself provide
the terminal cancellation: shifts divisible by $f$ carry neutral character
phase and remain inside the determinant bank.


% =============================================================================
% =============================================================================
\section{Frequency-resolved completion of the rectangular determinant current}
\label{sec:grh-frequency-resolved}
% =============================================================================

The primitive reciprocal congruences in \textup{(R.30)} are the first point
at which a genuine Kloosterman-fraction phase can enter the fixed-character
programme.  The operation must nevertheless be performed on the exact
determinant equation rather than on its congruence envelope.  This section
records the equality-preserving completion, separates its nonzero and zero
frequency channels, and identifies the finite source geometry which remains
after the already proved high-reader, same-product, and short-collision
closures.

\subsection{The low-reader circularity firewall}

Retain the decomposition
\[
 \mathcal Q_\chi(K;R_K,Y_K)
 =\mathcal Q_\chi^{\ge R_K}(K)
  +\mathcal Q_\chi^{<R_K}(K)
\]
from \Cref{thm:rectangular-high-reader}.

\begin{theorem}[The scalar low-reader corner is GRH-strength]
\label{thm:freq-low-corner}
One has the exact identity
\[
 \boxed{
 \mathcal Q_\chi^{<R_K}(K)
 =-M_\chi(K)+M_\chi(R_K)+M_\chi(Y_K)
  -\mathcal Q_\chi^{\ge R_K}(K).}
 \tag{F.1}
\]
Consequently,
\[
 \boxed{
 \mathcal Q_\chi^{<R_K}(K)
 =-M_\chi(K)+K^{1/2+o(1)}.}
 \tag{F.2}
\]
Thus a square-root estimate for the unsquared low-reader scalar is equivalent
to the twisted Mertens estimate at height $K$.
\end{theorem}

\begin{proof}
Subtract the high-reader part from the corner identity
\textup{(R.9)} and rearrange.  The two linear terms are
$K^{1/2+o(1)}$ by the trivial estimate, while the high-reader term is
$K^{1/2+o(1)}$ by \Cref{thm:rectangular-high-reader}.  This gives
\textup{(F.2)}.  The equivalence follows by applying the same identity in both
directions.
\end{proof}

\begin{obstruction}[No intermediate scalar-corner lemma]
\label{obs:freq-scalar-circular}
The estimate
\[
 |\mathcal Q_\chi^{<R_K}(K)|\le K^{1/2+o(1)}
\]
may be used as a final conditional closure, but not as a purportedly weaker
input to prove GRH.  By \Cref{thm:freq-low-corner}, it is already GRH-strength.
The admissible downstream targets are therefore the connected determinant
channels created only after the protected square and exact completion.
\end{obstruction}

\subsection{Exact equality-preserving determinant completion}

Write
\[
 \mathcal R_q:=\left(-\frac q2,\frac q2\right]\cap\Z
\]
for the symmetric complete residue system.  Let $u,v$ be coprime positive
integers, let $\ell\in\Z$, and let $(c_b)$ and $(d_n)$ be finitely supported
complex sequences.  For $0\le\theta\le1$, define
\[
 \widehat c_{v,\theta}(r)
 :=\sum_b c_b\,\ee{(r+\theta)b/v},
 \qquad r\in\mathcal R_v,
\]
\[
 \widehat d_{u,\theta}(s)
 :=\sum_n d_n\,\ee{(s-\theta)n/u},
 \qquad s\in\mathcal R_u.
 \tag{F.3}
\]
Here and below $\overline u$ denotes the inverse of $u$ modulo $v$, and
$\overline v$ denotes the inverse of $v$ modulo $u$, whenever these symbols
occur in a residue phase.

\begin{lemma}[Exact affine lattice of the determinant equation]
\label{lem:freq-affine-lattice}
Choose residues $b_0\pmod v$ and $n_0\pmod u$ satisfying
\[
 ub_0\equiv\ell\pmod v,
 \qquad
 vn_0\equiv-\ell\pmod u,
\]
and put
\[
 \tau_0=\frac{ub_0-vn_0-\ell}{uv}\in\Z.
\]
Every pair in the congruence envelope is uniquely of the form
$b=b_0+vy$, $n=n_0+uz$, and
\[
 \boxed{ub-vn-\ell=uv(\tau_0+y-z).}
 \tag{F.4}
\]
Consequently
\[
 \sum_{ub-vn=\ell}c_bd_n
 =\int_0^1\ee{\theta\tau_0}
   \left(\sum_yc_{b_0+vy}\ee{\theta y}\right)
   \left(\sum_zd_{n_0+uz}\ee{-\theta z}\right)\dd\theta.
 \tag{F.5}
\]
\end{lemma}

\begin{proof}
The residue parametrization is unique.  Substitution gives
\textup{(F.4)}.  Expanding the right side of \textup{(F.5)} and using
\[
 \int_0^1\ee{\theta m}\dd\theta=\one_{m=0}
 \qquad(m\in\Z)
\]
projects exactly onto $\tau_0+y-z=0$, which is equivalent to
$ub-vn=\ell$.
\end{proof}

\begin{theorem}[Exact double completion]
\label{thm:freq-double-completion}
For arbitrary finitely supported $c_b,d_n$ and coprime $u,v$,
\[
 \boxed{
 \begin{aligned}
 \sum_{ub-vn=\ell}c_bd_n
 ={}&\frac1{uv}\int_0^1
 \sum_{r\in\mathcal R_v}\sum_{s\in\mathcal R_u}
 \widehat c_{v,\theta}(r)\widehat d_{u,\theta}(s)\\
 &\times
 \ee{-(r+s)\ell\overline u/v}
 \ee{(s-\theta)\ell/(uv)}\dd\theta.
 \end{aligned}}
 \tag{F.6}
\]
The formula preserves the exact determinant equation; no congruence aliases
are introduced.
\end{theorem}

\begin{proof}
Apply finite Fourier inversion to the two residue-class sums in
\textup{(F.5)}.  The first contributes
\[
 \frac{\ee{-\theta b_0/v}}v
 \sum_{r\in\mathcal R_v}\ee{-rb_0/v}
 \widehat c_{v,\theta}(r),
\]
and the second contributes
\[
 \frac{\ee{\theta n_0/u}}u
 \sum_{s\in\mathcal R_u}\ee{-sn_0/u}
 \widehat d_{u,\theta}(s).
\]
The definition of $\tau_0$ reduces the continuous phase to
$\ee{-\theta\ell/(uv)}$.  The two canonical congruences and the reciprocity
identity
\[
 \frac{\overline u}{v}+\frac{\overline v}{u}
 \equiv\frac1{uv}\pmod1
\]
combine the discrete phases into the two factors displayed in
\textup{(F.6)}.
\end{proof}

The principal finite-field frequency is
\[
 \boxed{\nu:=r+s.}
 \tag{F.7}
\]
For $\nu\ne0$, the first phase in \textup{(F.6)} is a reciprocal-residue
phase.  For $\nu=0$, it disappears identically.

\begin{theorem}[Determinant displacement and completion frequency are independent]
\label{thm:freq-independent-labels}
Fix $\ell\ne0$.  The same determinant displacement occurs simultaneously in
the zero channel $\nu=0$ and in every nonzero frequency admitted by the two
completion ranges.  Hence a restriction $|\ell|\le H$ does not determine
whether the completed source is phase-free.
\end{theorem}

\begin{proof}
The displacement $\ell$ is fixed before the independent Fourier indices
$r,s$ are introduced.  Formula \textup{(F.6)} leaves $\ell$ unchanged while
summing over their total $\nu=r+s$.
\end{proof}

\begin{obstruction}[No phase-free radius threshold]
\label{obs:freq-no-radius-threshold}
There is no implication
\[
 1\le |\ell|<v
 \quad\Longrightarrow\quad
 \text{the completed determinant source is phase-free}.
\]
For example, with $\ell=1$ and $\nu=1$, the map
\[
 u\longmapsto\ee{-\overline u/v}
\]
is nonconstant on $(\Z/v\Z)^\times$ for every prime $v\ge3$.
\end{obstruction}

\begin{proof}
Inversion permutes the nonzero residue classes modulo a prime.  The displayed
phase therefore assumes nontrivial $v$th roots of unity although the physical
displacement equals one.
\end{proof}

Applied blockwise to the determinant lift \textup{(R.27)},
\Cref{thm:freq-double-completion} gives the exact decomposition
\[
 \boxed{
 \mathcal O_{\chi,K}
 =\mathcal O_{\chi,K}^{(0)}
  +\mathcal O_{\chi,K}^{(\ne0)},}
 \tag{F.8}
\]
where the superscript records the completion-frequency channel.  This is a
source decomposition, not an absolute-value split; all conductor, reader,
and dyadic states remain assembled inside each channel.

\subsection{The zero-frequency sinc sheet}

For coprime $u,v$, put
\[
 \mathcal J_{u,v}:=\mathcal R_u\cap(-\mathcal R_v),
 \qquad
 S_{u,v}:=|\mathcal J_{u,v}|.
\]
The set $\mathcal J_{u,v}$ is a consecutive interval; write
\[
 \mathcal J_{u,v}
 =\{m_{u,v},m_{u,v}+1,\ldots,m_{u,v}+S_{u,v}-1\}.
\]
One has
\[
 \min\{u,v\}-1\le S_{u,v}\le\min\{u,v\}.
 \tag{F.9}
\]
Define the resolution width
\[
 \boxed{W_{u,v}:=\frac{uv}{S_{u,v}}\asymp\max\{u,v\}.}
 \tag{F.10}
\]

\begin{definition}[Zero-frequency determinant kernel]
\label{def:freq-zero-kernel}
For $D\in\Z$, define
\[
 \boxed{
 \mathscr Z_{u,v}(D)
 :=\frac1{uv}\int_0^1
   \sum_{s\in\mathcal J_{u,v}}
   \ee{-(s-\theta)D/(uv)}\dd\theta.}
 \tag{F.11}
\]
\end{definition}

\begin{proposition}[Exact physical representation of the zero channel]
\label{prop:freq-zero-physical}
For finitely supported coefficient systems
$\rho_\ell,\alpha_u,\beta_v,c_b,d_n$, the zero-frequency part of the exact
completion is
\[
 \boxed{
 \mathscr D^{(0)}
 =\sum_{\ell}\rho_\ell
  \sum_{(u,v)=1}\alpha_u\beta_v
  \sum_{b,n}c_bd_n\,
  \mathscr Z_{u,v}(ub-vn-\ell).}
 \tag{F.11a}
\]
Any bounded conductor or reader weights may be incorporated into the displayed
coefficient systems without changing the kernel.
\end{proposition}

\begin{proof}
Restrict \textup{(F.6)} to $r+s=0$, expand the two shifted Fourier transforms,
and interchange the finite sums with the $\theta$ integral.  The combined
phase is
\[
 \ee{-(s-\theta)(ub-vn-\ell)/(uv)},
\]
which is exactly the kernel in \textup{(F.11)}.
\end{proof}

\begin{theorem}[Exact sinc formula]
\label{thm:freq-sinc}
Let
\[
 \operatorname{sinc}_\pi(x)
 :=\begin{cases}
 \dfrac{\sin(\pi x)}{\pi x},&x\ne0,\\[0.8em]
 1,&x=0.
 \end{cases}
\]
There exists $\gamma_{u,v}$ with $|\gamma_{u,v}|\le1$ such that
\[
 \boxed{
 \mathscr Z_{u,v}(D)
 =\ee{\gamma_{u,v}D/(uv)}
  \frac1{W_{u,v}}
  \operatorname{sinc}_\pi\!\left(\frac D{W_{u,v}}\right).}
 \tag{F.12}
\]
Consequently,
\[
 \boxed{
 |\mathscr Z_{u,v}(D)|
 \le\min\left\{\frac1{W_{u,v}},\frac1{\pi|D|}\right\}}
 \qquad(D\ne0),
 \tag{F.13}
\]
and $\mathscr Z_{u,v}(0)=W_{u,v}^{-1}$.
\end{theorem}

\begin{proof}
Put $x=D/(uv)$.  The consecutive geometric sum is
\[
 \sum_{s\in\mathcal J_{u,v}}\ee{-sx}
 =\ee{-(m_{u,v}+(S_{u,v}-1)/2)x}
   \frac{\sin(\pi S_{u,v}x)}{\sin(\pi x)}.
\]
Furthermore,
\[
 \int_0^1\ee{\theta x}\dd\theta
 =\ee{x/2}\frac{\sin(\pi x)}{\pi x}.
\]
The factors $\sin(\pi x)$ cancel.  Substituting
$W_{u,v}=uv/S_{u,v}$ yields \textup{(F.12)}, with
\[
 \gamma_{u,v}
 =\frac12-m_{u,v}-\frac{S_{u,v}-1}{2}.
\]
The symmetric endpoint convention gives $|\gamma_{u,v}|\le1$.  The bound
follows from $|\sin y|\le\min\{|y|,1\}$.
\end{proof}

The squared kernel is the Fej\'er-type profile
\[
 |\mathscr Z_{u,v}(D)|^2
 =\frac1{W_{u,v}^2}
  \operatorname{sinc}_\pi^2\!\left(\frac D{W_{u,v}}\right).
 \tag{F.14}
\]
At the protected-amplitude level, however, the kernel remains the signed sinc
in \textup{(F.12)}; replacing it by a positive window is a stronger problem.

\begin{theorem}[Core--annulus--tail compiler]
\label{thm:freq-core-tail}
Fix $\eta>0$.  Uniformly in $u,v$:
\begin{enumerate}[label=\textup{(\roman*)},leftmargin=3em]
\item if $|D|\le W_{u,v}K^{-\eta}$, then
\[
 \mathscr Z_{u,v}(D)
 =\frac1{W_{u,v}}\bigl(1+O(K^{-\eta})\bigr);
 \tag{F.15}
\]
\item if $|D|\ge W_{u,v}K^{\eta}$, then
\[
 |\mathscr Z_{u,v}(D)|
 \le\frac{K^{-\eta}}{W_{u,v}}.
 \tag{F.16}
\]
\end{enumerate}
\end{theorem}

\begin{proof}
In the core, $D/W_{u,v}=O(K^{-\eta})$, so the sinc factor is
$1+O(K^{-2\eta})$ and the harmless phase in \textup{(F.12)} is
$1+O(K^{-\eta})$.  In the tail, use \textup{(F.13)}.
\end{proof}

For a zero-frequency block with coefficient majorant
\[
 \mathfrak M_0
 :=\sum_{\ell,u,v,b,n}
   \frac{|\rho_\ell\alpha_u\beta_vc_bd_n|}{W_{u,v}},
 \tag{F.17}
\]
\Cref{thm:freq-core-tail} gives a relative $K^{-\eta}$ bound for the far tail
and for the error made by replacing the ultra-core by the constant strip
kernel $W_{u,v}^{-1}$.  The irreducible zero channel is therefore the
phase-free sinc-resolution strip
\[
 \boxed{|ub-vn-\ell|\le W_{u,v}K^{o(1)}.}
 \tag{F.18}
\]

\subsection{Product compression and one-dimensional entanglement}

Two further exact reductions remove coordinate artifacts before a spectral
estimate is attempted.

\begin{theorem}[Exact inverse-class product compression]
\label{thm:freq-product-compression}
Let $I_U,I_B$ be finite intervals, put
\[
 d_{\alpha,c}(m)
 :=\sum_{\substack{u\in I_U,\ b\in I_B\\ub=m}}\alpha_uc_b,
\]
and let $(A,q)=1$.  Then
\[
 \boxed{
 \begin{aligned}
 &\sum_{\substack{u\in I_U\\(u,q)=1}}\alpha_u
 \left(
  \sum_{\substack{b\in I_B\\b\equiv A\overline u\ ({\rm mod}\ q)}}c_b
  -\frac1{\varphi(q)}
   \sum_{\substack{b\in I_B\\(b,q)=1}}c_b
 \right)\\
 &\quad=
 \sum_{m\equiv A\ ({\rm mod}\ q)}d_{\alpha,c}(m)
 -\frac1{\varphi(q)}
  \sum_{(m,q)=1}d_{\alpha,c}(m).
 \end{aligned}}
 \tag{F.19}
\]
Thus the inverse map $u\mapsto A\overline u$ is only a coordinate chart for
the fixed product congruence $ub\equiv A\pmod q$.
\end{theorem}

\begin{proof}
Because $(A,q)=1$, the inverse-class congruence is equivalent to
$ub\equiv A\pmod q$.  Group both the progression term and the reduced-residue
mean by the product $m=ub$.
\end{proof}

\begin{corollary}[Imported logarithmic closure of admissible unbalanced leaves]
\label{cor:freq-FR-leaf}
Suppose the product coefficient in \textup{(F.19)} has a dyadic
factorization $d_{\alpha,c}=a*b$ with one divisor-bounded factor satisfying
the Siegel--Walfisz condition, and suppose its lengths and modulus fall in one
of the range systems of the Fouvry--Radziwi\l\l{} unbalanced-convolution
theorem \cite{FouvryRadziwill}.  Then the fixed-residue discrepancy in
\textup{(F.19)} has arbitrary logarithmic saving.  This removes the
inverse-class obstruction on every such separated leaf, but does not by
itself supply a fixed power at the boundary of the published range.
\end{corollary}

\begin{proof}
After \Cref{thm:freq-product-compression}, the leaf is exactly the product
progression discrepancy to which the cited theorem applies.  Its logarithmic
exponent may be chosen arbitrarily large.
\end{proof}

For the balanced generation-two coefficient geometry, fix integers $v$ and
$\Delta$ with $v(v+\Delta)\ne0$ and define
\[
 \mathsf T_{v,\Delta}(n,\ell,\ell')
 :=(vn+\ell,\,\Delta n+\ell'-\ell).
 \tag{F.20}
\]

\begin{theorem}[Primitive one-dimensional fibre kernel]
\label{thm:freq-lattice-kernel}
One has
\[
 \boxed{
 \Ker_{\Z}\mathsf T_{v,\Delta}
 =\Z(1,-v,-(v+\Delta)).}
 \tag{F.21}
\]
If $c_{n,\ell,\ell'}$ is finitely supported and
\[
 \Gamma_{v,\Delta}(w,a)
 :=\sum_{\substack{vn+\ell=w\\
                    \Delta n+\ell'-\ell=a}}
 c_{n,\ell,\ell'},
\]
then
\[
 \boxed{
 \sum_{w,a}|\Gamma_{v,\Delta}(w,a)|^2
 =\sum_{h\in\Z}\mathcal C_{v,\Delta}(h),}
 \tag{F.22}
\]
where
\[
 \mathcal C_{v,\Delta}(h)
 :=\sum_{n,\ell,\ell'}
 c_{n,\ell,\ell'}
 \overline{c_{n-h,\ell+vh,\ell'+(v+\Delta)h}}.
 \tag{F.23}
\]
The term $h=0$ is the literal source diagonal.
\end{theorem}

\begin{proof}
Solving $\mathsf T_{v,\Delta}(n,\ell,\ell')=(0,0)$ gives
$\ell=-vn$ and $\ell'=-(v+\Delta)n$, proving \textup{(F.21)}.  Expand the
square defining the left side of \textup{(F.22)}.  Two source triples have the
same image exactly when their difference is an integral multiple of the
primitive kernel generator; this gives \textup{(F.23)} and the stated sum.
\end{proof}

\begin{corollary}[Sharp deterministic short-fibre transfer]
\label{cor:freq-short-fibre}
Suppose the three source coordinates have interval diameters $N_0,L_1,L_2$
and put
\[
 H_{v,\Delta}
 :=\min\left\{N_0,\frac{L_1}{|v|},
                    \frac{L_2}{|v+\Delta|}\right\}.
\]
Then
\[
 \boxed{
 \sum_{w,a}|\Gamma_{v,\Delta}(w,a)|^2
 \le(2\lfloor H_{v,\Delta}\rfloor+3)
     \sum_{n,\ell,\ell'}|c_{n,\ell,\ell'}|^2.}
 \tag{F.24}
\]
If $L_1<|v|$ and $L_2<|v+\Delta|$, the map is injective on the support and
\textup{(F.24)} is the exact isometry with constant one.
\end{corollary}

\begin{proof}
Only $|h|\le H_{v,\Delta}+1$ can contribute to \textup{(F.22)}.  Each
correlation is at most the source norm squared by Cauchy--Schwarz.  Under the
strict short-fibre inequalities no nonzero integer $h$ survives.
\end{proof}

The factor in \textup{(F.24)} is sharp for arbitrary coefficients concentrated
on one fibre.  Arithmetic cancellation is therefore required only in the
long-fibre bank; paying an ambient three-index rank would be nonminimal.
Moreover, replacing the coherent $h$-average by a uniform estimate for every
fixed $h$ would demand a pointwise two-point M\"obius correlation theorem,
which is stronger than the source-native averaged target.

\begin{lemma}[Complete residue periods in a long reciprocal-phase bank]
\label{lem:freq-complete-periods}
Let $\rho_t$ be supported on a consecutive interval of length $L$.  For a
modulus $q$ and $r\pmod q$, put
\[
 \rho_{q,r}:=\sum_{t\equiv r\ ({\rm mod}\ q)}\rho_t.
\]
Then
\[
 \sum_t\rho_t\ee{-\nu t\overline p/q}
 =\sum_{r\ ({\rm mod}\ q)}\rho_{q,r}
   \ee{-\nu r\overline p/q},
 \tag{F.25}
\]
\[
 \sum_{r\ ({\rm mod}\ q)}|\rho_{q,r}|^2
 \le\left(\left\lfloor\frac Lq\right\rfloor+1\right)
      \sum_t|\rho_t|^2.
 \tag{F.26}
\]
The interval contains $\lfloor L/q\rfloor$ complete residue periods and an
edge of fewer than $q$ points.  If $|\rho_t|\le B$, the absolute edge
majorant is at most $Bq$, whereas the full cardinality majorant is $BL$.
Hence, when $L/q$ is a fixed power, the edge has a fixed-power saving relative
to the source majorant and the interior retains complete
Kloosterman-fraction phases for every primitive $\nu\ne0$.
\end{lemma}

\begin{proof}
The phase depends only on $t\pmod q$, proving \textup{(F.25)}.  Apply
Cauchy--Schwarz inside each residue class for \textup{(F.26)}.  Euclidean
division partitions the interval into complete $q$-blocks and one remainder
of fewer than $q$ points; the bounded-coefficient majorants then have the
stated sizes.
\end{proof}

\subsection{Frequency-resolved rectangular closure}

\begin{criterion}[Nonzero-phase and sinc-sheet closure]
\label{crit:freq-rectangular-closure}
Apply the exact completion \textup{(F.6)} to every remote determinant block
remaining after \Cref{thm:rectangular-high-reader,thm:rectangular-product-diagonal,thm:rectangular-short-shifts}.  Suppose, in the normalization of
\textup{(R.23)}, that the coherently assembled nonzero-frequency current and
the coherently assembled zero-frequency sinc-sheet current satisfy
\[
 |\mathcal O_{\chi,K}^{(\ne0)}|
 +|\mathcal O_{\chi,K}^{(0)}|
 \le K^{1+o(1)}.
 \tag{F.27}
\]
Then $M_\chi(K)=O(K^{1/2+o(1)})$ and GRH holds for $L(s,\chi)$.
\end{criterion}

\begin{proof}
By the exact source decomposition \textup{(F.8)}, condition
\textup{(F.27)} gives \textup{(R.23)}.  Apply
\Cref{crit:rectangular-determinant}.
\end{proof}

The two terms in \textup{(F.27)} require different analytic mechanisms.  The
nonzero channel retains reciprocal phases and is the legitimate landing point
for source-normalized Kloosterman-fraction technology such as
\cite{DFI97,BettinChandee}; any factorization or reader-matrix hypothesis of
the chosen theorem must remain visible.  The zero channel has no reciprocal
phase and is reduced instead to the phase-free strip \textup{(F.18)}.  A
radius-only theorem uniform across both channels is sufficient but is not
source minimal, because it suppresses the independent frequency label.

% =============================================================================
\section{Cofinal Mellin visibility and conditioned finite tomography}
\label{sec:cofinal-mellin}
% =============================================================================

The finite source geometry distinguishes three different requirements: a
source must see a spectral coefficient, its reconstruction map must be
conditioned, and its norm must be compared with the actual arithmetic
source.  These requirements remain distinct even when every finite
matrix is explicitly invertible.  This section gives a normalized
construction for the first requirement, quantitative bounds for the
second, and the precise limitation on passing to a cofinal closure.

\subsection{The regularity cost of a visible source}

For a compactly supported smooth function with values in a finite-dimensional
Hilbert space, define its Mellin reader by
\begin{equation}
 \Xi_G(s)=\int_\R e^{(s-1/2)t}G(t)\,\dd t.
 \label{eq:vis-mellin}
\end{equation}
The source is chosen independently of the zeros to be tested.

\begin{proposition}[Fixed smooth banks have no cofinal visibility floor]
\label{prop:vis-fixed-bank}
For every integer $r\ge0$ and every $\gamma\ne0$,
\begin{equation}
 \|\Xi_G(1/2+i\gamma)\|
 \le |\gamma|^{-r}\|G^{(r)}\|_{L^1}.
 \label{eq:vis-ibp}
\end{equation}
Consequently a fixed finite smooth bank, including any fixed finite
collection of translations, derivatives, or polynomially weighted readers,
cannot have a positive cutoff-independent lower bound on the critical
line.  If a cutoff-dependent bank satisfies
$\|\Xi_{G_T}(1/2+iT)\|\ge a_T$, then necessarily
$\|G_T^{(r)}\|_{L^1}\ge a_TT^r$.
\end{proposition}

\begin{proof}
Integrate \eqref{eq:vis-mellin} by parts $r$ times.  There are no boundary
terms.  Each fixed modification in the statement is again a compactly
supported smooth function, and a fixed finite linear combination satisfies
the same conclusion.  Apply the estimate at $\gamma=T$ for the last claim.
\end{proof}

\begin{theorem}[A normalized cofinal modulation bank]
\label{thm:vis-modulation}
Choose a nonzero $\phi\in C_c^\infty(\R)$ with $\phi\ge0$.  There are
constants $r_0,c_0>0$, depending only on $\phi$, such that, with
\begin{equation}
 \begin{split}
 L_T&=\lceil T/r_0\rceil+1,\qquad J_T=2L_T+1,\\
 G_T(t)&=J_T^{-1/2}\phi(t)
       \bigl(e^{-ijr_0t}\bigr)_{|j|\le L_T},
 \end{split}
 \label{eq:vis-bank}
\end{equation}
one has
\begin{equation}
 \|G_T(t)\|=|\phi(t)|,\qquad
 \inf_{\substack{0\le\beta\le1\\|\gamma|\le T}}
 \|\Xi_{G_T}(\beta+i\gamma)\|
 \ge \frac{c_0}{\sqrt{J_T}},\qquad
 J_T\le 2T/r_0+5.
 \label{eq:vis-bank-floor}
\end{equation}
For each fixed $r$, $\|G_T^{(r)}\|_{L^1}=O_r((1+T)^r)$.
\end{theorem}

\begin{proof}
The continuous function
$F(\beta,\xi)=\int e^{(\beta-1/2)t+i\xi t}\phi(t)\dd t$
is strictly positive at $\xi=0$, uniformly for $\beta\in[0,1]$.
Uniform continuity on a compact neighbourhood gives
$|F(\beta,\xi)|\ge c_0$ for $|\xi|\le r_0/2$ after reducing $r_0$.
For each $|\gamma|\le T$, some lattice point $jr_0$ in the bank lies
within $r_0/2$ of $\gamma$.  Its component alone gives the lower bound.
The pointwise norm identity is exact.  Leibniz's rule and
$|jr_0|=O(1+T)$ give the derivative estimate.
\end{proof}

The square-root factor in \eqref{eq:vis-bank-floor} is the normalization
cost of this particular bank, not a universal optimality assertion over
all source constructions.  Removing $J_T^{-1/2}$ increases the source
norm by $\sqrt{J_T}$ and its Gram by $J_T$; it does not provide a free
improvement in the relative metric of the Gran tensor.

\subsection{Translation reconstruction and its two conditioning factors}

For a translation $G_{T,k}(t)=G_T(t-k\tau)$,
\begin{equation}
 \Xi_{G_{T,k}}(s)=e^{k\tau(s-1/2)}\Xi_{G_T}(s).
 \label{eq:vis-translation}
\end{equation}
Suppose an independently justified completed explicit formula, including
its truncation and boundary terms, has produced a finite sequence
\begin{equation}
 Y_k=\sum_{j=1}^M z_j^kv_j,\qquad
 z_j=e^{\tau(\rho_j-1/2)},\qquad
 v_j=\mu(\rho_j)\Xi_{G_T}(\rho_j).
 \label{eq:vis-prony}
\end{equation}
Repeated zeros are combined into their nonzero vector weights.  This is
a finite algebraic model, not an assertion that the omitted zeros have
zero contribution to an actual arithmetic reader.

\begin{theorem}[Visibility, multiplier, and Vandermonde separation]
\label{thm:vis-hankel}
Assume the $z_j$ are distinct, the $\rho_j$ lie in the strip in
\eqref{eq:vis-bank-floor}, and $|\mu(\rho_j)|\ge\mu_T>0$.
Let $V=(z_j^k)_{0\le k<M,\,1\le j\le M}$ and define the vector-valued
Hankel map by
$(\mathsf H a)_k=\sum_{\ell=0}^{M-1}a_\ell Y_{k+\ell}$.
Then
\begin{equation}
 \sigma_{\min}(\mathsf H)
 \ge \frac{c_0\mu_T}{\sqrt{J_T}}\sigma_{\min}(V)^2.
 \label{eq:vis-hankel-floor}
\end{equation}
In particular, $0<\tau<\pi/T$ prevents aliasing among distinct strip
points, but does not provide a lower bound on $\sigma_{\min}(V)$.
\end{theorem}

\begin{proof}
Writing $D_v a=(a_jv_j)_j$, one has the exact factorization
$\mathsf H=(V\otimes I)D_vV^{\mathsf T}$.
Moreover $\sigma_{\min}(D_v)=\min_j\|v_j\|$.
Multiplication of the three lower bounds gives
\eqref{eq:vis-hankel-floor}.  Equality of two exponential nodes forces
equality of their real parts and an ordinate difference in
$(2\pi/\tau)\Z$; the strict upper bound on $\tau$ excludes a nonzero
such difference in $[-2T,2T]$.
\end{proof}

\begin{proposition}[Stable recovery of the finite recurrence]
\label{prop:vis-recurrence-stability}
Let $q$ solve $\mathsf Hq=-b$.  If the observed data are
$\widetilde{\mathsf H}=\mathsf H+E$ and $\widetilde b=b+e$ with
$\|E\|<\sigma_{\min}(\mathsf H)$, the least-squares solution
$\widetilde q=-\widetilde{\mathsf H}^{\dagger}\widetilde b$ satisfies
\begin{equation}
 \|\widetilde q-q\|
 \le\frac{\|e\|+\|E\|\|q\|}
 {\sigma_{\min}(\mathsf H)-\|E\|}.
 \label{eq:vis-recurrence-stability}
\end{equation}
\end{proposition}

\begin{proof}
The perturbed map has full column rank, so
$\widetilde q-q=-\widetilde{\mathsf H}^{\dagger}(e+Eq)$.
The smallest singular value changes by at most $\|E\|$.
This argument applies to the generally overdetermined vector system;
it does not assume exact consistency of its noisy equations.
\end{proof}

\begin{obstruction}[Proper-arc conditioning of monomial translation]
\label{obs:vis-proper-arc}
If $M\ge2$, $|z_j|=1$, and all nodes, after a common rotation, satisfy
$|\arg z_j|\le\alpha<\pi$, then
\begin{equation}
 \sigma_{\min}(V)
 \le\sqrt{M(2M-1)}\sin(\alpha/2)^{M-1}.
 \label{eq:vis-proper-arc}
\end{equation}
Thus the square monomial reconstruction is exponentially ill-conditioned
at a fixed proper-arc margin, even when every retained zero is on the
critical line.
\end{obstruction}

\begin{proof}
For $d=M-1$, use the coefficient vector of $(z-1)^d$.  Its squared norm
is $\binom{2d}{d}\ge4^d/(2d+1)$, whereas the modulus of its value at
each node is at most $(2\sin(\alpha/2))^d$.  The resulting Rayleigh
quotient for $V^{\mathsf T}$ gives \eqref{eq:vis-proper-arc}.
\end{proof}

If the translation step is $\tau\le(1-\delta)\pi/T$, the right side
contains $\cos(\pi\delta/2)^{M-1}$.  Avoiding this particular
superpolynomial obstruction requires a margin no larger than order
$\sqrt{\log M/M}$ for schemes demanding a fixed polynomial conditioning
bound.  Changing basis or using a redundant frame may change the
obstruction, but the analysis and synthesis norms of that replacement
must be included in the source comparison.

\subsection{The cofinal admission boundary}

The preceding construction removes the invisibility of a fixed smooth
bank.  It does not remove small completed multipliers, clustered spectral
nodes, truncation error, or a deficient lower comparison with the actual
arithmetic source.  A cofinal use of \eqref{eq:vis-prony} therefore requires
a declared family of finite cylinders, a complete Gram comparison for
$G_T$ and all translated readers, and an error bound below the product
in \eqref{eq:vis-hankel-floor}.  Without these quantitative hypotheses,
finite reconstructibility and cofinal observability are different
statements.  This is the spectral counterpart of the distinction between
carrier birth and signed retuning in Part~I.


\section{Completed theta heat flow and finite Jensen collisions}
\label{sec:grh-theta-heat}
% =============================================================================

The expansion--pullback theorem shows that an exact factorization of a
terminal square cannot be promoted to an independent positive source merely
by changing coordinates.  A genuinely independent GRH entrance is supplied
by the completed theta kernel.  It is fixed by the conductor, parity, gamma
factor, Gauss root number, and theta series before a terminal prime sum or a
zero location is read.

Throughout this section, $\chi$ is a primitive nonprincipal character of
conductor $f$, and
\[
 \chi(-1)=(-1)^{\kappa_\chi},
 \qquad \kappa_\chi\in\{0,1\}.
\]
Define
\[
 \mathfrak L_\chi(s)
 =\left(\frac f\pi\right)^{(s+\kappa_\chi)/2}
  \Gamma\!\left(\frac{s+\kappa_\chi}{2}\right)L(s,\chi).
 \tag{J.1}
\]
It is entire of order one and satisfies
\[
 \mathfrak L_\chi(s)
 =\varepsilon_\chi\mathfrak L_{\overline\chi}(1-s),
 \qquad |\varepsilon_\chi|=1.
 \tag{J.2}
\]
Fix once and for all $\omega_\chi$ with
$\omega_\chi^2=\varepsilon_\chi$ and $|\omega_\chi|=1$.

\subsection{The cutoff-independent completed source}

\begin{definition}[Root-number normalized critical entire function]
\label{def:theta-critical-entire}
Put
\[
 \boxed{
 \mathscr X_\chi(z)
 =\omega_\chi^{-1}
  \mathfrak L_\chi\!\left(\frac{1+\mathrm i z}{2}\right).}
 \tag{J.3}
\]
\end{definition}

\begin{theorem}[Reality and exact GRH equivalence]
\label{thm:theta-reality-GRH}
The function $\mathscr X_\chi$ is a real entire function of order one.  A
zero $\rho=\beta+\mathrm i\gamma$ of $L(s,\chi)$ corresponds to the zero
\[
 z_\rho=2\gamma+\mathrm i(1-2\beta)
 \tag{J.4}
\]
of $\mathscr X_\chi$.  Consequently,
\[
 \boxed{
 \text{GRH for }L(s,\chi)
 \iff \mathscr X_\chi\text{ has only real zeros}.}
 \tag{J.5}
\]
\end{theorem}

\begin{proof}
Conjugation and \textup{(J.2)} give, for real $x$,
\[
 \begin{aligned}
 \overline{\mathscr X_\chi(x)}
 &=\overline{\omega_\chi^{-1}}
   \mathfrak L_{\overline\chi}
   \!\left(\frac{1-\mathrm i x}{2}\right)\\
 &=\omega_\chi\varepsilon_\chi^{-1}
   \mathfrak L_\chi
   \!\left(\frac{1+\mathrm i x}{2}\right)
 =\mathscr X_\chi(x).
 \end{aligned}
\]
Solving $(1+\mathrm i z)/2=\rho$ gives \textup{(J.4)}, which is real exactly
when $\beta=1/2$.
\end{proof}

Define the primitive theta series
\[
 \vartheta_\chi(x)
 =\sum_{n\ge1}n^{\kappa_\chi}\chi(n)
  \exp\!\left(-\frac{\pi n^2x}{f}\right),
 \qquad x>0.
 \tag{J.6}
\]

\begin{lemma}[Primitive theta transformation]
\label{lem:theta-transform}
One has
\[
 \boxed{
 \vartheta_\chi(x)
 =\varepsilon_\chi x^{-\kappa_\chi-1/2}
  \vartheta_{\overline\chi}(1/x).}
 \tag{J.7}
\]
\end{lemma}

\begin{proof}
Apply Poisson summation to the parity-compatible full series
\[
 \Theta_\chi(x)
 =\sum_{n\in\Z}n^{\kappa_\chi}\chi(n)e^{-\pi n^2x/f}.
\]
The terms at $n$ and $-n$ agree, so $\Theta_\chi=2\vartheta_\chi$.  The
finite Fourier transform of a primitive character contributes its Gauss sum,
and the Fourier transform of
$t^{\kappa_\chi}e^{-\pi xt^2/f}$ contributes
$\mathrm i^{-\kappa_\chi}x^{-\kappa_\chi-1/2}$.  The resulting prefactor is
$\varepsilon_\chi$, which proves the formula after division by two.
\end{proof}

\begin{theorem}[Mellin and Fourier realization]
\label{thm:theta-Fourier}
For $\Re s>1$,
\[
 \mathfrak L_\chi(s)
 =\int_0^\infty\vartheta_\chi(x)
  x^{(s+\kappa_\chi)/2}\frac{\dd x}{x},
 \tag{J.8}
\]
and the theta transformation extends the identity to the entire completion.
With
\[
 \boxed{
 \Phi_\chi(u)
 =4\omega_\chi^{-1}e^{(1+2\kappa_\chi)u}
  \vartheta_\chi(e^{4u}),}
 \tag{J.9}
\]
one has
\[
 \boxed{
 \mathscr X_\chi(z)
 =\int_{-\infty}^{\infty}\Phi_\chi(u)e^{\mathrm i zu}\,\dd u.}
 \tag{J.10}
\]
Moreover,
\[
 \Phi_\chi(-u)=\overline{\Phi_\chi(u)},
 \tag{J.11}
\]
and $\Phi_\chi$ has super-exponential decay at both ends, up to fixed
exponential factors.
\end{theorem}

\begin{proof}
Termwise integration in $\Re s>1$ gives
\[
 \begin{aligned}
 &\int_0^\infty\vartheta_\chi(x)
  x^{(s+\kappa_\chi)/2}\frac{\dd x}{x}\\
 &\qquad=\sum_{n\ge1}n^{\kappa_\chi}\chi(n)
   \Gamma\!\left(\frac{s+\kappa_\chi}{2}\right)
   \left(\frac{f}{\pi n^2}\right)^{(s+\kappa_\chi)/2}
 =\mathfrak L_\chi(s).
 \end{aligned}
\]
The transformation \textup{(J.7)} controls the integral at zero.  Now set
$s=(1+\mathrm i z)/2$ and $x=e^{4u}$ to obtain \textup{(J.9)}--\textup{(J.10)}.
Finally \textup{(J.7)} and $\varepsilon_\chi=\omega_\chi^2$ give
\textup{(J.11)}.  Gaussian decay for $u\to+\infty$, followed by
\textup{(J.11)}, gives the stated decay at both ends.
\end{proof}

\begin{remark}[Principal channel]
For the primitive principal channel one replaces \textup{(J.1)} by the
classical pole-removed Riemann completion
$\xi(s)=\tfrac12s(s-1)\pi^{-s/2}\Gamma(s/2)\zeta(s)$.  The heat-flow and
Jensen constructions below then become the classical de Bruijn--Newman
programme.  The nonprincipal formulation is written out because it is the
new channel-specific component needed for Dirichlet GRH.
\end{remark}

\subsection{Heat deformation and the de Bruijn--Newman boundary}

\begin{definition}[Completed theta heat flow]
\label{def:theta-heat}
For $t\in\R$, define
\[
 \boxed{
 \mathscr X_{\chi,t}(z)
 =\int_{-\infty}^{\infty}
  e^{tu^2}\Phi_\chi(u)e^{\mathrm i zu}\,\dd u.}
 \tag{J.12}
\]
\end{definition}

\begin{theorem}[Entire heat flow]
\label{thm:theta-heat-equation}
For every real $t$, \textup{(J.12)} defines a real entire function and
\[
 \boxed{
 \partial_t\mathscr X_{\chi,t}
 =-\partial_z^2\mathscr X_{\chi,t},
 \qquad
 \mathscr X_{\chi,0}=\mathscr X_\chi.}
 \tag{J.13}
\]
\end{theorem}

\begin{proof}
The decay in \Cref{thm:theta-Fourier} dominates
$e^{tu^2+C|u|}$ for fixed $t$ on compact $z$-sets, so differentiation under
the integral is valid.  The identities
$\partial_te^{tu^2}=u^2e^{tu^2}$ and
$\partial_z^2e^{\mathrm i zu}=-u^2e^{\mathrm i zu}$ give the heat equation.
Conjugate symmetry makes the integral real on the real axis.
\end{proof}

\begin{theorem}[de Bruijn--Newman--Dobner threshold]
\label{thm:theta-DN}
For the normalization \textup{(J.12)}, there is a real number
$\Lambda_{\rm DN}(\chi)$ such that $\mathscr X_{\chi,t}$ has only real zeros
if and only if
\[
 t\ge\Lambda_{\rm DN}(\chi).
 \tag{J.14}
\]
For primitive Dirichlet $L$-functions,
\[
 \boxed{\Lambda_{\rm DN}(\chi)\ge0.}
 \tag{J.15}
\]
Consequently,
\[
 \boxed{
 \text{GRH for }L(s,\chi)
 \iff \Lambda_{\rm DN}(\chi)=0.}
 \tag{J.16}
\]
\end{theorem}

\begin{proof}
The threshold and forward propagation are the de Bruijn--Newman theorem for
this completed Fourier source \cite{deBruijn,Newman}.  Dobner proved the
existence and nonnegativity of the corresponding constant throughout the
extended Selberg class, which includes primitive Dirichlet $L$-functions
\cite{Dobner}.  By \Cref{thm:theta-reality-GRH}, GRH is equivalent to
real-rootedness at $t=0$, hence to $0\ge\Lambda_{\rm DN}(\chi)$.  Combine
this with \textup{(J.15)}.
\end{proof}

A real-rootedness theorem only at one fixed positive heat time gives
$\Lambda_{\rm DN}(\chi)\le t_0$, not GRH.  The protected boundary is exactly
$t=0$.

\subsection{Finite Jensen compiler}

Expand
\[
 \mathscr X_{\chi,t}(z)
 =\sum_{m=0}^\infty\gamma_{\chi,m}(t)\frac{z^m}{m!}.
 \tag{J.17}
\]

\begin{lemma}[Theta moments and heat transport]
\label{lem:theta-moments}
For every $m\ge0$,
\[
 \boxed{
 \gamma_{\chi,m}(t)
 =\mathrm i^m\int_{-\infty}^{\infty}
  u^m e^{tu^2}\Phi_\chi(u)\,\dd u\in\R,}
 \tag{J.18}
\]
and
\[
 \boxed{
 \frac{\dd}{\dd t}\gamma_{\chi,m}(t)
 =-\gamma_{\chi,m+2}(t).}
 \tag{J.19}
\]
\end{lemma}

\begin{proof}
Differentiate \textup{(J.12)} at $z=0$.  Reality follows because
$\mathscr X_{\chi,t}$ is real entire.  Equation \textup{(J.19)} is the
coefficient form of \textup{(J.13)}.
\end{proof}

\begin{definition}[Jensen packet]
\label{def:theta-Jensen}
For integers $d,n\ge0$, put
\[
 \boxed{
 J_{\chi,t}^{d,n}(Y)
 =\sum_{j=0}^d\binom dj\gamma_{\chi,n+j}(t)Y^j.}
 \tag{J.20}
\]
The packet is called hyperbolic when all its zeros are real, with the actual
degree used if leading coefficients vanish.
\end{definition}

\begin{theorem}[Exact Jensen recurrences]
\label{thm:theta-Jensen-recurrences}
For $d\ge1$,
\[
 \boxed{
 \partial_YJ_{\chi,t}^{d,n}
 =dJ_{\chi,t}^{d-1,n+1},}
 \tag{J.21}
\]
\[
 \boxed{
 J_{\chi,t}^{d,n}
 =J_{\chi,t}^{d-1,n}
  +YJ_{\chi,t}^{d-1,n+1},}
 \tag{J.22}
\]
and
\[
 \boxed{
 \partial_tJ_{\chi,t}^{d,n}
 =-J_{\chi,t}^{d,n+2}.}
 \tag{J.23}
\]
\end{theorem}

\begin{proof}
Differentiate \textup{(J.20)} and use
$j\binom dj=d\binom{d-1}{j-1}$ for \textup{(J.21)}.  Pascal's identity gives
\textup{(J.22)}, and \Cref{lem:theta-moments} gives \textup{(J.23)}.
\end{proof}

\begin{theorem}[Jensen--P\'olya criterion]
\label{thm:theta-Jensen-Polya}
For each fixed $t$, the following are equivalent:
\begin{enumerate}[label=\textup{(\roman*)},leftmargin=2.8em]
\item $\mathscr X_{\chi,t}$ has only real zeros;
\item $\mathscr X_{\chi,t}$ belongs to the Laguerre--P\'olya class;
\item every $J_{\chi,t}^{d,n}$ is hyperbolic for all $d,n\ge0$.
\end{enumerate}
In particular,
\[
 \boxed{
 \text{GRH for }L(s,\chi)
 \iff J_{\chi,0}^{d,n}\text{ is hyperbolic for all }d,n.}
 \tag{J.24}
\]
\end{theorem}

\begin{proof}
The equivalence between the Laguerre--P\'olya property and hyperbolicity of all
Jensen packets is the classical Jensen--P\'olya characterization; see
\cite{CravenCsordas}.  A real entire function of order one with only real
zeros belongs to that class.  Apply \Cref{thm:theta-reality-GRH} at $t=0$.
\end{proof}

Thus every obstruction is represented by a finite polynomial built from
source-fixed theta moments; no zero is used to choose its coefficients.

\subsection{First collision or cofinal escape}

\begin{theorem}[Finite collision, degree degeneration, or cofinal escape]
\label{thm:theta-collision-or-escape}
Assume GRH fails for $L(s,\chi)$, and put
$t_*=\Lambda_{\rm DN}(\chi)>0$.  At least one of the following occurs:
\begin{enumerate}[label=\textup{(C\arabic*)},leftmargin=3.2em]
\item there are fixed $d\ge2$, $n\ge0$ such that at $t_*$ the packet
      $J_{\chi,t_*}^{d,n}$ either drops degree or has a multiple real root;
\item there are $t_j\uparrow t_*$ and nonhyperbolic packets
      $J_{\chi,t_j}^{d_j,n_j}$ with $d_j+n_j\to\infty$.
\end{enumerate}
\end{theorem}

\begin{proof}
For every $t<t_*$, at least one Jensen packet is nonhyperbolic by
\Cref{thm:theta-DN,thm:theta-Jensen-Polya}.  Choose a witness along
$t_j\uparrow t_*$.  If $d_j+n_j$ is unbounded, branch~\textup{(C2)} holds.
Otherwise pass to a fixed pair $(d,n)$.  At $t_*$ the packet is hyperbolic.
If its degree were stable and all roots simple, continuity of polynomial
roots under coefficient perturbation would make it hyperbolic in a
neighbourhood below $t_*$, a contradiction.  Hence branch~\textup{(C1)}
holds.
\end{proof}

At a double collision define
\[
 \mathfrak K_{\chi,t}^{d,n}(y)
 =J_{\chi,t}^{d,n+2}(y)
  J_{\chi,t}^{d-2,n+2}(y),
 \qquad d\ge2.
 \tag{J.25}
\]

\begin{theorem}[Necessary signed curvature at a generic first collision]
\label{thm:theta-collision-sign}
Let $P(t,Y)=J_{\chi,t}^{d,n}(Y)$.  Suppose $P(t,\cdot)$ is hyperbolic for
$t_0<t\le t_0+\delta$ and has a nonreal conjugate pair for
$t_0-\delta\le t<t_0$.  Assume the transition occurs through one double root
$y_0$ with
\[
 P(t_0,y_0)=P_Y(t_0,y_0)=0,
 \qquad P_{YY}(t_0,y_0)\ne0,
 \qquad P_t(t_0,y_0)\ne0.
 \tag{J.26}
\]
Then
\[
 \boxed{\mathfrak K_{\chi,t_0}^{d,n}(y_0)>0.}
 \tag{J.27}
\]
Moreover,
\[
 J_{\chi,t_0}^{d-1,n}(y_0)
 =J_{\chi,t_0}^{d-1,n+1}(y_0)=0.
 \tag{J.28}
\]
\end{theorem}

\begin{proof}
Write $\tau=t-t_0$ and $v=Y-y_0$.  Taylor expansion gives
\[
 P(t,Y)=P_t(t_0,y_0)\tau
 +\frac12P_{YY}(t_0,y_0)v^2+o(|\tau|+v^2).
\]
For $\tau>0$ the two nearby roots are real, so
$-P_t/P_{YY}>0$.  By \Cref{thm:theta-Jensen-recurrences},
\[
 P_t=-J_{\chi,t}^{d,n+2},
 \qquad
 P_{YY}=d(d-1)J_{\chi,t}^{d-2,n+2}.
\]
The two factors in \textup{(J.25)} therefore have the same nonzero sign,
which proves \textup{(J.27)}.  Equation \textup{(J.21)} gives the second
zero in \textup{(J.28)}, and \textup{(J.22)} gives the first.
\end{proof}

\begin{criterion}[Completed theta--Jensen GRH closure]
\label{crit:theta-Jensen-GRH}
Suppose that for a primitive nonprincipal $\chi$ the following hold at every
positive heat boundary:
\begin{enumerate}[label=\textup{(T\arabic*)},leftmargin=3.2em]
\item the cofinal-escape branch in
      \Cref{thm:theta-collision-or-escape} does not occur;
\item every fixed first collision has stable degree and is a generic double
      collision of the form \textup{(J.26)};
\item at every such common real root,
      $\mathfrak K_{\chi,t}^{d,n}(y)\le0$.
\end{enumerate}
Then $\Lambda_{\rm DN}(\chi)=0$ and GRH holds for $L(s,\chi)$.
\end{criterion}

\begin{proof}
If GRH failed, \Cref{thm:theta-collision-or-escape} and assumptions
\textup{(T1)}--\textup{(T2)} would produce a generic fixed collision at
positive heat time.  \Cref{thm:theta-collision-sign} would force positive
curvature, contradicting \textup{(T3)}.  Thus the de Bruijn--Newman constant
cannot be positive; \Cref{thm:theta-DN} then gives the conclusion.
\end{proof}

The first nontrivial packet is the heat--Tur\'an current
\[
 \mathfrak T_{\chi,n}(t)
 =\gamma_{\chi,n+1}(t)^2
  -\gamma_{\chi,n}(t)\gamma_{\chi,n+2}(t).
 \tag{J.29}
\]

\begin{proposition}[Quadratic Jensen current]
\label{prop:theta-quadratic-Jensen}
The polynomial $J_{\chi,t}^{2,n}$ is hyperbolic if and only if
$\mathfrak T_{\chi,n}(t)\ge0$, and
\[
 \boxed{
 \mathfrak T_{\chi,n}'(t)
 =\gamma_{\chi,n+2}(t)^2
  +\gamma_{\chi,n}(t)\gamma_{\chi,n+4}(t)
  -2\gamma_{\chi,n+1}(t)\gamma_{\chi,n+3}(t).}
 \tag{J.30}
\]
At a simple zero of this determinant with $\gamma_{\chi,n+2}\ne0$, the
signed derivative is exactly the degree-two collision curvature after
substitution of the double root.
\end{proposition}

\begin{proof}
The discriminant of $J_{\chi,t}^{2,n}$ is
$4\mathfrak T_{\chi,n}(t)$.  Differentiate \textup{(J.29)} and use
\textup{(J.19)}.  At a zero, the double root is
$-\gamma_{\chi,n+1}/\gamma_{\chi,n+2}$; direct substitution identifies
\textup{(J.30)} with the corresponding factor in \textup{(J.25)}.
\end{proof}

The theta route, the frequency-resolved rectangular route, and the completed
Weil--Schur route below are complementary.  The first is a cutoff-independent
entire source whose unresolved issue is compactness of finite Jensen witnesses
and signed collision exclusion.  The second is a cutoff-dependent arithmetic
descent whose remaining source separates into a reciprocal-phase bank and a
phase-free sinc sheet.  The third is a localized spectral positivity problem
with a coercive high-mode complement and a finite protected low-mode
obstruction.  None is a coordinate repackaging of the Chebyshev
short-quotient packet.

% =============================================================================
% =============================================================================
\section{Finite multiplier blocks, heat collisions, and the padding boundary}
\label{sec:finite-heat-blocks}
% =============================================================================

The Jensen entrance admits a second, finite-scale analysis.  Multiplicative
coefficient blocks have an exact backward-heat limit, and the root dynamics
of that limit constrain how much imaginary height a prescribed collection
of real factors can remove.  This gives necessary conditions on a
factorwise certification architecture.  Such conditions must not be
confused with exclusion of a zero of the completed $L$-function.

Throughout this section $D$ denotes polynomial differentiation,
$\mathrm{He}_r=e^{-D^2/2}z^r$ is the monic probabilists' Hermite polynomial,
and
\begin{equation}
 \mathcal H_K P=e^{-KD^2/2}P
 =\sum_{j\ge0}\frac{(-K/2)^j}{j!}D^{2j}P.
 \label{eq:heatblock-operator}
\end{equation}
The sum terminates for a polynomial.  This parameterization has heat
velocity $-D^2/2$; a direct comparison with the theta heat time in
\Cref{sec:grh-theta-heat} uses $K=2t$.

\subsection{Finite multiplier tests and the fixed-block heat limit}

A real vector $(\mu_0,\ldots,\mu_M)$ acts diagonally on polynomials of
degree at most $M$ by multiplying their coefficients.  The finite
Grace--Szeg\H{o} characterization, in the form of the finite
multiplier criterion, tests this operation on
$\sum_{j=0}^M\binom Mj\mu_jq^j$: it preserves real-rootedness when this
polynomial has all zeros real and of one sign.  For the positive vectors
used below, that sign is negative.  This is the finite counterpart of
the Jensen criterion in \Cref{thm:theta-Jensen-Polya}; see
\cite{CravenCsordasFinite}.

\begin{lemma}[Linear and conjugate-quadratic coefficient blocks]
\label{lem:heatblock-quadratic-test}
For $u>0$, the vector $j+u$, $0\le j\le M$, is a finite multiplier
sequence.  For $M\ge2$, the positive vector $(j+u)^2+v^2$ is a finite
multiplier sequence exactly when
\begin{equation}
 M(1+4u-4v^2)+4(u^2+v^2)\ge0.
 \label{eq:heatblock-quadratic-discriminant}
\end{equation}
In particular $v^2\le u+1/4$ is sufficient uniformly in $M$.
\end{lemma}

\begin{proof}
The linear Jensen polynomial is
$(1+q)^{M-1}[u+(u+M)q]$.  For the quadratic vector it is
$(1+q)^{M-2}Q(q)$, where
\[
 Q(q)=(u^2+v^2)
 +\{M(1+2u)+2(u^2+v^2)\}q
 +\{(M+u)^2+v^2\}q^2.
\]
Its discriminant is $M$ times the left side of
\eqref{eq:heatblock-quadratic-discriminant}.  All coefficients are
positive, so its real roots, whenever they exist, are negative.
Apply the finite multiplier criterion.
\end{proof}

\begin{theorem}[Universal fixed-block Jensen limit]
\label{thm:heatblock-jensen-limit}
Fix an integer $r\ge1$ and $s>0$.  Suppose
$u_{\nu,M}=sM+\alpha_\nu\sqrt M+o(\sqrt M)$ for $1\le\nu\le r$,
with the multiset of factors invariant under conjugation.  Define
\begin{equation}
 \begin{split}
 \pi_M(j)&=\prod_{\nu=1}^r(j+u_{\nu,M}),\\
 \sum_{j=0}^M\binom Mj\pi_M(j)q^j
 &=(1+q)^{M-r}Q_M(q).
 \end{split}
 \label{eq:heatblock-finite-factor}
\end{equation}
Put $S=1+s$, $K=s(1+s)$, and
$\Pi_{\boldsymbol\alpha}(z)=\prod_\nu(z+\alpha_\nu)$.  Then,
coefficientwise and locally uniformly in $z$,
\begin{equation}
 \frac{S^r}{M^{r/2}}
 Q_M\!\left(-\frac{s}{S}+\frac{z}{S^2\sqrt M}\right)
 \longrightarrow \mathcal H_K\Pi_{\boldsymbol\alpha}(z).
 \label{eq:heatblock-jensen-limit}
\end{equation}
If the limit has simple real roots, the residual roots of $Q_M$ are
real and negative for all sufficiently large $M$.  A nonreal root of
the limit prevents eventual real-rootedness.  A multiple real root in
the limit alone decides neither conclusion.
\end{theorem}

\begin{proof}
After $k$ factors, write the Jensen polynomial as
$(1+q)^{M-k}Q_{k,M}(q)$.  Multiplication of the coefficient vector by
$j+u_{k+1,M}$ is $qD_q+u_{k+1,M}$, hence
\begin{equation}
 Q_{k+1,M}=
 \{u_{k+1,M}+(M-k+u_{k+1,M})q\}Q_{k,M}
 +q(1+q)Q_{k,M}'.
 \label{eq:heatblock-exact-recurrence}
\end{equation}
Use the affine change of variable in \eqref{eq:heatblock-jensen-limit}
and set $R_{k,M}=S^kM^{-k/2}Q_{k,M}$.  Since $k$ is fixed, the scaled
recurrence converges coefficientwise to
$R_{k+1}=(z+\alpha_{k+1})R_k-KR_k'$, with $R_0=1$.
The conjugation identity
$\mathcal H_K(z+\alpha)=(z+\alpha-KD)\mathcal H_K$
identifies its solution as $\mathcal H_K\Pi_{\boldsymbol\alpha}$.
Simple real roots of a real polynomial persist as real roots under
small real coefficient perturbations; their original $q$ coordinates
converge to $-s/S<0$.  A small disk around any nonreal limiting root
can be kept disjoint from the real axis, and root continuity supplies
nonreal roots of $Q_M$ inside it.  The final assertion is the absence
of either stability argument at a multiple real root.
\end{proof}

\subsection{Exact root dynamics and the total imaginary-height budget}

\begin{lemma}[Real roots and strip preservation]
\label{lem:heatblock-preservation}
For $K\ge0$, $\mathcal H_K$ preserves real-rootedness and every closed
horizontal root strip.  More generally it cannot decrease the number
of real roots of a real polynomial, counted with multiplicity.
\end{lemma}

\begin{proof}
For real $a$, the operator $I+aD$ preserves a horizontal strip.  Indeed,
outside the strip the imaginary part of $P'/P=\sum_j(z-z_j)^{-1}$ has
strictly one sign and cannot equal the real number $-1/a$.
For a real polynomial, Rolle's theorem applied between consecutive
real zeros gives the corresponding real-zero count for $P+aP'$;
there is also an exterior real zero, unless already accounted for by
multiplicity.  Thus the count is not decreased.  Now use
\[
 \mathcal H_KP=
 \lim_{N\to\infty}
 \left(1-\frac{K}{2N}D^2\right)^NP,
\]
factor each bracket into $I\pm\sqrt{K/(2N)}D$, and use coefficient
convergence.  The degree and leading coefficient are fixed, so real
roots cannot escape to infinity in this limit.
\end{proof}

\begin{theorem}[Root-flow identities and the total-height bound]
\label{thm:heatblock-total-height}
Let a monic real polynomial have $m\ge1$ nonreal conjugate pairs
$x_j\pm iy_j$ and $p$ real zeros, counted with multiplicity, and put
$n=2m+p$.  If $\mathcal H_KP$ is real-rooted, then
\begin{equation}
 \sum_{j=1}^m y_j^2\le m(2m+2p-1)K.
 \label{eq:heatblock-total-height}
\end{equation}
At simple roots $z_k=x_k+iy_k$ of $\mathcal H_\tau P$ one has
\begin{equation}
 \dot z_k=\sum_{\ell\ne k}\frac1{z_k-z_\ell},
 \label{eq:heatblock-root-flow}
\end{equation}
and the exact identities
\begin{equation}
 \begin{split}
 \frac{\dd}{\dd\tau}\sum_kx_k^2
 &=\sum_{k\ne\ell}\frac{(x_k-x_\ell)^2}{|z_k-z_\ell|^2},\\
 -\frac{\dd}{\dd\tau}\sum_ky_k^2
 &=\sum_{k\ne\ell}\frac{(y_k-y_\ell)^2}{|z_k-z_\ell|^2}.
 \end{split}
 \label{eq:heatblock-two-energies}
\end{equation}
The two displayed rates sum to $n(n-1)$.
\end{theorem}

\begin{proof}
Differentiating the root equation and using
$\partial_\tau P_\tau=-P_\tau''/2$ gives
$\dot z_k=P_\tau''(z_k)/(2P_\tau'(z_k))$, which is
\eqref{eq:heatblock-root-flow}.  Pair the ordered terms $k,\ell$ in
the derivatives of the two sums of squares to obtain
\eqref{eq:heatblock-two-energies}.
At least $p$ roots remain real by
\Cref{lem:heatblock-preservation}, so at least $p(p-1)$ ordered terms
in the imaginary-height rate vanish; every other term is at most one.
The total initial imaginary square is $2\sum_jy_j^2$ and the final
one is zero.  Integration gives
$2\sum_jy_j^2\le[n(n-1)-p(p-1)]K$.
The identities extend across collisions by continuity: the discriminant
is polynomial in $\tau$ and is not identically zero, since at large
$\tau$ the polynomial, after scaling, converges to $\mathrm{He}_n$
with simple real roots.  There are therefore only finitely many
collision times, and the displayed rates are bounded between them.
\end{proof}

\begin{theorem}[The shallowest-pair first-collision bound]
\label{thm:heatblock-shallow}
Under the same hypotheses,
\begin{equation}
 \min_j y_j^2\le(m+2p)K.
 \label{eq:heatblock-shallow}
\end{equation}
In particular, when $p\le m$,
$\min_jy_j^2\le(2m+p)K$.
\end{theorem}

\begin{proof}
Before the first pair reaches the real axis, track the least positive
height $y$.  Its conjugate contributes one to $-\dd(y^2)/\dd\tau$;
each real zero contributes at most two.  For any other pair at height
$b\ge y$ and horizontal separation $a$, its contribution is
\[
 2y\left\{\frac{y-b}{a^2+(y-b)^2}
          +\frac{y+b}{a^2+(y+b)^2}\right\}\le1.
\]
To verify the inequality put $A=a^2/y^2$, $B=b^2/y^2\ge1$; after
clearing the positive denominator it is
$(A+B-1)^2+4(B-1)\ge0$.
Thus the rate is at most $1+(m-1)+2p$.  At ties the same bound applies
to the one-sided derivatives of the minimum; isolated collisions are
handled by continuity as above.  If the first real collision occurs at
$T\le K$, integration from the initial minimum to zero proves
\eqref{eq:heatblock-shallow}.
\end{proof}

For three pairs and no real padding, \eqref{eq:heatblock-total-height}
gives $\sum_jy_j^2\le15K$.  The constant is attained: at $K=1$ the
inverse heat image of $z^6$ is
$z^6+15z^4+45z^2+15$, whose zeros form three purely imaginary pairs
and whose squared heights sum to $15$.  This exact total budget is not
a sharp classification of the smallest height, and it does not by
itself settle a mixed configuration with many real zeros.

\subsection{The sharp one-pair theorem with arbitrary real padding}

The shallow-pair bound gives $y^2\le(2p+1)K$ for one pair.  The complete
collision functional improves this to the degree bound, uniformly in
all positions of the real factors.

\begin{lemma}[A Hermite half-plane collision test]
\label{lem:heatblock-onepair-stability}
Fix $p\ge1$, $c\in\R$, $u\ge0$, and put $t=p+2+u$,
$H=\mathrm{He}_p$.  Define
\begin{equation}
 \begin{split}
 A&=(c^2+u+1)H+(z-2c)H',\\
 B&=(pz-2c(p+1))H-\bigl((z-c)^2+u\bigr)H'.
 \end{split}
 \label{eq:heatblock-AB}
\end{equation}
Then $A-iB$ has no zero in the open upper half-plane.  Its only
possible real zero occurs when $u=c=z=0$ and $p$ is odd.
\end{lemma}

\begin{proof}
For $z=x+iy$ with $y>0$, put $h=H'/H$.
The real simple zeros of $H$ and Cauchy--Schwarz imply
\begin{equation}
 -p\Im h-y|h|^2\ge0.
 \label{eq:heatblock-pick-bound}
\end{equation}
A zero of $A-iB$ would imply $\mathcal D h=\mathcal N$, where
\[
 \mathcal D=(z-2c)+i\bigl((z-c)^2+u\bigr),\qquad
 \mathcal N=-(c^2+u+1)+i(pz-2c(p+1)).
\]
Since $\Re\mathcal N=-(c^2+u+1)-py<0$, $\mathcal D$ cannot vanish.
Set $X=x-c$ and $M=c^2+p+u+1$.  Direct expansion yields
\begin{equation}
 |\mathcal D|^2(-p\Im h-y|h|^2)
 =-pMX^2+2cp(p+1)X-\mathcal C,
 \label{eq:heatblock-pick-contradiction}
\end{equation}
where
\begin{align*}
 \mathcal C={}&c^4y
 +c^2\{p^2y+p^2+pu+py^2+5py+2p+2uy+6y\}\\
 &+p^2uy+p^2y^2+pu^2+puy^2+puy+pu\\
 &+py^2+py+u^2y+2uy+y.
\end{align*}
In particular $\mathcal C\ge c^2p(p+2)$, and $\mathcal C>0$ when
$c=0$.  For $c\ne0$, completing the square bounds the right side by
\[
 -\frac{c^2p\{(p+2)c^2+p+1\}}{c^2+p+1}<0;
\]
for $c=0$ it is also strictly negative.  This contradicts
\eqref{eq:heatblock-pick-bound}.
For a real common zero of $A,B$ with $H\ne0$, elimination gives
\[
 MX^2-2c(p+1)X+c^2(p+u+2)+u(u+1)=0.
\]
Completing the square shows that its left side is positive unless
$c=u=X=0$; then $A=H\ne0$.  If $H=0$, its derivative is nonzero,
and \eqref{eq:heatblock-AB} first forces $z=2c$ and then $c=u=0$.
Finally $\mathrm{He}_p(0)=0$ exactly for odd $p$.
\end{proof}

\begin{theorem}[Arbitrary real padding of one conjugate pair]
\label{thm:heatblock-onepair}
Let
$P(z)=((z-a)^2+y^2)\prod_{j=1}^p(z-r_j)$, where $y>0$ and all
$a,r_j$ are real.  If $p\ge1$ and $\mathcal H_KP$ is real-rooted,
then
\begin{equation}
 y^2\le(p+2)K.
 \label{eq:heatblock-onepair}
\end{equation}
The inequality is strict for even $p$.  For every odd $p$ it is sharp.
For $p=0$ the exact condition is instead $y^2\le K$.
\end{theorem}

\begin{proof}
Let $T\le K$ be the first real collision of the nonreal pair.  Translate
the collision to zero and scale by $\sqrt T$.  A hypothetical
$y^2\ge(p+2)T$ gives an input
$((w-c)^2+t)\prod_j(w-r_j)$ with $t\ge p+2$.
The functional consisting of the heat value at zero minus $i$ times
its derivative there vanishes at a real collision.  As a function of
the real padding roots it is symmetric and multiaffine.  The
Grace--Walsh--Szeg\H{o} coincidence principle for the closed upper
half-plane therefore replaces those roots by a common
$\zeta$ in that half-plane without changing its zero value.
For the diagonal input $((w-c)^2+t)(w-\zeta)^p$, the functional,
up to the nonzero factor $(-1)^p$, is precisely $A(\zeta)-iB(\zeta)$
in \eqref{eq:heatblock-AB}.  This follows either by differentiating
the Hermite generating function or by using
$H''-zH'+pH=0$.  \Cref{lem:heatblock-onepair-stability} excludes
$t>p+2$, and excludes $t=p+2$ for even $p$.
For odd $p=2q+1$, the input
$P(z)=(z^2+p+2)z^p$ has heat image
\[
 \mathrm{He}_{p+2}(z)+(p+2)\mathrm{He}_p(z)
 =(-1)^q2^qq!\,z^3 L_q^{(3/2)}(z^2/2).
\]
The positive simple zeros of $L_q^{(3/2)}$ show that this image is
real-rooted.  Scaling proves sharpness for general $K$.
The unpadded quadratic is immediate from \eqref{eq:heatblock-operator}.
\end{proof}

The coincidence step retains the full conjugate pair and coalesces only
real factors.  Independent coalescence of upper and lower roots in a
multi-pair block changes the class of admissible inputs; it is not a
valid substitute for this argument.

\subsection{An exact quartic refinement}

For two real padding factors the degree bound is strict and can be
improved to an algebraic constant.

\begin{theorem}[Sharp quartic padding constant]
\label{thm:heatblock-quartic}
Let $P$ be a real quartic with one nonreal pair of imaginary height $y$
and two real roots, allowing the latter to coincide.  If
$\mathcal H_KP$ is real-rooted, then
\begin{equation}
 y^2\le c_\star K,
 \qquad
 4c_\star^3-21c_\star^2+24c_\star-8=0,
 \qquad 3<c_\star<4.
 \label{eq:heatblock-quartic-constant}
\end{equation}
Here $c_\star$ is the unique real zero of the displayed cubic
($c_\star\simeq3.814511867$), and equality is attained.
\end{theorem}

\begin{proof}
Normalize the first real collision to time one and location zero.
Write the input as
$\Pi(z)=((z-c)^2+t)(z^2+\alpha z+\beta)$,
where $t=y^2/T$ and $\alpha^2-4\beta\ge0$.
If $t\le3$ there is nothing to prove.  For $t>3$, put
$q=c^2+t$, $w=c^2$, $\Delta=4w+(q-1)(q-3)>0$.
The equations $(\mathcal H_1\Pi)(0)=(\mathcal H_1\Pi)'(0)=0$ give
\[
 \alpha=-\frac{4cq}{\Delta},\qquad
 \beta=\frac{12w+(q-3)^2}{\Delta},\qquad
 \alpha^2-4\beta=-\frac{4G(w,t)}{\Delta^2},
\]
where
\begin{align*}
 G(w,t)={}&w^4+(4t+2)w^3+6(t^2-t+2)w^2\\
 &+(4t^3-18t^2+18)w+(t-3)^3(t-1).
\end{align*}
Let $x>0$ be the unique root of $x^3+3x^2-8=0$, and set
\[
 c_\star=\frac{3x^2+12x+16}{x^2+3x+4},\qquad
 w_\star=\frac{x}{x^2+3x+4}.
\]
Polynomial division using the equation for $x$ gives
\[
 G(w,c_\star)=(w-w_\star)^2(w^2+A_\star w+B_\star),
\]
with
$A_\star=4c_\star+2+2w_\star>0$ and
$B_\star=6(c_\star^2-c_\star+2)+w_\star(8c_\star+4+3w_\star)>0$.
Moreover
\[
 \partial_tG=4w^3+(12t-6)w^2+12t(t-3)w
             +2(t-3)^2(2t-3)>0
 \quad(t>3,\ w\ge0).
\]
Thus $t>c_\star$ contradicts $\alpha^2-4\beta\ge0$.
For attainment choose
$r^2=3x/[(x+1)^2(x+4)]$, $a=xr$, and
$P(z)=(z+a)^2(z^2+c_\star)$.  Direct heat expansion gives
\[
 \mathcal H_1P(z)=(z-r)^3\bigl(z+(2x+3)r\bigr).
\]
Eliminating $x$ gives \eqref{eq:heatblock-quartic-constant}.
The discriminant of that cubic is negative, so it has one real root;
its values at $3$ and $4$ have opposite signs.
Finally $T\le K$ gives the asserted bound.
\end{proof}

\subsection{Real-axis exclusion for a general conjugate block}
\label{subsec:heatblock-appell}

Retain every conjugate factor and write
\begin{equation}
 \Pi_m(w)=\prod_{j=1}^m((w-c_j)^2+t_j),\qquad
 \widehat Q=\mathcal H_1\Pi_m,\qquad
 \mathcal P_r(x)=\widehat Q(D_x)\mathrm{He}_r(x).
 \label{eq:heatblock-appell}
\end{equation}
The notation $\mathcal P_r$ separates this Appell family from both the
Legendre basis and the original arithmetic Jensen packets.  Its exact
identities are
\begin{equation}
 \mathcal P_r'=r\mathcal P_{r-1},\qquad
 \sum_{r\ge0}\mathcal P_r(x)\frac{s^r}{r!}
 =e^{xs-s^2/2}\widehat Q(s).
 \label{eq:heatblock-appell-generating}
\end{equation}
Its leading coefficient is $\widehat Q(0)$, not generally one.

\begin{theorem}[The weakest real-axis collision criterion]
\label{thm:heatblock-real-axis}
Let
\[
 F_{m,p}(\zeta)=
 \left.\mathcal H_1[\Pi_m(w)(w-\zeta)^p]\right|_{w=0}
 -i\left.D_w\mathcal H_1[\Pi_m(w)(w-\zeta)^p]\right|_{w=0}.
\]
For real $x$, $F_{m,p}(x)=A(x)-iB(x)$ with
\begin{equation}
 A=(-1)^p\mathcal P_p,
 \qquad
 B=(-1)^p\bigl(\mathcal P_{p+1}-x\mathcal P_p\bigr).
 \label{eq:heatblock-real-axis-AB}
\end{equation}
The following are equivalent: $F_{m,p}$ has a real zero;
$\mathcal P_p$ and $\mathcal P_{p+1}$ share a real zero;
$\mathcal P_{p+1}$ has a real multiple zero.
\end{theorem}

\begin{proof}
Apply the Hermite generating function to
$\mathcal H_1[\Pi_m(w)(w-x)^p]$ at $w=0$.
Its value is $(-1)^p\widehat Q(D)\mathrm{He}_p$.
For the derivative use
$[\widehat Q(D),x]=\widehat Q'(D)$ and
$x\mathrm{He}_p=\mathrm{He}_{p+1}+p\mathrm{He}_{p-1}$;
the result is the second expression in
\eqref{eq:heatblock-real-axis-AB}.
Both expressions are real for real $x$.  Their simultaneous vanishing
is equivalent to $\mathcal P_p=\mathcal P_{p+1}=0$.
The derivative identity in
\eqref{eq:heatblock-appell-generating} proves the last equivalence.
\end{proof}

\begin{criterion}[Conjugation-preserving multi-pair exclusion]
\label{crit:heatblock-multipair}
Fix $m\ge2$, $p\ge0$, and $n=2m+p$.  Suppose that for every
$c_j\in\R$ and every $t_j\ge n$, the consecutive polynomials
$\mathcal P_p,\mathcal P_{p+1}$ in \eqref{eq:heatblock-appell}
have no common real zero.  Then every real polynomial with $m$ nonreal
pairs and $p$ real zeros whose heat image at time $K$ is real-rooted
satisfies
\begin{equation}
 \min_j y_j^2<nK.
 \label{eq:heatblock-conditional-band}
\end{equation}
\end{criterion}

\begin{proof}
The pure block $\widehat Q=\mathcal H_1\Pi_m$ has no real zero when
all $t_j\ge n>m$: the first-collision estimate in the proof of
\Cref{thm:heatblock-shallow}, with no real padding, prevents any pair
from reaching the axis by time one.  Hence $\widehat Q(0)\ne0$.
The leading coefficient of $F_{m,p}$ is therefore nonzero and its
degree stays $p$ on the connected parameter domain under consideration.
The hypothesis and \Cref{thm:heatblock-real-axis} prevent any of its
zeros from crossing the real axis.  Send all but one of the heights
to infinity, keeping their centers fixed, and divide by their product.
The collision polynomial tends to the one-pair polynomial, which has
no zero in the closed upper half-plane since $t_1\ge n>p+2$.
Consequently the upper-half-plane root count is zero everywhere on
the domain.  At a hypothetical first collision with
$\min_jy_j^2\ge nK$, translation and scaling give parameters in that
domain.  Coalescing only the real padding roots by the same coincidence
principle used in \Cref{thm:heatblock-onepair} contradicts this zero count.
\end{proof}

The criterion requires absence of a real common zero, not full
real-rootedness of the Appell sequence.  A nonzero resultant is a
stronger sufficient certificate, since it excludes common complex
zeros as well.  Positivity of an expanded determinant is stronger
still.  These sufficient conditions must not be identified with the
weak real-axis condition, and continuity at a symmetric configuration
does not establish the condition on the whole parameter domain.

\subsection{Pair couplings and coefficient determinants}
\label{subsec:heatblock-pairing}

The real-axis criterion can be made explicit without replacing conjugate
pairs by independent linear factors.  The relevant expansion records the
couplings created by the heat operator, rather than assigning a separate
positivity certificate to each pair.

\begin{proposition}[Exact pairing expansion]
\label{prop:heatblock-pairing}
Put $f_j(w)=(w-c_j)^2+t_j$ and $\widehat f_j=f_j-1$.  For integers
$a,b\ge0$ with $b$ even, define
\[
 \mathfrak m(a,b)=(-1)^{b/2}\E[Z^b(1-Z^2)^a],\qquad Z\sim N(0,1).
\]
Then
\begin{equation}
 \mathcal H_1\prod_{j=1}^m f_j
 =\sum_{\substack{A,B\subseteq\{1,\ldots,m\}\!,\ A\cap B=\varnothing\\
                   |B|\ {
m even}}}
 \mathfrak m(|A|,|B|)
 \prod_{j\notin A\cup B}\widehat f_j
 \prod_{j\in B}f_j'.
 \label{eq:heatblock-pairing}
\end{equation}
In particular, for two and three pairs,
\begin{align}
 \mathcal H_1(f_1f_2)
 &=\widehat f_1\widehat f_2-f_1'f_2'+2,
 \label{eq:heatblock-pairing-two}\\
 \mathcal H_1(f_1f_2f_3)
 &=\widehat f_1\widehat f_2\widehat f_3
   -\sum_{i<j}f_i'f_j'\widehat f_k
   +2\sum_k\widehat f_k+2\sum_{i<j}f_i'f_j'-8,
 \label{eq:heatblock-pairing-three}
\end{align}
where $k$ in the first sum of \eqref{eq:heatblock-pairing-three}
is the index complementary to $i,j$.
\end{proposition}

\begin{proof}
For every polynomial $F$, $\mathcal H_1F(w)=\E F(w+iZ)$, and
\[
 f_j(w+iZ)=\widehat f_j(w)+(1-Z^2)+iZf_j'(w).
\]
Expand the product, choosing its second term on $A$ and its third term
on $B$.  Odd Gaussian moments vanish.  The required low-order values are
$\mathfrak m(0,0)=1$, $\mathfrak m(1,0)=0$,
$\mathfrak m(2,0)=2$, $\mathfrak m(0,2)=-1$,
$\mathfrak m(1,2)=2$, and $\mathfrak m(3,0)=-8$.
\end{proof}

To state a determinant certificate with a fixed sign convention, set
\[
 A_p=\mathcal P_p,\qquad
 \widetilde B_p=x\mathcal P_p-\mathcal P_{p+1}.
\]
The sign change in the second expression does not change the common-zero
condition in \Cref{thm:heatblock-real-axis}.  When $p\ge2m+1$, expand
both polynomials in the Hermite functions of indices
$p-2m,\ldots,p+1$, and use
$\mathrm{He}_{r+1}=x\mathrm{He}_r-r\mathrm{He}_{r-1}$ to reduce them to
$U=\mathrm{He}_{p-2}$ and $V=\mathrm{He}_{p-3}$.  This specifies
rational-polynomial coefficients $A_U,A_V,B_U,B_V$ and the determinant
\begin{equation}
 \Delta_{m,p}(x)=A_U B_V-A_V B_U.
 \label{eq:heatblock-determinant-convention}
\end{equation}
Consecutive Hermite polynomials have no common real zero, as follows
by iterating their three-term recurrence to $\mathrm{He}_0=1$.
Consequently $\Delta_{m,p}(x)\ne0$ is a sufficient certificate against
a common zero at $x$.  Its vanishing is not a converse certificate.

\begin{proposition}[One-pair determinant and its exterior region]
\label{prop:heatblock-onepair-determinant}
For $m=1$, $p\ge3$, $c\in\R$, and $t>0$,
\begin{equation}
 \Delta_{1,p}(x)=p(p-1)(p-2)\Psi_1(x;c,t),
 \label{eq:heatblock-onepair-determinant}
\end{equation}
where
\begin{align}
 \Psi_1={}&(c^2+t-1)((x-c)^2+t)
        -2(p+1)\{c(x-c)+t\}+(p+1)(p+2).
 \label{eq:heatblock-onepair-psi}
\end{align}
If $c^2+t>1$ and $\rho=c^2+t-p-2$, then $\Psi_1$ is strictly
positive for every real $x$ precisely when
\begin{equation}
 \rho(t\rho+p+1)>0.
 \label{eq:heatblock-onepair-exact-region}
\end{equation}
In particular $c^2+t>p+2$ is a sufficient exterior condition.
On the closed height region $t\ge p+2$, the only possible zero of
$\Psi_1$ is $x=c=0$, $t=p+2$.
\end{proposition}

\begin{proof}
Here $\widehat Q(w)=w^2-2cw+c^2+t-1$, so
\[
 \mathcal P_p=(c^2+t-1)\mathrm{He}_p
 -2cp\mathrm{He}_{p-1}+p(p-1)\mathrm{He}_{p-2}.
\]
Reduction by the declared recurrence proves
\eqref{eq:heatblock-onepair-determinant}.  Write
$\Psi_1=Ax^2-Bx+C$.  Direct expansion gives
\[
 A=c^2+t-1,\qquad B=2c(c^2+t+p),\qquad
 4AC-B^2=4\rho(t\rho+p+1).
\]
The minimum of this quadratic is $(4AC-B^2)/(4A)$ when $A>0$.
If $t\ge p+2$, then $\rho\ge c^2\ge0$; equality can occur only
at $c=0,t=p+2$, where $\Psi_1=(p+1)x^2$.
\end{proof}

\begin{obstruction}[The one-pair modulus condition is not multiplicative]
\label{obs:heatblock-modulus-determinant}
A positive one-pair determinant does not imply a positive determinant
for a product of pairs.  For $m=2$, $p=10$,
\[
 c_1=c_2=\frac72,\qquad t_1=t_2=2,\qquad x=8,
\]
one has $c_j^2+t_j=57/4>2m+p$, but
\begin{equation}
 \Delta_{2,10}(8)=-\frac{968715}{16}<0.
 \label{eq:heatblock-modulus-witness}
\end{equation}
This disproves that proposed determinant-positivity extension, not the
weaker common-real-zero exclusion in \Cref{crit:heatblock-multipair}.
\end{obstruction}

\begin{proof}
Substitute $f_1=f_2=(w-7/2)^2+2$ in
\eqref{eq:heatblock-pairing-two}, form $\widehat Q(D)\mathrm{He}_{10}$
and $\widehat Q(D)\mathrm{He}_{11}$, and use the recurrence defining
\eqref{eq:heatblock-determinant-convention}.  The resulting determinant
at $x=8$ is \eqref{eq:heatblock-modulus-witness}.  Equivalently, after
division by $10\cdot9\cdot8=720$, its value is $-21527/256$.
The modulus hypothesis is different from the height hypothesis
$t_j\ge2m+p$ in \Cref{crit:heatblock-multipair}.
\end{proof}

\begin{criterion}[Certificate interface for the complete two-pair band]
\label{crit:heatblock-two-pair-certificates}
For each $p\ge3$, let $\mathcal P_r$ be constructed from two pairs as
in \eqref{eq:heatblock-appell}.  Suppose that throughout
$c_1,c_2\in\R$, $t_1,t_2\ge p+4$, one has
\[
 \begin{cases}
 \operatorname{Res}_x(\mathcal P_p,\mathcal P_{p+1})\ne0,
       &p=3,4,\\
 \Delta_{2,p}(x)\ne0\quad\text{for every }x\in\R,
       &p\ge5.
 \end{cases}
\]
Then every real polynomial with two nonreal pairs, $p\ge0$ real
zeros, and real-rooted heat image at time $K$ satisfies
\[
 \min(y_1^2,y_2^2)\le(p+4)K,
\]
strictly for $p\ge3$.  Under the ladder scaling
\eqref{eq:heatblock-gamma-scaling}, it follows that
\[
 \varkappa^2\ge
 \frac{R\min(\gamma_1^2,\gamma_2^2)}{(p+4)(1+\theta)}
 \ge\frac{\min(\gamma_1^2,\gamma_2^2)}{1+\theta}
 \quad\text{when }R\ge p+4.
\]
\end{criterion}

\begin{proof}
For $p\le2$, \Cref{thm:heatblock-shallow} gives
$\min y_j^2\le(2+2p)K\le(p+4)K$.  For $p=3,4$, the nonzero
resultant excludes every common zero.  For $p\ge5$, the determinant
and the nonvanishing Hermite basis vector exclude a common real zero.
Apply \Cref{crit:heatblock-multipair} and then substitute the ladder
scaling.
\end{proof}

The two families of algebraic certificates are explicit: their entries
are determined by \eqref{eq:heatblock-pairing-two} and the Hermite
recurrence.  A rational sum-of-squares or positive-semidefinite Gram
decomposition can verify them, but its complete coefficient identities
are required.  The reduction above does not infer a cone-wide
inequality from a positive sample or from the description of a
computational search.  In particular its arbitrary-padding conclusion
is retained with its certificate hypotheses.


\subsection{Signed Gaussian moments and the center-dependent symbol}

\begin{theorem}[Exact signed-moment and Fourier representations]
\label{thm:heatblock-signed-symbol}
For a standard normal random variable $Z$,
\begin{equation}
 \mathcal P_r(x)=\E[\Pi_m(iZ)(x-iZ)^r].
 \label{eq:heatblock-signed-moments}
\end{equation}
Let $\dd\nu(s)=\Pi_m(is)e^{-s^2/2}\dd s/\sqrt{2\pi}$ and use the
Fourier convention
$\widehat\nu(\xi)=\int e^{i\xi s}\dd\nu(s)$.
Then
\begin{equation}
 \begin{split}
 \widehat\nu(\xi)
 &=\left[\prod_{j=1}^m\{t_j+(D_\xi-c_j)^2\}\right]e^{-\xi^2/2}
   =e^{-\xi^2/2}G_{\boldsymbol c,\boldsymbol t}(\xi),\\
 G_{\boldsymbol c,\boldsymbol t}(\xi)&=\widehat Q(-\xi),\\
 \sum_{r\ge0}\mathcal P_r(x)\frac{s^r}{r!}
 &=e^{xs}\widehat\nu(-s)
   =e^{xs-s^2/2}G_{\boldsymbol c,\boldsymbol t}(-s).
 \end{split}
 \label{eq:heatblock-correct-fourier}
\end{equation}
The polynomial recurrence for $G$ consists of applying the mutually
commuting factors $t_j+(D_\xi-\xi-c_j)^2$ to $1$.
For centered equal heights,
\begin{equation}
 G_{\boldsymbol0,t}(\xi)
 =\sum_{k=0}^m\binom mk t^{m-k}\mathrm{He}_{2k}(\xi).
 \label{eq:heatblock-equal-height}
\end{equation}
\end{theorem}

\begin{proof}
The elementary Gaussian shift identity gives
$\E[\Pi_m(iZ)e^{-isZ}]=e^{-s^2/2}(\mathcal H_1\Pi_m)(s)$.
Multiplication by $e^{xs}$ and coefficient extraction proves
\eqref{eq:heatblock-signed-moments} and the last line of
\eqref{eq:heatblock-correct-fourier}.
Alternatively each factor $iZ-c_j$ corresponds under the declared
Fourier transform to $D_\xi-c_j$.  Conjugating by the Gaussian gives
$D_\xi-\xi-c_j$, proving the operator formula and
$G(\xi)=\widehat Q(-\xi)$.
When all centers vanish, $(D_\xi-\xi)^2$ raises
$\mathrm{He}_k$ to $\mathrm{He}_{k+2}$.  The binomial theorem proves
\eqref{eq:heatblock-equal-height}.
\end{proof}

The Fourier argument in the generating function is $-s$, with exactly
the convention in \eqref{eq:heatblock-correct-fourier}; it is not $is$.
The latter substitution would change the Gaussian sign and, already
for one centered pair, replace the correct symbol $s^2+t-1$ by
$t-1-s^2$.  Root-location conclusions based on that replacement do not
apply to the Appell family defined in \eqref{eq:heatblock-appell}.

For centered pairs $\nu$ is the real signed measure with density
$\prod_j(t_j-s^2)e^{-s^2/2}/\sqrt{2\pi}$; for nonzero centers it is
in general a complex measure with conjugate symmetry.  In the centered
case one has the exact even-moment expansion
\begin{equation}
 \mathcal P_r(x)=\sum_{2\ell\le r}\binom r{2\ell}
 (-1)^\ell\mu_{2\ell}x^{r-2\ell},\qquad
 \mu_{2\ell}=\int s^{2\ell}\dd\nu(s).
 \label{eq:heatblock-even-moments}
\end{equation}
Neither the sign-changing tail nor the degree dependence of these
moments can be discarded by a positivity assertion about the core.

\begin{obstruction}[Positive symmetric weights do not exclude double roots]
\label{obs:heatblock-positive-weight}
Positivity and symmetry of a measure alone do not imply the real-axis
condition in \Cref{thm:heatblock-real-axis}.  Indeed, the positive
probability measure
\[
 \nu_0=\frac89\delta_0+\frac1{18}(\delta_3+\delta_{-3})
\]
satisfies
\[
 \int(x-is)^4\dd\nu_0(s)=x^4-6x^2+9=(x^2-3)^2.
\]
Thus even this positive symmetric transform has real double roots.
\end{obstruction}

\begin{proof}
The second and fourth moments of $\nu_0$ are $1$ and $9$,
respectively, and its odd moments vanish.  Binomial expansion gives
the identity.
\end{proof}

For three pairs, \Cref{thm:heatblock-shallow} supplies the degree bound
when $p\le3$.  With more real padding, the weak common-real-zero
condition in \Cref{crit:heatblock-multipair} remains an independent
certificate.  In particular the ranges $p=4,5,6$ and a uniform
certificate for $p\ge7$ are not settled by the symbol identity.
Centered equal-height formulas, numerical separation at a symmetric
point, or a compactness bound for an unpadded configuration do not
establish uniform separation for a padded family with free centers.

\subsection{Critical points as heat-evolved real zeros}
\label{subsec:heatblock-critical-zero}

For the translated configuration relevant to gamma ladders, it is
convenient to write the centers with the opposite sign.  Let
$c_1,\ldots,c_M>0$, $t>0$, and define
\begin{equation}
 \begin{split}
 \Pi_M^+(\xi)&=\prod_{j=1}^M((\xi+c_j)^2+t),\qquad
 \mathscr G_M=\mathcal H_1\Pi_M^+,\\
 \mathscr F_M(\xi)&=e^{-\xi^2/2}\mathscr G_M(\xi),\qquad
 \mathscr R_M=\mathcal H_1[\xi\Pi_M^+(\xi)].
 \end{split}
 \label{eq:heatblock-critical-symbol}
\end{equation}
These are Gaussian symbols, not yet the finite arithmetic filters
constructed in \Cref{sec:admissible-spectral-filters}.

\begin{proposition}[Critical-point identity]
\label{prop:heatblock-critical-zero}
For every $M\ge1$,
\begin{equation}
 \mathscr F_M'(\xi)=-e^{-\xi^2/2}\mathscr R_M(\xi).
 \label{eq:heatblock-critical-identity}
\end{equation}
Thus the last real critical point of $\mathscr F_M$ is exactly the
largest real zero of the odd-degree polynomial $\mathscr R_M$.
Writing $\Pi_M^+(\xi)=\sum_{k=0}^{2M}g_k\xi^k$, one has
\[
 \mathscr R_M(\xi)=\sum_{k=0}^{2M}g_k\mathrm{He}_{k+1}(\xi).
\]
For one pair,
\[
 \mathscr R_1(\xi)=\xi^3+2c_1\xi^2+(c_1^2+t-3)\xi-2c_1,
\]
so a positive real critical point always exists.  For two pairs the
same conclusion holds whenever $c_1c_2+t>3$.
\end{proposition}

\begin{proof}
The heat commutator is
$\mathcal H_1(\xi P)=\xi\mathcal H_1P-(\mathcal H_1P)'$.
Differentiating \eqref{eq:heatblock-critical-symbol} proves
\eqref{eq:heatblock-critical-identity}; applying the heat operator to
monomials proves the Hermite expansion.  In the two-pair case,
\[
 \mathscr R_2(0)=-2(c_1+c_2)(c_1c_2+t-3).
\]
The polynomial is monic of odd degree and therefore tends to $+\infty$
on the positive axis.  Its negative value at zero gives a positive root.
\end{proof}

The initial real root at zero of $\xi\Pi_M^+(\xi)$ has positive
velocity $\sum_j2c_j/(c_j^2+t)$ under backward heat evolution.
This local velocity does not prove monotonicity throughout the flow:
conjugate roots can change their position relative to the tracked real
root.  The polynomial identity above avoids that unproved monotonicity.

\begin{proposition}[A finite large-height criterion]
\label{prop:heatblock-large-height}
Set $T_j=c_j^2+t$, $S=\sum_jT_j^{-1}$,
$c_- =\min_jc_j$, and $c_+=\max_jc_j$.  Suppose $S<1/2$, and put
\[
 \varepsilon_1=\E e^{SZ^2+2c_+S|Z|}-1,\qquad
 \varepsilon_2=\E\{|Z|(e^{SZ^2+2c_+S|Z|}-1)\}.
\]
If $c_-(1-\varepsilon_1)>\varepsilon_2$, then
$\mathscr G_M'(0)>0$ and $\mathscr F_M$ has a positive real
critical point.
\end{proposition}

\begin{proof}
Write
$w_j=(-Z^2+2ic_jZ)/T_j$.  For any omitted factor $j$,
\[
 \left|\prod_{\ell\ne j}(1+w_\ell)-1\right|
 \le e^{SZ^2+2c_+S|Z|}-1.
\]
The Gaussian representation and product differentiation give
\[
 \mathscr G_M'(0)=2\prod_jT_j\,
 \Re\sum_j\frac1{T_j}
 \E\left[(c_j+iZ)\prod_{\ell\ne j}(1+w_\ell)\right].
\]
The real part of each expectation is at least
$c_j(1-\varepsilon_1)-\varepsilon_2>0$.
Now $\mathscr R_M(0)=-\mathscr G_M'(0)<0$, and
\Cref{prop:heatblock-critical-zero} applies.
\end{proof}

The criterion is finite and quantitative.  It does not furnish a
uniform estimate when both the number of pairs and their largest
center grow.  For equispaced centers that regime has a separate
translation law.

\subsection{Growing gamma symbols and logarithmic translation}
\label{subsec:heatblock-growing-symbol}

Fix $s>0$, $t>0$, and $c_0>0$, and in
\eqref{eq:heatblock-critical-symbol} use the $M$ centers
$c_0,c_0+2s,\ldots,c_0+2s(M-1)$.  Define
\begin{equation}
 \begin{split}
 a_\pm(w)&=\frac{w+c_0\pm i\sqrt t}{2s},\qquad
 h(w)=\frac1{\Gamma(a_+(w))\Gamma(a_-(w))},\\
 A_M&=\frac{\log M}{s},\qquad
 C_M=(2s)^{2M}\Gamma(M)^2M^{c_0/s},\\
 q_M(w)&=C_M^{-1}e^{-A_Mw}\Pi_M^+(w).
 \end{split}
 \label{eq:heatblock-growing-normalization}
\end{equation}
For complex $w$, the two functions $a_\pm(w)$ are the declared shifts;
they are not interpreted as complex conjugates of one another.

In this subsection, on the nonpolynomial functions in question,
$\mathcal H_1q(\eta)$ is defined by the convergent Gaussian integral
$\E q(\eta+iZ)$.  It agrees with the polynomial heat operator; no
convergence of an infinite formal differential series is assumed.

\begin{theorem}[Translated limit of the growing Gaussian symbol]
\label{thm:heatblock-growing-symbol}
The functions
\begin{equation}
 \frac{\mathscr F_M(A_M+\eta)}{C_M}
 =e^{-\eta^2/2}(\mathcal H_1q_M)(\eta)
 \label{eq:heatblock-growing-exact-shift}
\end{equation}
converge locally uniformly in $\eta\in\C$, together with their fixed
order derivatives, to the nonzero entire function
\begin{equation}
 \mathscr B(\eta)=e^{-\eta^2/2}\E h(\eta+iZ)
 =\frac1{\sqrt{2\pi}}\int_\R e^{-v^2/2}h(iv)e^{-i\eta v}\dd v.
 \label{eq:heatblock-growing-limit}
\end{equation}
The function $\mathscr B$ is real on $\R$ and tends to zero at both
real ends.  There is a constant $C$, depending only on $s,t,c_0$, such
that the last real critical point $\xi_M^*$ of $\mathscr F_M$ obeys
\begin{equation}
 \xi_M^*\ge\frac{\log M}{s}-C
 \label{eq:heatblock-growing-last-critical}
\end{equation}
for all sufficiently large $M$.
\end{theorem}

\begin{proof}
The exact gamma product is
\[
 \Pi_M^+(w)=(2s)^{2M}
 \frac{\Gamma(M+a_+(w))\Gamma(M+a_-(w))}
      {\Gamma(a_+(w))\Gamma(a_-(w))}.
\]
The gamma-ratio asymptotic on compact sets gives $q_M\to h$.
For $\Re w$ in a fixed compact interval and $M$ sufficiently large,
\[
 |\Gamma(M+a_\pm(w))|
 \le\Gamma(M+\Re a_\pm(w)).
\]
After the normalization in \eqref{eq:heatblock-growing-normalization},
the remaining reciprocal gamma factors on a vertical line are bounded
by a polynomial times $e^{\pi|\Im w|/(2s)}$.  A Gaussian majorant
therefore proves locally uniform convergence of
$\E q_M(\eta+iZ)$; Cauchy's formula gives convergence of fixed
derivatives.

The exponential translation identity
\[
 \mathcal H_1[e^{Aw}q(w)](\xi)
 =e^{A\xi-A^2/2}(\mathcal H_1q)(\xi-A)
\]
proves \eqref{eq:heatblock-growing-exact-shift}.  It fixes the
normalization: there is no additional factor $e^{-A_M^2/2}$ in $C_M$.
To obtain the last integral in \eqref{eq:heatblock-growing-limit},
write the first integral on the vertical line $\Re w=\eta$ with
integrand $e^{w^2/2-\eta w}h(w)$ and shift that contour to $\Re w=0$.
The reciprocal gamma factors are entire and have subquadratic
exponential growth, so the Gaussian controls the horizontal ends.
The function $e^{-v^2/2}h(iv)$ is integrable and nonzero.  Fourier
uniqueness shows that $\mathscr B$ is not identically zero, and
Fourier decay gives its two real limits.  Conjugation symmetry makes
$\mathscr B$ real on $\R$.

Choose $\eta_1<\eta_*<\eta_2$ such that
$|\mathscr B(\eta_*)|$ exceeds both endpoint values.  Local uniform
convergence forces $|\mathscr F_M(A_M+\eta)|$ to attain an interior,
nonzero maximum on $[\eta_1,\eta_2]$ for all sufficiently large $M$.
That point is a real critical point.  Taking $C\ge-\eta_1$ proves
\eqref{eq:heatblock-growing-last-critical}.
\end{proof}

This is a growing-degree theorem for the Gaussian-polynomial symbol
itself.  It gives a lower bound, not an asymptotic equality for the
last critical point.  Its transfer to an arithmetic filter requires
control of a differential product whose order also grows with $M$;
a central Gaussian limit with only fixed-order derivative control
does not supply that transfer.  The exact finite kernel and the
remaining signed estimate are isolated in
\Cref{subsec:filter-beta-transfer}.


\subsection{Critical gamma-factor scaling and the remaining architecture}
\label{subsec:heatblock-gamma-boundary}

The finite multiplier test applies to products of real linear factors
and conjugate quadratic factors in a gamma-product coefficient
sequence.  At its critical square-root scaling, a fixed block has
\begin{equation}
 K=\frac{1+\theta}{\theta^2},\qquad
 y_j^2=\frac{\gamma_j^2R}{\theta^2\varkappa^2},
 \qquad\theta>0,
 \label{eq:heatblock-gamma-scaling}
\end{equation}
where $R$ is the number of gamma ladders available to the construction,
$\varkappa$ is its moving-line scale, and $\gamma_j$ is the ordinate
carried by the $j$th conjugate block.  Here $\varkappa$ is unrelated
to the character parity $\kappa$.  A block containing $m$ pairs and
$p$ real factors requires $R\ge2m+p$.  The conclusions below are
conditional on this particular scaling and allocation rule; no such
rule is asserted for an arbitrary spectral representation.

\begin{corollary}[Finite one-pair padding cannot cross the degree boundary]
\label{cor:heatblock-gamma-onepair}
For one pair with $p\ge1$ real factors in
\eqref{eq:heatblock-gamma-scaling}, a real-rooted heat certificate
requires
\begin{equation}
 \varkappa^2\ge
 \frac{\gamma^2R}{(p+2)(1+\theta)}
 \ge\frac{\gamma^2}{1+\theta}.
 \label{eq:heatblock-gamma-onepair}
\end{equation}
The first inequality is strict for even $p$.  For $p=2$ it sharpens
to $\varkappa^2\ge\gamma^2R/[c_\star(1+\theta)]$.
For a general fixed block the total-height estimate instead gives
\begin{equation}
 \varkappa^2\ge
 \frac{R\sum_{j=1}^m\gamma_j^2}
 {m(2m+2p-1)(1+\theta)}.
 \label{eq:heatblock-gamma-total}
\end{equation}
\end{corollary}

\begin{proof}
Substitute \eqref{eq:heatblock-gamma-scaling} in
\Cref{thm:heatblock-onepair,thm:heatblock-quartic,thm:heatblock-total-height}
and use the factor count.
\end{proof}

At minimal one-pair ladder count, odd real padding reaches the limiting
degree boundary but cannot enter its open interior; the optimized
quartic remains strictly outside it.  These statements close the
fixed one-pair, factorwise real-rootedness route through that interior.
They do not exclude growing block size, joint processing of several
conjugate factors, a nonuniform pole mesh, or a local variation-count
argument that does not require a globally real-rooted multiplier.
They also do not rule out the putative zero used to define the block.
The actual arithmetic closure still requires a source-valid estimate
of the signed current or a cofinal positivity certificate.
\Cref{thm:heatblock-growing-symbol} separately controls a growing
Gaussian symbol, and \Cref{sec:admissible-spectral-filters} makes the
finite filter, its admission conditions, and its completed compensation
explicit.  These connections do not replace the fixed-block limit by an
unproved increasing-order approximation.


\section{Completed Euler--Weil Schur renewal and finite-history coercivity}
\label{sec:grh-weil-schur}
% =============================================================================

The completed theta source of \Cref{sec:grh-theta-heat} encodes GRH through
real-rootedness.  A complementary spectral chart is obtained by keeping the
Euler cells as passive Schur systems and inserting the archimedean completion
before testing positivity.  The resulting construction is finite-support and
source native.  It produces exact positive local innovations, an orthogonal
Euler renewal, and a coercive high-mode complement for every fixed localized
Weil history.  It also supplies two exact no-go statements: the fully
completed scattering ratio is constant, and scalar normalization of the
finite Euler Hermite--Biehler functions cannot yield a nonzero infinite
limit.

Fix a primitive nonprincipal character $\chi$ of conductor $f$, with
$\chi(-1)=(-1)^{\kappa}$, $\kappa\in\{0,1\}$.  Put
\[
 s_z=\frac12-\mathrm i z,
 \qquad z\in\C.
 \tag{W.1}
\]
The critical line corresponds to real $z$.

\subsection{Local Euler scattering cells and exact Pick renewal}

Let $p\nmid f$ be prime and put
\[
 L_p=\log p,
 \qquad r_p=p^{-1/2},
 \qquad \alpha_{p,\chi}=\overline{\chi(p)}r_p,
 \qquad w_p(z)=\exp(\mathrm iL_pz).
\]
For $\im z>0$, one has $|w_p(z)|<1$.  Define
\[
 \boxed{
 b_{p,\chi}(z)
 :=\frac{w_p(z)-\alpha_{p,\chi}}
         {1-\overline{\alpha_{p,\chi}}w_p(z)}.}
 \tag{W.2}
\]

\begin{theorem}[The local Euler cell is strict Schur]
\label{thm:weil-local-Schur}
The function $b_{p,\chi}$ is analytic and satisfies
$|b_{p,\chi}(z)|<1$ on the upper half-plane and
$|b_{p,\chi}(t)|=1$ for real $t$.  Moreover,
\[
 \boxed{
 b_{p,\chi}(z)
 =p^{\mathrm iz}
  \frac{L_p(s_z,\chi)}{L_p(1-s_z,\overline\chi)},}
 \tag{W.3}
\]
where $L_p(s,\chi)=(1-\chi(p)p^{-s})^{-1}$.
\end{theorem}

\begin{proof}
The map $w\mapsto(w-\alpha)/(1-\overline\alpha w)$ is an automorphism of the
unit disk when $|\alpha|<1$.  Apply it to $w_p(z)$.  The boundary assertion
follows because disk automorphisms preserve the unit circle.  Finally,
$p^{-s_z}=p^{-1/2+\mathrm iz}=r_pw_p(z)$ and
$p^{-(1-s_z)}=r_pw_p(z)^{-1}$; substitution gives \textup{(W.3)}.
\end{proof}

Choose a continuous phase lift
$b_{p,\chi}(t)=\exp(\mathrm i\vartheta_{p,\chi}(t))$ on the real line.

\begin{theorem}[Positive phase density and prime-power Fourier modes]
\label{thm:weil-phase-density}
The phase is strictly increasing and
\[
 \boxed{
 \vartheta'_{p,\chi}(t)
 =L_p\frac{1-r_p^2}
 {|\exp(\mathrm iL_pt)-\alpha_{p,\chi}|^2}.}
 \tag{W.4}
\]
It has the absolutely and uniformly convergent expansion
\[
 \boxed{
 \vartheta'_{p,\chi}(t)
 =L_p\left(
 1+2\operatorname{Re}\sum_{k\ge1}
 \frac{\chi(p)^k}{p^{k/2}}
 \exp(\mathrm ikL_pt)
 \right).}
 \tag{W.5}
\]
Thus its nonconstant Fourier coefficients are exactly
$\chi(p)^k\Lambda(p^k)/\sqrt{p^k}$.
\end{theorem}

\begin{proof}
Logarithmic differentiation of the disk automorphism on the unit circle gives
its Poisson kernel, proving \textup{(W.4)}.  The standard Poisson expansion
\[
 \frac{1-|\alpha|^2}{|1-\overline\alpha e^{\mathrm i\theta}|^2}
 =1+2\operatorname{Re}\sum_{k\ge1}
   \overline\alpha^{\,k}e^{\mathrm ik\theta}
\]
with $\overline\alpha_{p,\chi}=\chi(p)p^{-1/2}$ gives \textup{(W.5)}.
\end{proof}

For an analytic function $F$ on the upper half-plane, define its Pick kernel
\[
 K_F(z,w)
 :=\frac{1-F(z)\overline{F(w)}}{-\mathrm i(z-\overline w)}.
 \tag{W.6}
\]
For $L>0$, the free delay $w_L(z)=\exp(\mathrm iLz)$ obeys
\[
 K_{w_L}(z,w)
 =\int_0^L\exp(\mathrm itz)
   \overline{\exp(\mathrm itw)}\dd t.
 \tag{W.7}
\]
Write $\mathcal H(K_F)$ for the reproducing-kernel Hilbert space generated by
$K_F$.

\begin{theorem}[Exact local Pick factorization]
\label{thm:weil-Pick-factorization}
Put
\[
 g_{p,\chi}(z)
 :=\frac{\sqrt{1-r_p^2}}
 {1-\chi(p)r_pw_p(z)}.
\]
Then
\[
 \boxed{
 K_{b_{p,\chi}}(z,w)
 =g_{p,\chi}(z)\overline{g_{p,\chi}(w)}K_{w_p}(z,w).}
 \tag{W.8}
\]
Hence the local Euler cell is unitarily equivalent, as a reproducing-kernel
space, to a filtered free delay interval of length $\log p$.
\end{theorem}

\begin{proof}
For $u,v$ in the disk and $|\alpha|<1$,
\[
 1-\frac{u-\alpha}{1-\overline\alpha u}
 \overline{\left(\frac{v-\alpha}{1-\overline\alpha v}\right)}
 =\frac{(1-|\alpha|^2)(1-u\overline v)}
 {(1-\overline\alpha u)(1-\alpha\overline v)}.
\]
Apply this identity with $u=w_p(z)$ and $v=w_p(w)$, divide by
$-\mathrm i(z-\overline w)$, and use \textup{(W.7)}.
\end{proof}

Let $\mathcal P=\{p_1<\cdots<p_m\}$ be a finite unramified prime bank and put
\[
 B_{0,\chi}=1,
 \qquad
 B_{j,\chi}=\prod_{\nu=1}^jb_{p_\nu,\chi}.
\]

\begin{theorem}[Exact finite Euler spectral renewal]
\label{thm:weil-Euler-renewal}
The finite cascade is Schur and
\[
 \boxed{
 \begin{aligned}
 K_{B_{m,\chi}}(z,w)
 =\sum_{j=1}^m&
 B_{j-1,\chi}(z)\overline{B_{j-1,\chi}(w)}\\
 &\times g_{p_j,\chi}(z)\overline{g_{p_j,\chi}(w)}
 K_{w_{p_j}}(z,w).
 \end{aligned}}
 \tag{W.9}
\]
Every summand is positive.  The associated model space has the orthogonal
decomposition
\[
 \boxed{
 \mathcal H(K_{B_{m,\chi}})
 =\bigoplus_{j=1}^m
 B_{j-1,\chi}g_{p_j,\chi}\mathcal H(K_{w_{p_j}}).}
 \tag{W.10}
\]
Thus adjoining a prime creates one positive filtered delay sector and does
not erase any previous sector.
\end{theorem}

\begin{proof}
The algebraic product identity
\[
 K_{FG}=K_F+F(z)\overline{F(w)}K_G
\]
iterated through the finite product, together with
\Cref{thm:weil-Pick-factorization}, gives \textup{(W.9)}.  The standard
model-space product decomposition gives the orthogonal direct sum
\textup{(W.10)}.
\end{proof}

\subsection{Finite Hermite--Biehler spectra and the completion audit}

Put
\[
 \Theta_{\mathcal P}:=\sum_{p\in\mathcal P}\log p,
 \qquad
 D_{\mathcal P,\chi}(z)
 :=\prod_{p\in\mathcal P}
  \left(1-\chi(p)p^{-1/2+\mathrm iz}\right),
\]
\[
 \boxed{
 E_{\mathcal P,\chi}(z)
 :=\exp\!\left(-\frac{\mathrm i}2\Theta_{\mathcal P}z\right)
 D_{\mathcal P,\chi}(z).}
 \tag{W.11}
\]
For an entire function $F$, write $F^\#(z)=\overline{F(\overline z)}$.

\begin{theorem}[Finite Euler Hermite--Biehler theorem]
\label{thm:weil-finite-HB}
The function $E_{\mathcal P,\chi}$ has no upper-half-plane zero and
\[
 \boxed{
 \frac{E_{\mathcal P,\chi}^\#(z)}
 {E_{\mathcal P,\chi}(z)}
 =B_{m,\chi}(z).}
 \tag{W.12}
\]
It is therefore Hermite--Biehler.  Its real and imaginary parts
\[
 A_{\mathcal P,\chi}=\frac{E+E^\#}{2},
 \qquad
 C_{\mathcal P,\chi}=\frac{E-E^\#}{2\mathrm i}
\]
have only real, simple, strictly interlacing zeros.
\end{theorem}

\begin{proof}
Every factor of $D_{\mathcal P,\chi}$ is nonzero in the upper half-plane.
Direct division gives \textup{(W.12)}, and the strict Schur inequality gives
$|E^\#|<|E|$ there.  A nonreal zero of $A$ or $C$ would force the quotient
$E^\#/E$ to equal $-1$ or $1$ in one half-plane, a contradiction.  On the
real axis, the strictly monotone phase from
\Cref{thm:weil-phase-density} gives simplicity and interlacing.
\end{proof}

The finite spectral systems do not pass to the completed function by a naive
limit.

\begin{theorem}[Functional-equation neutralization]
\label{thm:weil-completed-ratio}
For the completed function $\mathfrak L_\chi$ of \textup{(J.1)},
\[
 \boxed{
 \frac{\mathfrak L_\chi(s_z)}
 {\mathfrak L_{\overline\chi}(1-s_z)}
 =\varepsilon_\chi.}
 \tag{W.13}
\]
After cancellation of common zeros, the completed scattering ratio is the
constant root number.  It has zero Pick kernel and contains no zero-location
information.
\end{theorem}

\begin{proof}
This is the functional equation \textup{(J.2)} evaluated at $s=s_z$.
\end{proof}

\begin{theorem}[Minimal Hermite--Biehler clock and no scalar limit]
\label{thm:weil-no-scalar-limit}
For
\[
 E_{\mathcal P,\chi}^{(c)}(z)
 :=\exp(-\mathrm icz)D_{\mathcal P,\chi}(z),
\]
the Hermite--Biehler property holds exactly when
\[
 \boxed{c\ge\frac12\Theta_{\mathcal P}.}
 \tag{W.14}
\]
Let $\mathcal P_j$ increase with
$\Theta_{\mathcal P_j}\to\infty$, choose
$c_j\ge\Theta_{\mathcal P_j}/2$, and let $a_j\in\C$.  If
$a_jE_{\mathcal P_j,\chi}^{(c_j)}$ converges locally uniformly on $\C$, its
limit is identically zero.
\end{theorem}

\begin{proof}
One has
\[
 \frac{(E^{(c)})^\#}{E^{(c)}}
 =\exp\!\bigl(\mathrm i(2c-\Theta_{\mathcal P})z\bigr)
 B_{m,\chi}(z).
\]
The exponential is Schur precisely when
$2c-\Theta_{\mathcal P}\ge0$.  If the inequality fails, evaluation at
$z=\mathrm iy$ as $y\to\infty$ makes the quotient unbounded.

For the limit statement, suppose the limit $F$ is nonzero and choose
$1<y_1<y_2$ with $F(\mathrm iy_1)F(\mathrm iy_2)\ne0$.  The scalar $a_j$
cancels from the ratio of the two values, whereas
\[
 \left|
 \frac{E_{\mathcal P_j,\chi}^{(c_j)}(\mathrm iy_2)}
      {E_{\mathcal P_j,\chi}^{(c_j)}(\mathrm iy_1)}
 \right|
 =\exp(c_j(y_2-y_1))
  \left|
  \frac{D_{\mathcal P_j,\chi}(\mathrm iy_2)}
       {D_{\mathcal P_j,\chi}(\mathrm iy_1)}
  \right|.
\]
The Euler products in the second factor converge to finite nonzero limits
because $y_1>1$, while the exponential tends to infinity.  This contradicts
local uniform convergence to a finite nonzero ratio.
\end{proof}

\begin{principle}[Completion precedes positivity]
The finite Euler Schur positivity is genuine, but its free-delay clock has
divergent mean type.  The conductor and gamma counterterms must be installed
before an infinite-support positivity claim is formulated.  The constant
completed ratio itself cannot serve as the characteristic function; one must
retain a logarithmic derivative, a Weil form, or a nontrivial determinant
before functional-equation cancellation.
\end{principle}

\subsection{Localized Weil form and exact archimedean resummation}

For $g\in C_c^\infty(\R)$, put
\[
 \widehat g(t)=\int_\R g(x)\exp(-\mathrm itx)\dd x,
 \qquad
 C_g(u)=\int_\R g(x+u)\overline{g(x)}\dd x.
\]
The following standard input is the fixed-character Weil criterion; see
\cite{Bombieri2000,Lagarias2004}.

\begin{theorem}[Localized Guinand--Weil form; standard input]
\label{thm:weil-explicit-form}
There is a Hermitian archimedean--conductor form
\[
 \mathcal A_\chi[g]
 =\frac1{2\pi}\int_\R|\widehat g(t)|^2
 \left[
  \log\frac f\pi
  +\operatorname{Re}\frac{\Gamma'}{\Gamma}
   \!\left(\frac{1/2+\kappa+\mathrm it}{2}\right)
 \right]\dd t
 \tag{W.15}
\]
such that
\[
 \boxed{
 \mathcal W_\chi[g]
 =\mathcal A_\chi[g]
 -2\operatorname{Re}\sum_{n\ge2}
  \frac{\Lambda(n)\chi(n)}{\sqrt n}C_g(\log n).}
 \tag{W.16}
\]
GRH for $L(s,\chi)$ is equivalent to
$\mathcal W_\chi[g]\ge0$ for every such $g$.  If
$\supp g\subset(-A,A)$, only prime powers $n<\exp(2A)$ occur.
\end{theorem}

\begin{proof}[Source status]
The identity and positivity criterion are the classical Guinand--Weil
explicit formula and Weil criterion for a primitive Dirichlet $L$-function.
The support assertion follows because $C_g(u)=0$ for $|u|\ge2A$.
\end{proof}

\begin{remark}[Hermitian pairing for complex characters]
For nonreal $\chi$, the zero-side form pairs the $\chi$ channel with its
contragredient $\overline\chi$ under
$\rho\leftrightarrow1-\overline\rho$.  The arithmetic expression in
\textup{(W.16)} is real because the prime contribution is explicitly paired
with its complex conjugate.
\end{remark}

\subsection{Support renewal and compressed prime cells}

Let $I_A=(-A,A)$ and let $P_A$ be orthogonal projection from $L^2(\R)$ onto
$L^2(I_A)$.  Extend functions by zero outside $I_A$, put
\[
 (\tau_ug)(x)=g(x+u),
 \qquad
 T_{A,u}=P_A\tau_uP_A,
\]
and for $n\ge2$ define the self-adjoint prime-power cell
\[
 \boxed{
 \mathsf P_{\chi,A}(n)
 :=\chi(n)T_{A,\log n}
  +\overline{\chi(n)}T_{A,-\log n}.}
 \tag{W.16a}
\]

\begin{proposition}[Exact support activation and prime-cell operator]
\label{prop:weil-prime-cell}
If $g\in C_c^\infty(I_A)$, then $C_g(u)=0$ for $|u|\ge2A$ and
\[
 \boxed{
 2\operatorname{Re}\!\left(\chi(n)C_g(\log n)\right)
 =\langle\mathsf P_{\chi,A}(n)g,g\rangle.}
 \tag{W.16b}
\]
Consequently
\[
 \boxed{
 \mathcal W_{\chi,A}[g]
 =\mathcal A_{\chi,A}[g]
 -\sum_{2\le n<\exp(2A)}
  \frac{\Lambda(n)}{\sqrt n}
  \langle\mathsf P_{\chi,A}(n)g,g\rangle,}
 \tag{W.16c}
\]
and $\|\mathsf P_{\chi,A}(n)\|\le2$.
\end{proposition}

\begin{proof}
The two intervals supporting $g(x+u)$ and $g(x)$ are disjoint when
$|u|\ge2A$.  Moreover
$\langle T_{A,u}g,g\rangle=C_g(u)$ and
$C_g(-u)=\overline{C_g(u)}$, which proves \textup{(W.16b)} and hence
\textup{(W.16c)}.  Each $T_{A,u}$ is a compression of a unitary translation,
so it has norm at most one.
\end{proof}

\begin{theorem}[Ambient threshold norm jump]
\label{thm:weil-shift-jump}
For $u\ge0$,
\[
 \boxed{
 \|T_{A,u}\|=
 \begin{cases}
 1,&0\le u<2A,\\
 0,&u\ge2A.
 \end{cases}}
 \tag{W.16d}
\]
Thus a newly activated prime-power cell is not an operator-norm-small
perturbation merely because its overlap interval is short.
\end{theorem}

\begin{proof}
The operator is zero when the translated intervals are disjoint.  If
$u<2A$, choose a nonzero function supported in a subinterval on which the
translation remains inside $I_A$; on that vector the compression preserves
its full $L^2$ norm.  The reverse inequality follows from contraction of an
orthogonal compression.
\end{proof}

Put
\[
 \alpha=\frac{1+2\kappa}{4},
 \qquad
 b_n=n+\alpha,
\]
\[
 k_\alpha(t)
 :=\sum_{n\ge0}\exp(-2b_nt)
 =\frac{\exp(-2\alpha t)}{1-\exp(-2t)},
 \qquad
 s_\alpha(t):=\frac1{2t}-k_\alpha(t).
 \tag{W.17}
\]
The singularity of $s_\alpha$ at zero is removable.  On
$I_A=(-A,A)$, define
\[
 \mathfrak D_A[g]
 :=\frac14\int_{I_A}\int_{I_A}
 \frac{|g(x)-g(y)|^2}{|x-y|}\dd x\dd y
 \tag{W.18}
\]
and let $\mathsf S_{\alpha,A}$ be the bounded integral operator with kernel
$s_\alpha(|x-y|)$.

\begin{lemma}[One-resolvent difference energy]
\label{lem:weil-one-resolvent}
For $b,A>0$ and the integral operator
$\mathsf K_{b,A}(x,y)=\exp(-2b|x-y|)$,
\[
 \boxed{
 \begin{aligned}
 \frac1b\|g\|_2^2-\langle\mathsf K_{b,A}g,g\rangle
 ={}&\frac12\iint_{I_A^2}
 \exp(-2b|x-y|)|g(x)-g(y)|^2\dd x\dd y\\
 &+\int_{I_A}
 \frac{\exp(-2b(A-x))+\exp(-2b(A+x))}{2b}
 |g(x)|^2\dd x.
 \end{aligned}}
 \tag{W.19}
\]
In particular, every resolvent difference is nonnegative.
\end{lemma}

\begin{proof}
For fixed $x$,
\[
 \int_{I_A}\exp(-2b|x-y|)\dd y
 =\frac1b-\frac{\exp(-2b(A-x))+\exp(-2b(A+x))}{2b}.
\]
Expand the difference square and use symmetry.
\end{proof}

\begin{lemma}[Boundary writer constant]
\label{lem:weil-boundary-writer}
Define
\[
 F_\alpha(y):=\sum_{n\ge0}\frac{\exp(-2b_ny)}{b_n}.
\]
Then
\[
 \boxed{
 F_\alpha(y)+\log y
 \longrightarrow-\log2-\gamma-\psi(\alpha)
 \quad(y\downarrow0),}
 \tag{W.20}
\]
and $F_\alpha'(y)=-2k_\alpha(y)$.
\end{lemma}

\begin{proof}
Termwise differentiation is locally uniform for $y>0$.  Subtract the
comparison series
$\sum_{n\ge0}\exp(-2(n+1)y)/(n+1)=-\log(1-\exp(-2y))$.
The difference converges by dominated convergence to
$-\psi(\alpha)-\gamma$, and
$\log(1-\exp(-2y))=\log(2y)+o(1)$.
\end{proof}

\begin{theorem}[Exact parity-renormalized archimedean form]
\label{thm:weil-renormalization}
On the logarithmic form domain obtained by closing $C_c^\infty(I_A)$, the
localized archimedean form in \textup{(W.15)} is
\[
 \boxed{
 \begin{aligned}
 \mathcal A_{\chi,A}[g]
 ={}&\mathfrak D_A[g]\\
 &+\int_{-A}^A
 \left(
  \log\frac f{2\pi}-\gamma
  -\frac12\log(A^2-x^2)
 \right)|g(x)|^2\dd x\\
 &+\langle\mathsf S_{\alpha,A}g,g\rangle.
 \end{aligned}}
 \tag{W.21}
\]
The singular difference energy and the boundary logarithm are universal;
all parity dependence is confined to the smooth kernel $s_\alpha$.
\end{theorem}

\begin{proof}
The digamma partial-fraction expansion gives, with $b_n=n+\alpha$,
\[
 m_\chi(t)
 =m_\chi(0)+\sum_{n\ge0}
 \left(\frac1{b_n}-\frac{4b_n}{t^2+4b_n^2}\right),
\]
where
$m_\chi(t)$ is the multiplier in \textup{(W.15)}.  Fourier inversion turns
$4b_n/(t^2+4b_n^2)$ into the kernel
$\exp(-2b_n|x-y|)$.  Apply
\Cref{lem:weil-one-resolvent} and sum over $n$.

Now use $k_\alpha(t)=1/(2t)-s_\alpha(t)$.  The singular part is precisely
\textup{(W.18)}.  Expanding the smooth difference energy converts it into
$\langle\mathsf S_{\alpha,A}g,g\rangle$ plus a multiplication term.  The
primitive obtained from \Cref{lem:weil-boundary-writer} combines with the
endpoint writer and with
$m_\chi(0)=\log(f/\pi)+\psi(\alpha)$ to give
\[
 \log\frac f{2\pi}-\gamma-\frac12\log(A^2-x^2).
\]
All series converge first on a smooth core and then by closure of the
lower-bounded logarithmic form.
\end{proof}

\subsection{Legendre diagonalization and finite low-mode reduction}

Let $\mathsf A_{\chi,A}$ denote the lower-bounded self-adjoint operator
associated with the closed form in \textup{(W.21)}, and let
$U_A:L^2(-1,1)\to L^2(-A,A)$ be
$(U_Ah)(x)=A^{-1/2}h(x/A)$.  Under scaling, in the sense of closed
quadratic forms,
\[
 U_A^*\mathsf A_{\chi,A}U_A
 =\mathsf D+c_{f,A}I+\mathsf W
  +\mathsf S_{\alpha,A}^{\rm sc},
 \tag{W.22}
\]
where
\[
 c_{f,A}=\log\frac f{2\pi A}-\gamma,
 \qquad
 (\mathsf Wh)(t)=-\frac12\log(1-t^2)h(t),
\]
and the smooth kernel is
$A s_\alpha(A|t-v|)$.  The quadratic form of $\mathsf D$ is
\[
 \langle\mathsf Dh,h\rangle
 =\frac14\int_{-1}^1\int_{-1}^1
 \frac{|h(t)-h(v)|^2}{|t-v|}\dd t\dd v.
 \tag{W.23}
\]
Let $P_n$ be the Legendre polynomial and
$e_n(t)=\sqrt{(2n+1)/2}\,P_n(t)$.

\begin{theorem}[Exact harmonic diagonal]
\label{thm:weil-Legendre-diagonal}
One has
\[
 \boxed{\mathsf De_n=H_ne_n,}
 \qquad
 H_n=\sum_{j=1}^n\frac1j,
 \quad H_0=0.
 \tag{W.24}
\]
\end{theorem}

\begin{proof}
Define on polynomials
\[
 (\mathcal Lh)(t)
 :=\int_{-1}^1\frac{h(t)-h(v)}{|t-v|}\dd v.
\]
The operator is symmetric and preserves every polynomial-degree filtration.
For $t^n$, the leading coefficient contributed by each side of $v=t$ is
$H_n$, so the degree-$n$ quotient eigenvalue is $2H_n$.  Symmetry and
orthogonality force $\mathcal LP_n=2H_nP_n$.  Comparison of quadratic forms
gives $\mathsf D=\mathcal L/2$.
\end{proof}

At a fixed support wall, let $n_1,\ldots,n_r<\exp(2A)$ be the active
unramified prime powers and put
\[
 \lambda_j=\frac{\Lambda(n_j)}{\sqrt{n_j}}.
\]
Let
\[
 \mathsf P_{\chi,A}^{\rm sc}(n_j)
 :=U_A^*\mathsf P_{\chi,A}(n_j)U_A.
\]
By \Cref{prop:weil-prime-cell},
$\|\mathsf P_{\chi,A}^{\rm sc}(n_j)\|\le2$.
The complete finite-history operator is
\[
 \boxed{
 \mathsf H_{\chi,A}^{[r],\rm sc}
 =\mathsf D+c_{f,A}I+\mathsf W
  +\mathsf S_{\alpha,A}^{\rm sc}
  -\sum_{j=1}^r\lambda_j
   \mathsf P_{\chi,A}^{\rm sc}(n_j).}
 \tag{W.25}
\]

\begin{theorem}[Coercive high-mode complement for every finite history]
\label{thm:weil-high-mode}
Let $P_N$ project onto $\Span\{e_0,\ldots,e_{N-1}\}$ and put $Q_N=I-P_N$.
Then
\[
 \boxed{
 Q_N\mathsf H_{\chi,A}^{[r],\rm sc}Q_N
 \succeq d_NQ_N,}
 \tag{W.26}
\]
where one may take
\[
 \boxed{
 d_N=H_N+c_{f,A}
 -\|\mathsf S_{\alpha,A}^{\rm sc}\|
 -2\sum_{j=1}^r\lambda_j.}
 \tag{W.27}
\]
In particular $d_N>0$ for all sufficiently large $N$.
\end{theorem}

\begin{proof}
The harmonic diagonal contributes at least $H_N$ on $Q_N$,
$\mathsf W\succeq0$, the smooth operator is bounded, and each active prime
cell has norm at most two.  Since $H_N\to\infty$, the lower bound is eventually
positive.
\end{proof}

\begin{corollary}[Finite protected Schur obstruction]
\label{cor:weil-finite-Schur}
Choose $N$ with $d_N>0$ and write
$\mathsf H=\mathsf H_{\chi,A}^{[r],\rm sc}$.  Then
$\mathsf H\succeq0$ if and only if its finite Schur complement
\[
 \boxed{
 \mathsf F_{\chi,A,N}
 :=P_N\mathsf HP_N
  -P_N\mathsf HQ_N
   (Q_N\mathsf HQ_N)^{-1}
   Q_N\mathsf HP_N}
 \tag{W.28}
\]
is nonnegative on the $N$-dimensional protected low-mode space.  Thus every
fixed localized sign problem has an explicitly coercive complement and a
finite matrix obstruction.
\end{corollary}

\begin{proof}
Apply the finite/infinite block Feshbach formula with the strictly positive
$Q_N$ block supplied by \Cref{thm:weil-high-mode}.
\end{proof}

\begin{criterion}[Cofinal completed Weil--Schur closure]
\label{crit:weil-Schur-GRH}
Suppose that for every primitive character $\chi$ and every support wall
$A>0$, the finite Schur matrix \textup{(W.28)} is nonnegative for one
endpoint-independent cutoff $N=N(\chi,A)$, determined from the completed
active history and satisfying $d_N>0$.  Then every
localized Weil form is nonnegative, and GRH holds in every primitive channel.
The pole-removed zeta completion gives the principal analogue.
\end{criterion}

\begin{proof}
The Schur condition and \Cref{thm:weil-high-mode} imply positivity of the full
localized operator \textup{(W.25)}.  Undoing the scaling gives
$\mathcal W_\chi[g]\ge0$ for every compactly supported test function.  Apply
\Cref{thm:weil-explicit-form}.
\end{proof}

\begin{remark}[Scope of the spectral renewal]
The harmonic high-mode theorem is an unconditional coercivity result, not a
proof that every low-mode matrix is positive.  Its contribution is to replace
an infinite ambient sign problem at each fixed support wall by a finite,
source-fixed Schur obstruction.  The global difficulty is the cofinal control
of these protected low modes as the support and active prime-power history
grow.
\end{remark}


\subsection{Exact finite acquisition in the Legendre chart}
\label{subsec:weil-acquisition}

The finite Schur obstruction can be populated without an unspecified
matrix representation.  Work on $L^2(-1,1)$ with
$e_k(x)=\sqrt{(2k+1)/2}\,P_k(x)$, and put
\begin{equation}
 b_{mk}(\delta)=\int_{-1}^{1-\delta}e_m(x)e_k(x+\delta)\,\dd x,
 \qquad 0\le\delta\le2,
 \label{eq:acq-translation}
\end{equation}
with $B_N(\delta)=(b_{mk}(\delta))_{0\le m,k<N}$.
Here the translated functions are extended by zero at the support wall.

\begin{proposition}[Polynomial acquisition of every compressed shift]
\label{prop:acq-shift}
Define
$c_{j,a}=(-1)^a(2j-2a)!/[2^ja!(j-a)!(j-2a)!]$.
For $r=m-2a$ and $s=k-2b$,
\begin{equation}
\begin{split}
 b_{mk}(\delta)
 ={}&\frac{\sqrt{(2m+1)(2k+1)}}2
 \sum_{a=0}^{\lfloor m/2\rfloor}
 \sum_{b=0}^{\lfloor k/2\rfloor}c_{m,a}c_{k,b}\\
 &\quad\times\sum_{\ell=0}^{s}\binom{s}{\ell}\delta^{s-\ell}
 \frac{(1-\delta)^{r+\ell+1}-(-1)^{r+\ell+1}}
 {r+\ell+1}.
\end{split}
\label{eq:acq-shift-polynomial}
\end{equation}
Furthermore $B_N(0)=I$, $B_N(2)=0$, and
$b_{km}(\delta)=(-1)^{m+k}b_{mk}(\delta)$.
\end{proposition}

\begin{proof}
Expand the two Legendre polynomials in monomials, then expand
$(x+\delta)^s$ and integrate each monomial over $[-1,1-\delta]$.
The endpoint identities follow from orthonormality and the empty
integration interval.  The substitution $x\mapsto-x-\delta$, together
with $P_j(-x)=(-1)^jP_j(x)$, proves the transpose identity.
\end{proof}

For an active prime power let $\delta_n=(\log n)/A$.  Its low matrix is
\begin{equation}
 \mathcal P_{\chi,N}(n)
 =\chi(n)B_N(\delta_n)+\overline{\chi(n)}B_N(\delta_n)^{\mathsf T}.
 \label{eq:acq-prime-matrix}
\end{equation}
The same-parity entries contain $2\Re\chi(n)$ and the opposite-parity
entries contain $2i\Im\chi(n)$.  Multiplying every odd basis vector by
$i$ makes these matrices real symmetric.  The archimedean matrices
preserve parity, so the same fixed gauge works for the complete low
block.  For a real character the even and odd sectors already decouple.

\begin{proposition}[Wall and gamma matrices]
\label{prop:acq-arch}
For the logarithmic wall $W(x)=-\frac12\log(1-x^2)$, put
\begin{equation}
 j_q=\frac{2}{2q+1}\bigl(2\log2-2H_{2q+2}+H_{q+1}\bigr).
 \label{eq:acq-log-moment}
\end{equation}
Then $W_{mk}=0$ when $m+k$ is odd, and
\begin{equation}
 W_{mk}=\frac{\sqrt{(2m+1)(2k+1)}}{(k-m)(k+m+1)}
 \quad(m<k,\ m+k\text{ even}),
 \label{eq:acq-wall-offdiag}
\end{equation}
while
\begin{equation}
 W_{mm}=-\frac{2m+1}{4}
 \sum_{a,b=0}^{\lfloor m/2\rfloor}c_{m,a}c_{m,b}j_{m-a-b}.
 \label{eq:acq-wall-diag}
\end{equation}
For the smooth gamma kernel already defined in the localized Weil form,
\begin{equation}
 (\mathsf S_{\alpha,A}^{\rm sc})_{mk}
 =A\int_0^2s_\alpha(Ar)
       \bigl(b_{mk}(r)+b_{km}(r)\bigr)\,\dd r.
 \label{eq:acq-gamma-matrix}
\end{equation}
Thus its first $N$ modes require only the $2N$ scalar moments
$A\int_0^2r^q s_\alpha(Ar)\dd r$, $0\le q\le2N-1$.
\end{proposition}

\begin{proof}
Differentiating the beta integral at its second parameter equal to one
gives $\int_{-1}^1x^{2q}\log(1-x^2)\dd x=j_q$.
Monomial expansion proves \eqref{eq:acq-wall-diag}.  For completeness,
write $\lambda_k=k(k+1)$ and use the Legendre differential equation in
$\int\log(1-x^2)P_mP_k$.  The boundary term vanishes, and for $m<k$
of the same parity orthogonality reduces the resulting numerator to
$-2\int x(P_mP_k'-P_m'P_k)\dd x=-4$.
Here $\int P_m(x)xP_k'(x)\dd x=2$ follows by expanding $P_k'$ in
lower Legendre polynomials, whereas $xP_m'$ has degree at most $m$.
Division by $\lambda_k-\lambda_m$ gives
\eqref{eq:acq-wall-offdiag} after normalization.  Finally split the
double integral with kernel $A s_\alpha(A|x-y|)$ according to the sign
of $y-x$.  The two halves are precisely the two translation entries in
\eqref{eq:acq-gamma-matrix}.  Formula
\eqref{eq:acq-shift-polynomial} has degree at most $m+k+1$.
\end{proof}

\subsection{A finite-data bracket for the protected Schur operator}
\label{subsec:weil-schur-bracket}

Write the full scaled form as
$\mathsf H=\mathsf D+c_{f,A}I+\mathsf V$, where
\begin{equation}
 \mathsf V=\mathsf W+\mathsf S_{\alpha,A}^{\rm sc}
 -\mathsf E_{\chi,A},\qquad
 \mathsf E_{\chi,A}=\sum_{n<e^{2A}}\frac{\Lambda(n)}{\sqrt n}
                   \mathsf P_{\chi,A}^{\rm sc}(n).
 \label{eq:acq-complete-V}
\end{equation}
Retain the high-mode lower bound $d_N>0$ of
\Cref{thm:weil-high-mode}.  With $P=P_N$, $Q=Q_N$, define
\begin{equation}
 L_N=P\mathsf HP,\qquad
 R_N=(Q\mathsf VP)^*(Q\mathsf VP)\succeq0.
 \label{eq:acq-LR}
\end{equation}
The logarithmic wall sends each fixed polynomial into $L^2$, so the
finite-rank off-diagonal map in this formula is bounded.

\begin{theorem}[Certified Schur bracket]
\label{thm:acq-schur-bracket}
The exact low Schur operator $F_N$ satisfies
\begin{equation}
 \boxed{L_N-d_N^{-1}R_N\preceq F_N\preceq L_N.}
 \label{eq:acq-schur-bracket}
\end{equation}
All entries of $L_N$ are supplied by
\eqref{eq:acq-shift-polynomial}--\eqref{eq:acq-gamma-matrix} and the
harmonic diagonal.  If $v_k=\mathsf V e_k$, then
\begin{equation}
 (R_N)_{mk}=\langle v_m,v_k\rangle
 -\sum_{\ell<N}\langle v_m,e_\ell\rangle
                     \langle e_\ell,v_k\rangle,
 \label{eq:acq-tail-gram}
\end{equation}
using the inner product linear in its second argument.  Thus $R_N$ is
also given by finitely many explicit integrals, without evaluating an
infinite-dimensional inverse.
\end{theorem}

\begin{proof}
The operators $\mathsf D$ and $c_{f,A}I$ have no $P/Q$ cross block.
The exact Schur formula therefore subtracts
$(Q\mathsf VP)^*(Q\mathsf HQ)^{-1}(Q\mathsf VP)$ from $L_N$.
The high-mode estimate implies
$0\preceq(Q\mathsf HQ)^{-1}\preceq d_N^{-1}Q$.
Congruence gives \eqref{eq:acq-schur-bracket}.
The remaining formula is the orthogonal-projection identity for the
Gram of $Qv_k$.
\end{proof}

Nonnegativity of the lower endpoint in \eqref{eq:acq-schur-bracket}
certifies the complete form at this support.  A negative direction of
$L_N$ gives a polynomial witness, and hence a compact smooth witness by
form-domain approximation.  Positivity of $L_N$ alone does not certify
the Schur correction.  An interval straddling zero is inconclusive.
The GRH closure still requires the relevant nonnegative certificate on
a cofinal family of supports.

\subsection{Support activation, arithmetic resonance, and the cost of an absolute tail}
\label{subsec:weil-acquisition-cost}

\begin{proposition}[A support wall is strong, not norm, extinction]
\label{prop:acq-support-wall}
For $0\le\delta<2$, the zero-extended translation on $L^2(-1,1)$ has
operator norm one.  As $\delta\uparrow2$ it converges strongly to zero,
whereas every fixed matrix entry has the expansion
\begin{equation}
 b_{mk}(\delta)=
 \frac{(-1)^m\sqrt{(2m+1)(2k+1)}}2(2-\delta)
 +O_{m,k}((2-\delta)^2).
 \label{eq:acq-support-wall}
\end{equation}
\end{proposition}

\begin{proof}
A unit vector supported on the surviving input interval attains norm one.
For any fixed input its squared output norm is an integral over an
interval of length $2-\delta$, hence tends to zero.  Expand the integrand
in \eqref{eq:acq-translation} at $x=-1$, $x+\delta=1$ for the final formula.
\end{proof}

The distinction prevents a small fixed acquisition matrix at a newly
activated support wall from being used as an ambient operator-norm
estimate.

\begin{proposition}[Dimension cost of the absolute-prime-mass estimate]
\label{prop:acq-dimension-cost}
Let
$\mathcal L_f(A)=\sum_{n<e^{2A},(n,f)=1}\Lambda(n)/\sqrt n$.
For fixed conductor $f$, the prime number theorem and partial summation
give $\mathcal L_f(A)=2e^A+o(e^A)$.  In the particular high-mode estimate
\begin{equation}
 d_N=H_N+c_{f,A}-\|\mathsf S_{\alpha,A}^{\rm sc}\|
                         -2\mathcal L_f(A),
 \label{eq:acq-absolute-tail}
\end{equation}
positivity of $d_N$ necessarily entails
\begin{equation}
 \log N\ge (4+o(1))e^A\qquad(A\longrightarrow\infty).
 \label{eq:acq-dimension-cost}
\end{equation}
This is a cost of \eqref{eq:acq-absolute-tail}, not a lower bound on all
possible finite-history coercivity arguments.
\end{proposition}

\begin{proof}
Use $H_N\le1+\log N$, $c_{f,A}=O_f(1+\log A)$, and
$\|\mathsf S_{\alpha,A}^{\rm sc}\|\ge0$ in $d_N>0$.
The asserted asymptotic follows by partial summation of
$\sum_{n\le x,(n,f)=1}\Lambda(n)=x+o(x)$.
\end{proof}

\begin{proposition}[Finite-mode removal does not suppress prime-log resonance]
\label{prop:acq-resonance}
For every fixed $A$ and every finite Legendre projection $P_N$,
\begin{equation}
 \|Q_N\mathsf E_{\chi,A}Q_N\|
 \ge2\mathcal M_f(A),\qquad
 \mathcal M_f(A)=\sum_{\substack{n<e^{2A}\\(n,f)=1}}
 \frac{\Lambda(n)}{\sqrt n}
 \left(1-\frac{\log n}{2A}\right).
 \label{eq:acq-resonance}
\end{equation}
For fixed $f$, $\mathcal M_f(A)=2e^A/A+o(e^A/A)$.
\end{proposition}

\begin{proof}
Use $g_\tau(x)=2^{-1/2}e^{iA\tau x}$ on $(-1,1)$.  Direct integration
gives
\[
 \langle\mathsf E_{\chi,A}g_\tau,g_\tau\rangle
 =2\Re\sum_{n<e^{2A}}\frac{\Lambda(n)\chi(n)}{\sqrt n}
 \left(1-\frac{\log n}{2A}\right)e^{i\tau\log n}.
\]
The logarithms of the finitely many active primes not dividing $f$ are
rationally independent.  Their continuous torus flow therefore
approximates the phases $\overline{\chi(p)}$ simultaneously at an
unbounded sequence of times; all prime powers align at the same times.
The displayed expectation approaches $2\mathcal M_f(A)$.
Meanwhile $g_\tau\rightharpoonup0$, so $P_Ng_\tau\to0$.
Compression changes the expectation by $o(1)$ and proves the norm bound.
Partial summation gives the final asymptotic.
\end{proof}

This resonance is not a negative Weil witness.  The completed
archimedean energy also grows along the modulated sequence.  It rules
out a uniform improvement obtained merely by discarding finitely many
low modes, not a sign-sensitive comparison with that growing energy.


% =============================================================================
\section{Gaussian Weil evolution and coefficient-sensitive Pick tests}
\label{sec:gaussian-weil-pick}
% =============================================================================

The localized Schur formulation keeps support and polynomial degree
explicit.  A complementary formulation smooths the complete Weil form
rather than estimating its prime and archimedean parts separately.
It produces an exact scalar heat equation, a finite positivity
threshold, and a coefficient-sensitive continuation problem.  None of
these transformations supplies backward positivity automatically.

\subsection{A Gaussian partition of the complete form}

Retain the Fourier convention and the archimedean multiplier of
\Cref{thm:weil-explicit-form}, and abbreviate
\begin{equation}
 a_\chi(\xi)=\log\frac f\pi+
 \Re\psi\!\left(\frac{1/2+\kappa+i\xi}{2}\right),
 \qquad \lambda_n=\frac{\Lambda(n)}{\sqrt n}.
 \label{eq:gauss-arch-symbol}
\end{equation}
Here $\psi=\Gamma'/\Gamma$ and $\kappa$ is the character parity.
For $\sigma>0$, define normalized windows
\begin{equation}
 \eta_{\sigma,y}(x)=(2\pi\sigma^2)^{-1/4}
                 e^{-(x-y)^2/(4\sigma^2)}.
 \label{eq:gauss-window}
\end{equation}
They satisfy the exact overlap formulas
\begin{equation}
 \int_\R\eta_{\sigma,y}(x)^2\dd y=1,\qquad
 \int_\R\eta_{\sigma,y}(x)\eta_{\sigma,y}(x')\dd y
 =e^{-(x-x')^2/(8\sigma^2)}.
 \label{eq:gauss-overlap}
\end{equation}

\begin{theorem}[Gaussian averaging of the completed Weil form]
\label{thm:gauss-complete-average}
Set $t=(8\sigma^2)^{-1}$ and
$H_t(v)=(4\pi t)^{-1/2}e^{-v^2/(4t)}$.
For every $g\in C_c^\infty(\R)$,
\begin{equation}
 \begin{split}
 \mathcal W_{\chi,t}[g]
 &:={1\over2\pi}\int_\R|\widehat g(\xi)|^2
                         (a_\chi*H_t)(\xi)\,\dd\xi\\
 &\quad-2\Re\sum_{n\ge2}\lambda_n\chi(n)
                 e^{-t(\log n)^2}C_g(\log n)\\
 &=\int_\R\mathcal W_\chi[\eta_{\sigma,y}g]\,\dd y.
 \end{split}
 \label{eq:gauss-complete-average}
\end{equation}
For compactly supported $g$ the prime sum still contains only finitely
many active prime powers.  Moreover
$\mathcal W_{\chi,t}[g]\to\mathcal W_\chi[g]$ as $t\downarrow0$.
\end{theorem}

\begin{proof}
Completing a Gaussian square proves \eqref{eq:gauss-overlap}.
The second identity multiplies the complete spatial kernel by
$e^{-t(x-x')^2}$.  On its translation terms this gives
$e^{-t(\log n)^2}$; on the archimedean Fourier multiplier it gives
convolution by $H_t$.  Fubini is valid for compact smooth tests and the
logarithmically growing archimedean symbol.  The support assertion
follows from that of $C_g$.  Approximation by the heat identity on the
Fourier side and convergence of the finite prime sum prove the limit.
\end{proof}

For $M_vg(x)=e^{ivx}g(x)$ one also has
\begin{equation}
 \mathcal W_{\chi,t+s}[g]
 =\int_\R H_s(v)\mathcal W_{\chi,t}[M_vg]\,\dd v.
 \label{eq:gauss-dephasing}
\end{equation}
This follows from
$\widehat{M_vg}(\xi)=\widehat g(\xi-v)$ and
$C_{M_vg}(u)=e^{ivu}C_g(u)$.
It is a Gaussian mixture of unitary congruences of the whole form.
In particular, positivity propagates to larger heat times.

\subsection{The scalar symbol and the positivity threshold}

\begin{theorem}[Scalar heat-symbol formulation]
\label{thm:gauss-scalar-symbol}
For $t>0$, put
\begin{equation}
 m_{\chi,t}(\xi)=(a_\chi*H_t)(\xi)
 -2\Re\sum_{n\ge2}\lambda_n\chi(n)
                e^{-t(\log n)^2}e^{i\xi\log n}.
 \label{eq:gauss-scalar-symbol}
\end{equation}
The series and all its derivatives converge locally uniformly for
$t>0$.  The function is real analytic in $\xi$ and satisfies
\begin{equation}
 \mathcal W_{\chi,t}[g]={1\over2\pi}
       \int_\R m_{\chi,t}(\xi)|\widehat g(\xi)|^2\dd\xi,
 \qquad
 \partial_t m_{\chi,t}=\partial_\xi^2m_{\chi,t},
 \qquad m_{\chi,t+s}=H_s*m_{\chi,t}.
 \label{eq:gauss-scalar-identities}
\end{equation}
For each fixed $t$, positivity of the form on all compact smooth tests
is equivalent to $m_{\chi,t}(\xi)\ge0$ for every real $\xi$.
\end{theorem}

\begin{proof}
Gaussian decay in $\log n$ dominates every fixed power of $\log n$
and every exponential factor coming from $\xi$ in a compact complex
set.  This proves convergence and permits termwise differentiation.
The correlation identity
$C_g(u)=(2\pi)^{-1}\int|\widehat g(\xi)|^2e^{i\xi u}\dd\xi$
gives the form representation.  Each prime exponential and the
archimedean convolution obey the same heat equation, proving the
remaining identities.  A point with negative symbol is detected by
$g_R(x)=R^{-1/2}\phi(x/R)e^{i\xi_0x}$, where
$\phi\in C_c^\infty$, $\|\phi\|_2=1$: its spectral density converges
to a unit point mass at $\xi_0$.  This also proves necessity of the
pointwise inequality.
\end{proof}

\begin{theorem}[Eventual coercivity and exact heat-threshold closure]
\label{thm:gauss-threshold}
The set
\begin{equation}
 \mathfrak P_\chi=\{t>0:m_{\chi,t}(\xi)\ge0
                             \text{ for every }\xi\in\R\},\qquad
 \vartheta_\chi=\inf\mathfrak P_\chi
 \label{eq:gauss-threshold}
\end{equation}
is a nonempty upper interval.  If $\vartheta_\chi>0$, it contains its
left endpoint and the threshold symbol has minimum zero.  Furthermore
\begin{equation}
 \boxed{\mathrm{GRH}\text{ for }L(s,\chi)
                 \quad\Longleftrightarrow\quad\vartheta_\chi=0.}
 \label{eq:gauss-GRH}
\end{equation}
The principal channel requires its separate pole contributions and is
not included in this assertion.
\end{theorem}

\begin{proof}
The digamma asymptotic gives
$a_\chi(\xi)\ge\log(1+|\xi|)-C_\chi$.
An interval of length $\sqrt t$ has $H_t$ mass at most
$1/(2\sqrt\pi)$.  Thus, for $p_0=1-1/(2\sqrt\pi)>0$,
\begin{equation}
 \inf_\xi(a_\chi*H_t)(\xi)
 \ge p_0\log(1+\sqrt t/2)-C_\chi.
 \label{eq:gauss-large-time}
\end{equation}
The absolute prime bound
$R(t)=2\sum_{n\ge2}\lambda_ne^{-t(\log n)^2}$ tends to zero as
$t\to\infty$.  Therefore $m_{\chi,t}$ has a positive uniform floor
for all sufficiently large $t$.  The upper-interval property follows
from \eqref{eq:gauss-dephasing}.
On compact positive time intervals, symbol differences are uniformly
continuous in $\xi$: the gamma term has bounded second derivative,
and the differentiated prime sum is absolutely convergent.  Thus a
positive left endpoint belongs to the interval.  A positive uniform
floor there would persist at a smaller time, which is impossible.
For fixed $t$, the gamma symbol tends to infinity as $|\xi|\to\infty$
and the prime part is bounded, so the zero infimum is attained.
If GRH holds, \Cref{thm:weil-explicit-form} and
\eqref{eq:gauss-complete-average} give positivity at every $t>0$.
Conversely, threshold zero gives positivity at every $t>0$ by the
upper-interval property.  Letting $t\downarrow0$ on each fixed compact
test and using \Cref{thm:weil-explicit-form} proves GRH.
\end{proof}

The threshold $\vartheta_\chi$ belongs to the forward heat equation
for the Weil symbol.  It is not the de Bruijn--Newman parameter
$\Lambda_{\rm DN}(\chi)$ of the backward heat equation for the
completed theta transform.  Their vanishing encodes the same GRH
statement in different source geometries; no equality of their
positive values is asserted.

\begin{criterion}[Backward relative-form certificate]
\label{crit:gauss-relative-flow}
Suppose that $T$ is a time at which $\mathcal W_{\chi,T}$ is
nonnegative, and that a locally integrable real function $c$ on
$(0,T]$ satisfies
\begin{equation}
 \partial_t\mathcal W_{\chi,t}[g]
 \le c(t)\mathcal W_{\chi,t}[g]
 \quad\text{for every }g\in C_c^\infty(\R),\ 0<t\le T.
 \label{eq:gauss-relative-flow}
\end{equation}
Then $\vartheta_\chi=0$ and GRH holds for $L(s,\chi)$.
\end{criterion}

\begin{proof}
For a fixed test and $0<\varepsilon<T$, multiplication by
$\exp(-\int_\varepsilon^tc(r)\dd r)$ makes the form value nonincreasing
in $t$.  Its value at $\varepsilon$ is therefore nonnegative because
its value at $T$ is.  Let $\varepsilon\downarrow0$ and use
\Cref{thm:gauss-threshold}.
\end{proof}

The certificate is a relative inequality on the complete signed form.
Complete positivity of forward dephasing cannot supply it: the matrix
$\left(\begin{smallmatrix}1&be^{-t}\\be^{-t}&1\end{smallmatrix}\right)$,
with $b>1$, becomes positive exactly at $t=\log b$.
Eventual positivity and a positive smoothing semigroup therefore allow
a strictly positive threshold in the absence of an additional estimate.

\subsection{Signed contacts, negative mass, and the direction of evolution}

\begin{proposition}[Finite-order contacts at a positive threshold]
\label{prop:gauss-contact}
If $\vartheta_\chi>0$, the real zeros of $m_{\chi,\vartheta_\chi}$
form a nonempty finite set.  At each zero $\xi_0$, there is an integer
$r\ge1$ such that
\begin{equation}
 \partial_\xi^j m_{\chi,\vartheta_\chi}(\xi_0)=0
 \ (0\le j<2r),\qquad
 \partial_\xi^{2r}m_{\chi,\vartheta_\chi}(\xi_0)>0,
 \qquad
 \partial_t^r m_{\chi,t}(\xi_0)|_{t=\vartheta_\chi}>0.
 \label{eq:gauss-contact}
\end{equation}
\end{proposition}

\begin{proof}
The symbol is nonnegative, real analytic, and coercive at infinity.
Its zeros lie in a compact interval and cannot accumulate without the
symbol being identically zero.  Every isolated zero of a nonnegative
real analytic function has finite even order and positive first
nonzero coefficient.  Iterating the heat equation gives
$\partial_t^r m=\partial_\xi^{2r}m$.
\end{proof}

For a smooth curve of nondegenerate local minima $\xi(t)$, direct
differentiation of $m_\xi(t,\xi(t))=0$ gives
\begin{equation}
 \xi'(t)=-\frac{m_{\xi\xi\xi}}{m_{\xi\xi}},\qquad
 \frac{\dd}{\dd t}m(t,\xi(t))=m_{\xi\xi},\qquad
 \frac{\dd^2}{\dd t^2}m(t,\xi(t))
 =m_{\xi\xi\xi\xi}-\frac{m_{\xi\xi\xi}^2}{m_{\xi\xi}}.
 \label{eq:gauss-minimum-motion}
\end{equation}
These identities determine the kinematics of a contact; they do not
supply a sign for the last expression or a backward crossing exclusion.

\begin{proposition}[Dissipation of the negative part]
\label{prop:gauss-negative-entropy}
On a positive time interval, let $m=m_{\chi,t}$ and
$m_-=( -m)_+$.  For $q\ge2$,
\begin{equation}
 \frac{\dd}{\dd t}\int_\R m_-^q\dd\xi
 =-q(q-1)\int_{\{m<0\}}m_-^{q-2}|m_\xi|^2\dd\xi\le0.
 \label{eq:gauss-negative-entropy}
\end{equation}
At times when the endpoints of the negative intervals are simple,
\begin{equation}
 \frac{\dd}{\dd t}\int_\R m_-\dd\xi
 =-\sum_{\xi\in\partial\{m<0\}}|m_\xi(t,\xi)|.
 \label{eq:gauss-negative-flux}
\end{equation}
\end{proposition}

\begin{proof}
The negative set is compact, locally uniformly away from heat time
zero.  Differentiate the integral over its negative intervals, use
$m_t=m_{\xi\xi}$, and integrate by parts.  For $q\ge2$ the boundary
terms vanish because $m=0$.  For $q=1$ the endpoint velocities multiply
zero boundary values, while integrating $-m_{\xi\xi}$ gives precisely
the outgoing slope magnitudes.  The first identity also follows at
nongeneric contact times by a smooth convex approximation to the
negative-part power.
\end{proof}

Forward dissipation is consistent with disappearance of a negative
interval at a positive time.  It therefore sharpens the description of
a hypothetical obstruction but is not itself an exclusion theorem.
In particular, a small negative mass, finitely many contacts, or a
finite contact jet cannot replace the relative estimate
\eqref{eq:gauss-relative-flow}.

\subsection{The completed logarithmic derivative and the correct Pick rotation}
\label{subsec:pick-completed}

Use the completed primitive function $\mathfrak L_\chi$ from
\Cref{sec:grh-theta-heat}, and put
\begin{equation}
 F_\chi(z)=\frac{\mathfrak L_\chi'}{\mathfrak L_\chi}
                         (1/2-iz),\qquad z\in\C_+.
 \label{eq:pick-completed}
\end{equation}
This is initially meromorphic; holomorphy on all of $\C_+$ is part of
the desired conclusion, not an initial hypothesis.  The relevant
positive cone is $\Re F_\chi\ge0$, or equivalently the Pick property
of $iF_\chi$.  The unrotated function need not have an imaginary part
of one sign, even in the absolutely convergent Euler half-plane.

\begin{theorem}[Completed positive-real criterion]
\label{thm:pick-GRH}
GRH for the conjugate pair $\chi,\overline\chi$ is equivalent to the
assertion that both $F_\chi$ and $F_{\overline\chi}$ are holomorphic
on $\C_+$ and have nonnegative real part there.
\end{theorem}

\begin{proof}
Under GRH the nontrivial zeros are $\rho=1/2+i\gamma$.
In the regularized Hadamard logarithmic derivative of the completed
function they give
$im_\rho/(z+\gamma)$, with real part
$m_\rho\Im z/|z+\gamma|^2\ge0$.
The remaining constant is purely imaginary, as follows from the
functional equation on the critical line; the convergent real-part
sum is therefore nonnegative.
Conversely a zero with $\Re\rho>1/2$ produces a genuine pole at
$z=-\Im\rho+i(\Re\rho-1/2)$ in $\C_+$, contradicting holomorphy.
The functional equation transports a left-side zero to a right-side
zero in the conjugate channel.  Thus holomorphy for both channels
already forces all nontrivial zeros onto the critical line.
\end{proof}

This criterion must not be confused with the strict local Euler Schur
cells in \Cref{sec:grh-weil-schur}.  Their product has a positive local
phase measure, but the completed scattering ratio collapses to its
functional-equation constant.  Here the meromorphic logarithmic
derivative retains the zeros and hence retains the whole continuation
problem.

\subsection{Finite Pick matrices and exhaustive, not uniform, detection}

For distinct sample points $z_j\in\C_+$, set
\begin{equation}
 [\mathsf K_Z(F)]_{jk}
 =\frac{F(z_j)+\overline{F(z_k)}}{-i(z_j-\overline{z_k})}.
 \label{eq:pick-kernel}
\end{equation}
The Nevanlinna--Pick interpolation criterion says that a finite value
table admits a positive-real holomorphic interpolant on $\C_+$ exactly
when \eqref{eq:pick-kernel} is positive semidefinite.  One obtains this
form by applying the Cayley transform $(F-1)/(F+1)$ and a diagonal
congruence to the half-plane Schur kernel.  Compatibility of a finite
table does not identify its meromorphic continuation.

\begin{obstruction}[Exact finite-jet invisibility of a hidden pole]
\label{obs:pick-hidden-pole}
Fix finitely many sample points in $\Im z>1/2$, prescribed nonnegative
jet orders $r_j$, and a point $p$ with $0<\Im p<1/2$.
For any integer $L\ge0$, put
\begin{equation}
 Q(z)=\prod_j(z-z_j)^{r_j+1},\qquad
 R_p(z)=\frac{Q(z)}{(z-p)(z+i)^M},\qquad
 M>\deg Q+L+1.
 \label{eq:pick-hidden-pole}
\end{equation}
Then $R_p$ vanishes to all the prescribed orders, is $O(|z|^{-L-2})$
at infinity, and has a simple pole at $p$ and no other pole in
$\C_+$.  Thus these finite jets and asymptotic data alone cannot
exclude an interior pole.
\end{obstruction}

\begin{proof}
The factors in the numerator give the jet vanishing.  The denominator
has just one upper-half-plane zero, at which the numerator does not
vanish.  The degree inequality gives the stated decay.
\end{proof}

The rational perturbation in \eqref{eq:pick-hidden-pole} is an
obstruction to finite-data inference, not a modification of the actual
Dirichlet coefficients.

\begin{theorem}[Exhaustive Euler sampling detects every interior pole]
\label{thm:pick-exhaustive}
Enumerate a countable dense set $z_1,z_2,\ldots$ in the Euler
half-plane $\Im z>1/2$.  Let $F$ be holomorphic there and meromorphic
on $\C_+$.  If all initial Pick matrices
$\mathsf K_{\{z_1,\ldots,z_N\}}(F)$ are positive semidefinite, then
$F$ extends as a positive-real holomorphic function to all of $\C_+$.
Consequently any interior pole forces a strictly negative eigenvalue
at some finite stage.
\end{theorem}

\begin{proof}
For every $N$, finite Pick interpolation provides a positive-real
function with the first $N$ values.  Its Cayley transform belongs to
the unit ball of bounded holomorphic functions.  A subsequence
converges locally uniformly.  Its value at the first sample is fixed
and is not $1$, so its inverse Cayley transform is a holomorphic
positive-real function (the constant imaginary boundary case is
included).  It agrees with every prescribed value.  The identity
theorem gives agreement on the Euler half-plane, and uniqueness of
meromorphic continuation gives agreement on $\C_+$.  An interior
pole is impossible.  The last assertion is the contrapositive.
\end{proof}

For the arithmetic $F_\chi$, each Euler-region entry can be evaluated
from an absolutely convergent prime-power series and the gamma factor.
A negative eigenvalue is stable under sufficiently accurate interval
evaluation.  Thus an exhaustive, increasingly accurate computation
semidecides failure of the paired GRH criterion.  No fixed stage, nor
an unproved lower bound on all later stages, is a proof of GRH.  This
is the same finite/cofinal distinction as in
\Cref{sec:cofinal-mellin}, now in a coefficient-native chart.

\subsection{The exact directional Euler form and its first-prime obstruction}
\label{subsec:pick-directional}

For a column vector $\boldsymbol c$, define
\begin{equation}
 \begin{aligned}
 A(u)&=\sum_j\overline{c_j}e^{iz_ju},&
 B(u)&=\sum_j\overline{c_j}F(z_j)e^{iz_ju},\\
 E&=\int_0^\infty|A(u)|^2\dd u,&
 C(v)&=\int_0^\infty A(u+v)\overline{A(u)}\dd u.
 \end{aligned}
 \label{eq:pick-directional-readers}
\end{equation}
The coefficient conjugation here is fixed by the convention
$\boldsymbol c^*\mathsf K_Z(F)\boldsymbol c$.
The elementary Laplace identity for the denominator in
\eqref{eq:pick-kernel} gives
\begin{equation}
 \boldsymbol c^*\mathsf K_Z(F)\boldsymbol c
 =2\Re\int_0^\infty B(u)\overline{A(u)}\dd u.
 \label{eq:pick-directional-identity}
\end{equation}

\begin{theorem}[Completed Euler shift formula]
\label{thm:pick-euler-shifts}
If $\min_j\Im z_j>1/2$, then
\begin{equation}
 \begin{split}
 \boldsymbol c^*\mathsf K_Z(F_\chi)\boldsymbol c
 ={}&(\log(f/\pi)-\gamma_{\rm E})E\\
 &+2\int_0^\infty
 \frac{e^{-2v}E-e^{-(1/2+\kappa)v}\Re C(v)}{1-e^{-2v}}\dd v\\
 &-2\Re\sum_p\sum_{a\ge1}
 \frac{(\log p)\chi(p)^a}{p^{a/2}}C(a\log p).
 \end{split}
 \label{eq:pick-euler-shifts}
\end{equation}
The paired numerator of the integral must be retained at $v=0$.
The higher-prime-power function
\begin{equation}
 P_\chi^{(\ge2)}(z)=\sum_p\sum_{a\ge2}
       \frac{(\log p)\chi(p)^a}{p^{a/2}}e^{iza\log p}
 \label{eq:pick-higher-powers}
\end{equation}
is holomorphic throughout $\C_+$.  The gamma term is also holomorphic
there.  The only Euler layer whose absolute convergence stops at
$\Im z=1/2$ is
\begin{equation}
 P_\chi^{(1)}(z)=\sum_p\frac{(\log p)\chi(p)}{\sqrt p}
                                      e^{iz\log p}.
 \label{eq:pick-first-prime}
\end{equation}
Its meromorphic continuation has precisely the possible nontrivial-zero
poles of $-L'/L(1/2-iz,\chi)$ in $\C_+$.
\end{theorem}

\begin{proof}
In the Euler half-plane,
\[
 F_\chi(z)=\frac12\log\frac f\pi
 +\frac12\psi\!\left(\frac{1/2+\kappa-iz}{2}\right)
 -\sum_p\sum_{a\ge1}\frac{(\log p)\chi(p)^a}{p^{a/2}}
                                      e^{iza\log p}.
\]
The digamma integral equals
$\tfrac12\log(f/\pi)-\tfrac12\gamma_{\rm E}
+\int_0^\infty[e^{-2v}-e^{-(1/2+\kappa)v}e^{izv}]/(1-e^{-2v})\dd v$.
Substitution in \eqref{eq:pick-directional-identity} proves
\eqref{eq:pick-euler-shifts}.  Since $C(v)=E+O(v)$, the integral is
regular at zero.  If $y_* =\min_j\Im z_j>0$, then
$|C(v)|\le(\sum_j|c_j|)^2e^{-y_*v}/(2y_*)$.
This proves absolute convergence of the prime sum when $y_*>1/2$.
For $a\ge2$ the scalar majorant
$\sum_p\sum_{a\ge2}(\log p)p^{-a(1/2+\varepsilon)}$
converges for every $\varepsilon>0$, giving normal convergence on
$\Im z\ge\varepsilon$.  The digamma argument has positive real
part on $\C_+$.  Finally
$P_\chi^{(1)}=-L'/L(1/2-iz,\chi)-P_\chi^{(\ge2)}$
on the Euler half-plane and hence by meromorphic continuation.
The absolute first-prime majorant diverges when
$1/2+\Im z\le1$, by the prime number theorem.
\end{proof}

The first-prime localization is an analytic reduction, not an estimate
of its sign.  It identifies which coefficients must be controlled
while treating the gamma and higher-power layers without a false
continuation assumption.  In particular, positive finite
interpolation matrices, a small absolute tail, or a norm comparison
chosen after seeing a protected direction cannot replace a uniform
coefficient-sensitive lower bound on
\eqref{eq:pick-euler-shifts}.


% =============================================================================
\section{Admissible spectral filters and cofinal compensation}
\label{sec:admissible-spectral-filters}
% =============================================================================

The completed Euler and Pick identities distinguish a finite spectral
cancellation from an arithmetic positivity argument.  The same distinction
can be stated directly for finite radial filters.  A prescribed finite set
of poles, gamma nodes, or zero nodes can be interpolated exactly while
retaining positivity beyond a support wall.  What is not supplied by that
interpolation is a small normalized contribution from the unprescribed
spectrum.  The resulting compensation inequality gives a precise closure
criterion and a way to test proposed filter architectures before attempting
a cofinal limit.

\subsection{Arithmetic admission and the completed compensation identity}
\label{subsec:filter-compensation}

For a finite real exponential sum and an integer $k\ge1$, write
\begin{equation}
 \begin{split}
 W(u)&=\sum_{j=0}^J b_je^{-a_ju},\qquad b_j\in\R,\quad a_j\ge0,\\
 \mathcal L_y^{(k)}[W]
 &=\frac1{(k-1)!}\int_0^\infty u^{k-1}e^{-yu}W(u)\dd u
   =\sum_{j=0}^J\frac{b_j}{(y+a_j)^k},\qquad \Re y>0.
 \end{split}
 \label{eq:filter-response}
\end{equation}
The rational expression gives its continuation wherever its denominators
are nonzero.  For a modulus $f$, call $W$ \emph{arithmetically admitted} if
\begin{equation}
 W(\ell\log p)\ge0\qquad(p\nmid f,\ \ell\ge1).
 \label{eq:filter-arithmetic-admission}
\end{equation}
A convenient stronger sufficient condition is $W(u)\ge0$ for all $u\ge L$,
where $L\le\log\min\{p:p\nmid f\}$; the universal choice $L=\log2$
suffices.  Strict positivity on that half-line is stronger again.  A zero
at a nonarithmetic point contradicts only strict half-line positivity; a
negative interval contradicts nonnegative half-line positivity, but need
not by itself contradict \eqref{eq:filter-arithmetic-admission}.

Retain the full prime-power series, rather than removing its higher-power
terms.  For a character $\psi$ modulo $f$, define
\begin{equation}
 \mathfrak E_{\psi,k}^W(s)
 =\sum_{p\nmid f}\sum_{\ell\ge1}
 \frac{(\log p)(\ell\log p)^{k-1}}{(k-1)!}
 \psi(p)^\ell p^{-\ell s}W(\ell\log p),\qquad \Re s>1.
 \label{eq:filter-complete-euler}
\end{equation}
This is the finite radial combination of the normalized derivatives of
$-L'/L$ appearing in the completed Euler chart.  Let $\chi$ be primitive
of conductor $f$, and let $\chi_0$ be the principal character modulo $f$.
Fix a hypothetical zero $\rho_*=\beta+i\gamma$ with $\beta>1/2$,
put $\sigma>1$, $d=\sigma-1$, and $x=\sigma-\beta>d$.
For an admitted filter,
\begin{equation}
 \begin{split}
 \mathfrak S_{\chi,k}^W
 &=\mathfrak E_{\chi_0,k}^W(\sigma)
      +\Re\mathfrak E_{\chi,k}^W(\sigma+i\gamma)\ge0.
 \end{split}
 \label{eq:filter-two-channel-positive}
\end{equation}
Indeed, each prime-power coefficient is multiplied by
$1+\Re(\chi(p)^\ell e^{-i\gamma\ell\log p})\ge0$.

\begin{theorem}[Completed spectral compensation]
\label{thm:filter-compensation}
Assume $k\ge2$, or assume $k=1$ and $W(0)=\sum_jb_j=0$.
Suppose $W$ is arithmetically admitted and
\begin{equation}
 \mathcal L_d^{(k)}[W]=0,\qquad
 T=\mathcal L_x^{(k)}[W]>0.
 \label{eq:filter-normalized-target}
\end{equation}
Then the completed identity has the exact form
\begin{equation}
 \mathfrak S_{\chi,k}^W=-m_{\rho_*}T+
       \mathfrak R_{\chi,k}^W(\rho_*),\qquad
 \mathfrak R_{\chi,k}^W(\rho_*)\ge m_{\rho_*}T>0,
 \label{eq:filter-compensation}
\end{equation}
where $m_{\rho_*}$ is the zero multiplicity and $\mathfrak R$ is the
actual sum of the remaining completed contributions.  In particular it
is not an independently chosen absolute majorant.
\end{theorem}

\begin{proof}
For completeness, put $F(y)=\mathcal L_y^{(k)}[W]$, let
$\kappa_\chi\in\{0,1\}$ be the parity, and set $\delta_\chi=1$ for
the primitive principal character and $0$ otherwise.  Let $Z_\zeta$ and $Z_\chi$ denote the sets of distinct nontrivial
zeros, with multiplicities $m_\rho$.  An explicit
expression for the remainder is
\begin{align}
 \mathfrak R={}&F(\sigma)
 +\delta_\chi\Re\{F(\sigma+i\gamma)+F(d+i\gamma)\}
 -\sum_{\rho\in Z_\zeta}m_\rho F(\sigma-\rho)\notag\\
 &-\Re\sum_{\substack{\rho\in Z_\chi\\\rho\ne\rho_*}}
          m_\rho F(\sigma+i\gamma-\rho)
 -\sum_{m\ge0}F(\sigma+2m)\notag\\
 &-\Re\sum_{m\ge0}F(\sigma+\kappa_\chi+2m+i\gamma)
 -\mathfrak B_{f,k}^W(\sigma),
 \label{eq:filter-actual-remainder}
\end{align}
where the principal-character Euler correction is
\[
 \mathfrak B_{f,k}^W(\sigma)=
 \sum_{p\mid f}\sum_{\ell\ge1}
 \frac{(\log p)(\ell\log p)^{k-1}}{(k-1)!}
 p^{-\ell\sigma}W(\ell\log p).
\]
This follows by logarithmic differentiation of the completed functions,
with the factors $s(s-1)$ included in the zeta completion.  The pole
at $s=1$ in the unshifted principal channel contributes $F(d)$, which
vanishes.  The selected zero contributes $-m_{\rho_*}F(x)$.  The pole
at zero contributes $F(\sigma)$, and the gamma logarithmic derivatives
supply the two displayed negative ladder sums.  The factor
$\prod_{p\mid f}(1-p^{-s})$ supplies $-\mathfrak B$.

For $k\ge2$ the normalized derivative kills all constant terms.  For
$k=1$, $\sum_jb_j=0$ kills them in the radial combination and gives
$F(y)=O(|y|^{-2})$ at infinity.  In either case the zero and gamma sums
in \eqref{eq:filter-actual-remainder} converge absolutely, so the stated
regrouping is legitimate.  Equation \eqref{eq:filter-two-channel-positive}
then proves the lower bound in \eqref{eq:filter-compensation}.
\end{proof}

\begin{criterion}[Normalized cofinal compensation exclusion]
\label{crit:filter-cofinal-compensation}
Suppose that for every hypothetical off-line zero one can construct an
arithmetically admitted sequence satisfying the convergence hypotheses and
normalization of \Cref{thm:filter-compensation}, for which
\begin{equation}
 \limsup_{\nu\to\infty}
 \frac{\mathfrak R_{\chi,k_\nu}^{W_\nu}(\rho_*)}
      {\mathcal L_{x_\nu}^{(k_\nu)}[W_\nu]}<m_{\rho_*}.
 \label{eq:filter-cofinal-gap}
\end{equation}
Then that zero is impossible.  Establishing this condition for all
primitive characters excludes all zeros to the right of $1/2$ and,
by the functional equation, proves GRH.
\end{criterion}

\begin{proof}
Every quotient in \eqref{eq:filter-cofinal-gap} is at least
$m_{\rho_*}$ by \eqref{eq:filter-compensation}, a contradiction.
\end{proof}

This is a spectral closure criterion in the completed Euler entrance,
not an additional independent entrance.  When the filters are to be
implemented as renewal sources, their coefficients, acquisition costs,
and source declaration must also satisfy the finite-source conventions
of \Cref{sec:gran-sources,sec:gran-refinement}.  Finite interpolation
alone supplies none of those quantitative controls.

\subsection{Finite interpolation and the normalization cost}
\label{subsec:filter-interpolation-cost}

\begin{theorem}[Finite spectral interpolation with a positive tail]
\label{thm:filter-tail-interpolation}
Let $L>0$.  Prescribe finitely many complex pairs $(y,n)$ with
$\Re y>0$, $n\in\N\cup\{0\}$, and compatible target values
$c_{y,n}$.  The table is closed under conjugation, with conjugate values
at conjugate nodes and real values at real nodes; repeated identical
pairs have the same value.  There is a real polynomial $P$ such that
$W(u)=P(e^{-u})$ satisfies
\begin{equation}
 \frac1{n!}\int_0^\infty u^ne^{-yu}W(u)\dd u=c_{y,n}
 \quad\hbox{at every prescribed pair},\qquad
 W(u)>0\quad(u\ge L).
 \label{eq:filter-tail-interpolation}
\end{equation}
No bound on the degree or coefficient norm is asserted.
\end{theorem}

\begin{proof}
The distinct exponential-polynomial kernels $u^ne^{-yu}$ are linearly
independent on every nonempty interval.  This follows, for example, by
applying the differential annihilators for all exponents but one and
then using independence of the surviving monomials.  Their real and
imaginary parts, after imposing conjugation compatibility, therefore
induce a surjective moment map on $C_c(0,L;\R)$: otherwise a nonzero
linear combination of the kernels would annihilate every such function.

Choose $\varepsilon>0$.  Surjectivity gives $h\in C_c(0,L;\R)$ such
that $w=\varepsilon+h$ has all the prescribed moments.  Under
$q=e^{-u}$, the function $w$ extends continuously to $[0,1]$ and equals
$\varepsilon$ near both endpoints and on $[0,e^{-L}]$.
Polynomials approximate it uniformly.  The moment functionals are
continuous in this norm, since their absolute integrals are finite.
Moreover their restriction to polynomials has full rank: a linear
functional annihilating all polynomials would annihilate all continuous
functions by density and hence the original moment map.  Select a fixed
finite polynomial basis on which the moment map is invertible.
Correct the small moment error of a sufficiently accurate polynomial
approximant within that basis.  The correction is uniformly small, so
the resulting polynomial has the exact moments and remains larger than
$\varepsilon/2$ on $[0,e^{-L}]$.
\end{proof}

Consequently any finite collection of distinct pole, gamma, and zero
nodes can be cancelled while a distinct target response is set to one;
a fixed finite block of derivative orders can be prescribed as well.
This is a finite existence statement.  At order one the endpoint
condition $W(0)=0$ must also be retained when invoking the completed
identity.  Under a hypothetical zero,
\Cref{thm:filter-compensation} forces compensation into the remaining
unprescribed terms rather than making those terms small.

\begin{proposition}[Lower bounds forced by target normalization]
\label{prop:filter-normalization-cost}
Let $0<d<x$, let $W(u)\ge0$ for $u\ge L$, and suppose
\eqref{eq:filter-normalized-target} holds.  Then
\begin{equation}
 T\le\frac1{(k-1)!}\int_0^L u^{k-1}e^{-du}|W(u)|\dd u,
 \qquad
 \sup_{0\le u\le L}|W(u)|\ge\frac{k!T}{L^k}.
 \label{eq:filter-wall-cost}
\end{equation}
For the expansion \eqref{eq:filter-response}, this implies
$\sum_j|b_j|\ge k!T/L^k$.  In addition, for every $R>d$,
\begin{equation}
 \mathcal C_k(W;R):=\sum_j\frac{|b_j|}{(R+a_j)^k}
 \ge e^{-(R-d)L}T.
 \label{eq:filter-cauchy-cost}
\end{equation}
If all shifts satisfy $a_j\le A_k$, then also
\begin{equation}
 \mathcal C_k(W;R)\ge\frac{k!T}{L^k(R+A_k)^k}.
 \label{eq:filter-bounded-shift-cost}
\end{equation}
\end{proposition}

\begin{proof}
Subtract $e^{-(x-d)L}\mathcal L_d^{(k)}[W]=0$ from the target:
\[
 T=\frac1{(k-1)!}\int_0^\infty
 u^{k-1}e^{-du}\{e^{-(x-d)u}-e^{-(x-d)L}\}W(u)\dd u.
\]
The integrand on $[L,\infty)$ is nonpositive.  On $[0,L]$ the
braced factor lies between zero and one, proving the first inequality
in \eqref{eq:filter-wall-cost}.  Integrating $u^{k-1}$ proves the
second, and $|W(u)|\le\sum_j|b_j|$ proves the coefficient bound.
Finally
\[
 \frac1{(k-1)!}\int_0^L u^{k-1}e^{-(d+a_j)u}\dd u
 \le \frac{e^{(R-d)L}}{(R+a_j)^k}.
\]
Summing this inequality against $|b_j|$ proves
\eqref{eq:filter-cauchy-cost}; the shift cap gives
\eqref{eq:filter-bounded-shift-cost} directly from the coefficient bound.
\end{proof}

The cost is relative to $T$, so rescaling all coefficients cannot
remove it.  In particular, when $R$ and $L$ are fixed and $A_k=o(k)$,
the lower bound in \eqref{eq:filter-bounded-shift-cost} divided by $T$
grows faster than any fixed exponential.  Absolute gamma or Cauchy
bounds must retain this coefficient cost.  Their application to a
finite interpolant is not a proof of the signed estimate
\eqref{eq:filter-cofinal-gap}.

\subsection{A normalized equispaced gamma-prefix filter}
\label{subsec:filter-equispaced}

The finite interpolation theorem does not choose a stable family.
An explicit family exposes what must be controlled.  Fix
$1/2<\beta<\eta<1$, $\gamma\ne0$, $0<\theta<1$, and a parity
$\kappa_\chi\in\{0,1\}$.  For $M\ge1$, $d>0$, put
\begin{equation}
 N=3M+4,\qquad n=N-1,\qquad
 \varepsilon=\frac{\theta d}{N},\qquad
 \sigma=d+1,\quad x_d=\sigma-\beta,\quad\xi_d=\sigma-\eta.
 \label{eq:filter-equispaced-parameters}
\end{equation}
The three ladders correspond to the principal gamma sequence and the
two conjugate target gamma sequences.  Define
\begin{equation}
 \begin{split}
 P_{M,d}(y)={}&(y-d)(y-\xi_d)
 \prod_{m=0}^{M-1}(y-\sigma-2m)\\
 &\hspace{4mm}\cdot\prod_{m=0}^{M-1}
       \{(y-\sigma-\kappa_\chi-2m)^2+\gamma^2\},\\
 D_{M,d}(y)={}&\prod_{j=0}^{n}(y+j\varepsilon),\qquad
 F_{M,d}(y)=\frac{D_{M,d}(x_d)}{P_{M,d}(x_d)}
                     \frac{P_{M,d}(y)}{D_{M,d}(y)}.
 \end{split}
 \label{eq:filter-equispaced-rational}
\end{equation}
Here $\deg P_{M,d}=n-1$ and $\deg D_{M,d}=n+1$.
All factors at $x_d$ are nonzero.

\begin{proposition}[Exact residues and operator representation]
\label{prop:filter-equispaced-residues}
There are real residues $b_j$ such that
\begin{equation}
 F_{M,d}(y)=\sum_{j=0}^{n}\frac{b_j}{y+j\varepsilon},\qquad
 K_{M,d}(u)=\sum_{j=0}^{n}b_je^{-j\varepsilon u}.
 \label{eq:filter-equispaced-laplace}
\end{equation}
They satisfy $b_0>0$, $\sum_jb_j=0$, and
\begin{equation}
 \frac{K_{M,d}(u)}{b_0}=H_{M,d}(e^{-\varepsilon u}),\qquad
 H_{M,d}(q)=\sum_{j=0}^{n}(-1)^jc_jq^j,\qquad
 c_j=\binom nj\frac{P_{M,d}(-j\varepsilon)}{P_{M,d}(0)}>0.
 \label{eq:filter-alternating-polynomial}
\end{equation}
With $\lambda$ ranging over the roots of $P_{M,d}$, including
multiplicities,
\begin{equation}
 \begin{split}
 H_{M,d}(q)
 &=\prod_\lambda\left(1+\frac\varepsilon\lambda qD_q\right)(1-q)^n,\\
 \frac{K_{M,d}(u)}{b_0}
 &=\prod_\lambda\left(1-\frac1\lambda\partial_u\right)
                       (1-e^{-\varepsilon u})^n.
 \end{split}
 \label{eq:filter-equispaced-operators}
\end{equation}
The response cancels $d$, $\xi_d$, and the first $M$ nodes in each of
the three gamma ladders, while $F_{M,d}(x_d)=1$.
\end{proposition}

\begin{proof}
The poles $-j\varepsilon$ are simple and
\[
 D_{M,d}'(-j\varepsilon)
   =(-1)^j\varepsilon^n j!(n-j)!.
\]
Partial fractions yield \eqref{eq:filter-alternating-polynomial}.
Both $P_{M,d}(0)$ and $P_{M,d}(x_d)$ have sign $(-1)^M$,
so $b_0>0$.  The decay $F_{M,d}(y)=O(y^{-2})$ gives
$\sum_jb_j=0$.  Every real root of $P_{M,d}$ is positive and each
conjugate pair has positive real part, proving $c_j>0$.
Multiplication of the $j$th coefficient by $1+j\varepsilon/\lambda$
is $1+(\varepsilon/\lambda)qD_q$; under $q=e^{-\varepsilon u}$
this becomes $1-\lambda^{-1}\partial_u$.  The zero and target
identities follow directly from \eqref{eq:filter-equispaced-rational}.
\end{proof}

\begin{proposition}[A sufficient eventual-positivity wall]
\label{prop:filter-positive-wall}
For the family above,
\begin{equation}
 \frac{c_{j+1}}{c_j}\le e^\theta\frac{n-j}{j+1}\le e^\theta n.
 \label{eq:filter-coefficient-ratio}
\end{equation}
Consequently
\begin{equation}
 K_{M,d}(u)>0\quad\hbox{if}\quad
 u>\frac{\log n+\theta}{\varepsilon}.
 \label{eq:filter-positive-wall}
\end{equation}
In particular $d>N(\log n+\theta)/(\theta L)$ is sufficient for
strict positivity on $[L,\infty)$.  This is not a necessary condition.
\end{proposition}

\begin{proof}
All roots $\lambda$ have $\Re\lambda\ge d$.  Along the negative
real axis, each real factor contributes at most $\varepsilon/d$
to a logarithmic coefficient ratio, and each conjugate quadratic
contributes at most $2\varepsilon/d$.  Their total is at most
$(n-1)\varepsilon/d<\theta$, proving
\eqref{eq:filter-coefficient-ratio}.  When
$q<e^{-\theta}/n$, the terms $c_jq^j$ strictly decrease, so grouping
successive terms of the alternating polynomial proves positivity.
\end{proof}

Thus the rational cancellations are exact for every $M,d$, but admission
is a separate question.  The large-line sufficient condition above
cannot be reversed to decide the critical scaling $d\asymp\sqrt M$.

\subsection{Last-zero bounds and the role of real factors}
\label{subsec:filter-last-zero}

\begin{lemma}[A positive real factor moves the last zero outward]
\label{lem:filter-real-factor-zero}
Let $f$ be a nonzero real finite exponential sum with
$\lim_{u\to\infty}f(u)>0$ and last real zero $U$.  For $\lambda>0$,
$g=(1-\lambda^{-1}\partial_u)f$ is negative immediately to the right
of $U$ and has a sign-changing zero strictly larger than $U$.
\end{lemma}

\begin{proof}
Analyticity gives $f(u)=c(u-U)^r+O((u-U)^{r+1})$ near $U$,
where $r\ge1$ and $c>0$, because $f$ is positive to the right of its
last zero.  Hence $g(u)=-(cr/\lambda)(u-U)^{r-1}+O((u-U)^r)<0$
just to the right of $U$, whereas $g$ tends to the same positive limit
as $f$.  It therefore has a later sign-changing zero.  A nonzero real
finite exponential sum has only finitely many real zeros, by induction
and Rolle's theorem after removal of its smallest exponential factor.
\end{proof}

\begin{corollary}[An all-real lower bound]
\label{cor:filter-all-real-last-zero}
For $E(u)=(1-e^{-\varepsilon u})^n$, $n\ge1$, and
$\lambda_1,\ldots,\lambda_r>0$, $r\ge1$, the last zero $U$ of
$\prod_{i=1}^r(1-\lambda_i^{-1}\partial_u)E$ satisfies
\begin{equation}
 U\ge\max_i\frac1\varepsilon\log\left(1+\frac{n\varepsilon}{\lambda_i}\right)
 =\frac1\varepsilon\log\left(1+\frac{n\varepsilon}{\lambda_{\min}}\right).
 \label{eq:filter-all-real-last-zero}
\end{equation}
If $n\sim RM$, $d\sim\varkappa\sqrt M$,
$\varepsilon\sim\theta\varkappa/(R\sqrt M)$, and
$\lambda_{\min}\sim d$, then
\begin{equation}
 U\ge\left(\frac{R\log(1+\theta)}{\theta\varkappa}+o(1)\right)\sqrt M.
 \label{eq:filter-all-real-critical}
\end{equation}
\end{corollary}

\begin{proof}
The operators commute.  Apply the $i$th factor first; its last zero is
$\varepsilon^{-1}\log(1+n\varepsilon/\lambda_i)$ by direct
differentiation of $E$.  Every remaining real factor moves the last
zero outward by \Cref{lem:filter-real-factor-zero}.  The scaling follows
by substitution.
\end{proof}

Removing the pole factor $\lambda=d$ does not remove this obstruction
within the all-real class when the next real root is still asymptotic
to $d$.  This statement is not a theorem for the full conjugate-pair
filter: second-order conjugate factors need not preserve the same
last-zero ordering.

\begin{proposition}[An exact local test for a conjugate factor]
\label{prop:filter-first-pair}
Let a real polynomial $g$ have constant term one and smallest positive
simple root $r$, with $g'(r)<0$.  For $\lambda=\omega+i\gamma$,
$\omega>0$, and $a=\varepsilon/\lambda$, put
\[
 \Phi=(1+a qD_q)(1+\overline a qD_q).
\]
If
\begin{equation}
 \frac{r g''(r)}{|g'(r)|}<\frac{2\omega}{\varepsilon}+1,
 \label{eq:filter-pair-local-test}
\end{equation}
then $\Phi g$ has a positive root smaller than $r$.
For the pole-first polynomial
\[
 g_0(q)=(1-q)^{n-1}(1-q/q_1),\qquad
 q_1=(1+n\varepsilon/d)^{-1},
\]
this condition holds whenever $\omega>d(n-1)/n$.
\end{proposition}

\begin{proof}
At $r$,
\[
 (\Phi g)(r)=(2\Re a+|a|^2)r g'(r)+|a|^2r^2g''(r).
\]
Condition \eqref{eq:filter-pair-local-test} makes this value negative,
whereas $(\Phi g)(0)=1$.  For $g_0$ the left side is
$2(n-1)q_1/(1-q_1)=2d(n-1)/(n\varepsilon)$, giving the stated
sufficient condition.  Later conjugate factors require their own
curvature bounds; the first-step calculation cannot be iterated
without them.
\end{proof}

A different lower bound uses the product of roots and does not assume
that all roots are real.

\begin{theorem}[Root-product lower bound with exceptional roots]
\label{thm:filter-root-product}
For $H_{M,d}$ in \eqref{eq:filter-alternating-polynomial}, let $k'$ be
the number of roots, with multiplicity, outside $(0,1)$, and suppose
$k'<n$.  Let $U$ be the largest positive zero of $K_{M,d}$.  If
$S_{\rm real}=\sum_{\lambda\in\R}1/\lambda$, where the sum is over
the real roots of $P_{M,d}$, then
\begin{equation}
 U\ge(1-\theta/2)S_{\rm real}
       -\frac{k'(\theta+\log n)}{(n-k')\varepsilon}.
 \label{eq:filter-root-product-bound}
\end{equation}
In particular, at $d\sim\varkappa\sqrt M$, a bound
$k'=o(\sqrt M)$ would imply $U\to\infty$ and exclude strict
positivity on any fixed half-line.
\end{theorem}

\begin{proof}
The positive-coefficient polynomial $H_{M,d}(-q)$ has all roots of
modulus at least
$\rho_0=\min_{j<n}c_j/c_{j+1}\ge e^{-\theta}/n$ by the elementary
coefficient annulus bound.  For a direct proof, write
$H_{M,d}(-\rho_0 z)=\sum_j d_jz^j$, where
$d_0\ge d_1\ge\cdots\ge d_n>0$.  Multiplication by $1-z$
gives
\[
 d_0-\sum_{j=1}^n(d_{j-1}-d_j)z^j-d_nz^{n+1}.
\]
For $|z|<1$, the modulus of the sum subtracted from $d_0$ is
strictly smaller than $d_0$, so no root lies in that disk.
The product of all root moduli is $1/c_n$.  If $q_*$ is the smallest
root in $(0,1)$, then
\[
 q_*^{n-k'}\rho_0^{k'}\le\frac1{c_n},\qquad
 U=-\varepsilon^{-1}\log q_*
 \ge\frac{\log c_n+k'\log\rho_0}{(n-k')\varepsilon}.
\]
Every real factor contributes $\log(1+n\varepsilon/\lambda)$ to
$\log c_n$.  Each conjugate pair contributes a nonnegative logarithm.
Since $n\varepsilon/\lambda<\theta$ for real factors and
$\log(1+z)\ge(1-\theta/2)z$ on $[0,\theta]$,
\[
 \log c_n\ge(1-\theta/2)n\varepsilon S_{\rm real}.
\]
Substitution and $n/(n-k')\ge1$ give
\eqref{eq:filter-root-product-bound}.  Finally
\[
 S_{\rm real}\ge\sum_{m=0}^{M-1}\frac1{d+1+2m}
 \ge\frac12\log\left(1+\frac{2M}{d+1}\right).
\]
At the critical scaling this is of logarithmic size, whereas the
subtracted error is $o(\log M)$ under $k'=o(\sqrt M)$.
\end{proof}

The hypothesis on $k'$ is a separate root-count estimate; it is not
inferred from finite numerical samples.  Moreover a last zero alone
excludes strict positivity, not the discrete admission condition
\eqref{eq:filter-arithmetic-admission}.

\subsection{The finite beta kernel and the growing-order transfer}
\label{subsec:filter-beta-transfer}

The exact kernel underlying the filter is
\begin{equation}
 E_M(u)=(1-e^{-\varepsilon u})^n,\qquad
 V_M(u)=e^{-du}E_M(u),\qquad u\ge0.
 \label{eq:filter-beta-kernel}
\end{equation}
It is useful to separate its fixed-order central limit from the
increasing-order conjugate differential product.

\begin{proposition}[Exact beta-kernel geometry]
\label{prop:filter-beta-geometry}
The unique maximum of $V_M$ is at
\begin{equation}
 u_1=\frac1\varepsilon\log\left(1+\frac{n\varepsilon}{d}\right),
 \qquad q_1=e^{-\varepsilon u_1}.
 \label{eq:filter-beta-maximum}
\end{equation}
At this point,
\begin{equation}
 \begin{split}
 (\log V_M)''(u_1)&=-d(d/n+\varepsilon),\\
 (\log V_M)'''(u_1)&=d(1+q_1)(d/n+\varepsilon)^2.
 \end{split}
 \label{eq:filter-beta-derivatives}
\end{equation}
If $d/\sqrt M\to\varkappa>0$ and
\eqref{eq:filter-equispaced-parameters} holds, then
\begin{equation}
 s_M^2=\{d(d/n+\varepsilon)\}^{-1}
 \longrightarrow\frac{3}{\varkappa^2(1+\theta)},\qquad
 \frac{V_M(u_1+s_Mz)}{V_M(u_1)}\longrightarrow e^{-z^2/2}
 \label{eq:filter-beta-gaussian}
\end{equation}
locally uniformly with every fixed number of derivatives.
\end{proposition}

\begin{proof}
Differentiate $-du+n\log(1-e^{-\varepsilon u})$.
The first derivative vanishes exactly when
$n\varepsilon q/(1-q)=d$, giving \eqref{eq:filter-beta-maximum};
further differentiation gives \eqref{eq:filter-beta-derivatives}.
At critical scaling, $n\asymp M$, $d\asymp\sqrt M$,
$\varepsilon\asymp M^{-1/2}$, and $q_1\to(1+\theta)^{-1}<1$.
The logarithmic derivative of each fixed order $r\ge3$ on a fixed
neighborhood of $u_1$ is $O(M^{1-r/2})$.  Taylor expansion gives
\eqref{eq:filter-beta-gaussian} and its fixed-order differentiated
forms.  In particular the third derivative is $O(M^{-1/2})$;
no increasing-order derivative estimate is implied.
\end{proof}

Let $\omega_m=d+1+\kappa_\chi+2m$ and
$a_m=\omega_m-d=1+\kappa_\chi+2m$.  The conjugate-pair transform is
\begin{equation}
 \begin{split}
 \mathscr W_M(u)
 &=e^{-du}\prod_{m=0}^{M-1}
    \frac{(\partial_u-\omega_m)^2+\gamma^2}{\omega_m^2+\gamma^2}
           E_M(u)\\
 &=\prod_{m=0}^{M-1}
    \frac{(\partial_u-a_m)^2+\gamma^2}{(d+a_m)^2+\gamma^2}
           V_M(u).
 \end{split}
 \label{eq:filter-pair-beta-transform}
\end{equation}
The normalized pole-and-pairs subfilter, before the remaining real
factors are applied, is exactly
\begin{equation}
 \mathscr K_M^{\rm pp}(u)=-\frac{e^{du}}d\mathscr W_M'(u).
 \label{eq:filter-pole-pairs-derivative}
\end{equation}
Indeed, applying $1-d^{-1}\partial_u$ to
$e^{du}\mathscr W_M$ gives this expression.  The full filter is obtained
by subsequently applying the real factors at $\xi_d$ and
$\sigma,\sigma+2,\ldots,\sigma+2(M-1)$.

For an exact Gaussian $e^{-z^2/2}$ in the variable $z=(u-u_1)/s_M$,
each numerator in \eqref{eq:filter-pair-beta-transform}, after
multiplication by $s_M^2$, is
$(D_z-a_ms_M)^2+(\gamma s_M)^2$.  Its Gaussian symbol is therefore
precisely of the type in \Cref{subsec:heatblock-growing-symbol}, with
equispaced positive centers and equal heights.  This is an exact
identification for a Gaussian input.  It does not identify the output
of the actual beta kernel when $2M$ derivatives are taken.

\begin{obstruction}[No fixed-width finite-measure Gaussian mixture]
\label{obs:filter-gaussian-mixture}
Extend $V_M$ by zero to $u<0$.  With the Fourier convention
$\widehat V_M(\omega)=\int_\R V_M(u)e^{-i\omega u}\dd u$,
\begin{equation}
 \widehat V_M(\omega)
 =\frac1\varepsilon
   B\!\left(\frac{d+i\omega}{\varepsilon},n+1\right)
 =\frac{n!\,\varepsilon^n}
        {\prod_{j=0}^{n}(d+j\varepsilon+i\omega)}.
 \label{eq:filter-beta-fourier}
\end{equation}
For any fixed $s>0$, $V_M$ is not the convolution of a Gaussian of
variance $s^2$ with a finite measure, even a finite signed or complex
measure of bounded total variation.
\end{obstruction}

\begin{proof}
The change of variable $q=e^{-\varepsilon u}$ gives the beta integral;
the integer gamma-ratio formula gives the rational expression.
Thus $|\widehat V_M(\omega)|$ has nonzero polynomial order
$|\omega|^{-n-1}$ at infinity.  A convolution with a fixed-width
Gaussian and a finite measure would instead have Fourier transform
bounded in modulus by a constant times $e^{-s^2\omega^2/2}$.
\end{proof}

The obstruction concerns this specific mixture representation, not all
possible comparison kernels.  A coefficient comparison must also retain
its cancellation scale.  If a polynomial
$\sum_j(-1)^jA_jq^j$, $A_j>0$, is approximated by replacing $A_j$ with
$A_j(1+r_j)$, the resulting error is bounded by
\begin{equation}
 \sup_j|r_j|\sum_jA_jq^j,
 \label{eq:filter-cancellation-cost}
\end{equation}
not by $\sup_j|r_j|$ times the signed polynomial.  Already the binomial
example has relative signed size
$|(1-q)^n|/(1+q)^n=((1-q)/(1+q))^n$ for $0<q<1$.
A small coefficientwise error can therefore dominate the signed value.
Neither a Stirling approximation to the pair coefficients nor
\eqref{eq:filter-beta-gaussian} alone determines its sign.

\begin{criterion}[A signed transfer sufficient to exclude half-line admission]
\label{crit:filter-signed-transfer}
Along the critical scaling, suppose that for all sufficiently large
$M$ there is $u_M\ge u_1+\log M-C$, with $C$ independent of $M$,
such that $\mathscr W_M'(u_M)>0$.  Then the full filter
$K_{M,d}$ has a negative interval beyond every fixed wall for all
sufficiently large $M$.  In particular it is not nonnegative on any
fixed half-line.
\end{criterion}

\begin{proof}
Equation \eqref{eq:filter-pole-pairs-derivative} makes the
pole-and-pairs subfilter negative at $u_M$; its limit at infinity is
one.  Its last zero therefore exceeds $u_M$.  Apply the remaining
positive real factors and \Cref{lem:filter-real-factor-zero}.
At each step there is a later negative interval and the limit is
still one.  Since $u_M\to\infty$ and $b_0>0$, the full filter has the
asserted negative interval.
\end{proof}

The estimate in \Cref{crit:filter-signed-transfer} is a genuinely
signed, increasing-order statement about the finite beta kernel.
\Cref{thm:heatblock-growing-symbol} establishes translation and escape
for the Gaussian symbol, but does not prove this estimate.  Even a
proof of the latter would exclude this half-line-positive architecture,
not an off-line zero, and would not by itself exclude all filters
admitted only on prime-power logarithms.

\subsection{Exact finite sign certificates and the remaining closure}
\label{subsec:filter-finite-certificates}

Finite instances can be tested without a root finder or a floating-point
sign decision.  For the family \eqref{eq:filter-equispaced-rational},
take $\beta=3/4$, $\eta=9/10$, $\theta=1/2$, $\gamma=14$, and
$\kappa_\chi=0$.  The ordinate here is an algebraic parameter of the
filter; it is not asserted to be an ordinate of an actual zero.
Choose the rational values of $d$ in \Cref{tab:filter-exact-signs},
which are the nearest hundredths to $6\sqrt M$.

\begin{proposition}[Finite rational sign certificates]
\label{prop:filter-finite-signs}
For every row of \Cref{tab:filter-exact-signs},
$H_{M,d}(q)<0$.  Hence $K_{M,d}$ is negative at
$u=-\varepsilon^{-1}\log q$, has a later real zero, and fails to be
nonnegative on $[\log2,\infty)$ in each of these finite cases.
No assertion for all $M$ follows from this table.
\end{proposition}

\begin{table}[htbp]
\centering
\begin{tabular}{rrrr}
\toprule
$M$ & $d$ & $q$ & $\operatorname{sign}H_{M,d}(q)$\\
\midrule
6   & $147/10$   & $221/500$ & $-1$\\
10  & $1897/100$ & $209/500$ & $-1$\\
16  & $24$       & $101/250$ & $-1$\\
24  & $2939/100$ & $2/5$     & $-1$\\
48  & $4157/100$ & $407/1000$ & $-1$\\
96  & $5879/100$ & $213/500$ & $-1$\\
192 & $4157/50$  & $453/1000$ & $-1$\\
\bottomrule
\end{tabular}
\caption{Exact rational negative values for the full normalized
three-ladder filter, not only its pole-and-pairs subfilter.}
\label{tab:filter-exact-signs}
\end{table}

\begin{proof}
Here is an integer recurrence that specifies each certificate completely.
Choose a common denominator $Q$ for $d$, $\varepsilon$, $\xi_d$, and
$\sigma$, and set
\[
 e=Q\varepsilon,\quad D=Qd,\quad X=Q\xi_d,\quad S=Q\sigma,
 \quad G=14Q.
\]
For $0\le j\le n$, form the positive integers
\begin{equation}
 A_j=\binom nj(je+D)(je+X)
 \prod_{m=0}^{M-1}(je+S+2mQ)
 \prod_{m=0}^{M-1}\{(je+S+2mQ)^2+G^2\}.
 \label{eq:filter-integer-coefficients}
\end{equation}
They are a common positive multiple of the $c_j$ in
\eqref{eq:filter-alternating-polynomial}.  If $q=a/b$ in lowest
terms, set $T_n=(-1)^nA_n$ and, successively for $j=n-1,\ldots,0$,
\begin{equation}
 T_j=aT_{j+1}+(-1)^jA_jb^{n-j}.
 \label{eq:filter-integer-horner}
\end{equation}
Then $T_0=b^n\sum_j(-1)^jA_j(a/b)^j$.  Substitution of each row
gives $T_0<0$ by integer arithmetic; the recurrence uses only the
displayed rational parameters and finite products.  Also $H(0)=1$
and $H(1)=0$, independently checking the two endpoint values.

All listed $q$ are smaller than $1/2$ and all listed
$\varepsilon=d/(2N)$ are smaller than one.  Thus
$-\varepsilon^{-1}\log q>\log2$.  The positive limit
$K_{M,d}(u)\to b_0$ gives a later zero.  Continuity gives a negative
interval around the displayed point.
\end{proof}

The finite certificates, the all-real zero law, and the growing Gaussian
translation theorem answer different questions.  The first settles
specified finite filters; the second constrains real-factor architectures;
the third settles an increasing-degree Gaussian symbol.  None eliminates
the arithmetic compensation in \eqref{eq:filter-compensation}.  A
successful cofinal filter argument must retain actual prime-power
admission, normalize the target, and control the completed signed
remainder despite the interpolation cost.  Conversely, failure of a
particular half-line-positive architecture is not a failure of every
admitted spectral test.  This distinction is the spectral counterpart
of the source--endpoint separation used throughout the monograph.



\part{Conditional closure of affine prime correlations}
% =============================================================================

\section{Affine entrance and conventions}
\label{sec:affine-entrance}

The local theory is developed for a primitive nondegenerate pair
$\bm L=(L_1,L_2)$.  The explicit pressure closure is then specialized to the
twin shift.  Write
\[
 \mathfrak S(2)=2\prod_{p>2}\frac{p(p-2)}{(p-1)^2},
\]
and, with the shifted endpoint convention used in the outer compiler,
\[
 R_X(2)=\mathcal C_X(2)-\mathfrak S(2)X.
\]
Changing between $(X,2X]$ and the translated interval in the compiler alters
only $O((\log X)^2)$ endpoint weight and is absorbed into every $o(X)$ term.
The local comparator is
\[
 \gamma_2(d)=\frac d{\varphi(d)}\one_{(d,2)=1}.
\]
The square-root divisor cutoff is
\[
 D_X=X^{1/2}(\log X)^{-B},\qquad B>0
\]
with $B$ chosen sufficiently large for \Cref{thm:weighted-BV-input}, and
$Q_X=\lfloor(2X+2)/D_X\rfloor$.

\section{Primitive affine pairs and their local root geometry}
\label{sec:local-geometry}
% =============================================================================

\begin{definition}[Primitive nondegenerate affine pair]
\label{def:primitive-pair}
A primitive affine pair is
\[
 \bm L=(L_1,L_2),
 \qquad L_i(n)=a_i n+b_i,
 \qquad a_i,b_i\in\Z,
 \tag{3.1}
\]
with $a_i\ne0$, $\gcd(a_i,b_i)=1$, and
\[
 \boxed{\Delta(\bm L):=a_1b_2-a_2b_1\ne0.}
 \tag{3.2}
\]
At scale $X$, let $I_X(\bm L)$ be an interval of consecutive integers of
cardinality $M_X\asymp_{\bm L}X$, contained in
$[-C_{\bm L}X,C_{\bm L}X]$, on which both forms are at least two, and put
\[
 \mathcal C_{\bm L}(X)
 :=\sum_{n\in I_X(\bm L)}\Lambda(L_1(n))\Lambda(L_2(n)).
 \tag{3.3}
\]
We set $\Lambda(m)=0$ for $m\le1$.
\end{definition}

\begin{remark}[Degenerate diagonal branch]
If two primitive integer coefficient vectors $(a_i,b_i)$ are rationally
proportional, the proportionality factor is $\pm1$.  On a common positive
interval the nontrivial possibility is therefore $L_1=L_2$, which is a
one-form diagonal correlation rather than a prime-pair problem.  Condition
\textup{(3.2)} removes this branch.
\end{remark}

For a prime $p$, write
\[
 Z_{i,p}=\{r\in\F_p:L_i(r)=0\},
 \qquad
 \nu_p(\bm L)=\abs{Z_{1,p}\cup Z_{2,p}},
 \tag{3.4}
\]
and put
\[
 \varepsilon_i(p)=\one_{p\nmid a_i}.
 \tag{3.5}
\]

\begin{theorem}[Exact local root-count formula]
\label{thm:root-count}
For every prime $p$,
\[
 \boxed{
 \nu_p(\bm L)
 =\varepsilon_1(p)+\varepsilon_2(p)
 -\varepsilon_1(p)\varepsilon_2(p)\one_{p\mid\Delta(\bm L)}.}
 \tag{3.6}
\]
In particular $0\le\nu_p(\bm L)\le2$.  If
$p\nmid a_1a_2\Delta(\bm L)$, then $\nu_p(\bm L)=2$.
\end{theorem}

\begin{proof}
If $p\nmid a_i$, the equation $a_i r+b_i=0$ has the unique root
$r_i=-b_i a_i^{-1}$.  If $p\mid a_i$, primitivity gives $p\nmid b_i$, so the
form has no root.  When both roots exist, they coincide exactly when
\[
 -b_1a_1^{-1}=-b_2a_2^{-1}\pmod p,
\]
which is equivalent to $a_1b_2-a_2b_1=0\pmod p$.  Inclusion--exclusion gives
\textup{(3.6)}.
\end{proof}

\begin{definition}[Local factor and admissibility]
\label{def:local-factor}
Define
\[
 \boxed{
 \beta_p(\bm L)
 :=\left(1-\frac{\nu_p(\bm L)}p\right)
   \left(1-\frac1p\right)^{-2}.}
 \tag{3.7}
\]
The pair is locally admissible when $\beta_p(\bm L)>0$ for every prime $p$,
equivalently $\nu_p(\bm L)<p$ for every $p$.
\end{definition}

\begin{theorem}[For primitive pairs, the only local obstruction is parity]
\label{thm:parity-only-local}
A primitive nondegenerate one-variable affine pair is locally admissible if
and only if
\[
 \boxed{\nu_2(\bm L)\le1.}
 \tag{3.8}
\]
No odd prime can be a local obstruction.
\end{theorem}

\begin{proof}
By \Cref{thm:root-count}, $\nu_p\le2$.  Hence $\nu_p<p$ automatically for
$p\ge3$.  At $p=2$, admissibility is exactly $\nu_2<2$.
\end{proof}

\begin{corollary}[Concrete parity classification]
\label{cor:parity-classification}
The pair is inadmissible exactly when both $a_1,a_2$ are odd and the two roots
$-b_1a_1^{-1}$ and $-b_2a_2^{-1}$ are distinct modulo two.  Equivalently, the
two forms take opposite parities for every $n$.
\end{corollary}

\begin{proof}
A form has a root modulo two exactly when its slope is odd.  If both roots
exist, their union has size two exactly when they are distinct.  In that case
one of the two forms is even at each parity class.
\end{proof}

\begin{proposition}[Affine reparametrization invariance]
\label{prop:affine-invariance}
Let $p\nmid u$.  Replacing $n$ by $un+v$ permutes $\F_p$, and therefore
preserves $\nu_p$ and $\beta_p$.  More generally, for squarefree $W$ and
$\gcd(u,W)=1$, the same reparametrization gives a unitary equivalence of the
complete residue carriers modulo $W$.
\end{proposition}

\begin{proof}
The map $n\mapsto un+v$ is a bijection modulo $p$, and modulo $W$ by the
Chinese remainder theorem.  A bijective pullback preserves the union of the
zero sets and the uniform measure.  It also preserves every multiplication
operator generated by the residue indicators.
\end{proof}

% =============================================================================
\section{Singular-series convergence and the two benchmark formulae}
\label{sec:singular-series}
% =============================================================================

Put
\[
 C_2:=\prod_{p>2}\left(1-\frac1{(p-1)^2}\right).
 \tag{4.1}
\]

\begin{theorem}[Absolute convergence and positivity]
\label{thm:series-convergence}
For a primitive nondegenerate affine pair, the Euler product
\[
 \boxed{\mathfrak S(\bm L):=\prod_p\beta_p(\bm L)}
 \tag{4.2}
\]
converges absolutely after the factors are written as $1+u_p$.  It is positive
if and only if the pair is locally admissible.  More explicitly,
\[
 \boxed{
 \mathfrak S(\bm L)
 =\beta_2(\bm L)C_2
  \prod_{\substack{p>2\\p\mid a_1a_2\Delta(\bm L)}}
  \frac{\beta_p(\bm L)}{1-(p-1)^{-2}}.}
 \tag{4.3}
\]
The last product is finite.
\end{theorem}

\begin{proof}
For $p\nmid a_1a_2\Delta$, \Cref{thm:root-count} gives $\nu_p=2$, hence
\[
 \beta_p=\frac{1-2/p}{(1-1/p)^2}
 =\frac{p(p-2)}{(p-1)^2}
 =1-\frac1{(p-1)^2}.
\]
The sum of the absolute deviations is bounded by
$\sum_{m\ge2}m^{-2}<\infty$.  All remaining factors form a finite set.
An inadmissible factor is zero.  On the admissible branch every factor is
positive, and an absolutely convergent product of positive factors with
summable deviations from one has a positive limit.
\end{proof}

\begin{definition}[Fixed-shift and reflection packets]
For $h\ne0$ and $N\ge2$, put
\[
 \bm T_h=(n,n+h),
 \qquad
 \bm G_N=(n,N-n).
 \tag{4.4}
\]
Their determinants are respectively $h$ and $N$.
\end{definition}

\begin{theorem}[Explicit fixed-shift singular series]
\label{thm:shift-series}
The packet $\bm T_h$ is admissible exactly when $h$ is even.  On that branch,
\[
 \boxed{
 \mathfrak S(h)
 =2C_2\prod_{\substack{p\mid h\\p>2}}
       \frac{p-1}{p-2}.}
 \tag{4.5}
\]
\end{theorem}

\begin{proof}
Modulo $p$, the forbidden roots are $0$ and $-h$.  They coincide exactly when
$p\mid h$.  Thus $\nu_2=1$ for even $h$ and $\nu_2=2$ for odd $h$.  If
$p>2$ divides $h$, then
\[
 \beta_p=\frac{1-1/p}{(1-1/p)^2}=\frac p{p-1}.
\]
Dividing by the generic factor $p(p-2)/(p-1)^2$ gives
$(p-1)/(p-2)$.  The factor at two is $2$.
\end{proof}

\begin{theorem}[Explicit binary-Goldbach singular series]
\label{thm:goldbach-series}
The packet $\bm G_N$ is admissible exactly when $N$ is even.  On that branch,
\[
 \boxed{
 \mathfrak S_G(N)
 =2C_2\prod_{\substack{p\mid N\\p>2}}
       \frac{p-1}{p-2}.}
 \tag{4.6}
\]
\end{theorem}

\begin{proof}
The forbidden roots are $0$ and $N$ modulo $p$, so the root count is identical
to that of \Cref{thm:shift-series} with $h=N$.
\end{proof}

\begin{theorem}[Exact local reflection equivalence]
\label{thm:local-reflection}
Let $W$ be squarefree.  For $q\in\Z$, define
\[
 \mathcal R_{q,W}^{\rm sh}
 =\{r\pmod W:\gcd(r(r+q),W)=1\},
\]
\[
 \mathcal R_{q,W}^{\rm ref}
 =\{s\pmod W:\gcd(s(q-s),W)=1\}.
\]
Then
\[
 \boxed{r\longmapsto-r}
 \tag{4.7}
\]
is a bijection from $\mathcal R_{q,W}^{\rm sh}$ to
$\mathcal R_{q,W}^{\rm ref}$ and unitarily conjugates their complete local
indicator algebras.  In particular, fixed shift $h=q$ and moving sum $N=q$
have exactly the same finite local packet.
\end{theorem}

\begin{proof}
Under $s=-r$,
\[
 s(q-s)=(-r)(q+r).
\]
Multiplication by $-1$ does not change a gcd with $W$, so the admissibility
conditions agree.  Pullback by the bijection conjugates every residue
multiplication operator.
\end{proof}

\begin{remark}[Local equivalence is not global equivalence]
The map in \Cref{thm:local-reflection} acts on a finite residue ring.  It does
not identify the intervals $1\le n\le X$ and $1\le n<N$, nor does it remove the
external quantifier ``uniformly for every large even $N$''.  The two branches
must share their local theory and retain separate globalization theorems.
\end{remark}

% =============================================================================
\section{Local interaction closure and the two-axis grading}
\label{sec:local-closure}
% =============================================================================

The word ``grade'' is ambiguous unless one says what the slots are.  For a
prime pair there are two form slots, but the local model also has one coordinate
for every prime.  These two support filtrations behave differently.

\subsection{One-prime pair incidence}

For a prime $p$, let
\[
 d_{i,p}(r)=\one_{p\mid L_i(r)},
 \qquad
 q_{i,p}=1-d_{i,p},
 \qquad r\in\F_p.
 \tag{5.1}
\]
Let $\mathcal N_{p}^{<2}=\Span\{1,d_{1,p},d_{2,p}\}$ in
$L^2(\F_p)$ with uniform measure.

\begin{theorem}[One-prime pair reward has zero proper-support innovation]
\label{thm:one-prime-zero-innovation}
For every primitive affine pair and every prime $p$,
\[
 \boxed{d_{1,p}d_{2,p}\in\mathcal N_p^{<2},
 \qquad q_{1,p}q_{2,p}\in\mathcal N_p^{<2}.}
 \tag{5.2}
\]
Consequently the Schur residual of the one-prime pair-coprimality reward after
shorting constants and both one-form divisibility rewards is zero.
\end{theorem}

\begin{proof}
Each root set $Z_{i,p}$ is empty or a singleton.  Hence their intersection is
either empty or equals both singleton root sets.  In the first case
$d_{1,p}d_{2,p}=0$.  In the second case $d_{1,p}=d_{2,p}$ and
$d_{1,p}d_{2,p}=d_{1,p}$.  Thus the product belongs to the displayed span.
Finally,
\[
 q_{1,p}q_{2,p}=1-d_{1,p}-d_{2,p}+d_{1,p}d_{2,p}
\]
belongs to the same span.
\end{proof}

\begin{proposition}[Exact local covariance]
\label{prop:local-covariance}
If both forms have a root modulo $p$, then
\[
 \operatorname{Cov}(q_{1,p},q_{2,p})
 =\begin{cases}
  -p^{-2},&p\nmid\Delta,\\[0.3em]
  (p-1)p^{-2},&p\mid\Delta.
 \end{cases}
 \tag{5.3}
\]
If at least one form has no root, the covariance is zero.
\end{proposition}

\begin{proof}
When both roots exist,
$\mathbb E q_i=1-1/p$ and
$\mathbb E(q_1q_2)=1-\nu_p/p$.  Insert $\nu_p=2$ for distinct roots and
$\nu_p=1$ for coincident roots.  If one $q_i$ is constant one, the covariance
vanishes.
\end{proof}

\begin{remark}
A nonzero covariance is not the same thing as a new all-support source
direction.  Formula \textup{(5.2)} shows that the covariance is already encoded
by the deterministic root-collision relation inside the proper-support local
algebra.
\end{remark}

\subsection{Chinese-remainder tensorization}

For squarefree $W$, put
\[
 \mathcal H_W=L^2(\Z/W\Z),
 \qquad
 Q_{p,\bm L}=M_{\one_{p\nmid L_1L_2}},
 \qquad
 Q_{W,\bm L}=M_{\one_{\gcd(L_1L_2,W)=1}}.
 \tag{5.4}
\]

\begin{theorem}[Exact CRT tensor factorization]
\label{thm:CRT-tensor}
The Chinese remainder unitary gives
\[
 \boxed{
 \mathcal H_W\cong\bigotimes_{p\mid W}L^2(\F_p),
 \qquad
 Q_{W,\bm L}\cong\bigotimes_{p\mid W}Q_{p,\bm L}.}
 \tag{5.5}
\]
Moreover
\[
 \boxed{
 \frac1W\Tr Q_{W,\bm L}
 =\prod_{p\mid W}\left(1-\frac{\nu_p(\bm L)}p\right),}
 \tag{5.6}
\]
and
\[
 \boxed{
 \beta_W(\bm L)
 :=\frac{W\abs{\mathcal R_{\bm L,W}}}{\varphi(W)^2}
 =\prod_{p\mid W}\beta_p(\bm L),}
 \tag{5.7}
\]
where
$\mathcal R_{\bm L,W}=\{r\pmod W:\gcd(L_1(r)L_2(r),W)=1\}$.
\end{theorem}

\begin{proof}
The CRT identifies a residue modulo $W$ with its tuple of residues modulo the
prime divisors of $W$ and carries uniform measure to the product measure.  The
coprimality indicator is the product of the local indicators, proving
\textup{(5.5)}.  Taking expectations gives \textup{(5.6)}.  Finally
$W/\varphi(W)=\prod_{p\mid W}(1-1/p)^{-1}$, so multiplying
\textup{(5.6)} by $(W/\varphi(W))^2$ gives \textup{(5.7)}.
\end{proof}

Let
\[
 \alpha_p=\mathbb E q_{p,\bm L},
 \qquad
 f_p=q_{p,\bm L}-\alpha_p,
 \tag{5.8}
\]
where $q_{p,\bm L}=\one_{p\nmid L_1L_2}$.

\begin{theorem}[All prime-coordinate Hoeffding grades are monoidal]
\label{thm:prime-Hoeffding}
Under the CRT identification,
\[
 \boxed{
 q_{W,\bm L}
 =\sum_{S\subseteq\{p:p\mid W\}}
   \left(\prod_{p\notin S}\alpha_p\right)
   \bigotimes_{p\in S}f_p.}
 \tag{5.9}
\]
The summands are mutually orthogonal.  If every $f_p$ in a set $S$ is nonzero,
the $S$-grade is nonzero.  Hence no fixed cutoff in prime-coordinate tensor
order represents the exact hard-sieve projector as $W$ acquires more prime
factors.  Nevertheless every grade in \textup{(5.9)} is forced by the one-prime
sources and the CRT fusion; there is no nonmonoidal cross-prime source
innovation.
\end{theorem}

\begin{proof}
Expand
$\bigotimes_{p\mid W}(\alpha_p\one+f_p)$.  Distinct subsets differ at some
prime coordinate where one tensor factor is constant and the other has mean
zero, proving orthogonality.  A tensor product of nonzero vectors is nonzero.
The final statement follows because the CRT unitary is onto and the complete
source is exactly the displayed tensor product, with no orthogonal complement.
\end{proof}

\begin{corollary}[Cross-prime cumulants vanish]
\label{cor:prime-cumulants}
For distinct primes $p_1,\ldots,p_k\mid W$ and centred functions
$g_j\in L^2(\F_{p_j})$,
\[
 \boxed{
 \mathbb E\prod_{j=1}^k g_j=0.}
 \tag{5.10}
\]
Thus the high prime-coordinate grades in \textup{(5.9)} are product-connected
Hoeffding tensors, not irreducible cumulant interactions.
\end{corollary}

\begin{proof}
The CRT state is a product state, so the expectation factors into the product
of the zero means.
\end{proof}

\begin{theorem}[Bi-graded local closure verdict]
\label{thm:bigraded-verdict}
For a primitive affine pair, the exact local mechanism has the following
closure profile.
\begin{enumerate}[label=\textup{(B\arabic*)},leftmargin=3em]
\item Its form-support degree is at most two.
\item At one prime, the all-form local reward has zero innovation after the
      complete one-form root data are retained.
\item Across primes, arbitrarily high prime-coordinate Hoeffding grades occur,
      but they are exact monoidal CRT products with zero nonmonoidal innovation.
\item Therefore no genuinely new form-support grade $>2$ is required by local
      congruence admissibility or by the singular series.
\end{enumerate}
\end{theorem}

\begin{proof}
Items~\textup{(B1)}--\textup{(B3)} are respectively the definition of the pair
packet, \Cref{thm:one-prime-zero-innovation}, and
\Cref{thm:CRT-tensor,thm:prime-Hoeffding}.  Item~\textup{(B4)} follows because
every local source lies in the monoidal category generated by the two form
slots and the one-prime local factors.
\end{proof}

\begin{remark}[Why ``$m_{\rm cl}=2$'' needs a qualifier]
The local source innovation closes at form degree two.  The exact sieve
projector does not have bounded prime-coordinate tensor degree.  Confusing
these statements would either falsely claim a finite-grade local formula or
falsely interpret forced CRT products as new physical coordinates.
\end{remark}

% =============================================================================
\section{Actual degree-two occurrence and the nonmonoidal alternative}
\label{sec:occurrence}
% =============================================================================

Local tensorization does not prove that the actual two-copy von Mangoldt
occurrence is the product of two one-copy occurrences.  This is a distinct
same-history question.

\begin{definition}[Product and occurrence sectors]
Let $H_1=\C\Omega\oplus H_1^0$, let
$U_2:H_1^{\otimes2}\to H_2$ be the source-fusion isometry reconstructed from
the actual mixed occurrence Gram, and put
\[
 H_{\prodsec}=U_2(H_1^0\otimes H_1^0),
 \qquad
 H_{\occ}=(\Ran U_2)^\perp.
 \tag{6.1}
\]
For an actual pair source $Z:E\to H_2$, define
\[
 \Delta_{\occ}(Z)=Z^*(I-U_2U_2^*)Z\succeq0.
 \tag{6.2}
\]
\end{definition}

\begin{theorem}[Complete degree-two all-support decomposition]
\label{thm:degree-two-decomposition}
After every proper one-form support is shorted, the complete retained
projection is
\[
 \boxed{
 C_2^{\rm full}
 =U_2(Q\otimes Q)U_2^*+(I-U_2U_2^*).}
 \tag{6.3}
\]
Consequently
\[
 \boxed{
 Z^*C_2^{\rm full}Z
 =Z^*U_2(Q\otimes Q)U_2^*Z+\Delta_{\occ}(Z).}
 \tag{6.4}
\]
The two terms are orthogonal.  The first is the product-connected neutral pair
source; the second is a genuinely nonmonoidal source inside degree two.
\end{theorem}

\begin{proof}
Inside $\Ran U_2$, proper-support shorting leaves exactly the transported
Hoeffding top grade $U_2(H_1^0\otimes H_1^0)$.  The orthogonal complement of
the entire product occurrence range is not reached by any proper-support
product source and therefore also survives.  The two projections are
orthogonal, giving \textup{(6.3)} and \textup{(6.4)}.
\end{proof}

\begin{proposition}[Failure of product closure need not create grade three]
\label{cth:no-grade-three}
There are actual degree-two carriers with the prescribed one-copy source and
product degree-two source, but with an arbitrary additional orthogonal source
sector.  Thus a pair-closure failure may be entirely degree-two and
nonmonoidal; searching only for support size three can miss it.
\end{proposition}

\begin{proof}
Take
\[
 H_2=H_1^{\otimes2}\oplus K
\]
for any nonzero finite $K$, use the canonical inclusion as $U_2$, and give the
actual degree-two source an arbitrary nonzero $K$ component.  All one-copy and
product-copy data remain fixed, while \textup{(6.2)} has positive rank.
\end{proof}

\begin{theorem}[Occurrence Feshbach reduction]
\label{thm:pair-Feshbach}
Let the positive all-support action on
$H_{\prodsec}\oplus H_{\occ}$ be
\[
 L=\begin{pmatrix}A&B\\B^*&C\end{pmatrix}\succeq0,
 \qquad C\succeq\gamma I\quad(\gamma>0).
 \tag{6.5}
\]
Then the effective product-pair action is
\[
 \boxed{S_{\rm eff}=A-BC^{-1}B^*\succeq0.}
 \tag{6.6}
\]
For a source $z=(z_{\prodsec},z_{\occ})$, the effective retained source is
\[
 \boxed{z_{\rm eff}=z_{\prodsec}-BC^{-1}z_{\occ}.}
 \tag{6.7}
\]
Thus a nonzero occurrence sector modifies both the pair action and the pair
source; it cannot be discarded after computing only its marginal Gram.
\end{theorem}

\begin{proof}
Complete the square:
\[
 \left\langle\binom{x}{y},L\binom{x}{y}\right\rangle
 =\langle x,S_{\rm eff}x\rangle
 +\langle y+C^{-1}B^*x,
 C(y+C^{-1}B^*x)\rangle.
\]
This proves positivity and the action formula.  Solving the occurrence block
of $L(x,y)=z$ gives
$y=C^{-1}(z_{\occ}-B^*x)$; substitution into the product block gives
$S_{\rm eff}x=z_{\prodsec}-BC^{-1}z_{\occ}$.
\end{proof}

\begin{remark}[Local form-grade closure]
The local calculation shows that form grade three is not required locally.  It leaves two possible
survivors after complete local and proper-support shorting:
\[
 \boxed{
 \text{product-connected neutral degree two}
 \quad\oplus\quad
 \text{nonmonoidal occurrence degree two}.}
 \tag{6.8}
\]
The actual mixed two-copy panel decides the second summand.  Neither summand is
controlled by a degree-one GRH Ward identity.
\end{remark}

% =============================================================================
\section{The selector-stable parity square}
\label{sec:parity-square}
% =============================================================================

The classical parity barrier can now be stated as an exact finite quotient,
not as a slogan.  The canonical arithmetic loading supplies squarefree
factor-depth selectors.  Restrict each form to factor depth one or two and
apply the affine selector \emph{before} taking any signed current.

\begin{definition}[Pair factor-depth masses and currents]
\label{def:pair-parity}
Let $\mu$ be any finite nonnegative weight on accepted affine-pair occurrences
for which each selected value $L_i(n)$ is squarefree and has factor depth
$d_i(n)\in\{1,2\}$.  Put
\[
 \sigma_i(n)=(-1)^{d_i(n)}\in\{-1,+1\}.
 \tag{7.1}
\]
Let $P_{ab}$ be the total weight of the events with $(d_1,d_2)=(a,b)$, and
define
\[
 M=\sum_{a,b=1}^2P_{ab},
 \quad
 J_1=\sum_n\mu(n)\sigma_1(n),
 \quad
 J_2=\sum_n\mu(n)\sigma_2(n),
 \quad
 J_{12}=\sum_n\mu(n)\sigma_1(n)\sigma_2(n).
 \tag{7.2}
\]
The current $J_{12}$ is the neutral pair-parity current.
\end{definition}

\begin{theorem}[Exact Walsh inversion of the pair parity square]
\label{thm:Walsh-parity}
One has
\[
 \boxed{
 \begin{aligned}
 P_{11}&=\frac14(M-J_1-J_2+J_{12}),\\
 P_{12}&=\frac14(M-J_1+J_2-J_{12}),\\
 P_{21}&=\frac14(M+J_1-J_2-J_{12}),\\
 P_{22}&=\frac14(M+J_1+J_2+J_{12}).
 \end{aligned}}
 \tag{7.3}
\]
In particular,
\[
 \boxed{P_{11}>0
 \iff M-J_1-J_2+J_{12}>0.}
 \tag{7.4}
\]
\end{theorem}

\begin{proof}
The four functions
$1,\sigma_1,\sigma_2,\sigma_1\sigma_2$ are the character table of
$(\Z/2\Z)^2$.  The forward transform from the four masses to
$(M,J_1,J_2,J_{12})$ is the $4\times4$ Hadamard matrix.  Its square is four
times the identity, so its inverse is one quarter of itself, giving
\textup{(7.3)}.
\end{proof}

\begin{corollary}[A sufficient signed-current criterion]
\label{cor:current-sufficient}
If
\[
 \boxed{\abs{J_1}+\abs{J_2}+\abs{J_{12}}<M,}
 \tag{7.5}
\]
then $P_{11}>0$.  If along a growing family
$M\gg X$ and $J_1,J_2,J_{12}=o(X)$, then
$P_{11}=(1/4+o(1))M\gg X$.
\end{corollary}

\begin{proof}
By \textup{(7.3)},
$4P_{11}\ge M-|J_1|-|J_2|-|J_{12}|$.
\end{proof}

\begin{theorem}[The all-support parity quotient is one-dimensional]
\label{thm:parity-one-line}
On the four-point sign square, the proper-support function space is
\[
 \mathcal N_{<2}=\Span\{1,\sigma_1,\sigma_2\}.
 \tag{7.6}
\]
Its algebraic quotient has dimension one, generated by
$\sigma_1\sigma_2$.  With any full-support positive metric, the orthogonal
short of the pair sign by \textup{(7.6)} has rank one; if the sign support is
not full, its rank is at most one.
\end{theorem}

\begin{proof}
The four Walsh characters form a basis of all functions on the sign square.
The first three span the constant and the two proper one-slot grades, leaving
the product character as the one-dimensional quotient.  Changing to a
full-support positive metric changes the orthogonal representative but not the
quotient dimension.  Restricting the support can only lower dimension.
\end{proof}

\begin{proposition}[Degree-one parity data do not decide prime--prime mass]
\label{cth:degree-one-parity}
There are two positive pair carriers with the same $M,J_1,J_2$ and different
prime--prime mass.  One may have $P_{11}>0$ and the other $P_{11}=0$.
\end{proposition}

\begin{proof}
In the first carrier put
$P_{11}=P_{22}=M/2$ and $P_{12}=P_{21}=0$.  In the second put
$P_{12}=P_{21}=M/2$ and $P_{11}=P_{22}=0$.  Both have $J_1=J_2=0$.  Their
neutral currents are respectively $J_{12}=M$ and $J_{12}=-M$, and their
prime--prime masses are $M/2$ and zero.
\end{proof}

\begin{theorem}[Pair-parity closure theorem]
\label{thm:pair-parity-closure}
After the affine selector, squarefree depth-$1/2$ restriction, and simultaneous
short by the constant and the two one-form parity sources, every parity
obstruction to prime--prime mass is carried by one neutral degree-two line.
No form-support grade $>2$ is required for parity closure.  The remaining
possibilities are:
\begin{enumerate}[label=\textup{(P\arabic*)},leftmargin=3em]
\item a nonzero neutral product-sign current in the product occurrence sector;
\item a nonmonoidal degree-two occurrence innovation;
\item leakage from the selected depth-$1/2$ carrier into deeper factor-depth or
      prime-power sectors.
\end{enumerate}
\end{theorem}

\begin{proof}
The first assertion is \Cref{thm:parity-one-line}.  The product/nonmonoidal
split is \Cref{thm:degree-two-decomposition}.  Any event not represented by
the sign square lies outside the squarefree depth-$1/2$ selector and is exactly
the leakage in item~\textup{(P3)}.  These alternatives exhaust the source
space after the declared shorts.
\end{proof}

\begin{remark}[Selector stability]
The currents in \Cref{def:pair-parity} must be evaluated on the fixed-shift or
fixed-sum carrier itself.  A cancellation theorem for the unrestricted
one-variable Liouville or M\"obius source does not determine the same current
after imposing $L_2(n)-L_1(n)=h$ or $L_1(n)+L_2(n)=N$.
\end{remark}

% =============================================================================
\section{Finite congruence ledgers cannot break factor-depth parity}
\label{sec:parity-no-go}
% =============================================================================

The following elementary no-go identifies the exact limit of a fixed local
ledger.  It uses the classical Dirichlet theorem on primes in reduced residue
classes as an external established theorem.

\begin{theorem}[Prime and semiprime indistinguishability in every fixed residue ledger]
\label{thm:residue-parity-no-go}
Let $W\ge1$ and let $a$ be coprime to $W$.  There exist infinitely many primes
$q\equiv a\pmod W$ and infinitely many squarefree semiprimes
$m\equiv a\pmod W$.  Hence every observable measurable only with respect to
the residue modulo $W$ takes the same value on infinitely many factor-depth-one
and factor-depth-two integers.
\end{theorem}

\begin{proof}
Dirichlet's theorem gives infinitely many primes $q\equiv a\pmod W$.
It also gives infinitely many primes $r\equiv1\pmod W$ and infinitely many
primes $s\equiv a\pmod W$.  Choose $r\ne s$; then $m=rs$ is squarefree,
has factor depth two, and satisfies $m\equiv a\pmod W$.
\end{proof}

\begin{corollary}[Local factors cannot supply the neutral signed current]
\label{cor:local-no-sign}
No fixed finite collection of congruence indicators determines the Liouville
sign, the depth-one selector, or the pair-neutral current of
\Cref{def:pair-parity}.  Therefore the singular series and local admissibility
cannot by themselves close twins or Goldbach.
\end{corollary}

\begin{proof}
Every finite congruence collection factors through a residue modulo their
least common multiple $W$.  Apply \Cref{thm:residue-parity-no-go}.  The pair
statement follows because each one-form factor depth already fails to factor
through the local ledger.
\end{proof}

\begin{remark}[What the parity barrier means in the Renewal language]
The barrier is not evidence that a third form slot is missing.  It says that
the chronology-derived prime source has a signed factor-depth coordinate which
is invisible to every fixed local congruence quotient.  For a pair, the
proper-support parity coordinates leave one neutral all-support line.  That
line must be controlled by an actual pair action, current, or semiprime-profile
estimate.
\end{remark}

% =============================================================================
\section{Exact \texorpdfstring{$W$}{W}-fibre decomposition and the protected pair amplitude}
\label{sec:W-fibres}
% =============================================================================

We now separate the local model from the genuine global estimate without
assuming an asymptotic distribution theorem.

\begin{definition}[Admissible fibres and normalized pair amplitudes]
\label{def:fibre-amplitude}
Let $W$ be squarefree and divisible by $\rad(a_1a_2)$.  Put
\[
 \mathcal R_{\bm L,W}
 =\{r\pmod W:\gcd(L_1(r)L_2(r),W)=1\},
 \qquad R_{\bm L,W}=\abs{\mathcal R_{\bm L,W}}.
 \tag{9.1}
\]
For $r\in\mathcal R_{\bm L,W}$, define
\[
 S_{\bm L,X,W}(r)
 =\sum_{\substack{n\in I_X(\bm L)\\n\equiv r\ (W)}}
   \Lambda(L_1(n))\Lambda(L_2(n)),
 \tag{9.2}
\]
\[
 \boxed{
 A_{\bm L,X,W}(r)
 =\frac{\varphi(W)^2}{M_X W}
  S_{\bm L,X,W}(r).}
 \tag{9.3}
\]
The protected neutral amplitude and the fibre dispersion are
\[
 \mathfrak a_{\bm L}(X;W)
 =\frac1{R_{\bm L,W}}\sum_{r\in\mathcal R_{\bm L,W}}A(r),
 \tag{9.4}
\]
\[
 \mathfrak d_{\bm L}(X;W)^2
 =\frac1{R_{\bm L,W}}\sum_r
   \abs{A(r)-\mathfrak a_{\bm L}(X;W)}^2.
 \tag{9.5}
\]
Finally put
\[
 \mathfrak E_{\bm L}(X;W)
 =\frac1{R_{\bm L,W}}\sum_r\abs{A(r)-1}^2.
 \tag{9.6}
\]
\end{definition}

\begin{proposition}[Amplitude--dispersion Pythagoras]
\label{prop:amplitude-dispersion}
One has
\[
 \boxed{
 \mathfrak E_{\bm L}(X;W)
 =\abs{\mathfrak a_{\bm L}(X;W)-1}^2
  +\mathfrak d_{\bm L}(X;W)^2.}
 \tag{9.7}
\]
In particular,
\[
 \abs{\mathfrak a_{\bm L}(X;W)-1}
 \le\sqrt{\mathfrak E_{\bm L}(X;W)}.
 \tag{9.8}
\]
\end{proposition}

\begin{proof}
Write
$A(r)-1=(A(r)-\mathfrak a)+(\mathfrak a-1)$ and sum the squared norm.  The
cross term vanishes because $\sum_r(A(r)-\mathfrak a)=0$.
\end{proof}

\begin{definition}[Small-prime exceptional contribution]
Let
\[
 \mathcal E_{\bm L}^{\exc}(X;W)
 =\sum_{\substack{n\in I_X(\bm L)\\
 \gcd(L_1(n)L_2(n),W)>1}}
 \Lambda(L_1(n))\Lambda(L_2(n)).
 \tag{9.9}
\]
Let $Y_X=\max_{n\in I_X,i}L_i(n)$.
\end{definition}

\begin{theorem}[Exact fibre identity]
\label{thm:exact-fibre-identity}
For every admissible primitive pair and every squarefree $W$ as above,
\[
 \boxed{
 \mathcal C_{\bm L}(X)
 =M_X\beta_W(\bm L)\,
   \mathfrak a_{\bm L}(X;W)
  +\mathcal E_{\bm L}^{\exc}(X;W).}
 \tag{9.10}
\]
Moreover
\[
 \boxed{
 0\le\mathcal E_{\bm L}^{\exc}(X;W)
 \le2\,\omega(W)(\log Y_X)^2.}
 \tag{9.11}
\]
Consequently
\[
 \boxed{
 \begin{aligned}
 \abs{\mathcal C_{\bm L}(X)-M_X\beta_W(\bm L)}
 \le{}&M_X\beta_W(\bm L)
       \sqrt{\mathfrak E_{\bm L}(X;W)}\\
 &+2\omega(W)(\log Y_X)^2.
 \end{aligned}}
 \tag{9.12}
\]
\end{theorem}

\begin{proof}
Partition the sum according to residues modulo $W$.  The residues in
$\mathcal R_{\bm L,W}$ contribute
\[
 \sum_{r\in\mathcal R}S(r)
 =\frac{M_X W}{\varphi(W)^2}\sum_{r\in\mathcal R}A(r)
 =M_X\frac{WR_{\bm L,W}}{\varphi(W)^2}\mathfrak a,
\]
which is the first term in \textup{(9.10)} by \textup{(5.7)}.  The remaining
residues give \textup{(9.9)}.

If a summand in \textup{(9.9)} is nonzero, both form values are prime powers,
and for some $p\mid W$ at least one form value is $p^k$.  For fixed $i,p,k$,
the equation $L_i(n)=p^k$ has at most one solution.  Its weight is at most
$\log p\log Y_X$.  Summing over
$k\le\log Y_X/\log p$ gives at most $(\log Y_X)^2$ for each pair $(i,p)$.
Summing over two forms and the prime divisors of $W$ proves \textup{(9.11)}.
Finally use \textup{(9.8)} in \textup{(9.10)}.
\end{proof}

\begin{theorem}[The sharp scalar target is the protected amplitude]
\label{thm:scalar-target}
The total pair count in \textup{(9.10)} depends on the admissible-fibre vector
$A$ only through its projection onto the constant unit vector
\[
 u_W=R_{\bm L,W}^{-1/2}\one.
 \tag{9.13}
\]
The router dispersion $\mathfrak d^2$ is orthogonal to this observable.  Thus
an operator-norm connected-uniformity estimate is sufficient but not necessary
for the pair-counting theorem; the source-specific scalar amplitude is the
minimal protected target.
\end{theorem}

\begin{proof}
The admissible contribution is proportional to
$\sum_rA(r)=\sqrt{R}\,\ip{u_W}{A}$.  The centred vector
$A-\mathfrak a\one$ is orthogonal to $u_W$ and therefore cancels from the sum.
\end{proof}

\begin{theorem}[Protected Riesz estimate for the pair amplitude]
\label{thm:pair-Riesz}
Let $L_{X,W}^{\rm eff}\succeq0$ be an independently reconstructed effective
action on the admissible-fibre residual space after local, proper-support, and
occurrence shorting.  Put $r=A-\one$.  If
$u_W\in\Ran L_{X,W}^{\rm eff}$, then
\[
 \boxed{
 \abs{\mathfrak a_{\bm L}(X;W)-1}^2
 \le\frac1{R_{\bm L,W}}
 \ip{r}{L_{X,W}^{\rm eff}r}
 \ip{u_W}{(L_{X,W}^{\rm eff})^\dagger u_W}.}
 \tag{9.14}
\]
The first factor is the efficient connected energy; the second is the sharp
protected relation-source influence.
\end{theorem}

\begin{proof}
On $(\Ker L)^\perp$,
\[
 \ip{u_W}{r}
 =\ip{L^{\dagger/2}u_W}{L^{1/2}r}.
\]
Cauchy--Schwarz gives
$|\ip{u_W}{r}|^2\le
\ip{u_W}{L^\dagger u_W}\ip{r}{Lr}$.
Since
$\mathfrak a-1=R^{-1/2}\ip{u_W}{r}$, the result follows.
\end{proof}

\begin{remark}[Upstream refinement of the connected-uniformity target]
The broad affine theorem asks for small all-support connected energy.  The
exact fibre identity shows a sharper route: control only the matrix coefficient
seen by $u_W$.  A full operator-norm estimate remains valuable because it
handles a bank of affine systems simultaneously, but it is stronger than the
single twin or Goldbach amplitude needed for existence.
\end{remark}

% =============================================================================
\section{General fixed-packet globalization}
\label{sec:fixed-globalization}
% =============================================================================

\begin{lemma}[Elementary size of the primorial]
\label{lem:primorial}
Let $W(z)=\prod_{p\le z}p$.  There is an absolute constant $C$ such that
\[
 \boxed{\log W(z)\le Cz\qquad(z\ge2).}
 \tag{10.1}
\]
Consequently, if $z=o(\log X)$, then $W(z)=X^{o(1)}$.
\end{lemma}

\begin{proof}
The product of the primes in $(m,2m]$ divides $\binom{2m}{m}$, so
\[
 \sum_{m<p\le2m}\log p\le\log\binom{2m}{m}\le2m\log2.
\]
Sum this estimate over dyadic intervals ending at $z$.  The resulting geometric
series is $O(z)$, proving the first assertion.  Exponentiation gives the
second.
\end{proof}

\begin{lemma}[Fixed-packet singular tail]
\label{lem:fixed-tail}
For a fixed primitive admissible pair, if $W(z)$ also contains
$\rad(a_1a_2)$, then
\[
 \boxed{\beta_{W(z)}(\bm L)=\mathfrak S(\bm L)(1+O_{\bm L}(z^{-1}))}
 \tag{10.2}
\]
as $z\to\infty$.
\end{lemma}

\begin{proof}
Once $z$ exceeds every prime divisor of $a_1a_2\Delta$, all omitted factors are
$1-(p-1)^{-2}$.  The logarithm of their product is
$O(\sum_{p>z}p^{-2})=O(z^{-1})$.  Exponentiating gives the result.
\end{proof}

\begin{theorem}[Conditional Hardy--Littlewood bridge for a fixed affine pair]
\label{thm:fixed-HL-bridge}
Let $\bm L$ be a fixed primitive admissible pair and let $z=z(X)$ satisfy
\[
 z(X)\to\infty,
 \qquad z(X)=o(\log X),
 \tag{10.3}
\]
with $W_X=W(z(X))$ enlarged by the fixed factor $\rad(a_1a_2)$.  Suppose
\[
 \boxed{\mathfrak E_{\bm L}(X;W_X)\longrightarrow0.}
 \tag{10.4}
\]
Then
\[
 \boxed{
 \mathcal C_{\bm L}(X)
 =M_X\mathfrak S(\bm L)+o(M_X).}
 \tag{10.5}
\]
\end{theorem}

\begin{proof}
By \Cref{lem:primorial}, $\omega(W_X)=O(z(X))$ and
$W_X=X^{o(1)}$.  Since $Y_X=O_{\bm L}(X)$ and $M_X\asymp X$, the exceptional
bound in \textup{(9.11)} is $o(M_X)$.  Equation \textup{(9.12)} and
\textup{(10.4)} give
$\mathcal C_{\bm L}(X)=M_X\beta_{W_X}(\bm L)+o(M_X)$.  Apply
\Cref{lem:fixed-tail}.
\end{proof}

\begin{theorem}[Existence bridge needs only a positive amplitude]
\label{thm:fixed-lower-bridge}
In the setting of \Cref{thm:fixed-HL-bridge}, assume instead that for some
$\eta>0$ and all sufficiently large $X$,
\[
 \boxed{\mathfrak a_{\bm L}(X;W_X)\ge\eta.}
 \tag{10.6}
\]
Then
\[
 \mathcal C_{\bm L}(X)\gg_{\bm L,\eta}X.
 \tag{10.7}
\]
Thus a positive source-specific amplitude is enough for infinitude of the
corresponding prime pairs after the prime-power cleanup of
\Cref{sec:prime-power-cleanup}; fibre equidistribution is not required.
\end{theorem}

\begin{proof}
The exceptional term in \textup{(9.10)} is nonnegative.  The finite local
products converge to the positive number $\mathfrak S(\bm L)$, so they are
bounded below for large $X$.  Insert \textup{(10.6)} into \textup{(9.10)}.
\end{proof}

% =============================================================================
\section{Fixed shifts and the twin/Polignac branch}
\label{sec:fixed-shift}
% =============================================================================

For fixed even $h$, take
\[
 I_X(\bm T_h)=\{1,\ldots,X\},
 \qquad M_X=X,
 \tag{11.1}
\]
and write
\[
 \mathfrak E_h(X;W)=\mathfrak E_{\bm T_h}(X;W),
 \qquad
 \mathfrak a_h(X;W)=\mathfrak a_{\bm T_h}(X;W).
 \tag{11.2}
\]

\begin{theorem}[Fixed-shift closure criterion]
\label{thm:shift-closure}
Let $h$ be a fixed nonzero even integer and choose any
$z(X)\to\infty$ with $z(X)=o(\log X)$.  If
\[
 \boxed{\mathfrak E_h(X;W(z(X)))\to0,}
 \tag{11.3}
\]
then
\[
 \boxed{
 \sum_{n\le X}\Lambda(n)\Lambda(n+h)
 =\mathfrak S(h)X+o_h(X).}
 \tag{11.4}
\]
If only
$\liminf_X\mathfrak a_h(X;W(z(X)))>0$, then the weighted correlation is
$\gg_hX$.
\end{theorem}

\begin{proof}
Apply \Cref{thm:fixed-HL-bridge,thm:fixed-lower-bridge} and
\Cref{thm:shift-series}.
\end{proof}

\begin{corollary}[Twin-prime specialization]
\label{cor:twin-specialization}
For $h=2$, the singular series is $2C_2$.  Either condition in
\Cref{thm:shift-closure} implies infinitely many twin-prime pairs after
\Cref{thm:weighted-to-primes} below.
\end{corollary}

\begin{remark}[Quantifier structure]
The fixed-shift branch asks for one cofinal estimate on one fixed pair source.
Uniformity in a finite bank of even $h$ is useful but not logically necessary
for the twin-prime specialization.  This is the first downstream point at
which the fixed-shift and Goldbach frameworks differ.
\end{remark}

% =============================================================================
\section{Moving sums and a subpolynomial Goldbach local regulator}
\label{sec:Goldbach}
% =============================================================================

For even $N$, take
\[
 I_N(\bm G_N)=\{1,\ldots,N-1\},
 \qquad M_N=N-1.
 \tag{12.1}
\]
The local packet now changes with $N$.

\begin{lemma}[Uniform moving singular-series tail]
\label{lem:Goldbach-tail}
Let $z\ge5$, $W(z)=\prod_{p\le z}p$, and $N$ be even.  Then
\[
 \boxed{
 \left|\log\frac{\mathfrak S_G(N)}{\beta_{W(z)}(\bm G_N)}\right|
 \le C\left(\frac1z+
 \frac{\log N}{z\log z}\right)}
 \tag{12.2}
\]
for an absolute constant $C$.
\end{lemma}

\begin{proof}
For $p>2$, write
\[
 c_p=1-\frac1{(p-1)^2},
 \qquad d_p=\frac p{p-1}.
\]
The local factor is $c_p$ when $p\nmid N$ and $d_p$ when $p\mid N$.  Therefore
\[
 \frac{\mathfrak S_G(N)}{\beta_{W(z)}(\bm G_N)}
 =\prod_{\substack{p>z\\p\nmid N}}c_p
  \prod_{\substack{p>z\\p\mid N}}d_p.
\]
For $p\ge5$,
$|\log c_p|\ll p^{-2}$ and $\log d_p\ll p^{-1}$.  The first tail is
$O(z^{-1})$.  If $k$ distinct primes greater than $z$ divide $N$, then
$z^k<N$, so $k\le\log N/\log z$.  Hence
\[
 \sum_{\substack{p>z\\p\mid N}}\frac1p
 \le\frac{k}{z}
 \le\frac{\log N}{z\log z}.
\]
Combining the two estimates proves the claim.
\end{proof}

\begin{definition}[Goldbach local scale]
For sufficiently large $N$, put
\[
 \boxed{
 z_N=\frac{\log N}{\sqrt{\log\log N}},
 \qquad W_N=\prod_{p\le z_N}p.}
 \tag{12.3}
\]
\end{definition}

\begin{theorem}[Subpolynomial moving local regulator]
\label{thm:Goldbach-regulator}
Along even $N\to\infty$,
\[
 \boxed{W_N=N^{o(1)},}
 \tag{12.4}
\]
and
\[
 \boxed{
 \beta_{W_N}(\bm G_N)
 =\mathfrak S_G(N)
  \left(1+O\left((\log\log N)^{-1/2}\right)\right)}
 \tag{12.5}
\]
uniformly in $N$.
\end{theorem}

\begin{proof}
By \Cref{lem:primorial},
$\log W_N=O(z_N)=o(\log N)$, proving \textup{(12.4)}.  In
\textup{(12.2)}, the first term tends to zero, while
\[
 \frac{\log N}{z_N\log z_N}
 =O\left((\log\log N)^{-1/2}\right).
\]
The logarithmic ratio is therefore of that order; exponentiation gives
\textup{(12.5)}.
\end{proof}

\begin{theorem}[Moving-sum Hardy--Littlewood bridge]
\label{thm:Goldbach-HL-bridge}
Suppose
\[
 \boxed{
 \mathfrak E_{\bm G_N}(N;W_N)\longrightarrow0}
 \tag{12.6}
\]
uniformly along the even integers.  Then
\[
 \boxed{
 \mathcal G(N)
 =(N-1)\mathfrak S_G(N)\bigl(1+o(1)\bigr)}
 \tag{12.7}
\]
along the even integers.  Equivalently, the error is
$o\!\left(N\mathfrak S_G(N)\right)$.
\end{theorem}

\begin{proof}
The exceptional contribution is at most
$2\omega(W_N)(\log N)^2$.  Since
$\omega(W_N)\le z_N$ and $z_N(\log N)^2=o(N)$, it is $o(N)$.  Apply
\textup{(9.12)}, \textup{(12.6)}, and \Cref{thm:Goldbach-regulator}.
\end{proof}

\begin{theorem}[Moving-sum existence bridge]
\label{thm:Goldbach-lower-bridge}
If there is an $\eta>0$ such that
\[
 \boxed{
 \mathfrak a_{\bm G_N}(N;W_N)\ge\eta}
 \tag{12.8}
\]
for every sufficiently large even $N$, then
\[
 \boxed{\mathcal G(N)\gg_\eta N}
 \tag{12.9}
\]
uniformly for large even $N$, and every sufficiently large even $N$ is a sum
of two primes.
\end{theorem}

\begin{proof}
For even $N$, \textup{(4.6)} gives
$\mathfrak S_G(N)\ge2C_2>0$.  By \textup{(12.5)}, the finite local products
are uniformly bounded below.  Insert \textup{(12.8)} into
\textup{(9.10)}.  The final prime-prime conclusion is
\Cref{thm:Goldbach-prime-cleanup}.
\end{proof}

\begin{remark}[The exact moving-boundary difference]
The local singular series can be tracked with a modulus $W_N=N^{o(1)}$, so the
moving local factors are not the central obstruction.  What changes is the
quantifier on the neutral amplitude:
\[
 \boxed{
 \text{fixed shift: one source as }X\to\infty,
 \qquad
 \text{Goldbach: every source }\bm G_N.}
 \tag{12.10}
\]
The Goldbach theorem therefore requires regulator-uniform control in the
external sum parameter.
\end{remark}

% =============================================================================
\section{Prime-power cleanup and genuine prime-pair consequences}
\label{sec:prime-power-cleanup}
% =============================================================================

The von Mangoldt observable counts prime powers.  The following estimates show
that a linear pair lower bound cannot be carried by proper prime powers.

\begin{lemma}[Number of proper prime powers]
\label{lem:prime-power-count}
The number of integers $m\le Y$ of the form $p^k$ with $p$ prime and $k\ge2$
is $O(\sqrt Y)$.
\end{lemma}

\begin{proof}
The squares contribute at most $\sqrt Y$.  For $k\ge3$, the number is at most
\[
 \sum_{3\le k\le\log_2Y}Y^{1/k}
 \le(\log_2Y)Y^{1/3}=O(\sqrt Y).
\]
\end{proof}

\begin{theorem}[Proper-prime-power contamination for a fixed affine pair]
\label{thm:fixed-prime-power-noise}
Let $\bm L$ be fixed and primitive, with $Y_X=O_{\bm L}(X)$.  The contribution
to $\mathcal C_{\bm L}(X)$ from occurrences for which at least one form value
is a proper prime power is
\[
 \boxed{O_{\bm L}(X^{1/2}(\log X)^2).}
 \tag{13.1}
\]
\end{theorem}

\begin{proof}
For each $i$, the equation $L_i(n)=m$ has at most one solution.  By
\Cref{lem:prime-power-count}, there are $O(\sqrt{Y_X})$ possible proper prime
powers $m$.  Each pair weight is at most $(\log Y_X)^2$.  Sum over the two
forms.
\end{proof}

\begin{theorem}[Weighted affine lower bound forces genuine prime pairs]
\label{thm:weighted-to-primes}
For a fixed primitive affine pair, if
\[
 \mathcal C_{\bm L}(X)\ge cX
 \tag{13.2}
\]
for all sufficiently large $X$ and some $c>0$, then there are infinitely many
$n$ for which both $L_1(n)$ and $L_2(n)$ are prime.  More quantitatively, the
number of such $n\le X$ is $\gg_{\bm L,c}X/(\log X)^2$.
\end{theorem}

\begin{proof}
By \Cref{thm:fixed-prime-power-noise}, proper prime powers contribute $o(X)$.
The genuine prime--prime weighted contribution is therefore at least $cX/2$
for large $X$.  Each genuine pair has weight at most $O_{\bm L}((\log X)^2)$,
which gives the counting lower bound.
\end{proof}

\begin{theorem}[Goldbach prime-power contamination]
\label{thm:Goldbach-prime-cleanup}
Let
\[
 \mathcal G_{\PP}(N)
 =\sum_{\substack{p+q=N\\p,q\ \mathrm{prime}}}\log p\log q.
 \tag{13.3}
\]
Then
\[
 \boxed{
 \mathcal G(N)-\mathcal G_{\PP}(N)
 =O(\sqrt N(\log N)^2).}
 \tag{13.4}
\]
Consequently, if $\mathcal G(N)\ge cN$ for all sufficiently large even $N$,
then every such $N$ is a sum of two primes; indeed the number of ordered prime
representations is $\gg_cN/(\log N)^2$.
\end{theorem}

\begin{proof}
If a non-prime summand has nonzero von Mangoldt weight, it is a proper prime
power.  There are $O(\sqrt N)$ such integers by
\Cref{lem:prime-power-count}; for each one the complementary summand is fixed,
and the weight is at most $(\log N)^2$.  This proves \textup{(13.4)}.  A
linear lower bound dominates the error, and each genuine representation has
weight at most $(\log N)^2$.
\end{proof}

\begin{corollary}[Conditional twin-prime and strong-Goldbach conclusions]
\label{cor:headline-conditional}
The fixed-shift amplitude condition \textup{(10.6)} with $h=2$ implies the
twin-prime conjecture.  The moving-sum amplitude condition \textup{(12.8)}
implies the strong Goldbach conjecture for all sufficiently large even
integers; the finitely many remaining cases are a finite verification.
\end{corollary}

\begin{proof}
Combine \Cref{thm:fixed-lower-bridge,thm:weighted-to-primes} for twins and
\Cref{thm:Goldbach-lower-bridge,thm:Goldbach-prime-cleanup} for Goldbach.
\end{proof}

% =============================================================================


% =============================================================================
\section{The outer M\"obius--log switch for twin correlations}
\label{sec:affine-outer}
% =============================================================================

\subsection{Exact outer M\"obius--log compiler}
\label{subsec:affouter-outer-compiler}
% -----------------------------------------------------------------------------

Put
\[
 \mathscr I_X^+
 :=\{m\in\mathbb Z:X+2<m\le2X+2\}.
\]
Then the sharp twin correlation is exactly
\begin{equation}
 \label{eq:affouter-shifted-pair}
 \boxed{
 \mathcal C_X(2)
 =\sum_{m\in\mathscr I_X^+}\Lambda(m-2)\Lambda(m).}
 \tag{A.1}
\end{equation}
Use the exact local density,
\[
 g_2(d)=
 \begin{cases}
  \varphi(d)^{-1},&(d,2)=1,\\
  0,&(d,2)>1,
 \end{cases}
\]
and put
\[
 \gamma_2(d):=d\,g_2(d).
\]
For every $d\ge1$ define the actual, local, and completed outer rows
\begin{equation}
 \label{eq:affouter-outer-rows}
 \boxed{
 \begin{aligned}
 \mathcal A_X^{\rm out}(d)
 &:=\sum_{\substack{q\ge1\\dq\in\mathscr I_X^+}}
      (\log q)\Lambda(dq-2),\\
 \mathcal L_X^{\rm out}(d)
 &:=\gamma_2(d)
      \sum_{\substack{q\ge1\\dq\in\mathscr I_X^+}}
      \log q,\\
 \mathcal E_X^{\rm out}(d)
 &:=\mathcal A_X^{\rm out}(d)
   -\mathcal L_X^{\rm out}(d).
 \end{aligned}}
 \tag{A.2}
\end{equation}

\begin{theorem}[Exact outer M\"obius--log assembly]
\label{thm:affouter-outer-assembly}
One has
\begin{equation}
 \label{eq:affouter-actual-outer-assembly}
 \boxed{
 \mathcal C_X(2)
 =\sum_{d\le2X+2}\mu(d)\mathcal A_X^{\rm out}(d).}
 \tag{A.3}
\end{equation}
Moreover,
\begin{equation}
 \label{eq:affouter-local-outer-assembly}
 \boxed{
 \sum_{d\le2X+2}\mu(d)\mathcal L_X^{\rm out}(d)
 =\mathfrak S(2)X+o(X).}
 \tag{A.4}
\end{equation}
Consequently,
\begin{equation}
 \label{eq:affouter-residual-outer-assembly}
 \boxed{
 R_X(2)
 =\sum_{d\le2X+2}\mu(d)\mathcal E_X^{\rm out}(d)
  +o(X).}
 \tag{A.5}
\end{equation}
\end{theorem}

\begin{proof}
The finite convolution identity
\[
 \Lambda(m)
 =\sum_{d\mid m}\mu(d)\log\frac md
\]
gives
\[
 \Lambda(m)
 =\sum_{dq=m}\mu(d)\log q.
\]
Insert this identity in \textup{(A.1)} and interchange the finite sums to
obtain \textup{(A.3)}.

For the local statement define
\[
 F_2(n)
 :=\sum_{dq=n}\mu(d)\gamma_2(d)\log q.
\]
For $\Re s>1$ its Dirichlet series is
\[
 \begin{aligned}
 \sum_{n\ge1}\frac{F_2(n)}{n^s}
 &=-\zeta'(s)
   \sum_{d\ge1}\frac{\mu(d)\gamma_2(d)}{d^s}\\
 &=-\frac{\zeta'(s)}{\zeta(s)}\,\mathscr H_2(s),
 \end{aligned}
\]
where
\begin{equation}
 \label{eq:affouter-local-Euler-factor}
 \boxed{
 \mathscr H_2(s)
 =\frac1{1-2^{-s}}
  \prod_{p>2}
  \frac{1-p^{1-s}/(p-1)}{1-p^{-s}}.}
 \tag{A.6}
\end{equation}
The Euler product multiplying $-\zeta'/\zeta$ is analytic and bounded in the
classical zero-free region of the classical de la Vall\'ee Poussin type, because each odd-prime factor
is $1+O(p^{-1-\Re s})$ there.  At $s=1$,
\[
 \mathscr H_2(1)
 =2\prod_{p>2}\frac{p(p-2)}{(p-1)^2}
 =\mathfrak S(2).
\]
The same Perron shift and zero-free-region estimate used for
\Cref{thm:zero-free-input} gives some $c>0$ such that
\[
 \sum_{n\le y}F_2(n)
 =\mathfrak S(2)y
  +O\!\left(ye^{-c\sqrt{\log y}}\right).
\]
Subtract this formula at $2X+2$ and $X+2$ to obtain
\textup{(A.4)}.  Subtracting \textup{(A.4)} from
\textup{(A.3)} proves \textup{(A.5)}.
\end{proof}

\begin{remark}[Relation to the complete source]
A full Vaughan decomposition uses two auxiliary cutoffs and leaves a completed
three-factor balanced packet.  Identity \textup{(A.5)} starts from the same
symmetric von-Mangoldt source, but first absorbs the \emph{entire}
square-root progression head.  The resulting tail has only the large divisor
and its complementary quotient.  It is an exact reassembly of the Type-I and
balanced coordinates, not a new arithmetic assumption.
\end{remark}

% -----------------------------------------------------------------------------
\subsection{The complete square-root head closes with the logarithmic quotient weight}
\label{subsec:affouter-head}
% -----------------------------------------------------------------------------

Let
\begin{equation}
 \label{eq:affouter-DX}
 \boxed{
 D_X=X^{1/2}(\log X)^{-B}}
 \tag{A.7}
\end{equation}
with $B$ chosen sufficiently large for the maximal von-Mangoldt
maximal Bombieri--Vinogradov input of \Cref{thm:weighted-BV-input}.

\begin{theorem}[Weighted square-root head closure]
\label{thm:affouter-head-closure}
For every fixed $A>0$, after increasing $B$ if necessary,
\begin{equation}
 \label{eq:affouter-head-closure}
 \boxed{
 \sum_{d\le D_X}\mu(d)\mathcal E_X^{\rm out}(d)
 \ll_A\frac{X}{(\log X)^A}.}
 \tag{A.8}
\end{equation}
\end{theorem}

\begin{proof}
Suppose first that $d$ is odd.  Put
\[
 w_d(t)=\log\frac{t+2}{d}
 \qquad(X<t\le2X).
\]
Then
\[
 \mathcal A_X^{\rm out}(d)
 =\sum_{\substack{X<n\le2X\\n\equiv-2\ (d)}}
   \Lambda(n)w_d(n).
\]
Let
\[
 E_d(t)
 =\sum_{\substack{X<n\le t\\n\equiv-2\ (d)}}\Lambda(n)
  -\frac{t-X}{\varphi(d)}.
\]
Stieltjes summation by parts gives
\[
 \mathcal A_X^{\rm out}(d)
 -\frac1{\varphi(d)}\int_X^{2X}w_d(t)\,\dd t
 \ll
 (\log X)\max_{X\le t\le2X}|E_d(t)|.
\]
The integral and the discrete local row differ by the Riemann-sum error
\[
 \frac1{\varphi(d)}\int_X^{2X}w_d(t)\,\dd t
 -\gamma_2(d)
  \sum_{\substack{q\ge1\\dq\in\mathscr I_X^+}}\log q
 \ll\frac{d\log X}{\varphi(d)}.
\]
Therefore
\[
 |\mathcal E_X^{\rm out}(d)|
 \ll
 (\log X)\max_{X\le t\le2X}|E_d(t)|
 +\frac{d\log X}{\varphi(d)}.
\]
Summing the first term over $d\le D_X$ uses \Cref{thm:weighted-BV-input} with an arbitrarily large
logarithmic exponent.  The elementary estimate
\[
 \sum_{d\le D_X}\frac{d}{\varphi(d)}\ll D_X\log\log(3D_X)
\]
makes the second contribution $X^{1/2+o(1)}=o(X(\log X)^{-A})$.

If $d$ is even, the local row vanishes.  A nonzero actual term requires
$dq-2$ to be a power of two.  The crude bound
\[
 \sum_{\substack{d\le D_X\\2\mid d}}
 |\mathcal A_X^{\rm out}(d)|
 \ll D_X(\log X)^3
\]
is already $o(X(\log X)^{-A})$ for each fixed $A$.  Combining the odd and
even rows proves \textup{(A.8)}.
\end{proof}

\begin{remark}[Why the local row is part of the head]
The actual weighted progression and the local density
$\gamma_2(d)=d/\varphi(d)$ must be subtracted before the modulus is declared
small or large.  Estimating
$\mathcal A_X^{\rm out}(d)$ and
$\mathcal L_X^{\rm out}(d)$ separately would discard the singular-series
orientation already assembled in \textup{(A.5)}.
\end{remark}

% -----------------------------------------------------------------------------
\subsection{Exact target-free short-quotient tail}
\label{subsec:affouter-tail}
% -----------------------------------------------------------------------------

Put
\[
 Q_X:=\left\lfloor\frac{2X+2}{D_X}\right\rfloor
 =X^{1/2+o(1)}.
\]
For $2\le q\le Q_X$ define the completed quotient row
\begin{equation}
 \label{eq:affouter-quotient-row}
 \boxed{
 \begin{aligned}
 \mathfrak r_X(q)
 :={}&
 \sum_{\substack{d>D_X\\dq\in\mathscr I_X^+}}
 \mu(d)\Lambda(dq-2)
 \\
 &-
 \sum_{\substack{d>D_X\\dq\in\mathscr I_X^+}}
 \mu(d)\gamma_2(d).
 \end{aligned}}
 \tag{A.9}
\end{equation}
The second sum is automatically supported on odd $d$.

\begin{definition}[Affine M\"obius--prime bilinear packet]
\label{def:affouter-AMP-packet}
Define
\begin{equation}
 \label{eq:affouter-AMP-packet}
 \boxed{
 \mathfrak T_X^{\rm AMP}
 :=\sum_{q=2}^{Q_X}(\log q)\mathfrak r_X(q).}
 \tag{A.10}
\end{equation}
\end{definition}

\begin{theorem}[Source-minimal short-quotient reduction]
\label{thm:affouter-tail-reduction}
One has
\begin{equation}
 \label{eq:affouter-tail-reduction}
 \boxed{
 R_X(2)=\mathfrak T_X^{\rm AMP}+o(X).}
 \tag{A.11}
\end{equation}
\end{theorem}

\begin{proof}
Split the exact sum in \textup{(A.5)} at $D_X$.  The head is negligible by
\Cref{thm:affouter-head-closure}.  In the tail, interchange the finite $d$- and
$q$-sums.  The coefficient of $q=1$ is $\log1=0$, so the tail begins at
$q=2$ and is exactly \textup{(A.10)}.  This proves \textup{(A.11)}.\end{proof}

\begin{theorem}[No target-return atom]
\label{thm:affouter-no-target-return}
Every nonzero actual occurrence in $\mathfrak T_X^{\rm AMP}$ satisfies
\[
 m=dq\in\mathscr I_X^+,
 \qquad d>D_X>1,
 \qquad q\ge2.
\]
Hence the left endpoint $m$ is composite.  The section $q=1$, which would
contain the prime left endpoint, vanishes identically before any estimate is
made.

After removing proper prime powers at the right endpoint,
\begin{equation}
 \label{eq:affouter-prime-cleanup}
 \boxed{
 \mathfrak T_X^{\rm AMP}
 =\mathfrak T_X^{\rm AMP,pr}+o(X),}
 \tag{A.13}
\end{equation}
where every actual occurrence of
$\mathfrak T_X^{\rm AMP,pr}$ has
\[
 dq\text{ composite},
 \qquad dq-2\text{ prime}.
\]
Thus the outer M\"obius--log packet is target free in the literal occurrence sense.
\end{theorem}

\begin{proof}
The first statement follows directly from the support in \textup{(A.9)} and the exact
factor $\log q$.

For the proper-prime-power contribution write $dq-2=p^k$ with $k\ge2$.
For one such endpoint $m=p^k+2$, the total absolute factorization weight is
bounded by
\[
 \sum_{d\mid m}\mu(d)^2
 \log\frac md\,\log p
 \ll \tau(m)(\log X)^2.
\]
There are $O(X^{1/2})$ proper prime powers in $(X,2X]$, and
$\tau(m)=X^{o(1)}$ uniformly.  Their total contribution is
$X^{1/2+o(1)}=o(X)$, proving \textup{(A.13)}.
\end{proof}

\begin{remark}[The first genuine analytic boundary]
The outer source has exactly the geometry sought at the end of the preceding parity analysis:
\[
 \boxed{
 \mu(d)
 \times
 \bigl[\Lambda(dq-2)-\gamma_2(d)\bigr]
 \times
 \log q,
 \qquad
 d>X^{1/2-o(1)},
 \quad
 2\le q\le X^{1/2+o(1)}.}
\]
It is a signed M\"obius--prime bilinear packet, with the local affine row
inside the same bracket.  Neither factor coordinate reproduces the twin
source.
\end{remark}

% -----------------------------------------------------------------------------
\subsection{The first literal affine bilinear certificate}
\label{subsec:affouter-bilinear-certificate}
% -----------------------------------------------------------------------------

\begin{definition}[Affine M\"obius--prime bilinear certificate]
\label{def:affouter-bilinear-certificate}
The source-native bilinear certificate at the square-root head is
\begin{equation}
 \label{eq:affouter-bilinear-certificate}
 \boxed{
 \mathsf{AMPB}_X:
 \qquad
 \sum_{q=2}^{Q_X}(\log q)
 \left[
  \sum_{\substack{d>D_X\\dq\in\mathscr I_X^+}}
  \mu(d)
  \bigl(\Lambda(dq-2)-\gamma_2(d)\bigr)
 \right]
 =o(X).}
 \tag{A.14}
\end{equation}
The cutoff $D_X$, quotient range, M\"obius orientation, logarithmic writer,
and local density are fixed before the protected coefficient is evaluated.
No coefficient may be chosen from the signs of the returned arithmetic row.
\end{definition}

\begin{theorem}[Exact equivalence of the first bilinear certificate]
\label{thm:affouter-bilinear-equivalence}
The following statements are equivalent:
\begin{enumerate}[label=\textup{(E\arabic*)},leftmargin=3em]
\item the affine M\"obius--prime certificate $\mathsf{AMPB}_X$ holds;
\item $\mathfrak T_X^{\rm AMP}=o(X)$;
\item $R_X(2)=o(X)$.
\end{enumerate}
Thus \textup{(A.14)} is a literal signed asymptotic-sieve certificate on a
target-free occurrence carrier whose verification closes the twin shift.
\end{theorem}

\begin{proof}
The first two statements are identical by
\Cref{def:affouter-bilinear-certificate,def:affouter-AMP-packet}.
Their equivalence with the third follows from the exact reduction
\textup{(A.11)}.
\end{proof}

\begin{remark}[Meaning of the word ``axiom'']
\textup{(A.14)} is not asserted as a new assumption of the Renewal Gran tensor and is
not proved by finite arithmetic loading.  It is the explicitly isolated
conjecture-specific quantitative certificate.  Its value is that it names the
precise signed bilinear row to which an asymptotic-sieve, dispersion,
Kloosterman-fraction, or other analytic theorem must apply, without a target
atom or an unassembled local counterterm.
\end{remark}

% -----------------------------------------------------------------------------
\subsection{Positive fallbacks and the complete two-occurrence Gram}
\label{subsec:affouter-dispersion}
% -----------------------------------------------------------------------------

Let
\[
 h_X(q)=\log q,
 \qquad
 r_X(q)=\mathfrak r_X(q),
 \qquad
 2\le q\le Q_X.
\]
Then
\[
 \mathfrak T_X^{\rm AMP}=\langle h_X,r_X\rangle_{\ell^2}.
\]

\begin{corollary}[A sufficient diagonal quotient-dispersion scale]
\label{cor:affouter-L2-sufficient}
Since
\[
 \sum_{q=2}^{Q_X}(\log q)^2
 =Q_X(\log Q_X)^2+O(Q_X\log Q_X),
\]
the estimate
\begin{equation}
 \label{eq:affouter-L2-sufficient}
 \boxed{
 \sum_{q=2}^{Q_X}|\mathfrak r_X(q)|^2
 =o\!\left(
 \frac{X^2}{Q_X(\log Q_X)^2}
 \right)}
 \tag{A.15}
\end{equation}
is sufficient for \textup{(A.14)}.
\end{corollary}

\begin{proof}
Apply Cauchy--Schwarz to
$\langle h_X,r_X\rangle$ and use the displayed norm of $h_X$.
\end{proof}

For $(d,q)$ in the outer tail put
\[
 \xi_X(d,q)
 :=\one_{\mathscr I_X^+}(dq)
   \bigl[\Lambda(dq-2)-\gamma_2(d)\bigr].
\]
Then
\[
 r_X(q)=\sum_{d>D_X}\mu(d)\xi_X(d,q).
\]

\begin{theorem}[Complete two-occurrence quotient-dispersion Gram]
\label{thm:affouter-complete-Gram}
The positive quantity in \textup{(A.15)} is exactly
\begin{equation}
 \label{eq:affouter-complete-Gram}
 \boxed{
 \sum_{q=2}^{Q_X}|\mathfrak r_X(q)|^2
 =\sum_{d,d'>D_X}\mu(d)\mu(d')
   \mathbb G_X^{\rm quot}(d,d'),}
 \tag{A.16}
\end{equation}
where
\begin{equation}
 \label{eq:affouter-Gram-block}
 \boxed{
 \mathbb G_X^{\rm quot}(d,d')
 :=\sum_{q=2}^{Q_X}
   \xi_X(d,q)\xi_X(d',q)
 \succeq0}
 \tag{A.17}
\end{equation}
is the Gram of the complete completed-row vectors
$q\mapsto\xi_X(d,q)$.

Expanding \textup{(A.17)} produces all four blocks
\[
 \boxed{
 \text{actual--actual}
 -\text{actual--local}
 -\text{local--actual}
 +\text{local--local}.}
\]
Every common quotient, common endpoint, common factor, nonreduced, and local
row is therefore present before the square is taken.
\end{theorem}

\begin{proof}
Expand
\[
 \sum_q
 \left|
  \sum_d\mu(d)\xi_X(d,q)
 \right|^2
\]
and interchange the finite sums.  Positivity follows because
$\mathbb G_X^{\rm quot}$ is the Gram matrix of the row vectors
$\xi_X(d,\cdot)$.
\end{proof}

\begin{proposition}[Per-divisor or per-quotient energies are not source minimal]
\label{cth:affouter-diagonal-no-go}
Neither
\[
 \sum_d\mathbb G_X^{\rm quot}(d,d)
\]
nor
\[
 \sum_q(\log q)^2|\mathfrak r_X(q)|^2
\]
is determined by the protected signed coefficient \textup{(A.10)}.  In finite
dimension two nonzero rows can cancel exactly in the protected sum while
both diagonal energies remain positive.  Consequently a dispersion proof
must retain the complete $d$--$d'$ Gram or prove a stronger theorem explicitly.
\end{proposition}

\begin{proof}
Take two row vectors $v_1=v$ and $v_2=-v$ with equal M\"obius orientation.
Their protected sum vanishes, whereas
$\|v_1\|^2+\|v_2\|^2=2\|v\|^2>0$.
The same example applies after grouping by $q$.
\end{proof}

% -----------------------------------------------------------------------------
\subsection{Positive realization of the signed bilinear coefficient}
\label{subsec:affouter-positive-realization}
% -----------------------------------------------------------------------------

On the occurrence set
\[
 \Omega_X^{\rm tail}
 =\{(d,q):d>D_X,\ q\ge2,\ dq\in\mathscr I_X^+\},
\]
put
\[
 a_{d,q}=\mu(d)^2(\log q)\Lambda(dq-2),
 \qquad
 \ell_{d,q}=\mu(d)^2(\log q)\gamma_2(d).
\]
Both are nonnegative.  Define
\[
 x_X=(\sqrt a)\oplus(\sqrt\ell),
\]
\[
 y_X=\bigl(\mu(d)\sqrt a\bigr)
     \oplus\bigl(-\mu(d)\sqrt\ell\bigr).
\]

\begin{theorem}[One positive polarization packet]
\label{thm:affouter-polarization}
The outer coefficient is the mixed entry
\begin{equation}
 \label{eq:affouter-polarization}
 \boxed{
 \langle x_X,y_X\rangle
 =\mathfrak T_X^{\rm AMP}.}
 \tag{A.18}
\end{equation}
Hence
\begin{equation}
 \label{eq:affouter-positive-2x2}
 \boxed{
 \begin{pmatrix}
  \|x_X\|^2&\mathfrak T_X^{\rm AMP}\\
  \overline{\mathfrak T_X^{\rm AMP}}&\|y_X\|^2
 \end{pmatrix}
 \succeq0.}
 \tag{A.19}
\end{equation}
The complete signed coefficient is recoverable from the four positive
polarization energies of $x_X\pm y_X$ and $x_X\pm iy_X$.
\end{theorem}

\begin{proof}
Direct multiplication gives
\[
 \langle x_X,y_X\rangle
 =\sum_{d,q}\mu(d)(\log q)
  \bigl[\Lambda(dq-2)-\gamma_2(d)\bigr],
\]
which is exactly \textup{(A.10)}.  The block in \textup{(A.19)} is the Gram of $x_X,y_X$ and is
therefore positive.  The polarization identity gives the last statement.
\end{proof}

\begin{remark}[Why positivity still does not close the packet]
The norms of $x_X$ and $y_X$ retain total actual and local mass.  They do not
determine the M\"obius orientation in the mixed entry.  A first-kernel
crossing or a common action becomes informative only if it is populated on
this occurrence carrier independently of the returned value of \textup{(A.18)}.
Choosing a rank-one action from $x_X,y_X$ after reading the coefficient would
again rewrite the target rather than estimate it.
\end{remark}

% -----------------------------------------------------------------------------
\subsection{Conditional closure theorem and finite obstruction}
\label{subsec:affouter-closure}
% -----------------------------------------------------------------------------

\begin{theorem}[Affine bilinear closure]
\label{thm:affouter-closure}
Choose $D_X$ as in \textup{(A.7)}.  Each of the following implies
$R_X(2)=o(X)$ and hence, through the previously constructed prime-power handoff, the
Hardy--Littlewood twin-prime asymptotic:
\begin{enumerate}[label=\textup{(C\arabic*)},leftmargin=3em]
\item the signed affine M\"obius--prime certificate \textup{(A.14)};
\item the quotient-dispersion estimate \textup{(A.15)};
\item an independently populated positive action $L_X$ on the complete
      quotient packet satisfies
\[
 \boxed{
 \langle r_X,L_Xr_X\rangle
 \left\langle h_X,L_X^{\dagger}h_X\right\rangle
 =o(X^2),}
\]
      with $h_X(q)=\log q$ and the action fixed before $r_X$ is read;
\item a stronger asymptotic-sieve, bilinear-dispersion, or
      Kloosterman-fraction theorem proves the same completed mixed coefficient
      after preserving every block in \textup{(A.16)}--\textup{(A.17)}.
\end{enumerate}
\end{theorem}

\begin{proof}
The first implication is
\Cref{thm:affouter-bilinear-equivalence}.  The second is
\Cref{cor:affouter-L2-sufficient}.  The third is the protected-observable Riesz
inequality applied to $\langle h_X,r_X\rangle$.  The fourth is a restatement
of the first condition with a stronger independently proved analytic input.
\end{proof}

\begin{theorem}[Finite obstruction returned by failure]
\label{thm:affouter-obstruction}
If
\[
 |R_X(2)|\ge\eta X
\]
along an unbounded sequence, then for all sufficiently large $X$ on that
sequence,
\begin{equation}
 \label{eq:affouter-obstruction}
 \boxed{
 \left|
 \sum_{q=2}^{Q_X}(\log q)
 \sum_{\substack{d>D_X\\dq\in\mathscr I_X^+}}
 \mu(d)
 \bigl[\Lambda(dq-2)-\gamma_2(d)\bigr]
 \right|
 \ge\frac\eta2X.}
 \tag{A.20}
\end{equation}
The returned witness contains:
\begin{enumerate}[label=\textup{(W\arabic*)},leftmargin=3em]
\item one fixed square-root cutoff $D_X$;
\item one quotient bank $2\le q\le X^{1/2+o(1)}$;
\item one long M\"obius coordinate $d>D_X$;
\item one prime right endpoint $dq-2$, after an $o(X)$ cleanup;
\item the complete local density $\gamma_2(d)$ inside the same bracket;
\item no prime left endpoint and no target-return section.
\end{enumerate}
\end{theorem}

\begin{proof}
Use \textup{(A.11)} and absorb its $o(X)$ error into $\eta X/2$.  The six structural
properties are
\Cref{thm:affouter-tail-reduction,thm:affouter-no-target-return} and the definitions of
$D_X,Q_X,\gamma_2$.
\end{proof}

The finite tensor identities prove none of the quantitative conditions in
\Cref{thm:affouter-closure} and therefore prove neither twin primes, general
Polignac, nor strong Goldbach.



% =============================================================================
\section{Endpoint fusion and the retuning frontier}
\label{sec:affine-endpoint}
% =============================================================================

\subsection{Exact endpoint fusion of the outer M\"obius--log packet}
\label{subsec:affend-endpoint-fusion}
% -----------------------------------------------------------------------------

For a real threshold $D\ge1$ and $m\ge2$, define the actual and local
endpoint writers
\begin{equation}
 \label{eq:affend-endpoint-writers}
 \boxed{
 \begin{aligned}
 \Theta_D(m)
 &:=\sum_{\substack{d\mid m\\d>D}}
     \mu(d)\log\frac md,\\
 \Gamma_D(m)
 &:=\sum_{\substack{d\mid m\\d>D}}
     \mu(d)\gamma_2(d)\log\frac md.
 \end{aligned}}
 \tag{E.1}
\end{equation}
Put
\begin{equation}
 \label{eq:affend-endpoint-row}
 \boxed{
 \mathfrak e_{X,D}(m)
 :=\Theta_D(m)\Lambda(m-2)-\Gamma_D(m),}
 \tag{E.2}
\end{equation}
and
\begin{equation}
 \label{eq:affend-endpoint-packet}
 \boxed{
 \mathfrak T_{X,D}^{\rm end}
 :=\sum_{m\in\mathscr I_X^+}\mathfrak e_{X,D}(m).}
 \tag{E.3}
\end{equation}

\begin{theorem}[Exact quotient-to-endpoint fusion]
\label{thm:affend-endpoint-fusion}
For every $D\ge1$,
\begin{equation}
 \label{eq:affend-fusion-identity}
 \boxed{
 \mathfrak T_{X,D}^{\rm end}
 =\sum_{\substack{d>D,\ q\ge1\\dq\in\mathscr I_X^+}}
   \mu(d)(\log q)
   \bigl[\Lambda(dq-2)-\gamma_2(d)\bigr].}
 \tag{E.4}
\end{equation}
In particular, at the square-root threshold,
\begin{equation}
 \label{eq:affend-affouter-equivalence}
 \boxed{
 \mathfrak T_{X,D_X}^{\rm end}
 =\mathfrak T_X^{\rm AMP}.}
 \tag{E.5}
\end{equation}
Consequently,
\begin{equation}
 \label{eq:affend-residual-endpoint}
 \boxed{
 R_X(2)
 =\mathfrak T_{X,D_X}^{\rm end}+o(X).}
 \tag{E.6}
\end{equation}
\end{theorem}

\begin{proof}
In the finite sum on the right of \textup{(E.4)}, set $m=dq$ and group all
terms with the same product.  Since $q=m/d$, the actual coefficient at $m$ is
\[
 \Lambda(m-2)
 \sum_{\substack{d\mid m\\d>D}}
 \mu(d)\log\frac md
 =\Theta_D(m)\Lambda(m-2),
\]
and the local coefficient is $\Gamma_D(m)$.  This proves
\textup{(E.4)}.  At $D=D_X$, the $q=1$ summands vanish because
$\log1=0$, so the remaining range is exactly \textup{(A.10)}.  Finally use \textup{(A.11)}.
\end{proof}

\begin{theorem}[Prime-null endpoint support]
\label{thm:affend-prime-null}
For every prime $p$ and every $D\ge1$,
\begin{equation}
 \label{eq:affend-prime-null}
 \boxed{
 \Theta_D(p)=\Gamma_D(p)=\mathfrak e_{X,D}(p)=0.}
 \tag{E.7}
\end{equation}
More generally,
\begin{equation}
 \label{eq:affend-rad-support}
 \boxed{
 \rad(m)\le D
 \quad\Longrightarrow\quad
 \Theta_D(m)=\Gamma_D(m)=0.}
 \tag{E.8}
\end{equation}
Thus every nonzero actual endpoint in \textup{(E.6)} has a composite left
endpoint with radical exceeding $D_X$.
\end{theorem}

\begin{proof}
For a prime $p$, the only divisor exceeding $D$ can be $d=p$; its complement
is one and therefore its logarithmic coefficient is zero.  If
$\rad(m)\le D$, every divisor with nonzero M\"obius coefficient is at most
$D$, so both sums in \textup{(E.1)} are empty.
\end{proof}

\begin{remark}[Target-free status after fusion]
The disappearance of the prime left endpoint is now an endpoint theorem, not
an accident of the quotient coordinate.  No factorization section of
\textup{(E.4)} can reproduce the twin target after the fibre has been
coherently fused.
\end{remark}

% -----------------------------------------------------------------------------
\subsection{Complementary-divisor normal form and explicit shallow rows}
\label{subsec:affend-complementary-normal-form}
% -----------------------------------------------------------------------------

For a squarefree integer $r$ and $Y>0$, put
\[
 M_r(Y):=\sum_{\substack{e\mid r\\e<Y}}\mu(e),
 \qquad
 L_r(Y):=\sum_{\substack{e\mid r\\e<Y}}\mu(e)\log e.
\]

\begin{theorem}[Radical complementary-divisor formula]
\label{thm:affend-radical-formula}
Let $r=\rad(m)$ and $c=m/r$.  Then
\begin{equation}
 \label{eq:affend-radical-theta}
 \boxed{
 \Theta_D(m)
 =\mu(r)
 \left[
  (\log c)M_r\!\left(\frac rD\right)
  +L_r\!\left(\frac rD\right)
 \right].}
 \tag{E.9}
\end{equation}
The local endpoint writer is
\begin{equation}
 \label{eq:affend-radical-gamma}
 \boxed{
 \Gamma_D(m)
 =\mu(r)
 \sum_{\substack{e\mid r\\e<r/D}}
 \mu(e)\gamma_2(r/e)
 \left(\log c+\log e\right).}
 \tag{E.10}
\end{equation}
\end{theorem}

\begin{proof}
Every divisor contributing to \textup{(E.1)} is squarefree and hence divides
$r$.  Set $e=r/d$.  Since $e$ and $d$ are complementary divisors of the
squarefree integer $r$,
\[
 \mu(d)=\mu(r/e)=\mu(r)\mu(e).
\]
Moreover
\[
 d>D\iff e<r/D,
 \qquad
 \log\frac md=\log\frac mr+\log e=\log c+\log e.
\]
Substitution gives both formulas.
\end{proof}

\begin{corollary}[Exact semiprime and prime-power rows]
\label{cor:affend-shallow-rows}
Let $p<q$ be distinct odd primes.  Then
\begin{equation}
 \label{eq:affend-semiprime-theta}
 \boxed{
 \Theta_D(pq)=
 \begin{cases}
  0,&q\le D,\\
  -\log p,&p\le D<q,\\
  -\log(pq),&D<p.
 \end{cases}}
 \tag{E.11}
\end{equation}
and
\begin{equation}
 \label{eq:affend-semiprime-gamma}
 \boxed{
 \Gamma_D(pq)=
 \begin{cases}
  0,&q\le D,\\
  -\dfrac{q}{q-1}\log p,&p\le D<q,\\[0.8em]
  -\dfrac{p}{p-1}\log q
  -\dfrac{q}{q-1}\log p,&D<p.
 \end{cases}}
 \tag{E.12}
\end{equation}
For a prime power $p^a$, $a\ge2$,
\begin{equation}
 \label{eq:affend-prime-power-row}
 \boxed{
 \Theta_D(p^a)=
 \begin{cases}
  -(a-1)\log p,&p>D,\\
  0,&p\le D,
 \end{cases}}
 \tag{E.13}
\end{equation}
with the analogous local coefficient obtained by multiplying the nonzero row
by $p/(p-1)$.
\end{corollary}

\begin{proof}
List the divisors with nonzero M\"obius coefficient.  For $pq$ the term
$d=pq$ always has logarithmic complement zero.  The remaining terms are
$d=p$ and $d=q$, included exactly in the stated threshold ranges.  For $p^a$
only $d=1,p$ contribute to the M\"obius sum; when $p>D$, the $d=p$ term is
$-\log(p^{a-1})$.
\end{proof}

\begin{remark}[No factor-depth short is declared]
The formulas above are diagnostics of the fused endpoint coefficient.  The endpoint analysis
does not estimate the semiprime, prime-power, or higher-depth rows separately.
Their cross-depth cancellation remains part of the completed endpoint packet.
Proper prime powers remain $o(X)$ by the previously constructed cleanup.
\end{remark}

% -----------------------------------------------------------------------------
\subsection{Exact quotient-fibre short and read-null innovation}
\label{subsec:affend-fibre-short}
% -----------------------------------------------------------------------------

For $m\in\mathscr I_X^+$ put
\[
 \Omega_D(m)=\{d:d\mid m,\,d>D\},
 \qquad n_D(m)=|\Omega_D(m)|.
\]
On a nonempty fibre define
\[
 a_{X,D;m}(d)
 :=\mu(d)\log\frac md
   \bigl[\Lambda(m-2)-\gamma_2(d)\bigr].
\]
Let $F_m$ be the fibre-sum map and let $P_m$ be orthogonal projection onto the
constant line in $\ell^2(\Omega_D(m))$.

\begin{theorem}[Endpoint-fusion Pythagoras]
\label{thm:affend-fibre-Pythagoras}
For every nonempty fibre,
\begin{equation}
 \label{eq:affend-fibre-sum}
 \boxed{
 F_ma_{X,D;m}
 =\mathfrak e_{X,D}(m).}
 \tag{E.14}
\end{equation}
Moreover,
\begin{equation}
 \label{eq:affend-fibre-Pythagoras}
 \boxed{
 \sum_{d\in\Omega_D(m)}|a_{X,D;m}(d)|^2
 =\frac{|\mathfrak e_{X,D}(m)|^2}{n_D(m)}
  +\|(I-P_m)a_{X,D;m}\|_2^2.}
 \tag{E.15}
\end{equation}
The second term is positive factorization innovation invisible to the endpoint
read.
\end{theorem}

\begin{proof}
The first identity is the definition of the endpoint-fused row.  The
projection of $a$ onto the constant line is
\[
 P_ma=\frac{F_ma}{n_D(m)}\one_{\Omega_D(m)}.
\]
Orthogonality of $P_m$ and $I-P_m$ gives \textup{(E.15)}.
\end{proof}

\begin{corollary}[Minimum factor-history cost at fixed endpoint]
\label{cor:affend-minimum-fibre-cost}
Among all coefficient vectors $b$ on the fibre satisfying
$F_mb=\mathfrak e_{X,D}(m)$, the minimum Euclidean cost is
\begin{equation}
 \label{eq:affend-minimum-fibre-cost}
 \boxed{
 \inf_{F_mb=\mathfrak e_{X,D}(m)}\|b\|_2^2
 =\frac{|\mathfrak e_{X,D}(m)|^2}{n_D(m)}.}
 \tag{E.16}
\end{equation}
Every quotient-resolved realization pays this endpoint cost plus a nonnegative
within-fibre innovation.
\end{corollary}

\begin{proof}
The affine constraint fixes the constant projection.  The orthogonal fibre
contrast is arbitrary and has minimum norm zero.
\end{proof}

\begin{remark}[Per-quotient dispersion may charge a pressure-null sector]
The outer quotient square and the $d$--$d'$ Gram remain valid sufficient
coordinates.  They are not source minimal unless their within-endpoint fibre
innovation is either retained as a named stronger cost or shorted after the
complete endpoint fusion.  Large factorization dispersion can coexist with a
zero fused endpoint coefficient.
\end{remark}

% -----------------------------------------------------------------------------
\subsection{Carrier-birth calculus on the quotient bank}
\label{subsec:affend-quotient-birth}
% -----------------------------------------------------------------------------

Fix $X$.  For $2\le Q\le Q_X$, let an independently populated positive action
on the first quotient coordinates be denoted by $G_{X,Q}$.  After quotienting
its null space, assume it is positive definite.  Put
\[
 h_Q=(\log2,\ldots,\log Q)^{\mathsf T},
 \qquad
 \Lambda_{X,Q}=\langle h_Q,G_{X,Q}^{-1}h_Q\rangle.
\]
Suppose the fine action after adding $Q+1$ has block form
\[
 G_{X,Q+1}^{+}
 =\begin{pmatrix}A_Q&B_Q\\B_Q^*&D_Q\end{pmatrix}\succeq0,
 \qquad A_Q\succ0.
\]
Define
\begin{equation}
 \label{eq:affend-birth-data}
 \boxed{
 \begin{aligned}
 S_Q&=D_Q-B_Q^*A_Q^{-1}B_Q,\\
 \zeta_Q&=\log(Q+1)-B_Q^*A_Q^{-1}h_Q,\\
 \eta_Q^{\rm birth}
 &=\langle\zeta_Q,S_Q^{\dagger}\zeta_Q\rangle,\\
 \mathcal C_Q^{\rm ret}
 &=\Lambda_{X,Q}-\langle h_Q,A_Q^{-1}h_Q\rangle.
 \end{aligned}}
 \tag{E.17}
\end{equation}

\begin{theorem}[Exact quotient birth--retuning balance]
\label{thm:affend-birth-retuning}
If $\zeta_Q\in\Ran S_Q$, then
\begin{equation}
 \label{eq:affend-birth-retuning}
 \boxed{
 \Lambda_{X,Q+1}
 =\Lambda_{X,Q}
  -\mathcal C_Q^{\rm ret}
  +\eta_Q^{\rm birth}.}
 \tag{E.18}
\end{equation}
If $\zeta_Q\notin\Ran S_Q$, the sharp influence at the fine cutoff is infinite.
The birth term is nonnegative; only old-sector retuning can give a favorable
decrement.
\end{theorem}

\begin{proof}
Use the congruence factorization
\[
 \begin{pmatrix}A&B\\B^*&D\end{pmatrix}
 =\begin{pmatrix}I&0\\B^*A^{-1}&I\end{pmatrix}
  \begin{pmatrix}A&0\\0&S\end{pmatrix}
  \begin{pmatrix}I&A^{-1}B\\0&I\end{pmatrix}.
\]
The inverse left factor sends $h_Q\oplus\log(Q+1)$ to
$h_Q\oplus\zeta_Q$.  The Moore--Penrose value is therefore
$\langle h_Q,A_Q^{-1}h_Q\rangle+\eta_Q^{\rm birth}$.  Insert the definition of
$\mathcal C_Q^{\rm ret}$.
\end{proof}

\begin{corollary}[Pure quotient extension cannot lower influence]
\label{cor:affend-pure-extension}
If the old action is strictly nested,
\[
 \boxed{A_Q=G_{X,Q},}
\]
then
\begin{equation}
 \label{eq:affend-pure-extension}
 \boxed{
 \Lambda_{X,Q+1}
 =\Lambda_{X,Q}+\eta_Q^{\rm birth}
 \ge\Lambda_{X,Q}.}
 \tag{E.19}
\end{equation}
Equality holds exactly when the new logarithmic read is predicted by the old
minimum decoder on the faithful detail quotient.
\end{corollary}

\begin{proof}
The retuning term vanishes in \textup{(E.18)}.
\end{proof}

\begin{corollary}[Cumulative quotient ledger]
\label{cor:affend-quotient-telescope}
For any predeclared quotient chain $Q_0<Q_1<\cdots<Q_J=Q_X$ built by adjacent
coordinate additions,
\begin{equation}
 \label{eq:affend-quotient-telescope}
 \boxed{
 \Lambda_{X,Q_X}
 =\Lambda_{X,Q_0}
  -\sum_{Q=Q_0}^{Q_X-1}\mathcal C_Q^{\rm ret}
  +\sum_{Q=Q_0}^{Q_X-1}\eta_Q^{\rm birth}.}
 \tag{E.20}
\end{equation}
\end{corollary}

\begin{proof}
Telescope \textup{(E.18)}.
\end{proof}

\begin{proposition}[Scalar quotient rows do not determine birth cost]
\label{cth:affend-marginals-no-birth}
The positive Grams
\[
 G_0=\begin{pmatrix}1&0\\0&2\end{pmatrix},
 \qquad
 G_1=\begin{pmatrix}1&1\\1&2\end{pmatrix}
\]
have the same old and new marginal actions.  For old and detail reads both
equal to one, however,
\begin{equation}
 \label{eq:affend-marginals-no-birth}
 \boxed{
 \eta_0^{\rm birth}=\frac12,
 \qquad
 \eta_1^{\rm birth}=0.}
 \tag{E.21}
\end{equation}
Thus the quotient coefficients $\mathfrak r_X(q)$, their separate energies,
and the new-row norm do not determine the Christoffel payment.  The mixed
old--detail polarization is load bearing.
\end{proposition}

\begin{proof}
For $G_0$, the Schur value is $2$ and the residual read is one.  For $G_1$,
the Schur value is one but the residual read vanishes.
\end{proof}

% -----------------------------------------------------------------------------
\subsection{Diagonal and scalar quotient actions are pure-birth compilers}
\label{subsec:affend-diagonal-no-go}
% -----------------------------------------------------------------------------

\begin{proposition}[Exact diagonal influence ledger]
\label{prop:affend-diagonal-ledger}
Let
\[
 G_{X,Q}=\operatorname{diag}(w_2,\ldots,w_Q),
 \qquad w_q>0,
\]
and retain the old weights when $Q$ is enlarged.  Then
\begin{equation}
 \label{eq:affend-diagonal-influence}
 \boxed{
 \Lambda_{X,Q}
 =\sum_{q=2}^{Q}\frac{(\log q)^2}{w_q},}
 \tag{E.22}
\end{equation}
and every new quotient is a pure birth with cost
\begin{equation}
 \label{eq:affend-diagonal-birth}
 \boxed{
 \eta_Q^{\rm birth}
 =\frac{(\log(Q+1))^2}{w_{Q+1}}.}
 \tag{E.23}
\end{equation}
For the ordinary quotient $L^2$ action, $w_q=1$ and
\[
 \boxed{
 \Lambda_{X,Q}\sim Q(\log Q)^2.}
\]
\end{proposition}

\begin{proof}
The inverse action is diagonal.  The fine mixed block is zero, so the Schur
value is $w_{Q+1}$ and the residual read is $\log(Q+1)$.
\end{proof}

\begin{proposition}[Scalar retuning preserves the old action--influence product]
\label{prop:affend-scalar-retuning}
If $A_Q=\lambda_QG_{X,Q}$ with $\lambda_Q>0$, then for every old quotient
source vector $x$,
\begin{equation}
 \label{eq:affend-scalar-retuning}
 \boxed{
 \langle x,A_Qx\rangle
 \langle h_Q,A_Q^{-1}h_Q\rangle
 =\langle x,G_{X,Q}x\rangle
  \langle h_Q,G_{X,Q}^{-1}h_Q\rangle.}
 \tag{E.24}
\end{equation}
Useful retuning must therefore be selective: it must be heavy on the endpoint
decoder while remaining light on the protected arithmetic source.
\end{proposition}

\begin{proof}
The two factors scale by $\lambda_Q$ and $\lambda_Q^{-1}$.
\end{proof}

\begin{remark}[Effect on the outer compiler]
The quotient $L^2$ dispersion of \textup{(A.15)} is a valid stronger certificate, but
its action is a pure orthogonal-birth action.  Adding quotient coordinates or
rescaling all old rows cannot manufacture the needed cancellation.  A
successful non-diagonal action must come from a predeclared same-history
arithmetic panel and must expose the mixed old--detail blocks entering
\textup{(E.17)}.
\end{remark}

% -----------------------------------------------------------------------------
\subsection{Hard-threshold motion is signed endpoint retuning}
\label{subsec:affend-threshold-retuning}
% -----------------------------------------------------------------------------

For $1\le D_1<D_2$, define the completed threshold-shell current
\begin{equation}
 \label{eq:affend-threshold-shell}
 \boxed{
 \begin{aligned}
 \mathcal C_X^{\rm cut}(D_1,D_2)
 :=\sum_{m\in\mathscr I_X^+}
 \sum_{\substack{d\mid m\\D_1<d\le D_2}}
 \mu(d)\log\frac md
 \bigl[\Lambda(m-2)-\gamma_2(d)\bigr].
 \end{aligned}}
 \tag{E.25}
\end{equation}

\begin{theorem}[Exact endpoint retuning identity]
\label{thm:affend-threshold-retuning}
For every $D_1<D_2$,
\begin{equation}
 \label{eq:affend-threshold-retuning}
 \boxed{
 \mathfrak T_{X,D_1}^{\rm end}
 -\mathfrak T_{X,D_2}^{\rm end}
 =\mathcal C_X^{\rm cut}(D_1,D_2).}
 \tag{E.26}
\end{equation}
If
\[
 D_X=D_0<D_1<\cdots<D_J\ge2X+2
\]
is any source-independent threshold partition, then
\begin{equation}
 \label{eq:affend-threshold-telescope}
 \boxed{
 \mathfrak T_X^{\rm AMP}
 =\sum_{j=0}^{J-1}
   \mathcal C_X^{\rm cut}(D_j,D_{j+1}).}
 \tag{E.27}
\end{equation}
\end{theorem}

\begin{proof}
Subtract the two endpoint divisor sums.  The surviving divisors are exactly
those in the shell $(D_1,D_2]$.  At $D_J\ge2X+2$ every tail sum is empty, so
\textup{(E.27)} follows from \textup{(E.5)}.
\end{proof}

\begin{corollary}[Retuning, not endpoint birth]
\label{cor:affend-retuning-not-birth}
At fixed $X$, changing $D$ alters the coefficient underlying to each already
visible endpoint $m$ but introduces no new endpoint coordinate.  The shell
current \textup{(E.25)} is therefore a signed old-sector retuning.  Any
positive quotient-carrier birth payment appearing before endpoint fusion is a
factorization-provenance cost, not a favorable endpoint decrement.
\end{corollary}

\begin{proof}
The endpoint index set $\mathscr I_X^+$ is independent of $D$.  Formula
\textup{(E.26)} changes only the endpoint writer.  The quotient-fibre
Pythagoras identifies every additional factorization contrast as a
nonnegative read-null innovation.
\end{proof}

\begin{remark}[Threshold shells must be reassembled]
The exact telescope does not authorize separate absolute estimates of the
shells.  Two nonzero shell currents may cancel in \textup{(E.27)}.  A
shellwise analytic theorem is sufficient only when its signed estimates sum
to $o(X)$ or when the complete cross-shell Gram is retained.
\end{remark}

% -----------------------------------------------------------------------------
\subsection{Birth-corrected action criterion}
\label{subsec:affend-action-criterion}
% -----------------------------------------------------------------------------

Let
\[
 r_{X,Q_X}
 =\bigl(\mathfrak r_X(2),\ldots,\mathfrak r_X(Q_X)\bigr)^{\mathsf T}.
\]
Suppose an independently populated positive quotient action $G_{X,Q_X}$ has
been built through the adjacent chain above and that all range conditions
hold.

\begin{theorem}[Birth-corrected quotient-action closure]
\label{thm:affend-birth-corrected-closure}
Put
\[
 \mathfrak A_X^{\rm quot}
 =\langle r_{X,Q_X},G_{X,Q_X}r_{X,Q_X}\rangle.
\]
If
\begin{equation}
 \label{eq:affend-birth-corrected-criterion}
 \boxed{
 \mathfrak A_X^{\rm quot}
 \left[
  \Lambda_{X,Q_0}
  -\sum_{Q=Q_0}^{Q_X-1}\mathcal C_Q^{\rm ret}
  +\sum_{Q=Q_0}^{Q_X-1}\eta_Q^{\rm birth}
 \right]
 =o(X^2),}
 \tag{E.28}
\end{equation}
then
\[
 R_X(2)=o(X).
\]
\end{theorem}

\begin{proof}
The bracket in \textup{(E.28)} is the exact final protected influence by
\textup{(E.20)}.  The protected-observable Riesz inequality gives
\[
 |\langle h_{Q_X},r_{X,Q_X}\rangle|^2
 \le
 \mathfrak A_X^{\rm quot}\Lambda_{X,Q_X}.
\]
The left side is $|\mathfrak T_X^{\rm AMP}|^2$.  Apply \textup{(A.11)}.
\end{proof}

\begin{corollary}[Endpoint-fused action closure]
\label{cor:affend-endpoint-action-closure}
Let $z_{X,D_X}$ be the endpoint vector with entries
$\mathfrak e_{X,D_X}(m)$ and let $\one_X$ be the constant endpoint Read.  If an
independently populated positive endpoint action $L_X$ satisfies
\begin{equation}
 \label{eq:affend-endpoint-action-closure}
 \boxed{
 \langle z_{X,D_X},L_Xz_{X,D_X}\rangle
 \langle\one_X,L_X^{\dagger}\one_X\rangle
 =o(X^2),}
 \tag{E.29}
\end{equation}
then $R_X(2)=o(X)$.
\end{corollary}

\begin{proof}
Apply the Riesz inequality to
$\mathfrak T_{X,D_X}^{\rm end}=\langle\one_X,z_{X,D_X}\rangle$ and use
\textup{(E.6)}.
\end{proof}

\begin{remark}[Consequences of complete birth accounting]
A quotient action is not certified by its diagonal rows.  It must carry the
mixed old--detail polarization needed to compute every Christoffel payment.
A pure nested action only increases influence, and a scalar retuning merely
moves cost between energy and influence.  The first potentially useful action
is therefore either:
\begin{enumerate}[label=\textup{(\roman*)},leftmargin=2.8em]
\item an endpoint-fused arithmetic action satisfying \textup{(E.29)}; or
\item a selectively retuned quotient action satisfying the complete
      birth-corrected ledger \textup{(E.28)}.
\end{enumerate}
\end{remark}


% =============================================================================
\section{Complement pressure and the odd--odd endpoint current}
\label{sec:affine-pressure}
% =============================================================================

\subsection{Base-two complement pressure on the fused endpoint carrier}
\label{subsec:affpress-pressure-definition}
% -----------------------------------------------------------------------------

For $q\ge2$ put
\[
 \alpha(q):=\log_2q=\frac{\log q}{\log2}\ge1.
\]
For $0\le t\le1$, a real threshold $D\ge1$, and $m\ge2$, define
\begin{equation}
 \label{eq:affpress-pressure-writers}
 \boxed{
 \begin{aligned}
 \Theta_{D,t}(m)
 &:=\sum_{\substack{d\mid m\\d>D}}
     \mu(d)\log\frac md\,
     t^{\alpha(m/d)},\\
 \Gamma_{D,t}(m)
 &:=\sum_{\substack{d\mid m\\d>D}}
     \mu(d)\gamma_2(d)\log\frac md\,
     t^{\alpha(m/d)}.
 \end{aligned}}
 \tag{P.1}
\end{equation}
The term with $m/d=1$ is zero because its logarithmic coefficient vanishes.
At $t=0$ every remaining term is zero.  Put
\begin{equation}
 \label{eq:affpress-pressure-row}
 \boxed{
 \mathfrak T_{X,D}(t)
 :=\sum_{m\in\mathscr I_X^+}
 \left[
  \Theta_{D,t}(m)\Lambda(m-2)-\Gamma_{D,t}(m)
 \right].}
 \tag{P.2}
\end{equation}
At the square-root threshold write
\[
 \mathfrak T_X(t):=\mathfrak T_{X,D_X}(t).
\]

\begin{theorem}[Exact complement-pressure compiler]
\label{thm:affpress-pressure-compiler}
For every $0\le t\le1$,
\begin{equation}
 \label{eq:affpress-pressure-quotient}
 \boxed{
 \mathfrak T_X(t)
 =\sum_{q=2}^{Q_X}
   (\log q)t^{\alpha(q)}\mathfrak r_X(q),}
 \tag{P.3}
\end{equation}
where $\mathfrak r_X(q)$ is the completed outer quotient row.  The path obeys
\begin{equation}
 \label{eq:affpress-pressure-boundary}
 \boxed{
 \mathfrak T_X(0)=0,
 \qquad
 \mathfrak T_X(1)=\mathfrak T_X^{\rm AMP}.}
 \tag{P.4}
\end{equation}
It is continuously differentiable on $[0,1]$, with
\begin{equation}
 \label{eq:affpress-pressure-current}
 \boxed{
 \begin{aligned}
 \mathfrak J_X(t)
 :=\frac{\dd}{\dd t}\mathfrak T_X(t)
 =\frac1{\log2}
 \sum_{q=2}^{Q_X}
 (\log q)^2t^{\alpha(q)-1}\mathfrak r_X(q).
 \end{aligned}}
 \tag{P.5}
\end{equation}
Consequently,
\begin{equation}
 \label{eq:affpress-pressure-integral}
 \boxed{
 R_X(2)
 =\int_0^1\mathfrak J_X(t)\,\dd t+o(X).}
 \tag{P.6}
\end{equation}
\end{theorem}

\begin{proof}
In \textup{(P.2)} write $m=dq$ and interchange the finite sums.  This gives
\textup{(P.3)}.  Since $q\ge2$, every exponent $\alpha(q)$ is at least one;
hence each summand is $C^1$ at zero.  Differentiation gives
\textup{(P.5)}.  Integrating
\[
 \frac{(\log q)^2}{\log2}t^{\alpha(q)-1}
\]
from zero to one gives $\log q$.  Thus
$\int_0^1\mathfrak J_X=\mathfrak T_X^{\rm AMP}$, and \textup{(A.11)} proves
\textup{(P.6)}.
\end{proof}

\begin{corollary}[Exact pressure entrance]
\label{cor:affpress-pressure-entrance}
At the zero-pressure endpoint,
\begin{equation}
 \label{eq:affpress-pressure-entrance}
 \boxed{
 \mathfrak J_X(0)
 =(\log2)\mathfrak r_X(2).}
 \tag{P.7}
\end{equation}
Only the complementary quotient $q=2$ is visible to the first pressure
current.
\end{corollary}

\begin{proof}
For $q=2$, $\alpha(q)-1=0$.  For $q>2$, the exponent is positive and the
corresponding term vanishes at $t=0$.
\end{proof}

\begin{theorem}[Prime-nullity and radical support along the complete pressure path]
\label{thm:affpress-pressure-prime-null}
For every $0\le t\le1$ and every prime $p$,
\begin{equation}
 \label{eq:affpress-pressure-prime-null}
 \boxed{
 \Theta_{D,t}(p)=\Gamma_{D,t}(p)=0.}
 \tag{P.8}
\end{equation}
More generally,
\begin{equation}
 \label{eq:affpress-pressure-radical-null}
 \boxed{
 \rad(m)\le D
 \quad\Longrightarrow\quad
 \Theta_{D,t}(m)=\Gamma_{D,t}(m)=0.}
 \tag{P.9}
\end{equation}
Thus every nonzero actual endpoint remains composite throughout the pressure
path.
\end{theorem}

\begin{proof}
For a prime endpoint the only possible tail divisor with nonzero M\"obius
weight is $d=p$, whose complementary quotient is one and whose logarithmic
coefficient is zero.  If $\rad(m)\le D$, every divisor with nonzero M\"obius
coefficient is at most $D$.
\end{proof}

\begin{theorem}[Pressure-deformed radical complementary-divisor formula]
\label{thm:affpress-radical-pressure}
Write
\[
 r=\rad(m),
 \qquad
 c=\frac mr.
\]
Define
\[
 \begin{aligned}
 M_{r,D}(t)
 &:=\sum_{\substack{e\mid r\\e<r/D}}
     \mu(e)t^{\alpha(e)},
 &
 L_{r,D}(t)
 &:=\sum_{\substack{e\mid r\\e<r/D}}
     \mu(e)(\log e)t^{\alpha(e)},\\
 M_{r,D}^{\gamma}(t)
 &:=\sum_{\substack{e\mid r\\e<r/D}}
     \mu(e)\gamma_2(r/e)t^{\alpha(e)},
 &
 L_{r,D}^{\gamma}(t)
 &:=\sum_{\substack{e\mid r\\e<r/D}}
     \mu(e)\gamma_2(r/e)(\log e)t^{\alpha(e)}.
 \end{aligned}
\]
Then
\begin{equation}
 \label{eq:affpress-radical-pressure-actual}
 \boxed{
 \Theta_{D,t}(m)
 =\mu(r)t^{\alpha(c)}
 \left[
  (\log c)M_{r,D}(t)+L_{r,D}(t)
 \right],}
 \tag{P.10}
\end{equation}
while
\begin{equation}
 \label{eq:affpress-radical-pressure-local}
 \boxed{
 \Gamma_{D,t}(m)
 =\mu(r)t^{\alpha(c)}
 \left[
  (\log c)M_{r,D}^{\gamma}(t)
  +L_{r,D}^{\gamma}(t)
 \right].}
 \tag{P.11}
\end{equation}
The complete factor-depth hierarchy therefore remains fused inside one
finite radical boundary polynomial for every~$t$.
\end{theorem}

\begin{proof}
Every divisor with nonzero M\"obius coefficient is a divisor of $r$.  Put
$d=r/e$.  Then $d>D$ is equivalent to $e<r/D$,
$\mu(d)=\mu(r)\mu(e)$, and $m/d=ce$.  Finally
$t^{\alpha(ce)}=t^{\alpha(c)}t^{\alpha(e)}$.  Substitution gives both
formulas.
\end{proof}

% -----------------------------------------------------------------------------
\subsection{Complete positive source panel and the zero-birth theorem}
\label{subsec:affpress-zero-birth}
% -----------------------------------------------------------------------------

Let
\[
 \Omega_X
 =\left\{(d,q):d>D_X,\ q\ge2,\ dq\in\mathscr I_X^+,\ \mu(d)^2=1\right\}.
\]
Use two copies of $\ell^2(\Omega_X)$, one for the actual row and one for the
local row.  For $0\le t\le1$ define diagonal source maps
\begin{equation}
 \label{eq:affpress-positive-sources}
 \boxed{
 \begin{aligned}
 [\mathscr A_t x]_{d,q}
 &:=t^{\alpha(q)/2}
    \sqrt{(\log q)\Lambda(dq-2)}\,x_{d,q},\\
 [\mathscr L_t y]_{d,q}
 &:=t^{\alpha(q)/2}
    \sqrt{(\log q)\gamma_2(d)}\,y_{d,q}.
 \end{aligned}}
 \tag{P.12}
\end{equation}
Put
\[
 \mathscr S_t=\mathscr A_t\oplus\mathscr L_t,
 \qquad
 \mathsf J_X=M_{\mu}\oplus(-M_{\mu}),
\]
where $M_{\mu}$ multiplies the $(d,q)$ coordinate by $\mu(d)$.  If
$\mathbf1_X$ denotes the constant coefficient vector on both copies, then
\begin{equation}
 \label{eq:affpress-positive-signed-read}
 \boxed{
 \mathfrak T_X(t)
 =\left\langle
   \mathscr S_t\mathbf1_X,
   \mathsf J_X\mathscr S_t\mathbf1_X
  \right\rangle.}
 \tag{P.13}
\end{equation}
The complete untwisted Gram of the two positive source families is positive;
$\mathsf J_X$ retains the M\"obius and actual-minus-local orientation as a
signed mixed Read.

\begin{theorem}[Finite pressure transport has zero carrier-birth innovation]
\label{thm:affpress-zero-birth}
For $0<s,t\le1$, let $M_{t/s}$ be the positive diagonal multiplier
\[
 [M_{t/s}f]_{d,q}
 =\left(\frac ts\right)^{\alpha(q)/2}f_{d,q}
\]
on both source copies.  Then
\begin{equation}
 \label{eq:affpress-source-transport}
 \boxed{
 \mathscr S_t=M_{t/s}\mathscr S_s.}
 \tag{P.14}
\end{equation}
The multiplier is invertible, and
\begin{equation}
 \label{eq:affpress-source-range}
 \boxed{
 \Ran\mathscr S_t=\Ran\mathscr S_s.}
 \tag{P.15}
\end{equation}
Consequently both mutual Schur innovations vanish.  The source derivative is
\begin{equation}
 \label{eq:affpress-source-derivative}
 \boxed{
 \dot{\mathscr S}_t
 =\mathsf K_t\mathscr S_t,
 \qquad
 [\mathsf K_tf]_{d,q}
 =\frac{\alpha(q)}{2t}f_{d,q},}
 \tag{P.16}
\end{equation}
and hence its range-leakage Gram is zero:
\begin{equation}
 \label{eq:affpress-zero-leakage}
 \boxed{
 \mathscr S_t^*\mathsf K_t
 (I-P_{\Ran\mathscr S_t})
 \mathsf K_t\mathscr S_t=0.}
 \tag{P.17}
\end{equation}
Thus the whole finite pressure path is source retuning and contains no genuine
carrier birth.
\end{theorem}

\begin{proof}
The first identity follows directly from \textup{(P.12)}.  On the finite
coordinate support the multiplier has strictly positive diagonal entries and
is invertible.  Each source range is the same coordinate subspace determined
by the nonzero actual or local amplitudes, proving \textup{(P.15)}.  The
mutual Schur residuals are therefore zero.  Differentiating the diagonal
multiplier gives \textup{(P.16)}, whose range lies in the same coordinate
subspace.  Orthogonal projection onto its complement gives
\textup{(P.17)}.
\end{proof}

\begin{corollary}[Exact birth--retuning balance on the pressure path]
\label{cor:affpress-retuning-only}
On every finite interval $0<s\le t\le1$, the carrier-birth term of the
complement-pressure transport is
\[
 \boxed{\eta_{s\to t}^{\rm birth}=0.}
\]
Every change of a protected influence along this path is therefore an
old-sector retuning effect.  Any positive birth payment obtained by revealing
quotient or divisor labels sequentially belongs to that finer provenance
filtration, not to the source-minimal pressure deformation.
\end{corollary}

\begin{proof}
The new source range is contained in, and in fact equals, the old range by
\textup{(P.15)}.  The detail residual and its Christoffel cost therefore
vanish.
\end{proof}

\begin{proposition}[The zero-pressure endpoint is not a positive action anchor]
\label{cth:affpress-zero-anchor}
At $t=0$,
\[
 \boxed{\mathscr S_0=0.}
\]
Hence every nonzero endpoint or quotient Read lies outside the source range
and has infinite Moore--Penrose influence.  The scalar identity
\textup{(P.6)} may be integrated from zero, but a positive action or kernel
crossing argument must begin at a strictly positive pressure.
\end{proposition}

\begin{proof}
Every complement exponent is positive, so all source amplitudes vanish at
zero.  The range of the zero Gram is the zero space.
\end{proof}

\begin{remark}[No regular first crossing is created]
For every $t>0$ the positive source support is constant.  The only singular
endpoint is $t=0$, where the entire source collapses simultaneously.  Thus the
GRH first-kernel-crossing mechanism does not return a new affine arithmetic
row here.  The nontrivial datum is exactly the signed retuning derivative
$\mathfrak J_X(t)$ already computed in \textup{(P.5)}.
\end{remark}

% -----------------------------------------------------------------------------
\subsection{Weighted Mertens control and parity cleanup}
\label{subsec:affpress-parity-cleanup}
% -----------------------------------------------------------------------------

\begin{lemma}[Weighted odd Mertens estimate]
\label{lem:affpress-weighted-Mertens}
There is a constant $c>0$ such that
\begin{equation}
 \label{eq:affpress-weighted-Mertens}
 \boxed{
 \sum_{n\le y}\mu(n)\gamma_2(n)
 \ll ye^{-c\sqrt{\log y}}
 \qquad(y\ge3).}
 \tag{P.18}
\end{equation}
The same bound holds for every interval by subtraction.
\end{lemma}

\begin{proof}
For $\Re z>1$,
\[
 \sum_{n\ge1}\frac{\mu(n)\gamma_2(n)}{n^z}
 =\prod_{p>2}
  \left(1-\frac{p^{1-z}}{p-1}\right)
 =\frac{\mathscr H_2(z)}{\zeta(z)},
\]
where $\mathscr H_2$ is the analytic Euler factor of \textup{(A.6)}.  It is bounded in
the classical zero-free region for $\zeta$.  Perron inversion and the standard
zero-free contour shift used in the local singular-series analysis and the outer-compiler estimates give \textup{(P.18)}.
\end{proof}

Split the pressure packet according to the parity of $q$ and $d$.  Let
$\mathfrak T_X^{\rm evq}(t)$ and $\mathfrak J_X^{\rm evq}(t)$ denote the
subpackets with $2\mid q$.  Let
$\mathfrak T_X^{\rm ode}(t)$ and $\mathfrak J_X^{\rm ode}(t)$ denote the
subpackets with $q$ odd and $d$ even.  The remaining packet is denoted by
$\mathfrak T_X^{\rm oo}(t)$ and $\mathfrak J_X^{\rm oo}(t)$.

\begin{theorem}[Uniform closure of the non-odd--odd sectors]
\label{thm:affpress-parity-closure}
Uniformly for $0\le t\le1$,
\begin{equation}
 \label{eq:affpress-even-closure}
 \boxed{
 \begin{aligned}
 |\mathfrak T_X^{\rm evq}(t)|
 +|\mathfrak J_X^{\rm evq}(t)|
 +|\mathfrak T_X^{\rm ode}(t)|
 +|\mathfrak J_X^{\rm ode}(t)|
 =o(X).
 \end{aligned}}
 \tag{P.19}
\end{equation}
Consequently,
\begin{equation}
 \label{eq:affpress-odd-odd-reduction}
 \boxed{
 \mathfrak T_X(t)=\mathfrak T_X^{\rm oo}(t)+o(X),
 \qquad
 \mathfrak J_X(t)=\mathfrak J_X^{\rm oo}(t)+o(X),}
 \tag{P.20}
\end{equation}
uniformly on the full pressure interval.
\end{theorem}

\begin{proof}
If $q$ is even, then $dq-2$ is even.  A nonzero actual von Mangoldt row must
therefore satisfy
\[
 dq-2=2^a.
\]
There are $O(\log X)$ such endpoints, and each has at most $X^{o(1)}$
divisor factorizations.  The logarithmic weights in both
$\mathfrak T$ and $\mathfrak J$ are polynomial in $\log X$, so the complete
actual even-$q$ contribution is $X^{o(1)}=o(X)$ uniformly in~$t$.

For the local even-$q$ contribution use
\textup{(P.18)} on the interval
\[
 d>D_X,
 \qquad
 dq\in\mathscr I_X^+.
\]
Its endpoints are at least $D_X$ whenever the interval is nonempty.  Since
$t^{\alpha(q)}\le1$ and $t^{\alpha(q)-1}\le1$,
\[
 \begin{aligned}
 &\sum_{\substack{q\le Q_X\\2\mid q}}
  (\log q)^2
  \left|
   \sum_{\substack{d>D_X\\dq\in\mathscr I_X^+}}
   \mu(d)\gamma_2(d)
  \right|\\
 &\qquad\ll
 Xe^{-c'\sqrt{\log X}}(\log X)^4
 =o(X).
 \end{aligned}
\]
The pressure row has one fewer logarithm and obeys the same conclusion.

If $q$ is odd and $d$ is even, the local coefficient is zero by definition of
$\gamma_2$.  The actual endpoint $dq-2$ is again even and is handled by the
same power-of-two argument.  This proves \textup{(P.19)} and hence
\textup{(P.20)}.
\end{proof}

\begin{corollary}[Literal odd--odd current]
\label{cor:affpress-odd-odd-current}
Uniformly for $0\le t\le1$,
\begin{equation}
 \label{eq:affpress-odd-odd-current}
 \boxed{
 \begin{aligned}
 \mathfrak J_X(t)
 ={}&\frac1{\log2}
 \sum_{\substack{3\le q\le Q_X\\q\ {\rm odd}}}
 (\log q)^2t^{\alpha(q)-1}
 \\[-0.2em]
 &\times
 \sum_{\substack{d>D_X\\d\ {\rm odd}\\dq\in\mathscr I_X^+}}
 \mu(d)
 \left[
  \Lambda(dq-2)-\frac d{\varphi(d)}
 \right]
 +o(X).
 \end{aligned}}
 \tag{P.21}
\end{equation}
After removing proper prime powers at the right endpoint, every actual
occurrence in the displayed row has
\[
 \boxed{
 d,q\text{ odd},
 \quad d>D_X>1,
 \quad q\ge3,
 \quad dq\text{ composite},
 \quad dq-2\text{ prime}.}
\]
The cleanup is uniform in~$t$.
\end{corollary}

\begin{proof}
The formula is \textup{(P.5)} after
\Cref{thm:affpress-parity-closure}.  Proper prime powers contribute at most
\[
 X^{1/2+o(1)}(\log X)^3=o(X)
\]
uniformly because the pressure weights are at most their endpoint logarithmic
weights.  The support statements then follow from $d>D_X$ and $q\ge3$.
\end{proof}

% -----------------------------------------------------------------------------
\subsection{Closed pressure entrance and a constant-scale odd--odd witness}
\label{subsec:affpress-pressure-obstruction}
% -----------------------------------------------------------------------------

\begin{lemma}[Crude complete pressure mass]
\label{lem:affpress-pressure-mass}
One has
\begin{equation}
 \label{eq:affpress-pressure-mass}
 \boxed{
 \sum_{\substack{d>D_X,\ q\ge2\\dq\in\mathscr I_X^+}}
 (\log q)^2
 \left[
  \Lambda(dq-2)+\gamma_2(d)
 \right]
 \ll X(\log X)^4.}
 \tag{P.22}
\end{equation}
\end{lemma}

\begin{proof}
Group by the endpoint $m=dq$.  Each endpoint has at most $\tau(m)$ divisor
histories.  Use
\[
 \Lambda(m-2)\le\log(2X+2),
 \qquad
 \gamma_2(d)\ll\log\log(3X),
\]
$\log q\ll\log X$, and the classical divisor-sum bound
\[
 \sum_{m\le2X+2}\tau(m)\ll X\log X.
\]
The displayed logarithmic power is a harmless uniform envelope.
\end{proof}

Put
\[
 \delta_0:=\log_23-1>0,
 \qquad
 t_{0,X}:=(\log X)^{-20}.
\]

\begin{theorem}[Unconditional pressure-entrance closure]
\label{thm:affpress-entrance-closure}
Uniformly for $0\le t\le t_{0,X}$,
\begin{equation}
 \label{eq:affpress-entrance-current}
 \boxed{
 \mathfrak J_X(t)=o(X).}
 \tag{P.23}
\end{equation}
Moreover,
\begin{equation}
 \label{eq:affpress-entrance-pressure}
 \boxed{
 \mathfrak T_X(t)=o(X).}
 \tag{P.24}
\end{equation}
In particular, the first pressure row $q=2$ is unconditionally closed.
\end{theorem}

\begin{proof}
The $q=2$ row belongs to the even-quotient sector and is $o(X)$ uniformly by
\Cref{thm:affpress-parity-closure}.  For $q\ge3$ and $t\le t_{0,X}$,
\[
 t^{\alpha(q)-1}
 \le t_{0,X}^{\delta_0}.
\]
Apply \Cref{lem:affpress-pressure-mass} to obtain
\[
 |\mathfrak J_X^{q\ge3}(t)|
 \ll X(\log X)^4t_{0,X}^{\delta_0}
 =o(X).
\]
This proves \textup{(P.23)}.  For the pressure itself,
$t^{\alpha(q)}\le t_{0,X}$ for every $q\ge2$; the analogous bound with one
fewer logarithm proves \textup{(P.24)}.
\end{proof}

\begin{theorem}[Constant-scale pressure obstruction]
\label{thm:affpress-pressure-obstruction}
Assume that for some $\eta>0$,
\[
 |R_X(2)|\ge\eta X
\]
along an unbounded sequence.  Then for all sufficiently large $X$ on that
sequence there exists a source-independent pressure parameter
\[
 t_X^*\in[t_{0,X},1]
\]
such that
\begin{equation}
 \label{eq:affpress-pressure-obstruction}
 \boxed{
 \left|
 \frac1{\log2}
 \sum_{\substack{3\le q\le Q_X\\q\ {\rm odd}}}
 (\log q)^2(t_X^*)^{\alpha(q)-1}
 \sum_{\substack{d>D_X\\d\ {\rm odd}\\dq\in\mathscr I_X^+}}
 \mu(d)
 \left[
  \Lambda(dq-2)-\frac d{\varphi(d)}
 \right]
 \right|
 \ge\frac\eta4X.}
 \tag{P.25}
\end{equation}
The actual row may be restricted to prime values of $dq-2$ at the cost of an
$o(X)$ error.
\end{theorem}

\begin{proof}
By \textup{(P.6)} and \Cref{thm:affpress-entrance-closure},
\[
 R_X(2)
 =\int_{t_{0,X}}^1\mathfrak J_X(t)\,\dd t+o(X).
\]
Use the uniform odd--odd reduction and proper-prime-power cleanup from
\Cref{cor:affpress-odd-odd-current}.  For sufficiently large $X$ the absolute value
of the remaining integral is at least $\eta X/2$.  Since the interval has
length at most one, one parameter has absolute current at least $\eta X/2$;
absorbing the uniform $o(X)$ remainder gives the displayed conservative
constant $\eta/4$.
\end{proof}

\begin{remark}[Role of complement pressure]
The pressure family gives a finite, source-independent route from the closed
$q=2$ entrance to the complete outer M\"obius--log packet.  It does not prove cancellation,
but it removes three artificial losses at once:
\begin{enumerate}[label=\textup{(\roman*)},leftmargin=2.8em]
\item no quotient interval is selected;
\item no carrier-birth payment occurs for $t>0$;
\item a macroscopic residual returns a constant-scale odd--odd current rather
      than a current divided by $1+\log\log X$.
\end{enumerate}
\end{remark}

% -----------------------------------------------------------------------------
\subsection{Analytic certificate and maximal unconditional range}
\label{subsec:affpress-certificate}
% -----------------------------------------------------------------------------

\begin{definition}[Complement-pressure M\"obius--prime certificate]
\label{def:affpress-CPMP-certificate}
Let $\mathfrak J_X^{\rm oo,pr}(t)$ denote the prime-cleaned odd--odd row in
\textup{(P.25)}.  The regulator-uniform certificate is
\begin{equation}
 \label{eq:affpress-CPMP-certificate}
 \boxed{
 \mathsf{CPMP}_X:\qquad
 \sup_{t_{0,X}\le t\le1}
 \left|\mathfrak J_X^{\rm oo,pr}(t)\right|
 =o(X).}
 \tag{P.26}
\end{equation}
\end{definition}

\begin{theorem}[Complement-pressure closure]
\label{thm:affpress-CPMP-closure}
If \textup{(P.26)} holds, then
\[
 \boxed{R_X(2)=o(X).}
\]
The same conclusion follows from the weaker signed integral estimate
\begin{equation}
 \label{eq:affpress-integrated-certificate}
 \boxed{
 \int_{t_{0,X}}^1
 \mathfrak J_X^{\rm oo,pr}(t)\,\dd t=o(X).}
 \tag{P.27}
\end{equation}
\end{theorem}

\begin{proof}
Combine the exact integral compiler, the entrance closure, the parity
reduction, and the prime-power cleanup.
\end{proof}


% =============================================================================

% =============================================================================
\section{Reflected prime readers and the rigid-modulus frontier}
\label{sec:affine-hard-modulus}
% =============================================================================

The complement-pressure theorem leaves one global signed current.  A
different sharpening is obtained by returning to the complete prime reader
before opening a second M\"obius--log factor.  This change of order allows a
published fixed-residue distribution theorem to remove a large,
source-declared modulus sector absolutely.  The remainder has a canonical
divisor gap and a finite-depth rough cofactor.

\subsection{The reflected outer compiler}

Let
\[
 \mathscr I_X^-:=\{n\in\Z:X<n\le2X\}.
\]
For $a\ge1$, define
\[
 \begin{aligned}
 \mathcal A_X^+(a)
 &:=\sum_{\substack{b\ge1\\ab\in\mathscr I_X^-}}
       (\log b)\Lambda(ab+2),\\
 \mathcal L_X^+(a)
 &:=\gamma_2(a)
     \sum_{\substack{b\ge1\\ab\in\mathscr I_X^-}}\log b,\\
 \mathcal E_X^+(a)
 &:=\mathcal A_X^+(a)-\mathcal L_X^+(a),
 \end{aligned}
 \tag{M.1}
\]
where $\gamma_2(a)=a/\varphi(a)$ for odd $a$ and zero for even $a$.

\begin{theorem}[Exact reflected outer assembly]
\label{thm:hard-reflected-assembly}
One has
\[
 \boxed{
 \sum_{X<n\le2X}\Lambda(n)\Lambda(n+2)
 =\sum_{a\le2X}\mu(a)\mathcal A_X^+(a).}
 \tag{M.2}
\]
Moreover,
\[
 \boxed{
 \sum_{a\le2X}\mu(a)\mathcal L_X^+(a)
 =\mathfrak S(2)X+o(X),}
 \tag{M.3}
\]
and, with the square-root cutoff $D_X$ of \Cref{sec:affine-entrance},
\[
 \boxed{
 R_X(2)=\sum_{a>D_X}\mu(a)\mathcal E_X^+(a)+o(X).}
 \tag{M.4}
\]
Thus the twin residual has two source-native orientations, obtained by
expanding respectively the right or left von Mangoldt factor:
\[
 \begin{aligned}
 R_X(2)
 &=\sum_{d>D_X}\mu(d)
   \sum_{q\ge2}(\log q)
   [\Lambda(dq-2)-\gamma_2(d)]+o(X),\\
 &=\sum_{a>D_X}\mu(a)
   \sum_{b\ge2}(\log b)
   [\Lambda(ab+2)-\gamma_2(a)]+o(X),
 \end{aligned}
 \tag{M.5}
\]
with the corresponding endpoint interval restrictions.
\end{theorem}

\begin{proof}
Insert
$\Lambda(n)=\sum_{ab=n}\mu(a)\log b$
into the left endpoint of the pair and interchange finite sums to obtain
\textup{(M.2)}.  If
\[
 F_2(n)=\sum_{ab=n}\mu(a)\gamma_2(a)\log b,
\]
then
\[
 \sum_{n\ge1}\frac{F_2(n)}{n^s}
 =-\frac{\zeta'(s)}{\zeta(s)}\mathscr H_2(s),
\]
with the same analytic Euler factor $\mathscr H_2$ as in
\Cref{thm:affouter-outer-assembly} and
$\mathscr H_2(1)=\mathfrak S(2)$.  The zero-free-region input gives
\textup{(M.3)} on $(X,2X]$.  For $a\le D_X$, maximal
Bombieri--Vinogradov and partial summation close the reduced progression
$2\pmod a$.  Even $a$ can contribute to the actual row only when $ab+2$ is
a power of two; their total is $X^{o(1)}$.  Subtract the local assembly and
the small-divisor head to obtain \textup{(M.4)}.  The first line of
\textup{(M.5)} is \Cref{thm:affouter-tail-reduction}; the second is
\textup{(M.4)}.
\end{proof}

The two orientations are equal only after complete source assembly.  Their
individual divisor or quotient slices are not canonically paired.

\subsection{Reader convolution is invisible before dispersion}

For a modulus $m$ and a reduced residue $h$, put
\[
 \mathscr B_m(x;h)
 =\one_{x\equiv h\pmod m}
  -\frac{\one_{(x,m)=1}}{\varphi(m)}.
 \tag{M.6}
\]
For finitely supported sequences $\alpha,\beta$ and modulus coefficients
$\Gamma$, define
\[
 \mathfrak D_h[\alpha,\beta;\Gamma]
 =\sum_m\Gamma(m)
   \sum_u\alpha_u\sum_v\beta_v\mathscr B_m(uv;h).
 \tag{M.7}
\]

\begin{lemma}[Exact absorption of a reader convolution]
\label{lem:hard-convolution-absorption}
For finitely supported $a,b,\beta$,
\[
 \boxed{
 \mathfrak D_h[a*b,\beta;\Gamma]
 =\mathfrak D_h[b,a*\beta;\Gamma].}
 \tag{M.8}
\]
The identity remains exact when independent smooth dyadic weights are
incorporated into the three sequences.
\end{lemma}

\begin{proof}
Expand the left side and write the first reader variable as $u v$:
\[
 \sum_m\Gamma(m)\sum_{u,v,w}
 a_u b_v\beta_w\mathscr B_m(uvw;h).
\]
Grouping $u$ and $w$ by their product gives the coefficient
$(a*\beta)$ and yields the right side.
\end{proof}

\begin{corollary}[Linear factorization does not create a new estimate]
\label{cor:hard-linear-invisibility}
Before the first Cauchy, dispersion, or spectral step, displaying a reader as
$a*b$ gives no new linear information: one factor can be absorbed exactly
into the complementary reader.  A gain from multiplicative structure must
keep the factors separate through the operation which creates the shifted
product relation.
\end{corollary}

This identity is the affine counterpart of the endpoint pullback audit.  It
explains why a prime-distribution theorem must be applied after restoring the
complete von Mangoldt reader, while a convolution theorem must preserve its
factors beyond the linear stage.

\subsection{Smooth prime-progression rows}

Let $\sigma\in\{-1,+1\}$.  Let
$\mathscr W=\mathscr W_{M,Q}$ be a smooth function supported where
\[
 m\asymp M,
 \qquad q\asymp Q,
 \qquad mq\asymp X,
\]
and suppose
\[
 |\partial_m^j\partial_q^k\mathscr W(m,q)|
 \ll_{j,k}M^{-j}Q^{-k}(\log X)^C
 \tag{M.9}
\]
for every fixed $j,k$.  We require the prime reader itself to lie in the
fixed dyadic interval,
\[
 X\le mq+2\sigma\le2X
\]
on the support.  Define
\[
 \boxed{
 \mathcal E_{\sigma,\mathscr W}(m)
 =\sum_{q\ge1}(\log q)\mathscr W(m,q)
  [\Lambda(mq+2\sigma)-\gamma_2(m)].}
 \tag{M.10}
\]
The sign $-1$ is the original orientation and $+1$ the reflected one.

For fixed $m$, set
\[
 G_{\sigma,m}(t)
 =\log\!\left(\frac{t-2\sigma}{m}\right)
  \mathscr W\!\left(m,\frac{t-2\sigma}{m}\right).
 \tag{M.11}
\]
On the support,
\[
 \norm{G_{\sigma,m}}_\infty
 +X\norm{G'_{\sigma,m}}_\infty
 \ll(\log X)^C.
 \tag{M.12}
\]

\begin{lemma}[Discrete local row and prime-density integral]
\label{lem:hard-local-integral}
Uniformly for odd $m\asymp M$,
\[
 \boxed{
 \gamma_2(m)\sum_q(\log q)\mathscr W(m,q)
 =\frac1{\varphi(m)}\int_\R G_{\sigma,m}(t)\,\dd t
 +O\!\left(\frac{m}{\varphi(m)}(\log X)^C\right).}
 \tag{M.13}
\]
If $M\le X^{17/33}$, the total lattice error after summing over
$m\asymp M$ is $O_A(X(\log X)^{-A})$ for every fixed $A$.
\end{lemma}

\begin{proof}
Euler summation gives
\[
 \sum_q(\log q)\mathscr W(m,q)
 =\int_0^\infty(\log u)\mathscr W(m,u)\,\dd u
  +O((\log X)^C).
\]
Use $\gamma_2(m)=m/\varphi(m)$ and the change of variables
$t=mu+2\sigma$.  Finally $m/\varphi(m)\ll\log\log(3m)$ and
$X^{17/33}(\log X)^{C+1}=o_A(X(\log X)^{-A})$.
\end{proof}

\subsection{Convenient-divisor closure}

Write $\pi(x;m,h)$ for the number of primes $p\le x$ with
$p\equiv h\pmod m$.

\begin{theorem}[Maynard fixed-residue convenient-divisor input]
\label{thm:hard-Maynard-input}
Let $h\ne0$, $0<\delta<1/42$, and
\[
 0<\eta<\frac{1-42\delta}{4}.
\]
Let $\mathscr Q_{\delta,\eta}(x)$ be the moduli $m\le x^{1/2+\delta}$
which possess a divisor $e\mid m$ satisfying
\[
 x^{2\delta+\eta}<e<
 \min\{x^{1/10-7\delta/5-\eta},
       x^{1/2-19\delta-\eta}\}.
 \tag{M.14}
\]
Then, for every $A>0$,
\[
 \boxed{
 \sum_{\substack{m\in\mathscr Q_{\delta,\eta}(x)\\(m,h)=1}}
 \left|\pi(x;m,h)-\frac{\pi(x)}{\varphi(m)}\right|
 \ll_{h,\delta,\eta,A}\frac{x}{(\log x)^A}.}
 \tag{M.15}
\]
The estimate may be subtracted at two endpoints whenever the modulus bank
belongs to the convenient-divisor set at both endpoints.
\end{theorem}

\begin{proof}
This is the convenient-divisor corollary of Maynard's fixed-residue theorem
\cite{MaynardFixed}.  The endpoint-subtraction statement is immediate.
\end{proof}

\begin{corollary}[Weighted von Mangoldt form]
\label{cor:hard-Maynard-Lambda}
Let $F_m$ be smooth functions supported in $[x,2x]$ with
\[
 \norm{F_m}_\infty+x\norm{F_m'}_\infty\ll(\log x)^C.
\]
Let $\mathscr A$ be a modulus bank satisfying
\[
 \mathscr A\subseteq
 \bigcap_{x\le y\le2x}\mathscr Q_{\delta,\eta}(y).
 \tag{M.16a}
\]
Then
\[
 \boxed{
 \sum_{\substack{m\in\mathscr A\\(m,h)=1}}
 \left|
  \sum_{n\equiv h\ (m)}\Lambda(n)F_m(n)
  -\frac1{\varphi(m)}\int F_m(t)\,\dd t
 \right|
 \ll_{h,\delta,\eta,A}\frac{x}{(\log x)^A}.}
 \tag{M.16}
\]
\end{corollary}

\begin{proof}
Apply \textup{(M.15)} with endpoint $y$ for every $y\in[x,2x]$; the
uniform membership condition \textup{(M.16a)} keeps the same modulus bank
admissible throughout the partial-summation integral.  Partial summation
then gives the logarithmically weighted prime sum, and the prime number
theorem converts its global main term to the Lebesgue integral.  Proper
prime powers contribute $x^{1/2+o(1)}$ after summing over all divisor moduli
of the shifted endpoint.  Increasing the input logarithmic exponent absorbs
the smooth variation and this error.
\end{proof}

Take
\[
 \delta=\frac1{66},
 \qquad
 \eta_*=\frac1{10000},
 \tag{M.17}
\]
and work on an upper balanced dyadic block satisfying
\[
 M=X^{17/33}(\log X)^{-O(1)},
 \qquad
 Q\asymp X/M=X^{16/33}(\log X)^{O(1)},
 \qquad
 m\le X^{17/33}\text{ on the support}.
 \tag{M.18}
\]
The smaller upper exponent in \textup{(M.14)} is
\[
 \frac1{10}-\frac7{5}\frac1{66}=\frac{13}{165}.
\]
Define the strict interior window
\[
 \boxed{
 L_X=X^\lambda,
 \qquad U_X=X^\upsilon,
 \qquad
 \lambda=\frac1{33}+2\eta_*,
 \qquad
 \upsilon=\frac{13}{165}-2\eta_*.}
 \tag{M.19}
\]
Every divisor in $[L_X,U_X]$ lies inside Maynard's admissible interval.  Put
\[
 \boxed{
 \kappa=\upsilon-\lambda
 =\frac8{165}-4\eta_*.}
 \tag{M.20}
\]
Then
\[
 11\kappa>\frac{17}{33}.
 \tag{M.21}
\]

\begin{theorem}[Convenient-divisor closure of a balanced physical row]
\label{thm:hard-convenient-row}
Let $\mathscr C_X(M)$ be any set of odd moduli $m\asymp M$ possessing a
divisor in $[L_X,U_X]$.  At the scale \textup{(M.18)}, for either
$\sigma\in\{-1,+1\}$,
\[
 \boxed{
 \sum_{m\in\mathscr C_X(M)}
 \abs{\mathcal E_{\sigma,\mathscr W}(m)}
 \ll_A\frac{X}{(\log X)^A}.}
 \tag{M.22}
\]
Bounded source-fixed outer coefficients, including $\mu(m)$, may be inserted.
\end{theorem}

\begin{proof}
For fixed $m$,
\[
 \sum_q(\log q)\mathscr W(m,q)\Lambda(mq+2\sigma)
 =\sum_{n\equiv2\sigma\ (m)}\Lambda(n)G_{\sigma,m}(n).
\]
The residue is reduced for odd $m$.  If $X\le y\le2X$, then for all
sufficiently large $X$ the margin in \textup{(M.19)} gives
\[
 y^{2\delta+\eta_*}<L_X\le e\le U_X
 <y^{13/165-\eta_*}
 \le
 \min\{y^{1/10-7\delta/5-\eta_*},
          y^{1/2-19\delta-\eta_*}\}.
\]
Moreover $m\le X^{17/33}\le y^{17/33}$.  Hence every modulus in the
stated bank belongs to
$\bigcap_{X\le y\le2X}\mathscr Q_{\delta,\eta_*}(y)$.
Apply \Cref{cor:hard-Maynard-Lambda}, and use
\Cref{lem:hard-local-integral} to replace the integral by the discrete local
row.  The total replacement error has size $X^{17/33}(\log X)^{O(1)}$ and is smaller
than every $X(\log X)^{-A}$.
\end{proof}

The order of operations is load bearing: the theorem controls the complete
von Mangoldt occurrence in one residue class.  It does not control an
arbitrary dyadic component obtained after reopening that occurrence by
$\Lambda=\mu*\log$.

\subsection{Canonical hard moduli}

\begin{definition}[Maynard-hard critical moduli]
\label{def:hard-moduli}
Let
\[
 \boxed{
 \mathscr H_X(M)
 =\{m\asymp M:
   m\text{ is odd and squarefree, and no }e\mid m
   \text{ lies in }[L_X,U_X]\}.}
 \tag{M.23}
\]
Define the hard prime-progression current
\[
 \boxed{
 \mathsf H_{\sigma}(M,Q;\mathscr W)
 =\sum_{m\in\mathscr H_X(M)}
   \mu(m)\mathcal E_{\sigma,\mathscr W}(m).}
 \tag{M.24}
\]
\end{definition}

\begin{theorem}[Convenient rows disappear before positivity]
\label{thm:hard-prime-reduction}
At the balanced scale,
\[
 \boxed{
 \sum_{m\asymp M}\mu(m)\mathcal E_{\sigma,\mathscr W}(m)
 =\mathsf H_{\sigma}(M,Q;\mathscr W)
  +O_A\!\left(\frac{X}{(\log X)^A}\right).}
 \tag{M.25}
\]
The removed modulus bank is controlled absolutely.
\end{theorem}

\begin{proof}
Partition the odd squarefree moduli according as they do or do not possess a
divisor in $[L_X,U_X]$, and apply
\Cref{thm:hard-convenient-row} to the first bank.
\end{proof}

For $n>1$, let $P^-(n)$ denote its least prime factor.

\begin{definition}[Maximal subwindow divisor]
\label{def:hard-maximal-divisor}
For $m\in\mathscr H_X(M)$, put
\[
 \boxed{
 r(m)=\max\{r:r\mid m,\ r<L_X\},
 \qquad s(m)=m/r(m).}
 \tag{M.26}
\]
\end{definition}

\begin{theorem}[Rigid divisor-gap compiler]
\label{thm:hard-gap}
For every $m\in\mathscr H_X(M)$:
\begin{enumerate}[label=\textup{(G\arabic*)},leftmargin=3.2em]
\item $r(m)<L_X$, $s(m)>1$, $(r(m),s(m))=1$, and
      $\mu(m)=\mu(r(m))\mu(s(m))$;
\item every divisor $e\mid s(m)$ with $e>1$ satisfies
      $r(m)e>U_X$;
\item
      \[
       \boxed{
       P^-(s(m))>\frac{U_X}{r(m)}
       >\frac{U_X}{L_X}=X^\kappa;}
       \tag{M.27}
      \]
\item for all sufficiently large $X$,
      \[
       \boxed{1\le\omega(s(m))\le10.}
       \tag{M.28}
      \]
\end{enumerate}
The map $m\mapsto(r(m),s(m))$ is canonical and injective.
\end{theorem}

\begin{proof}
Squarefreeness gives the first statement.  If $e>1$ divides $s(m)$, then
$r(m)e$ is a divisor of $m$ larger than $r(m)$.  It cannot be below $L_X$ by
maximality, and it cannot lie in $[L_X,U_X]$ by hardness; hence it exceeds
$U_X$.  Taking $e=P^-(s(m))$ proves \textup{(M.27)}.  If
$j=\omega(s(m))$, then $m\ge s(m)>X^{\kappa j}$.  Since
$m=X^{17/33+o(1)}$ and \textup{(M.21)} holds, $j\le10$ for large $X$.
The maximum in \textup{(M.26)} is unique.
\end{proof}

\begin{corollary}[Finite-depth hard current]
\label{cor:hard-depth}
One has
\[
 \boxed{
 \mathsf H_{\sigma}(M,Q;\mathscr W)
 =\sum_{j=1}^{10}
  \sum_{\substack{m\in\mathscr H_X(M)\\\omega(s(m))=j}}
  \mu(r(m))\mu(s(m))
  \mathcal E_{\sigma,\mathscr W}(m).}
 \tag{M.29}
\]
The depth-one layer $s(m)=p$ contains all prime moduli, through $r(m)=1$,
and the canonical near-prime moduli $m=rp$ with $r<L_X$.
\end{corollary}

\begin{proof}
Partition \textup{(M.24)} by $j=\omega(s(m))$ and use
\Cref{thm:hard-gap}.
\end{proof}

On the prime-modulus base layer, the $+2$ orientation is
\[
 \boxed{
 -\sum_{\substack{p\asymp M\\ p\ {\rm prime}}}
  \sum_{q\asymp Q}(\log q)\mathscr W(p,q)
  \left[\Lambda(pq+2)-\frac{p}{p-1}\right].}
 \tag{M.30}
\]
Its actual part is a literal bilinear prime--prime occurrence: the outer
prime is the hard modulus and the inner von Mangoldt factor reads $pq+2$.
Opening either prime occurrence returns to a deeper factorization chart; it
does not remove the parity core.

\begin{criterion}[Hard-modulus affine closure]
\label{crit:hard-affine-closure}
Let a source-independent smooth partition decompose one reflected outer
residual as
\[
 R_X(2)
 =\sum_{\nu\in\mathcal B_X}
   \sum_{m\asymp M_\nu}
    \mu(m)\mathcal E_{+,\mathscr W_\nu}(m)
  +\mathfrak R_X^{\rm rem}+o(X),
 \tag{M.31}
\]
where $|\mathcal B_X|=X^{o(1)}$, every displayed block has
$M_\nu=X^{17/33}(\log X)^{-O(1)}$ and is supported on
$m\le X^{17/33}$; the remaining dyadic ranges are included in
$\mathfrak R_X^{\rm rem}$.  If
\[
 \mathfrak R_X^{\rm rem}=o(X)
 \tag{M.32}
\]
and
\[
 \sum_{\nu\in\mathcal B_X}
 \mathsf H_+(M_\nu,X/M_\nu;\mathscr W_\nu)=o(X),
 \tag{M.33}
\]
then $R_X(2)=o(X)$, the Hardy--Littlewood twin asymptotic holds, and there
are infinitely many twin primes.
\end{criterion}

\begin{proof}
Apply \Cref{thm:hard-prime-reduction} to every balanced block.  The absolute
convenient-modulus errors sum to $o(X)$ because the block count is
subpower.  Insert \textup{(M.32)}--\textup{(M.33)} into \textup{(M.31)},
then use the prime-power cleanup already proved in
\Cref{sec:prime-power-cleanup}.
\end{proof}

This criterion deliberately leaves the nonbalanced smooth ranges explicit.
The rigid-modulus theorem is an analytic pruning result for the balanced
prime-reader frontier, not a hidden proof that every other dyadic range is
small.

\subsection{Exact separation from the fixed-character GRH reader}

\begin{theorem}[No source-native Maynard transfer to GRH]
\label{thm:hard-no-transfer-GRH}
The convenient-divisor theorem used above has no source-preserving
specialization to either the short-quotient GRH packet of
\Cref{sec:grh-hyperbola} or the rectangular packet of
\Cref{sec:grh-rectangular}.  Its essential data are a varying divisor modulus
$m$, a fixed progression residue $\pm2\pmod m$, a complete prime reader, and
summation over $m$ before absolute values.  A fixed primitive character
supplies none of these progression-modulus occurrences.
\end{theorem}

\begin{proof}
In the affine compiler, one endpoint is opened into a divisor $m$ while the
other endpoint remains a von Mangoldt reader.  This produces the congruence
$n\equiv\pm2\pmod m$ and the modulus average in
\Cref{thm:hard-Maynard-input}.  The GRH packets contain one fixed character
and no second affine endpoint whose divisor could supply a varying modulus.
Adjoining such a modulus would enlarge the arithmetic source rather than
reparametrize it.
\end{proof}

The common M\"obius--log coefficient therefore survives, but the first
noncircular analytic readers differ.  In the GRH branch it is a fixed-character
Mertens current, a frequency-resolved two-M\"obius rectangle, a completed
theta packet, or a localized Weil form.  In the affine branch it can be
restored to a prime progression and pruned to the finite-depth hard-modulus
current \textup{(M.29)}.

\part{The common arithmetic spine}
% =============================================================================

\section{Exact linkage of the GRH and affine-prime compilers}
\label{sec:common-spine}

\subsection{One universal factorization identity}

\begin{theorem}[Common outer-compiler theorem]
\label{thm:common-outer}
Let $I_X$ be a finite endpoint interval and let $B_X(m)$ be any predeclared
arithmetic occurrence row.  Define
\[
 \mathcal O_X(B)=\sum_{m\in I_X}B_X(m)\Lambda(m).
\]
Then
\[
 \mathcal O_X(B)
 =\sum_{q\ge2}(\log q)
   \sum_{\substack{d\ge1\\dq\in I_X}}\mu(d)B_X(dq).
\]
The formula factors through the quotient carrier $q\ge2$ and therefore has
zero projection onto the direct $q=1$ endpoint section.
\end{theorem}

\begin{proof}
This is \Cref{thm:universal-outer,cor:q1-annihilation}.
\end{proof}

Two choices of $B_X$ recover the two branches:
\[
 B_X(m)=\chi(m)\one_{m\le X}
\]
produces the Dirichlet Chebyshev compiler after complete multiplicativity and
summation in $d$, whereas
\[
 B_X(m)=\Lambda(m-2)\one_{m\in\mathscr I_X^+}
\]
produces the twin compiler.  Subtracting the appropriate local expectation is
a linear operation on the same occurrence carrier.

\begin{corollary}[Shared vanishing top-birth section]
Any source assembly, nuisance short, endpoint fusion, conductor-block short,
or pressure localization acting entirely on $q\ge2$ cannot recreate the
direct terminal coordinate.  A nonzero top-coordinate birth can occur only if
a later construction explicitly adjoins $q=1$ or an equivalent terminal read.
\end{corollary}

\begin{proof}
Every admitted operation factors through the fixed quotient space
$\ell^2\{2,\ldots,Q_X\}$.  Linear maps and orthogonal shorts cannot enlarge
the set of endpoint coordinates through which the source factors.
\end{proof}

\subsection{Common factorization-to-endpoint short}

Let $\mathcal O_X$ be the divisor--quotient occurrence set and let
$F_X:\ell^2(\mathcal O_X)\to\ell^2(I_X)$ be endpoint summation.  Every
protected GRH or affine read in the outer compiler factors through $F_X$ after
its public character or occurrence multiplier is installed.

\begin{theorem}[Read-null innovation theorem]
\label{thm:common-read-null}
For every factorization vector $z$,
\[
 \norm z^2
 =\norm{P_{\rm end}z}^2+\norm{(I-P_{\rm end})z}^2,
\]
while every endpoint read satisfies
\[
 \ip{F_X^*c}{z}=\ip{F_X^*c}{P_{\rm end}z}.
\]
Consequently, a raw quotient norm may have a positive kinematic floor even
when the protected signed endpoint is zero.  Replacing the signed closure
certificate by a raw energy estimate is valid only if the stronger estimate
is actually proved with the complete cross-fibre Gram retained.
\end{theorem}

\begin{proof}
Apply the orthogonal endpoint projection constructed in
\Cref{sec:arithmetic-loading}.  The second summand belongs to $\Ker F_X$ and
is invisible to every endpoint read.
\end{proof}

This theorem explains both the hyperbola--Volterra floor in the GRH quotient
profile and the quotient-fibre innovation in the affine packet.  In each case
the positive floor is genuine source information but is not the protected
signed scalar.

\subsection{Common birth--retuning chronology}

\begin{theorem}[Retuning-first principle for both branches]
\label{thm:common-retuning}
Consider any cutoff or parameter deformation of either outer packet.
\begin{enumerate}[label=\textup{(R\arabic*)},leftmargin=3.2em]
\item If the deformation enlarges the source range, its protected influence
      changes by signed old-sector retuning minus a nonnegative carrier-birth
      innovation as in \Cref{thm:gran-birth-retuning}.
\item If the source range is constant and the deformation is implemented by
      invertible diagonal transport, the carrier-birth innovation is zero and
      the deformation is pure retuning.
\item If the deformation begins from the zero source, every nonzero protected
      read has infinite influence at the endpoint and no regular positive
      action is anchored there.
\end{enumerate}
\end{theorem}

\begin{proof}
The first statement is the old/detail Schur formula.  For the second, an
invertible transport maps the source range bijectively onto itself, so there is
no new detail quotient.  For the third, the zero Gram has range $\{0\}$, and
the influence definition is infinite off that range.
\end{proof}

In the Dirichlet compiler, removing $q=1$ makes the circular top coordinate a
zero-birth section; conductor blocks then separate deterministic transport
from a mean-zero signed current.  In the affine compiler, hard-threshold
motion retunes the endpoint writer, while positive complement pressure
transports the full quotient source by invertible diagonal multipliers and
therefore has zero birth.

\subsection{Structural comparison}

\begin{center}
\small
\begin{tabularx}{\textwidth}{@{}L{0.19\textwidth}L{0.36\textwidth}L{0.36\textwidth}@{}}
\toprule
\textbf{Layer} & \textbf{Dirichlet/GRH branch} & \textbf{Affine-prime branch}\\
\midrule
Protected endpoint
& $\Psi_\chi(X)$, $M_\chi(K)$, localized Weil forms, or completed critical zeros
& $R_X(2)$, a fixed affine amplitude, or a moving Goldbach amplitude\\[0.35em]
Common outer row
& $\chi(q)\log q\,[M_\chi(X/q)-M_\chi(D_X)]$
& $\log q\sum_{d>D_X}\mu(d)[\Lambda(dq-2)-\gamma_2(d)]$\\[0.35em]
Direct section
& $q=1$ absent
& $q=1$ absent in both reflected orientations\\[0.35em]
Independent entrance
& phase action; two-M\"obius rectangle; theta heat; completed Weil and Pick charts
& $W$-fibre amplitude; endpoint pressure; restored fixed-residue prime reader\\[0.35em]
Positive nonminimal sector
& hyperbola floor; quotient-one block; divergent raw Euler clock
& quotient-fibre, occurrence, and sequential provenance innovation\\[0.35em]
Exact analytic pruning
& divisor head; high reader; product diagonal; short collisions; frequency split; high Legendre modes
& square-root divisor head; non-odd--odd pressure sectors; convenient-divisor moduli\\[0.35em]
First remaining frontier
& short quotient; reciprocal-phase/sinc currents; Jensen escape; Schur sign; backward heat; first-prime continuation; normalized compensation
& odd--odd pressure current; protected amplitude; finite-depth hard prime progression\\[0.35em]
Non-transfer point
& fixed character has no varying prime-progression modulus
& Maynard pruning requires the restored prime endpoint and modulus average\\
\bottomrule
\end{tabularx}
\end{center}

\subsection{Combined closure theorem}

\begin{theorem}[Arithmetic Renewal Gran closure theorem]
\label{thm:combined-closure}
Assume that all source maps, shorts, occurrence fusions, reader restorations,
and cutoff transports are constructed on one compatible arithmetic Renewal
Gran tensor as in \Cref{sec:gran-cylinder,sec:arithmetic-loading}.  Then:
\begin{enumerate}[label=\textup{(C\arabic*)},leftmargin=3.2em]
\item subpower extensive source-fixed phase action in every primitive channel
      implies GRH;
\item a square-root calibrated short-quotient head implies GRH in every fixed
      nonprincipal channel; full GRH follows after an independent principal
      closure or a square-root estimate for its explicit head and tail;
\item a square-root protected rectangular current implies GRH in every fixed
      nonprincipal channel.  After the proved high-reader, diagonal, and
      short-collision removals, it is enough to control the coherently
      assembled nonzero reciprocal-phase current and zero-frequency sinc
      sheet;
\item exclusion of cofinal Jensen escape together with generic nonpositive
      collision curvature implies vanishing of the channel
      de Bruijn--Newman constant and hence GRH; the pole-removed zeta
      completion gives the principal analogue;
\item nonnegativity of the finite protected Weil--Schur matrix at every
      compact support wall, after choosing an endpoint-independent cutoff with a
      positive harmonic high-mode complement, implies GRH; the principal
      completion has the same form;
\item convergence of the fixed-packet protected affine amplitude to one at a
      subpolynomial regulator gives the Hardy--Littlewood asymptotic, while a
      positive lower bound gives genuine prime pairs;
\item the corresponding regulator-uniform moving reflection amplitude gives
      the strong Goldbach conclusion;
\item the prime-cleaned odd--odd complement-pressure certificate gives the
      twin Hardy--Littlewood asymptotic and infinitely many twin primes;
\item after a source-independent smooth decomposition, an $o(X)$ nonbalanced
      remainder and $o(X)$ total finite-depth hard prime-progression current
      give the same twin conclusion, because every convenient-divisor
      balanced modulus has already been removed absolutely;
\item a backward relative-form estimate for the complete Gaussian Weil
      evolution implies a zero heat-positivity threshold and GRH in each
      primitive nonprincipal channel;
\item positivity of every finite Pick matrix on a fixed dense Euler-region
      sample, for both actual completed conjugate channels, implies paired
      GRH.  Positivity at finitely many stages alone is not this hypothesis;
\item arithmetically admitted target-normalized filters satisfying
      \Cref{crit:filter-cofinal-compensation} exclude every hypothetical
      off-line zero and hence imply GRH.  Cancellation of a finite spectral
      prefix without the completed remainder estimate is not sufficient.
\end{enumerate}
None of these implications requires a positive action to control every
ambient quotient, factorization, matrix, or spectral direction.  It is
sufficient to control the actual protected read, provided the source is fixed
independently, the complete mixed Gram is retained, and an analytic theorem is
applied to the reader it actually assumes.
\end{theorem}

\begin{proof}
Items~\textup{(C1)} and \textup{(C2)} are
\Cref{thm:intro-entropic-GRH,thm:intro-short-GRH}.  Item~\textup{(C3)} is
\Cref{thm:rectangular-corner,crit:freq-rectangular-closure}.  Item~\textup{(C4)}
is \Cref{thm:theta-DN,crit:theta-Jensen-GRH}.  Item~\textup{(C5)} is
\Cref{thm:weil-high-mode,crit:weil-Schur-GRH}.  Items~\textup{(C6)} and
\textup{(C7)} follow from the exact $W$-fibre identity, singular-tail control,
and prime-power cleanup.  Item~\textup{(C8)} is the pressure integral
identity, entrance closure, parity reduction, and prime-power cleanup in
\Cref{sec:affine-pressure}.  Item~\textup{(C9)} is
\Cref{thm:hard-prime-reduction,crit:hard-affine-closure}.
Item~\textup{(C10)} is \Cref{crit:gauss-relative-flow}; item~\textup{(C11)}
is \Cref{thm:pick-GRH,thm:pick-exhaustive}.  Item~\textup{(C12)}
is \Cref{thm:filter-compensation,crit:filter-cofinal-compensation}.
In item~\textup{(C5)},
\Cref{thm:acq-schur-bracket} supplies the finite-data sufficient test
$L_N-d_N^{-1}R_N\succeq0$ at every relevant support.
The final statement
is the Riesz protection inequality together with
\Cref{thm:arithmetic-expansion-pullback,thm:common-read-null,lem:hard-convolution-absorption}.
\end{proof}

\begin{remark}[What remains analytic]
The tensor construction removes ambiguity about source choice, occurrence,
reader restoration, factorization fusion, cutoff birth, completion frequency,
and the order in which external estimates may be used.  The unresolved
statements are the displayed signed short-quotient, reciprocal-phase,
sinc-sheet, protected-amplitude, finite low-mode Schur, hard-prime,
collision-curvature, backward relative-heat, first-prime continuation,
normalized completed-compensation, or cofinal-compactness estimates.  The one-pair padding and total-height bounds
restrict finite certification architectures; the mixed-block Appell condition
is an additional finite-algebraic certificate, not a GRH theorem.
The arbitrary-padding two-pair reduction likewise retains its full
coefficient-certificate hypotheses.  Finite negative values of a
particular filter and logarithmic escape of a Gaussian symbol do not
establish a uniform signed transfer to the finite beta kernel.
Positivity of a
complete Gram, finiteness of a reader alphabet, and coercivity of the high
Legendre modes are representation and reduction theorems, not cancellation or
low-mode positivity theorems.
\end{remark}

\section{Conclusion}
\label{sec:conclusion}

The arithmetic Renewal Gran tensor supplies a common finite language for two
apparently different prime-distribution problems.  Its decisive operations are
complete mixed-source assembly, target-free M\"obius--log compilation,
reader restoration, endpoint fusion, and the separation of signed retuning
from nonnegative carrier birth.  The expansion--pullback theorem sharpens this
language: an exact factorization or determinant expansion of a terminal square
is still the same protected endpoint after fusion and cannot be declared an
independent coercive source.

For GRH, five complementary entrances are now explicit.  The phase route
identifies a source-fixed entropy at the exact extensive square-root scale.
The Chebyshev route removes the direct target, closes the nonprincipal divisor
head, exposes the raw quotient energy floor, and removes the terminal quotient
range.  Direct two-threshold M\"obius descent creates a finite-state
rectangular current in which the high reader, full same-product diagonal, and
short product collisions close.  Exact completion then shows that its remote
frontier is not one undifferentiated determinant norm: it is the sum of a live
nonzero reciprocal-frequency current and a phase-free zero-frequency sinc
sheet.  The completed theta route is cutoff independent and converts a
positive de Bruijn--Newman boundary into either a finite Jensen collision with
signed curvature or a cofinal escape of finite witnesses.  The completed
Euler--Weil route supplies strict finite Schur cells and Hermite--Biehler
spectra, while exact archimedean resummation makes the logarithmic high modes
coercive and reduces every fixed support wall to a finite protected Schur
matrix.

These reductions now retain the quantitative price of acquisition.  A
normalized Mellin bank is visible on a growing strip, but the complete
multiplier and Vandermonde conditioning remain separate.  A state-preserving
predictor bounds the upper score defect only through its actual mixed-Gram
residual.  Exact Legendre entries make the Schur obstruction accessible to
finite integration, while support-wall concentration and prime-log resonance
explain why a small fixed block is not an ambient norm estimate.  Gaussian
averaging provides an independent description of the same completed sign
problem: a finite heat threshold, explicit contact geometry, and a backward
relative inequality sufficient to force positivity at zero time.  The Pick
chart identifies the first-prime continuation as the remaining coefficient
layer and makes the distinction between finite interpolation and exhaustive
positivity exact.

The finite multiplier analysis has a different role.  Its universal Hermite
limit gives sharp one-pair and quartic padding bounds and exact multi-pair
height budgets.  The remaining mixed-block test is the absence of a real
common zero of consecutive Appell polynomials.  The signed Gaussian symbol
makes this condition explicit without turning numerical root separation,
positivity of a core weight, or a symmetric special case into a uniform
exclusion theorem.  Exact pairing expansions show why separate one-pair
certificates do not multiply.  The critical-point identity converts a
Gaussian symbol into one heat-evolved real-root problem, and the growing
gamma-product formula gives a logarithmically translating limit.  This
last statement concerns the Gaussian symbol itself; an increasing-order
transfer to the arithmetic beta kernel remains distinct.

The spectral-filter analysis makes the quantitative obstruction sharper.
Finite pole, gamma, and zero data can be interpolated with a strictly
positive tail, but target normalization forces a coefficient cost.
Arithmetic admission then forces a positive completed compensation of
at least the target multiplicity.  Exact residues, real-factor last-zero
bounds, and finite rational negative values restrict specific equispaced
architectures without confusing their failure with exclusion of the
hypothetical zero.  The remaining criterion controls the actual signed
remainder, not the size of a discarded spectral tail or an absolute norm
that ignores the normalization cost.

For affine prime pairs, the local tensor closes at two form slots, the
selector-stable parity quotient is one-dimensional, and the $W$-fibre
identity separates local density from the protected pair amplitude.  In the
twin specialization, the target-free packet fuses to endpoints and admits a
zero-birth complement pressure.  Restoring the reflected prime reader adds a
second analytic chart: Maynard's fixed-residue theorem removes every balanced
modulus with a divisor in an explicit interval, and the surviving modulus has
a canonical prefix and a rough cofactor of depth at most ten.  The prime
modulus is not a nuisance endpoint; it is the base bilinear prime--prime
occurrence of the hard current.

The exact link is not an implication between GRH and twin primes.  It is a
shared arithmetic compiler and a shared proof discipline.  Both branches
annihilate the direct quotient, short read-null factorization energy, and
require a source fixed before the protected coefficient is known.  They then
separate at their first genuinely analytic readers: fixed-character Mertens,
reciprocal-phase, completed-theta, and localized Weil currents on the GRH
side, and prime-progression, pressure, or protected-amplitude currents on the
affine side.  The absence of a
source-native Maynard transfer to GRH is part of this separation, not a defect
of the common tensor.

All finite identities and all conditional implications stated here are exact.
The remaining programme is therefore expressed by explicit quantitative
certificates, several of which must hold cofinally: prove a protected GRH
short-quotient, reciprocal-phase, sinc-sheet, Jensen-collision, low-mode
Schur, backward relative-heat, coefficient-continuation, or normalized
spectral-compensation estimate, and prove an affine pressure, amplitude,
or hard-prime certificate.  None may
replace the protected scalar by an ambient norm, erase the acquisition cost,
or change the reader after the analytic theorem has been chosen.

\small
\begin{thebibliography}{99}

\bibitem{CravenCsordasFinite}
T.~Craven and G.~Csordas,
\newblock Location of zeros. Part I: Real polynomials and entire functions,
\newblock \emph{Illinois J. Math.} \textbf{27} (1983), Theorem~3.1.

\bibitem{CravenCsordas}
T.~Craven and G.~Csordas,
\newblock Jensen polynomials and the Tur\'an and Laguerre inequalities,
\newblock \emph{Pacific J. Math.} \textbf{136} (1989), 241--260.

\bibitem{deBruijn}
N.~G.~de Bruijn,
\newblock The roots of trigonometric integrals,
\newblock \emph{Duke Math. J.} \textbf{17} (1950), 197--226.

\bibitem{Dobner}
A.~Dobner,
\newblock A proof of Newman's conjecture for the extended Selberg class,
\newblock \emph{Acta Arith.} \textbf{201} (2021), 29--62.

\bibitem{MaynardFixed}
J.~Maynard,
\newblock Primes in arithmetic progressions to large moduli I: Fixed residue classes,
\newblock \emph{Mem. Amer. Math. Soc.} \textbf{306} (2025), no.~1542.

\bibitem{Newman}
C.~M.~Newman,
\newblock Fourier transforms with only real zeros,
\newblock \emph{Proc. Amer. Math. Soc.} \textbf{61} (1976), 245--251.


\bibitem{Bombieri2000}
E.~Bombieri,
\newblock Remarks on Weil's quadratic functional in the theory of prime numbers. I,
\newblock \emph{Atti Accad. Naz. Lincei Cl. Sci. Fis. Mat. Natur. Rend.
Lincei (9) Mat. Appl.} \textbf{11} (2000), no.~3, 183--233.

\bibitem{Lagarias2004}
J.~C.~Lagarias,
\newblock Li coefficients for automorphic $L$-functions,
\newblock \emph{Ann. Inst. Fourier (Grenoble)} \textbf{57} (2007),
1689--1740.

\bibitem{DFI97}
W.~Duke, J.~Friedlander, and H.~Iwaniec,
\newblock Bilinear forms with Kloosterman fractions,
\newblock \emph{Invent. Math.} \textbf{128} (1997), no.~1, 23--43.

\bibitem{BettinChandee}
S.~Bettin and V.~Chandee,
\newblock Trilinear forms with Kloosterman fractions,
\newblock \emph{Adv. Math.} \textbf{328} (2018), 1234--1262.

\bibitem{FouvryRadziwill}
E.~Fouvry and M.~Radziwi\l\l{},
\newblock Level of distribution of unbalanced convolutions,
\newblock \emph{Ann. Sci. \'Ec. Norm. Sup\'er.} \textbf{55} (2022),
no.~2, 537--568.

\bibitem{Bombieri1965}
E.~Bombieri,
\newblock On the large sieve,
\newblock \emph{Mathematika} \textbf{12} (1965), 201--225.

\bibitem{Davenport}
H.~Davenport,
\newblock \emph{Multiplicative Number Theory},
\newblock 3rd ed., revised by H.~L.~Montgomery, Springer, 2000.

\bibitem{HardyLittlewood}
G.~H.~Hardy and J.~E.~Littlewood,
\newblock Some problems of ``Partitio Numerorum'' III: On the expression of a
number as a sum of primes,
\newblock \emph{Acta Math.} \textbf{44} (1923), 1--70.

\bibitem{IK}
H.~Iwaniec and E.~Kowalski,
\newblock \emph{Analytic Number Theory},
\newblock American Mathematical Society, 2004.

\bibitem{MV}
H.~L.~Montgomery and R.~C.~Vaughan,
\newblock \emph{Multiplicative Number Theory I: Classical Theory},
\newblock Cambridge University Press, 2007.

\bibitem{Vinogradov}
A.~I.~Vinogradov,
\newblock The density hypothesis for Dirichlet $L$-series,
\newblock \emph{Izv. Akad. Nauk SSSR Ser. Mat.} \textbf{29} (1965), 903--934.

\end{thebibliography}

\end{document}
